\documentclass{situationist-preprint}
\usepackage{longtable,array}

\title{Gödel's Ontological Argument Without Modal Collapse}
\msid{REL-01-c}
\msfamily{REL-01}
\msdiscipline{Philosophy of Religion}
\mstype{Alternative proof}
\mskeywords{ontological argument; Gödel; positive properties; modal collapse; Anderson's emendation; proper filter; finite models; natural deduction}

% generated by gen/generate.py -- do not edit
\newcommand{\csModels}{2{,}025}
\newcommand{\csSym}{816}
\newcommand{\csNon}{1{,}209}
\newcommand{\csSymColl}{231}
\newcommand{\csSymNoColl}{585}
\newcommand{\csSymT}{317}
\newcommand{\csSymTColl}{95}
\newcommand{\csNonFail}{407}
\newcommand{\csNonC}{366}
\newcommand{\csNonCFail}{0}
\newcommand{\csNonNotC}{843}
\newcommand{\csNonNotCFail}{407}
\newcommand{\csNonNotCOk}{436}
\newcommand{\csNonOk}{802}
\newcommand{\csNonFailT}{338}
\newcommand{\csNonFailFour}{203}
\newcommand{\csNonFailTFour}{175}
\newcommand{\csUltra}{1{,}365}
\newcommand{\csScottModels}{1{,}365}
\newcommand{\csScottSym}{510}
\newcommand{\csScottSymColl}{0}
\newcommand{\csScottNonColl}{464}
\newcommand{\csUltraNotScott}{0}
\newcommand{\csManyGods}{342}
\newcommand{\csRejClosure}{839{,}019}
\newcommand{\csRejFTwo}{47{,}723}
\newcommand{\csRejGcomp}{0}
\newcommand{\csRejAThree}{3{,}832}
\newcommand{\csRejNEcomp}{0}
\newcommand{\csRejAFive}{0}
\newcommand{\csXModels}{2{,}022}
\newcommand{\csXFormulas}{23{,}001}
\newcommand{\csCandidates}{53{,}580}
\newcommand{\csLinesAf}{76}
\newcommand{\csRulesAf}{22}
\newcommand{\csLinesScott}{92}
\newcommand{\csRulesScott}{23}
\newcommand{\csLinesPositive}{21}
\newcommand{\csRulesPositive}{14}
\newcommand{\csLinesAnderson}{25}
\newcommand{\csRulesAnderson}{12}
\newcommand{\csLinesIndisc}{20}
\newcommand{\csRulesIndisc}{8}
\newcommand{\csLinesConverse}{17}
\newcommand{\csRulesConverse}{13}

\newcommand{\Pos}{\mathcal{P}}
\newcommand{\NE}{\mathit{NE}}
\newcommand{\Ess}{\mathrm{Ess}}
\newcommand{\GS}{G_{\!S}}

\begin{document}
\maketitle

\begin{abstract}
Gödel's ontological argument derives the necessary existence of a godlike being from axioms about positive properties, but in Scott's standard version its axioms entail modal collapse, the thesis that every truth is necessary. We give a variant that proves the conclusion without the collapse. Its positivity axioms say that the positive properties form a proper filter, closed under conjunction and necessary entailment and excluding the empty property, which is the hypothesis of Lemma~2417 of \emph{On the Existence of God}; its definitions of godlikeness, essence and necessary existence are Anderson's. We prove that necessarily something is godlike, and that some being is godlike, has godlikeness as its essence and exists necessarily, using of the accessibility conditions only symmetry. Relative to the main argument (REL-01-a), the variant avoids the possibility premise, which becomes a lemma. Modal collapse fails in models of the variant, including a full model with two worlds and one being. The computed part enumerates \csModels{} finite models up to isomorphism, checks the axioms and the conclusion in each and lists the countermodels to collapse, and gives the main derivations as natural-deduction proofs checked line by line by machine.
\end{abstract}

\section{Introduction}

\subsection{The problem}

Among ontological arguments, Gödel's is the one most often called a proof. It is stated in a second-order modal logic, its premises are a handful of axioms about a primitive notion of positive property, and its conclusion, that necessarily something is godlike, follows from them by steps that have now been checked by machine \citep{Benzmuller2014}. Its weakness lies elsewhere. \citet{Sobel1987} showed that the axioms, in the form Scott gave them, entail that whatever is true is necessarily true. This consequence, modal collapse, is unacceptable to almost everyone and in particular to theists, for whom the existence of the world and the choices of God are contingent. An argument whose premises entail modal collapse cannot be sound, since collapse is false; and if collapse were true, the argument's conclusion would be trivial, since necessary existence would differ from existence in nothing.

The question of this manuscript is whether the collapse is the price of Gödel's conclusion, or an accident of a particular formulation. We argue that it is an accident. The variant we give keeps Gödel's conclusion and Gödel's method, the derivation of a godlike being's necessary existence from axioms about positive properties, and does not entail collapse.

\subsection{Thesis}

The variant, which we call AF, has two components. Its axioms about positivity say that the positive properties form a proper filter: they are closed under conjunction and under necessary entailment, and the empty property is not positive. That is exactly the hypothesis of Lemma~2417 of \emph{On the Existence of God}, stated in Volume~X and analysed in REL-01-f. Its definitions of godlikeness, essence and necessary existence are those of \citet{Anderson1990}: a being is godlike if the properties it has necessarily are exactly the positive ones, and an essence of a being is a property that necessarily entails exactly the properties the being has necessarily. Godlikeness and necessary existence are positive. We prove:

\begin{enumerate}[label=(\roman*)]
\item In every model of AF on a symmetric frame, at every world, necessarily something is godlike, and some being is godlike, has godlikeness as an essence and has necessary existence (Theorem~\ref{t:main}).
\item AF does not entail modal collapse: there are models of AF, including full models, in which a proposition is true but not necessary (Proposition~\ref{p:nocollapse}).
\item Relative to the main argument of REL-01-a, AF avoids the possibility premise (Assumption~3.1 of REL-01-a), which it derives as Lemma~\ref{l:poss}, and of the S5 conditions (Assumption~2.1) it uses only symmetry.
\end{enumerate}

\subsection{Method}

Both claims of the thesis that admit mechanical checking are checked mechanically. The main theorem, and Sobel's collapse theorem for Scott's system, are given in Appendix~\ref{app:derivations} as derivations in a labelled natural-deduction system, written as data and checked line by line by a program. The absence of collapse, and the dependence of the theorem on symmetry, are exhibited by an exhaustive enumeration of finite models in Appendix~\ref{app:models}: all models of AF up to isomorphism with at most three worlds and at most six points (world--individual pairs) on arbitrary frames, together with larger models on symmetric and on reflexive frames, \csModels{} models in all. For each model the program checks the axioms, the conclusion, the corresponding Scott axioms and the collapse of every proposition of the form \enquote{something is $Y$}.

\section{Gödel's argument and its emendations}\label{s:history}

\subsection{Gödel and Scott}

Gödel's argument survives in a two-page note dated 1970, published in the third volume of his \emph{Collected Works} \citep{Godel1995}. Dana Scott reformulated it in 1972, and Scott's version, with an essence clause that Gödel's note omits, is the one most later discussions take as standard \citep{Sobel1987, Sobel2004, Fitting2002}. Scott's system has five axioms, about a second-order predicate $\Pos$ of properties (\enquote{$Y$ is positive}):

\begin{center}\small
\begin{tabular}{@{}ll@{}}
S1 & $\Pos(\lambda x.\neg Y(x)) \leftrightarrow \neg\Pos(Y)$ \\
S2 & $\Pos(Y) \wedge \nec\forall x(Y(x) \to Z(x)) \to \Pos(Z)$ \\
S3 & $\Pos(\GS)$ \\
S4 & $\Pos(Y) \to \nec\Pos(Y)$ \\
S5 & $\Pos(\NE_S)$
\end{tabular}
\end{center}

\noindent and three definitions: $x$ is godlike, $\GS(x)$, if it has every positive property; $Y$ is an essence of $x$, $\Ess_S(Y,x)$, if $Y(x)$ and every property of $x$ is necessarily entailed by $Y$; and $x$ has necessary existence, $\NE_S(x)$, if every essence of $x$ is necessarily instantiated. From these, in a modal logic in which $\pos\nec q \to q$ holds, it follows that necessarily something is godlike. Gödel's note has a conjunction axiom in place of S3 and an essence clause without the conjunct $Y(x)$; \citet{Benzmuller2016} showed by automated reasoning that the note's axioms, so read, are inconsistent, which is a reason for taking Scott's version as the standard one.

\subsection{Positivity}

Everything in Gödel's argument depends on the primitive notion of a positive property. Gödel's note glosses it in two ways: positive in a moral and aesthetic sense, independently of the accidental structure of the world, and, alternatively, as pure attribution in contrast to privation, that is, as a property whose logical form contains no negation \citep{Godel1995}. Neither gloss determines the axioms. The moral-aesthetic reading makes it plausible that positive properties are compatible with one another and that what a positive property entails is not bad; the attribution reading makes it plausible that exactly one of a property and its negation is positive, since one of them contains a privation. Scott's S1 expresses the second thought in its full strength, and Anderson's weakening of it to one direction expresses the first. \citet{Oppy1995} and \citet{Sobel2004} press the objection that, on any reading on which the axioms are plausible, the notion is either too weak to yield the conclusion or too strong to be credited independently of it.

The filter axioms of the variant below take a definite position on this question. They say that positive properties are jointly compatible (properness and conjunction) and that whatever a positive property necessarily entails is positive (entailment), and they say nothing about negations. They are therefore the axioms suited to the moral-aesthetic reading, and they are the conditions that Volume~III of \emph{On the Existence of God} places on the great-making properties when it appeals to Lemma~2417.

\subsection{Collapse}

\citet{Sobel1987} proved that Scott's axioms entail $q \to \nec q$ for every proposition $q$. The proof is short. A godlike being has every positive property, and by S1 it has a property or its negation positive; so every property it has is positive, and by S4 necessarily positive, and hence necessarily entailed by godlikeness. Godlikeness is therefore an essence of the godlike being, and the being's essence entails every property it has, including the property of being such that $q$, whenever $q$ is true. Since necessarily something is godlike, necessarily something is such that $q$; that is, necessarily $q$. Appendix~\ref{app:derivations} gives the proof as a checked derivation (Derivation~B.2).

\subsection{Emendations}

\citet{Anderson1990} located the source of collapse in the combination of S1, which makes the positive properties maximal, with the definitions of godlikeness and essence, which make every property of a godlike being, contingent ones included, flow from its essence. He weakened S1 to its left-to-right half, that if a property is positive its negation is not, and redefined godlikeness and essence in terms of the properties a being has necessarily. Anderson's version proves the necessary existence of a godlike being, and it does not entail collapse. \citet{Hajek2002} studied the relations between Anderson's axioms and Gödel's, and proposed a further small emendation. \citet{Fitting2002} gave a careful treatment of the argument in intensional type theory, distinguishing extensional and intensional readings of positivity, and \citet{Sobel2004} placed the collapse problem within a general assessment of ontological arguments. \citet{Benzmuller2014} verified Scott's version and Sobel's collapse with higher-order automated theorem provers and reported that the modal logic KB suffices for the main theorem.

The variant we study belongs to the Anderson line. It differs from Anderson's version in replacing his first axiom by the filter axioms below, which entail it, and it is chosen so that its positivity axioms coincide with the hypothesis of Lemma~2417 of \emph{On the Existence of God}. The choice is not idle: it shows that the filter hypothesis on which Volume~III relies is consistent with denying modal collapse, and it makes the comparison with REL-01-a and REL-01-f exact.

\section{Framework}\label{s:framework}

\subsection{Language and models}

The language has individual variables $x, y, \ldots$, property variables $Y, Z, \ldots$, the positivity predicate $\Pos$, the operators $\nec$ and $\pos$, and $\lambda$-abstraction for forming properties. Formulas are evaluated at worlds of a model.

\begin{definition}\label{d:model}
A \emph{model} consists of a non-empty set $W$ of worlds with an accessibility relation $R$, a non-empty domain $D$ of individuals, the same at every world, a set $\Pi$ of admissible properties, each a function from worlds to subsets of $D$, and a set $\Pos \subseteq \Pi$ of positive properties, the same at every world. Property quantifiers range over $\Pi$. $\Pi$ is required to be closed under the Boolean operations, under $\nec$ (the property of being necessarily $Y$), and under existential projection (the property of being such that something is $Y$), and to contain the properties defined below. A model is \emph{full} if $\Pi$ contains every function from worlds to subsets of $D$.
\end{definition}

Models with restricted $\Pi$ are general models in the sense of \citet{Henkin1950}. The closure conditions guarantee that every property term used in the derivations of Appendix~\ref{app:derivations} denotes a member of $\Pi$; the checker of Appendix~\ref{app:derivations} admits exactly such terms. The positive properties are taken to be the same at every world, so that Scott's axiom S4 holds automatically; we state it as an axiom below so that the derivations show where it is used.

\subsection{Definitions}

\begin{definition}[Godlikeness]\label{d:godlike}
$G(x) \;:\leftrightarrow\; \forall Y\,(\Pos(Y) \leftrightarrow \nec Y(x))$: the properties $x$ has necessarily are exactly the positive properties.
\end{definition}

\begin{definition}[Essence]\label{d:essence}
$\Ess(Y,x) \;:\leftrightarrow\; \forall Z\,(\nec Z(x) \leftrightarrow \nec\forall y\,(Y(y) \to Z(y)))$: the properties $Y$ necessarily entails are exactly the properties $x$ has necessarily.
\end{definition}

\begin{definition}[Necessary existence]\label{d:ne}
$\NE(x) \;:\leftrightarrow\; \forall Y\,(\Ess(Y,x) \to \nec\exists y\,Y(y))$: every essence of $x$ is necessarily instantiated.
\end{definition}

Definitions~\ref{d:godlike} and~\ref{d:essence} are Anderson's. They differ from Scott's in two ways. Godlikeness requires that the positive properties be had necessarily, and that only positive properties be had necessarily; a godlike being may have contingent properties, provided none of them is positive and none of its non-positive properties is had necessarily. An essence is tied to the necessary properties of its bearer, not to all its properties, so that a being's essence does not entail its contingent properties.

\subsection{Axioms}

The axioms of AF are the following. The labels in brackets are those used in the derivations of Appendix~\ref{app:derivations}; F stands for filter, and A3--A5 keep Scott's numbering for the axioms that AF shares with his system in form.

\begin{assumption}[F1, conjunction]\label{a:f1}
$\Pos(Y) \wedge \Pos(Z) \to \Pos(\lambda x.\,Y(x) \wedge Z(x))$.
\end{assumption}

\begin{assumption}[F2, entailment]\label{a:f2}
$\Pos(Y) \wedge \nec\forall x\,(Y(x) \to Z(x)) \to \Pos(Z)$.
\end{assumption}

\begin{assumption}[F3, properness]\label{a:f3}
$\neg\Pos(\lambda x.\,\bot)$: the empty property is not positive.
\end{assumption}

\begin{assumption}[A3]\label{a:a3}
$\Pos(G)$: godlikeness is positive.
\end{assumption}

\begin{assumption}[A4]\label{a:a4}
$\Pos(Y) \to \nec\Pos(Y)$.
\end{assumption}

\begin{assumption}[A5]\label{a:a5}
$\Pos(\NE)$: necessary existence is positive.
\end{assumption}

Assumptions~\ref{a:f1} to~\ref{a:f3} say that the positive properties form a proper filter, which is the hypothesis of Lemma~2417 of Volume~X as REL-01-f states it. They entail Anderson's first axiom, that the negation of a positive property is not positive: if $Y$ and $\lambda x.\neg Y(x)$ were both positive, so would be their conjunction, which necessarily entails the empty property (Derivation~B.4). They do not entail Scott's S1, since a proper filter need not contain, for every property, either it or its negation. Every axiom is assumed to hold at every world.

In a finite model, a family of admissible properties satisfying Assumptions~\ref{a:f1} and~\ref{a:f2} is determined by its least member, which we call its \emph{core}: the positive properties are exactly the admissible properties that include the core. This is how the models of Appendix~\ref{app:models} are generated.

\begin{definition}[Modal collapse]\label{d:collapse}
A model exhibits \emph{modal collapse} if $q \to \nec q$ holds at every world for every sentence $q$. A \emph{collapse witness} is an admissible property $Y$ and a world $w$ at which $\exists x\,Y(x)$ is true and $\nec\exists x\,Y(x)$ false.
\end{definition}

A model with a collapse witness does not exhibit modal collapse. The converse need not hold, since collapse may fail for sentences of other forms; the witnesses are what the enumeration searches for.

\section{The argument}\label{s:argument}

\begin{lemma}[Possibility]\label{l:poss}
Assumptions~\ref{a:f2}, \ref{a:f3} and~\ref{a:a3} entail, at every world, that possibly something is godlike: $\pos\exists x\,G(x)$.
\end{lemma}

\begin{proof}
Suppose that at a world $w$ nothing is godlike at any accessible world. Then at $w$ it is necessary that everything that is godlike has the empty property, since nothing is godlike. Godlikeness is positive (Assumption~\ref{a:a3}), so by Assumption~\ref{a:f2} the empty property is positive, contrary to Assumption~\ref{a:f3}.
\end{proof}

The same argument shows that every positive property is possibly instantiated (Derivation~B.3). It is the argument of Lemma~2417 in the form REL-01-f gives it, and it uses no condition on accessibility. In particular, it shows that the axioms force every world to have an accessible world.

\begin{lemma}[Rigidity]\label{l:rigid}
Assumption~\ref{a:a3} entails that whatever is godlike is necessarily godlike: $G(x) \to \nec G(x)$.
\end{lemma}

\begin{proof}
If $x$ is godlike, every positive property is had by $x$ necessarily (Definition~\ref{d:godlike}, left to right). Godlikeness is positive.
\end{proof}

Lemma~\ref{l:rigid} is the point at which Anderson's definition does the work that Scott's system assigns to essence and necessary existence. A Scott-godlike being has all positive properties, and showing that it has them necessarily requires S1, S4 and the essence clause. An Anderson-godlike being has its positive properties necessarily by definition, so the persistence of godlikeness along accessibility is immediate.

\begin{lemma}[Essence]\label{l:essence}
On a symmetric frame, if $x$ is godlike at $w$, then godlikeness is an essence of $x$ at $w$: $G(x) \to \Ess(G,x)$.
\end{lemma}

\begin{proof}
Let $x$ be godlike at $w$. We must show, for every admissible $Z$, that $\nec Z(x)$ holds at $w$ if and only if $\nec\forall y(G(y) \to Z(y))$ does.

Suppose first that $\nec Z(x)$ holds at $w$. Since $x$ is godlike, $Z$ is positive. Let $u$ be accessible from $w$ and $y$ godlike at $u$. By Lemma~\ref{l:rigid}, $y$ is necessarily godlike at $u$, and by symmetry $w$ is accessible from $u$; so $y$ is godlike at $w$. Since $Z$ is positive and $y$ is godlike at $w$, $\nec Z(y)$ holds at $w$, and so $Z(y)$ holds at $u$. Hence $\nec\forall y(G(y) \to Z(y))$ holds at $w$.

Suppose conversely that $\nec\forall y(G(y) \to Z(y))$ holds at $w$. By Lemma~\ref{l:rigid}, $x$ is godlike at every world accessible from $w$, and so $Z(x)$ holds at each of them. Hence $\nec Z(x)$ holds at $w$.
\end{proof}

The positivity of $Z$ is used at $w$, where it was established, rather than at $u$; that is why the proof needs no appeal to Assumption~\ref{a:a4}. Symmetry is used to bring the godlike being $y$ back from $u$ to $w$.

\begin{lemma}[Necessary existence]\label{l:ne}
On a symmetric frame, if $x$ is godlike at $w$ and some world is accessible from $w$, then $x$ has necessary existence at $w$.
\end{lemma}

\begin{proof}
Let $v$ be accessible from $w$. By Lemma~\ref{l:rigid}, $x$ is godlike at $v$. Necessary existence is positive (Assumption~\ref{a:a5}), so $\nec\NE(x)$ holds at $v$. By symmetry $w$ is accessible from $v$, so $\NE(x)$ holds at $w$.
\end{proof}

\begin{theorem}\label{t:main}
In every model of AF on a symmetric frame, at every world $w$:
\[
\nec\exists x\,G(x) \;\wedge\; \exists x\,(G(x) \wedge \Ess(G,x) \wedge \NE(x)).
\]
\end{theorem}

\begin{proof}
By Lemma~\ref{l:poss}, there are a world $v$ accessible from $w$ and a being $g$ godlike at $v$. By Lemma~\ref{l:rigid}, $g$ is necessarily godlike at $v$, and by symmetry $w$ is accessible from $v$; so $g$ is godlike at $w$. By Lemma~\ref{l:rigid} again, $g$ is godlike at every world accessible from $w$, which gives the first conjunct. Lemma~\ref{l:essence} gives $\Ess(G,g)$ at $w$, and Lemma~\ref{l:ne}, with $v$ as the accessible world, gives $\NE(g)$ at $w$.
\end{proof}

\begin{corollary}\label{c:actual}
At the actual world some being is godlike and exists necessarily, in the strong sense that its essence, godlikeness, is necessarily instantiated.
\end{corollary}

\begin{proof}
By Theorem~\ref{t:main} for $w = @$; the essence $G$ of the godlike being $g$ is necessarily instantiated by the first conjunct, and also by $\NE(g)$ with Definition~\ref{d:ne}.
\end{proof}

The derivation of Theorem~\ref{t:main} in Appendix~\ref{app:derivations} (Derivation~B.1) has \csLinesAf{} lines. The checker reports that it depends on no open assumption, that the only relational rule it uses is symmetry, and that the axioms it cites are F2, F3, A3 and A5. Two of the six axioms of AF are therefore idle for the main theorem: the conjunction axiom F1, as REL-01-f found for Lemma~2417, and the axiom A4 that positivity is necessary. They belong to AF because they belong to the filter hypothesis of Volume~X and to the Gödel--Scott tradition respectively, and they do no harm: the enumeration verifies them in every model.

\subsection*{Beyond symmetry}

Symmetry is sufficient for Theorem~\ref{t:main}, and it is the only relational rule in Derivation~B.1, but it is not necessary. The persistence of godlikeness along accessibility (Lemma~\ref{l:rigid}) iterates: whatever is godlike at $v$ is godlike at every world reachable from $v$ by a chain of accessibility steps. Each use of symmetry in the proofs of Section~\ref{s:argument} serves to return from a world $v$ accessible from $w$ to $w$ itself, and a chain of steps from $v$ back to $w$ serves as well.

\begin{proposition}\label{p:retrace}
Theorem~\ref{t:main} holds on every retraceable frame, that is, every frame on which $wRv$ implies that $w$ is reachable from $v$ by a chain of accessibility steps.
\end{proposition}

\begin{proof}
In the proof of Theorem~\ref{t:main}, $g$ is godlike at $v$ and hence, iterating Lemma~\ref{l:rigid}, at every world on a chain from $v$ to $w$, including $w$. In the proof of Lemma~\ref{l:essence}, the godlike $y$ at $u$ is brought back to $w$ along a chain from $u$ to $w$ in the same way. In the proof of Lemma~\ref{l:ne}, let $v'$ be the last world before $w$ on a chain from $v$ to $w$; $x$ is godlike at $v'$, so $\nec\NE(x)$ holds at $v'$, and $v' R w$.
\end{proof}

Retraceability is the condition that REL-01-b identifies as exact for its version of the modal argument without the S5 axiom. The enumeration confirms Proposition~\ref{p:retrace}: among the \csNon{} models on non-symmetric frames, \csNonC{} lie on retraceable frames, and the conclusion and all four lemmas hold in each of them; of the \csNonNotC{} models on frames that are not retraceable, the first conjunct of the conclusion fails in \csNonNotCFail{}. In the remaining \csNonNotCOk{} the conclusion holds for reasons particular to the model. We have not determined whether retraceability is necessary for the theorem in the sense in which REL-01-b shows it necessary for its schemas; the countermodels show only that it cannot be dropped altogether.

\section{Modal collapse}\label{s:collapse}

\begin{proposition}[Sobel]\label{p:sobel}
On a symmetric frame, Scott's axioms entail, at every world, that whatever is instantiated is necessarily instantiated: $\forall Y(\exists x\,Y(x) \to \nec\exists x\,Y(x))$. More generally, they entail $q \to \nec q$ for every sentence $q$ such that the property $\lambda z.\,q$ is admissible.
\end{proposition}

\begin{proof}
The argument of Section~\ref{s:history}. Derivation~B.2 proves the first claim; the general claim has the same proof with $\lambda z.\,q$ in place of $\lambda z.\,\exists x\,Y(x)$ at line~80 of that derivation.
\end{proof}

\begin{proposition}\label{p:nocollapse}
AF does not entail modal collapse. In particular, there is a full model of AF with two worlds, universal accessibility and one individual, in which some proposition is true at one world and false at the other.
\end{proposition}

\begin{proof}
Model S18 of the listing in Appendix~\ref{app:models}. Let $W = \{0, 1\}$ with every world accessible from every world, $D = \{d\}$, and let $\Pi$ contain all four properties. Let the core of positivity be the property that $d$ has at both worlds, so that a property is positive exactly when $d$ has it at both worlds. Then $d$ is godlike at both worlds: the properties it has necessarily are exactly those it has at both worlds, which are exactly the positive ones. One checks the remaining axioms directly; the program checks them all. The property $Y$ that $d$ has at world~0 only is admissible, and $\exists x\,Y(x)$ is true at world~0 and false at world~1.
\end{proof}

\begin{proposition}\label{p:consistent}
AF is consistent with modal collapse.
\end{proposition}

\begin{proof}
Model S1 of the listing has one world, accessible from itself, and one individual; every axiom of AF holds, and every truth is necessary, since the only accessible world is the world of evaluation.
\end{proof}

More informatively, in the enumeration each of the \csUltra{} models of AF whose positive properties form an ultrafilter of the admissible properties is also a model of Scott's axioms, and none of the \csScottSym{} among them that lie on symmetric frames has a collapse witness, as Proposition~\ref{p:sobel} requires, since the admissible properties are closed under existential projection.

Propositions~\ref{p:nocollapse} and~\ref{p:consistent} together locate the collapse precisely. AF neither entails nor excludes it; what entails it, given AF's filter, is the further requirement that the filter be an ultrafilter, which is Scott's S1. In Derivation~B.2 the right-to-left half of S1 is used at lines~25 to~32, to infer from the fact that a godlike being has a property $Z$ that $Z$ is positive. That inference is the step through which every contingent property of a godlike being acquires the status of a positive, and therefore necessary, property. AF has no counterpart of it. In AF a godlike being has as positive properties exactly the properties it has necessarily, and its contingent properties, such as the property of being such that world~0 obtains in model S18, are neither positive nor negatively positive.

Scott's essence clause is the second ingredient. Under Definition~\ref{d:essence} an essence entails only the necessary properties of its bearer, so that even if godlikeness is an essence of a being, the essence does not transmit the being's contingent properties to every world. In Derivation~B.2 the collapse is reached at line~80 by instantiating Scott's essence clause with the property $\lambda z.\,\exists x\,Y(x)$, which a godlike being has whenever something is $Y$. The corresponding instance of Definition~\ref{d:essence} requires that the godlike being have that property necessarily, which is what collapse would assert and what the countermodel denies.

\subsection{The four systems compared}

Table~\ref{tab:systems} summarizes how the four systems discussed in Section~\ref{s:history} stand to the two properties at issue. The entries for AF and for Scott's system are proved in this manuscript and checked in its appendices; the entries for Gödel's and Anderson's systems report the results of \citet{Benzmuller2016} and \citet{Anderson1990}.

\begin{table}[ht]
\caption{Four versions of Gödel's argument.}\label{tab:systems}
\begin{center}\small
\begin{tabular}{@{}lllll@{}}
\toprule
& positivity axioms & godlikeness, essence & necessary existence & collapse \\
\midrule
Gödel (note) & conjunction, S1, S2 & Scott's, no $Y(x)$ in essence & axioms inconsistent & -- \\
Scott & S1, S2, S3 & Scott's & proved (KB) & entailed \\
Anderson & half of S1, S2, \ldots & Anderson's & proved & not entailed \\
AF & F1--F3, A3 & Anderson's & proved (KB) & not entailed \\
\bottomrule
\end{tabular}
\end{center}
\end{table}

\section{The premise avoided}\label{s:avoided}

The main argument of REL-01-a rests on two assumptions: the S5 conditions on accessibility (Assumption~2.1) and the possibility premise (Assumption~3.1), that possibly some being is maximally great. AF avoids the second. In its place stand the positivity axioms, from which Lemma~\ref{l:poss} derives the possibility of a godlike being. This is exactly the use to which Volume~III puts Lemma~2417, transferring support from the coherence of a family of perfections to the possibility of their joint instantiation; AF makes the transfer part of the argument rather than a separate defence of its premise. Of the S5 conditions, AF uses symmetry alone. Reflexivity and transitivity are not needed for Theorem~\ref{t:main}, in agreement with the observation of \citet{Benzmuller2014} that KB suffices for Scott's version, and in contrast with REL-01-a, whose possibility premise yields necessary maximal greatness only on uniform frames, a condition that symmetry does not provide (REL-01-b).

Avoiding the possibility premise does not mean avoiding a substantive premise. The work is done by A3, that godlikeness is positive, together with the properness of the filter. If godlikeness is positive and the positive properties are not jointly impossible, godlikeness is possibly instantiated, and then, by Lemma~\ref{l:rigid} and symmetry, actually and necessarily instantiated. Someone who rejects the conclusion must reject A3 or the filter axioms. The converse argument of REL-01-a (Proposition~4.1) has an analogue here: from Lemma~\ref{l:rigid}, the possibility that nothing is godlike entails, on any frame, that nothing is godlike (Derivation~B.6). The analogue is not a symmetric rival within AF, since AF refutes its premise, but it shows where the burden lies: the choice between the possibility of a godlike being and the possibility of its absence is made, in AF, by the positivity axioms. Whether they are true is a question about the notion of positivity, and it is the question that Volume~III addresses for the great-making properties.

\section{The computation}\label{s:computation}

\subsection{Summary}

The enumeration of Appendix~\ref{app:models} examines \csCandidates{} candidate structures (a frame, a domain, an algebra of admissible properties closed as Definition~\ref{d:model} requires, and a core) and finds \csModels{} models of AF up to isomorphism: \csSym{} on symmetric frames and \csNon{} on frames that are not symmetric. Of the candidates rejected, \csRejFTwo{} violate the entailment axiom F2 at some world and \csRejAThree{} make godlikeness non-positive; in no candidate does necessary existence fail to be positive once the other axioms hold, and in none do godlikeness or necessary existence fall outside the admissible properties. For each model the program verifies Lemmas~\ref{l:poss}--\ref{l:ne} semantically, evaluates the conclusion of Theorem~\ref{t:main} at every world, and searches for a collapse witness. As an independent check, the axioms and conclusions written in the syntax of the derivations are evaluated directly in the \csXModels{} models with at most 64 admissible properties, \csXFormulas{} evaluations in all, and agree with the hand-coded checks.

The results are these. In all \csSym{} models on symmetric frames the conclusion of Theorem~\ref{t:main} holds at every world, as the theorem requires. \csSymColl{} of them have a collapse witness, \csSymTColl{} of which are on reflexive symmetric frames. Of the \csNon{} models on non-symmetric frames, the first conjunct of the conclusion fails at some world in \csNonFail{}, including \csNonFailT{} on reflexive frames and \csNonFailTFour{} on reflexive transitive frames: symmetry cannot be replaced by reflexivity, or by reflexivity and transitivity together. Of the \csModels{} models, \csScottModels{} have an ultrafilter of positive properties and are models of Scott's axioms; on symmetric frames none of these has a collapse witness, while on non-symmetric frames \csScottNonColl{} do, which shows that Sobel's collapse theorem, like the main theorem, depends on symmetry. In \csManyGods{} models on symmetric frames more than one individual is godlike at world~0. Table~\ref{tab:counts} breaks the counts down by size.

\begin{table}[ht]
\caption{Models of AF by number of worlds $n$ and individuals $m$. Scope \enquote{all}: every frame; \enquote{B}: symmetric frames only; \enquote{B or T}: symmetric or reflexive frames. \enquote{Collapse fails}: models with a collapse witness.}\label{tab:counts}
% generated by gen/generate.py -- do not edit
\begin{center}\footnotesize\begin{tabular}{@{}rrrrrrrrr@{}}\toprule
$n$ & $m$ & scope & models & symmetric & collapse fails & $\nec\exists G$ fails & Scott models & Scott, collapse fails\\\midrule
1 & 1 & all & 1 & 1 & 0 & 0 & 1 & 0\\
1 & 2 & all & 2 & 2 & 0 & 0 & 2 & 0\\
1 & 3 & all & 4 & 4 & 0 & 0 & 4 & 0\\
2 & 1 & all & 9 & 5 & 3 & 1 & 8 & 2\\
2 & 2 & all & 22 & 12 & 8 & 3 & 19 & 6\\
2 & 3 & all & 55 & 31 & 20 & 8 & 49 & 16\\
3 & 1 & all & 135 & 19 & 65 & 30 & 117 & 47\\
3 & 2 & all & 365 & 47 & 204 & 105 & 305 & 159\\
1 & 4 & B & 7 & 7 & 0 & 0 & 7 & 0\\
2 & 4 & B & 85 & 85 & 7 & 0 & 72 & 0\\
3 & 3 & B & 143 & 143 & 25 & 0 & 106 & 0\\
4 & 2 & B & 335 & 335 & 121 & 0 & 164 & 0\\
4 & 1 & B or T & 862 & 125 & 585 & 260 & 511 & 234\\
\bottomrule\end{tabular}\end{center}

\end{table}

\subsection{Six models}

\emph{S18} is the countermodel of Proposition~\ref{p:nocollapse}: frame F2.10 (two worlds, universal accessibility), one individual, all four properties admissible, the core equal to the individual's whole profile. The listing writes its algebra as \texttt{a|b}: the individual's point at world~0 is the atom \texttt{a} and its point at world~1 the atom \texttt{b}. The collapse witness \texttt{a@0} is the property of being such that world~0 obtains, instantiated at world~0 and not at world~1. S18 is the smallest model of AF with a collapse witness, and it is full, so the failure of collapse does not depend on restricting the admissible properties.

\emph{S19} has the same frame and individual but only the two trivial properties admissible. Its positive properties form an ultrafilter, Scott's axioms hold, and there is no collapse witness. S19 and S18 differ only in the richness of the property algebra, and the comparison shows that collapse in Scott's system is a matter of what the positivity axioms say about properties that distinguish worlds.

\emph{S158} lies on frame F3.96: world~0 sees all three worlds, and worlds~1 and~2 see each other and world~0 but not themselves. The frame is symmetric but neither reflexive nor transitive. The conclusion holds at every world, and the property of being such that world~0 obtains is a collapse witness. S158 shows that Theorem~\ref{t:main} holds, without collapse, in a model that S5, S4 and KTB all exclude.

\emph{S55} has two worlds, three individuals and an algebra (\texttt{aaa|bbb}) that cannot distinguish individuals. All three individuals are godlike at every world, and the proposition that world~0 obtains is again a collapse witness. AF does not prove that there is exactly one godlike being; in models whose admissible properties cannot tell beings apart it cannot.

\emph{N3} lies on frame F2.9, the smallest reflexive transitive frame that is not an equivalence relation: world~0 sees both worlds, world~1 only itself. The one individual is godlike at world~1 only. All the axioms hold, possibly something is godlike at world~0, and nothing is godlike there, so the conclusion fails at world~0. The same model satisfies Scott's axioms, whose conclusion fails there too (\texttt{s} in the listing). N3 is the counterexample that shows symmetry to be needed in both systems; it has the shape of the S4 countermodel F2.9 to the main argument in REL-01-b.

\emph{N1} lies on frame F2.5, in which both worlds see world~0 only. It satisfies Scott's axioms, Scott's conclusion holds at both worlds, and yet the proposition that world~1 obtains is a collapse witness at world~1: it is true there, and false at world~0, the only world accessible from world~1. Sobel's derivation fails at the step at which it needs a godlike being at the world of evaluation; at world~1 there is none. N1 is the smallest model in which Scott's axioms hold without collapse.

\section{Objections and replies}\label{s:objections}

\begin{objection}
Anderson's godlikeness makes A3 a disguised possibility premise. To say that godlikeness is positive, when godlikeness is defined as having the positive properties necessarily, is to say that it is possible to have them necessarily; the proof merely moves the possibility premise of REL-01-a into an axiom.
\end{objection}

\begin{reply}
The objection is half right, and the half that is right is the point of the manuscript. A3 and the filter axioms together entail the possibility of a godlike being (Lemma~\ref{l:poss}), and nothing weaker would do, since the possibility is necessary for the conclusion. What changes is the form of the premise. The possibility premise of REL-01-a is a bare modal claim, and REL-01-a records that its converse has the same logical standing. A3 and the filter axioms are claims about a family of properties: that godlikeness belongs to it, and that the family is closed and contains no impossibility. These are the claims that REL-01-f identifies as the content Volume~III must establish, and they are claims of a kind that a theory of perfections can support or refute. The proof does not make them easier to believe; it makes clear what believing them commits one to.
\end{reply}

\begin{objection}
A godlike being that has contingent properties is not God. The tradition holds that God is simple and that all of God's attributes are necessary; collapse is the price of taking that seriously, and avoiding it gives up divine perfection.
\end{objection}

\begin{reply}
In AF every positive property of a godlike being is had necessarily, so every perfection is necessary. What may be contingent are properties that are not positive, such as being such that a particular contingent world obtains, or having created this world rather than another. A tradition that holds that God's free choices are contingent needs exactly this. Simplicity is a further thesis, which Volume~VIII takes up; it is consistent with AF, which says nothing about how a godlike being's properties are related to one another.
\end{reply}

\begin{objection}
The countermodels are finite and, mostly, general rather than full models. Collapse might fail in them only because they lack the properties needed to derive it.
\end{objection}

\begin{reply}
That is true of Scott's system on finite models, where collapse is derived through the admissibility of $\lambda z.\,q$, and our closure conditions are designed to make every instance used in the derivations admissible. For AF the point does not arise, because the countermodel S18 is full: every function from worlds to sets of individuals is admissible. The other countermodels show how common the failure of collapse is; S18 alone shows that it does not depend on restricting $\Pi$.
\end{reply}

\begin{objection}
The entailment axiom F2 makes every property entailed by a perfection positive, including trivial and disjunctive ones; this is the objection that REL-01-f records against closure under entailment, and AF inherits it.
\end{objection}

\begin{reply}
AF uses F2 once, in Lemma~\ref{l:poss}, to infer from the necessary emptiness of godlikeness that the empty property is positive. REL-01-f shows that the finite intersection property suffices for Lemma~2417; in the present setting, the finite intersection property applied to the family $\{G\}$ gives Lemma~\ref{l:poss} directly, and Derivation~B.1 then uses no axiom about positivity except A3 and A5. Anyone moved by the objection can replace F1 to F3 by the finite intersection property without loss. We state AF with the filter axioms because they are the hypothesis of Lemma~2417 as Volume~X states it, and the aim is to show that that hypothesis is compatible with denying collapse.
\end{reply}

\begin{objection}
Without uniqueness, the conclusion is not monotheism; S55 has three godlike beings.
\end{objection}

\begin{reply}
The conclusion asserts existence. In AF any two godlike beings have exactly the same necessary properties, namely the positive ones (Derivation~B.5, which uses no axiom and no frame condition), so they are indiscernible in every respect that matters to godlikeness, and in S55 the admissible properties cannot distinguish them at all. Whether indiscernible godlike beings are identical is the question of Volume~VIII, on uniqueness and simplicity, and it requires principles about identity that no version of Gödel's axioms supplies.
\end{reply}

\section{What the result does not establish}\label{s:not}

The result is a validity result and a consistency result. It does not show that the axioms of AF are true: that godlikeness and necessary existence are positive, and that the positive properties form a proper filter. It does not show that symmetry holds for metaphysical accessibility, though it shows that nothing more is needed. It does not show that a godlike being in the sense of Definition~\ref{d:godlike} is maximally great in the sense of Definition~2.4 of REL-01-a; the two notions are related through the positivity of necessary existence and of the components of maximal excellence, a connection that belongs to Volume~IV. Finally, the countermodels to collapse show that AF does not entail it; they do not show that collapse is false, which is a premise of the motivation of this manuscript, not a conclusion.

\section{Conclusion}

Modal collapse is not the price of Gödel's conclusion. With Anderson's definitions and positivity axioms that say only that the positive properties form a proper filter, which is the hypothesis of Lemma~2417, the necessary existence of a godlike being follows on every symmetric frame, and there are models, among them a full model with two worlds and one being, in which some truths are contingent. The variant avoids the possibility premise of REL-01-a by deriving it, and it uses only the symmetry of accessibility. Its decisive premises are that godlikeness is positive and that the positive properties are jointly possible.

\begin{thebibliography}{99}\small
\bibitem[Anderson(1990)]{Anderson1990} Anderson, C.~A. (1990). Some emendations of Gödel's ontological proof. \emph{Faith and Philosophy}, 7(3), 291--303.
\bibitem[Benzmüller and Woltzenlogel~Paleo(2014)]{Benzmuller2014} Benzmüller, C. and Woltzenlogel Paleo, B. (2014). Automating Gödel's ontological proof of God's existence with higher-order automated theorem provers. In \emph{Proceedings of the 21st European Conference on Artificial Intelligence (ECAI 2014)} (pp.~93--98). IOS Press.
\bibitem[Benzmüller and Woltzenlogel~Paleo(2016)]{Benzmuller2016} Benzmüller, C. and Woltzenlogel Paleo, B. (2016). The inconsistency in Gödel's ontological argument: A success story for AI in metaphysics. In \emph{Proceedings of the Twenty-Fifth International Joint Conference on Artificial Intelligence (IJCAI 2016)}. AAAI Press.
\bibitem[Fitting(2002)]{Fitting2002} Fitting, M. (2002). \emph{Types, Tableaus, and Gödel's God}. Kluwer Academic Publishers.
\bibitem[Gödel(1995)]{Godel1995} Gödel, K. (1995). Ontological proof. In S.~Feferman et al. (Eds.), \emph{Collected Works, Volume III: Unpublished Essays and Lectures} (pp.~403--404). Oxford University Press.
\bibitem[Hájek(2002)]{Hajek2002} Hájek, P. (2002). A new small emendation of Gödel's ontological proof. \emph{Studia Logica}, 71, 149--164.
\bibitem[Henkin(1950)]{Henkin1950} Henkin, L. (1950). Completeness in the theory of types. \emph{The Journal of Symbolic Logic}, 15(2), 81--91.
\bibitem[Lemmon(1965)]{Lemmon1965} Lemmon, E.~J. (1965). \emph{Beginning Logic}. Nelson.
\bibitem[Oppy(1995)]{Oppy1995} Oppy, G. (1995). \emph{Ontological Arguments and Belief in God}. Cambridge University Press.
\bibitem[Sobel(1987)]{Sobel1987} Sobel, J.~H. (1987). Gödel's ontological proof. In J.~J. Thomson (Ed.), \emph{On Being and Saying: Essays for Richard Cartwright} (pp.~241--261). MIT Press.
\bibitem[Sobel(2004)]{Sobel2004} Sobel, J.~H. (2004). \emph{Logic and Theism: Arguments for and against Beliefs in God}. Cambridge University Press.
\bibitem[Viganò(2000)]{Vigano2000} Viganò, L. (2000). \emph{Labelled Non-Classical Logics}. Kluwer Academic Publishers.
\end{thebibliography}

\clearpage
\appendix
\section{Part A: finite models}\label{app:models}

\subsection{What is enumerated}

A candidate structure consists of a frame, a domain size $m$, an algebra $\Pi$ of admissible properties and a core $c$. Worlds are numbered $0, \ldots, n-1$ and individuals $0, \ldots, m-1$; a property is a set of points $(w,d)$, encoded as an $nm$-bit mask. Frames are those of the enumeration of REL-01-b and carry its identifiers F$n$.$k$. Every Boolean algebra of properties is generated by a partition of the $nm$ points into atoms, and the program generates all set partitions of the points, keeping those whose algebra is closed under $\nec$ and existential projection on the given frame (Definition~\ref{d:model}); \csRejClosure{} partitions fail the closure test across all frames. For each surviving algebra, every non-empty element is tried as the core $c$, with the positive properties those admissible properties that include $c$. As explained in Section~\ref{s:framework}, every family of positive properties satisfying F1 and F2 with a non-empty least member arises in this way, so for a given frame and algebra the enumeration of models of AF is exhaustive. Candidates are reduced up to isomorphism by applying every automorphism of the frame and every permutation of the individuals and keeping the least image.

The scope is: (I) every frame with one to three worlds and every $m$ with $nm \leq 6$; (II) every symmetric frame with $(n,m)$ equal to $(1,4)$, $(2,4)$, $(3,3)$ or $(4,2)$; (III) every frame with four worlds that is symmetric or reflexive, with $m = 1$.

\subsection{How each entry is checked}

For each candidate the program computes, by brute force over $\Pi$, the extension of Anderson's godlikeness and necessary existence (Definitions~\ref{d:godlike}--\ref{d:ne}) and checks that both are admissible and positive (A3, A5), and that F2 holds at every world for every pair of admissible properties. F1, F3 and A4 hold by construction. For each model it then checks Lemma~\ref{l:poss} (every positive property is possibly instantiated at every world), Lemma~\ref{l:rigid}, and, at every world and for every godlike individual, Lemmas~\ref{l:essence} and~\ref{l:ne}; it records the worlds at which something is godlike, the godlike individuals at world~0, and the worlds at which $\nec\exists x\,G(x)$ fails. It searches all admissible properties and worlds for a collapse witness (Definition~\ref{d:collapse}) and records the one with fewest points. If the core is an atom, so that the positive properties form an ultrafilter, it computes Scott's godlikeness and necessary existence and checks Scott's axioms and his conclusion. The program halts with an error if, on a symmetric frame, the conclusion or any lemma fails, if on a retraceable frame the conclusion fails, or a model of Scott's axioms whose conclusion holds has a collapse witness; and if any of the direct evaluations of the formulas of Appendix~\ref{app:derivations} disagrees with these checks. It completes without error.

\subsection{How to read the listing}

Each row gives: the model's identifier (S for symmetric frames, N otherwise); the frame; its standard conditions T, B, 4 (a dash for each that fails); $m$; the algebra $\Pi$, written by giving, for each world in turn and each individual in order, the letter of the atom containing that point, worlds separated by bars; the core; the extensions of $G$ and $\NE$, as sets of atoms; the worlds at which something is godlike; the godlike individuals at world~0; \texttt{+} if $\nec\exists x\,G(x)$ holds at every world, or a dash followed by the first world at which it fails; Scott's verdict (\texttt{S}: Scott's axioms and conclusion hold; \texttt{s}: the axioms hold and the conclusion fails at some world; \texttt{.}: the positive properties do not form an ultrafilter); and the collapse witness with fewest points, written as a set of atoms followed by \texttt{@} and the world, or \texttt{none}.

\subsection{Models on symmetric frames}

% generated by gen/generate.py -- do not edit
{\scriptsize\setlength{\tabcolsep}{2.6pt}
\begin{longtable}{@{}lllrlllllllll@{}}\toprule
id & frame & TB4 & $m$ & $\Pi$ (atoms by point) & core & $G$ & $\mathit{NE}$ & $\exists G$ at & gods@0 & $\nec\exists G$ & Sc & collapse wit.\\
\midrule\endfirsthead
\toprule
id & frame & TB4 & $m$ & $\Pi$ (atoms by point) & core & $G$ & $\mathit{NE}$ & $\exists G$ at & gods@0 & $\nec\exists G$ & Sc & collapse wit.\\
\midrule\endhead
\multicolumn{13}{@{}l}{\rule{0pt}{2.8ex}\footnotesize\bfseries $n=1$ worlds, $m=1$ individuals: 1 model}\\*
S1 & F1.2 & \texttt{TB4} & 1 & \texttt{a} & \texttt{a} & \texttt{a} & \texttt{a} & \texttt{0} & \texttt{0} & \texttt{+} & \texttt{S} & \texttt{none}\\
\multicolumn{13}{@{}l}{\rule{0pt}{2.8ex}\footnotesize\bfseries $n=1$ worlds, $m=2$ individuals: 2 models}\\*
S2 & F1.2 & \texttt{TB4} & 2 & \texttt{ab} & \texttt{a} & \texttt{a} & \texttt{ab} & \texttt{0} & \texttt{0} & \texttt{+} & \texttt{S} & \texttt{none}\\
S3 & F1.2 & \texttt{TB4} & 2 & \texttt{aa} & \texttt{a} & \texttt{a} & \texttt{a} & \texttt{0} & \texttt{01} & \texttt{+} & \texttt{S} & \texttt{none}\\
\multicolumn{13}{@{}l}{\rule{0pt}{2.8ex}\footnotesize\bfseries $n=1$ worlds, $m=3$ individuals: 4 models}\\*
S4 & F1.2 & \texttt{TB4} & 3 & \texttt{abc} & \texttt{a} & \texttt{a} & \texttt{abc} & \texttt{0} & \texttt{0} & \texttt{+} & \texttt{S} & \texttt{none}\\
S5 & F1.2 & \texttt{TB4} & 3 & \texttt{aab} & \texttt{b} & \texttt{b} & \texttt{ab} & \texttt{0} & \texttt{2} & \texttt{+} & \texttt{S} & \texttt{none}\\
S6 & F1.2 & \texttt{TB4} & 3 & \texttt{aab} & \texttt{a} & \texttt{a} & \texttt{ab} & \texttt{0} & \texttt{01} & \texttt{+} & \texttt{S} & \texttt{none}\\
S7 & F1.2 & \texttt{TB4} & 3 & \texttt{aaa} & \texttt{a} & \texttt{a} & \texttt{a} & \texttt{0} & \texttt{012} & \texttt{+} & \texttt{S} & \texttt{none}\\
\multicolumn{13}{@{}l}{\rule{0pt}{2.8ex}\footnotesize\bfseries $n=1$ worlds, $m=4$ individuals: 7 models}\\*
S8 & F1.2 & \texttt{TB4} & 4 & \texttt{abcd} & \texttt{a} & \texttt{a} & \texttt{abcd} & \texttt{0} & \texttt{0} & \texttt{+} & \texttt{S} & \texttt{none}\\
S9 & F1.2 & \texttt{TB4} & 4 & \texttt{aabc} & \texttt{b} & \texttt{b} & \texttt{abc} & \texttt{0} & \texttt{2} & \texttt{+} & \texttt{S} & \texttt{none}\\
S10 & F1.2 & \texttt{TB4} & 4 & \texttt{aabc} & \texttt{a} & \texttt{a} & \texttt{abc} & \texttt{0} & \texttt{01} & \texttt{+} & \texttt{S} & \texttt{none}\\
S11 & F1.2 & \texttt{TB4} & 4 & \texttt{aaab} & \texttt{b} & \texttt{b} & \texttt{ab} & \texttt{0} & \texttt{3} & \texttt{+} & \texttt{S} & \texttt{none}\\
S12 & F1.2 & \texttt{TB4} & 4 & \texttt{aaab} & \texttt{a} & \texttt{a} & \texttt{ab} & \texttt{0} & \texttt{012} & \texttt{+} & \texttt{S} & \texttt{none}\\
S13 & F1.2 & \texttt{TB4} & 4 & \texttt{abba} & \texttt{b} & \texttt{b} & \texttt{ab} & \texttt{0} & \texttt{12} & \texttt{+} & \texttt{S} & \texttt{none}\\
S14 & F1.2 & \texttt{TB4} & 4 & \texttt{aaaa} & \texttt{a} & \texttt{a} & \texttt{a} & \texttt{0} & \texttt{0123} & \texttt{+} & \texttt{S} & \texttt{none}\\
\multicolumn{13}{@{}l}{\rule{0pt}{2.8ex}\footnotesize\bfseries $n=2$ worlds, $m=1$ individuals: 5 models}\\*
S15 & F2.6 & \texttt{-{}B-{}} & 1 & \texttt{a\textbar{}a} & \texttt{a} & \texttt{a} & \texttt{a} & \texttt{01} & \texttt{0} & \texttt{+} & \texttt{S} & \texttt{none}\\
S16 & F2.7 & \texttt{TB4} & 1 & \texttt{a\textbar{}a} & \texttt{a} & \texttt{a} & \texttt{a} & \texttt{01} & \texttt{0} & \texttt{+} & \texttt{S} & \texttt{none}\\
S17 & F2.8 & \texttt{-{}B-{}} & 1 & \texttt{a\textbar{}a} & \texttt{a} & \texttt{a} & \texttt{a} & \texttt{01} & \texttt{0} & \texttt{+} & \texttt{S} & \texttt{none}\\
S18 & F2.10 & \texttt{TB4} & 1 & \texttt{a\textbar{}b} & \texttt{ab} & \texttt{ab} & \texttt{ab} & \texttt{01} & \texttt{0} & \texttt{+} & \texttt{.} & \texttt{a@0}\\
S19 & F2.10 & \texttt{TB4} & 1 & \texttt{a\textbar{}a} & \texttt{a} & \texttt{a} & \texttt{a} & \texttt{01} & \texttt{0} & \texttt{+} & \texttt{S} & \texttt{none}\\
\multicolumn{13}{@{}l}{\rule{0pt}{2.8ex}\footnotesize\bfseries $n=2$ worlds, $m=2$ individuals: 12 models}\\*
S20 & F2.6 & \texttt{-{}B-{}} & 2 & \texttt{ab\textbar{}ab} & \texttt{a} & \texttt{a} & \texttt{ab} & \texttt{01} & \texttt{0} & \texttt{+} & \texttt{S} & \texttt{none}\\
S21 & F2.6 & \texttt{-{}B-{}} & 2 & \texttt{aa\textbar{}aa} & \texttt{a} & \texttt{a} & \texttt{a} & \texttt{01} & \texttt{01} & \texttt{+} & \texttt{S} & \texttt{none}\\
S22 & F2.7 & \texttt{TB4} & 2 & \texttt{ab\textbar{}ab} & \texttt{a} & \texttt{a} & \texttt{ab} & \texttt{01} & \texttt{0} & \texttt{+} & \texttt{S} & \texttt{none}\\
S23 & F2.7 & \texttt{TB4} & 2 & \texttt{ab\textbar{}ba} & \texttt{b} & \texttt{b} & \texttt{ab} & \texttt{01} & \texttt{1} & \texttt{+} & \texttt{S} & \texttt{none}\\
S24 & F2.7 & \texttt{TB4} & 2 & \texttt{aa\textbar{}aa} & \texttt{a} & \texttt{a} & \texttt{a} & \texttt{01} & \texttt{01} & \texttt{+} & \texttt{S} & \texttt{none}\\
S25 & F2.8 & \texttt{-{}B-{}} & 2 & \texttt{ab\textbar{}ab} & \texttt{a} & \texttt{a} & \texttt{ab} & \texttt{01} & \texttt{0} & \texttt{+} & \texttt{S} & \texttt{none}\\
S26 & F2.8 & \texttt{-{}B-{}} & 2 & \texttt{aa\textbar{}aa} & \texttt{a} & \texttt{a} & \texttt{a} & \texttt{01} & \texttt{01} & \texttt{+} & \texttt{S} & \texttt{none}\\
S27 & F2.10 & \texttt{TB4} & 2 & \texttt{ab\textbar{}cd} & \texttt{ac} & \texttt{ac} & \texttt{abcd} & \texttt{01} & \texttt{0} & \texttt{+} & \texttt{.} & \texttt{a@0}\\
S28 & F2.10 & \texttt{TB4} & 2 & \texttt{aa\textbar{}bb} & \texttt{ab} & \texttt{ab} & \texttt{ab} & \texttt{01} & \texttt{01} & \texttt{+} & \texttt{.} & \texttt{a@0}\\
S29 & F2.10 & \texttt{TB4} & 2 & \texttt{ab\textbar{}ab} & \texttt{a} & \texttt{a} & \texttt{ab} & \texttt{01} & \texttt{0} & \texttt{+} & \texttt{S} & \texttt{none}\\
S30 & F2.10 & \texttt{TB4} & 2 & \texttt{ab\textbar{}ba} & \texttt{ab} & \texttt{ab} & \texttt{ab} & \texttt{01} & \texttt{01} & \texttt{+} & \texttt{.} & \texttt{none}\\
S31 & F2.10 & \texttt{TB4} & 2 & \texttt{aa\textbar{}aa} & \texttt{a} & \texttt{a} & \texttt{a} & \texttt{01} & \texttt{01} & \texttt{+} & \texttt{S} & \texttt{none}\\
\multicolumn{13}{@{}l}{\rule{0pt}{2.8ex}\footnotesize\bfseries $n=2$ worlds, $m=3$ individuals: 31 models}\\*
S32 & F2.6 & \texttt{-{}B-{}} & 3 & \texttt{abc\textbar{}abc} & \texttt{a} & \texttt{a} & \texttt{abc} & \texttt{01} & \texttt{0} & \texttt{+} & \texttt{S} & \texttt{none}\\
S33 & F2.6 & \texttt{-{}B-{}} & 3 & \texttt{abc\textbar{}cba} & \texttt{b} & \texttt{b} & \texttt{abc} & \texttt{01} & \texttt{1} & \texttt{+} & \texttt{S} & \texttt{none}\\
S34 & F2.6 & \texttt{-{}B-{}} & 3 & \texttt{aab\textbar{}aab} & \texttt{b} & \texttt{b} & \texttt{ab} & \texttt{01} & \texttt{2} & \texttt{+} & \texttt{S} & \texttt{none}\\
S35 & F2.6 & \texttt{-{}B-{}} & 3 & \texttt{aab\textbar{}aab} & \texttt{a} & \texttt{a} & \texttt{ab} & \texttt{01} & \texttt{01} & \texttt{+} & \texttt{S} & \texttt{none}\\
S36 & F2.6 & \texttt{-{}B-{}} & 3 & \texttt{aaa\textbar{}aaa} & \texttt{a} & \texttt{a} & \texttt{a} & \texttt{01} & \texttt{012} & \texttt{+} & \texttt{S} & \texttt{none}\\
S37 & F2.7 & \texttt{TB4} & 3 & \texttt{abc\textbar{}abc} & \texttt{a} & \texttt{a} & \texttt{abc} & \texttt{01} & \texttt{0} & \texttt{+} & \texttt{S} & \texttt{none}\\
S38 & F2.7 & \texttt{TB4} & 3 & \texttt{abc\textbar{}cba} & \texttt{b} & \texttt{b} & \texttt{abc} & \texttt{01} & \texttt{1} & \texttt{+} & \texttt{S} & \texttt{none}\\
S39 & F2.7 & \texttt{TB4} & 3 & \texttt{abc\textbar{}cba} & \texttt{c} & \texttt{c} & \texttt{abc} & \texttt{01} & \texttt{2} & \texttt{+} & \texttt{S} & \texttt{none}\\
S40 & F2.7 & \texttt{TB4} & 3 & \texttt{aab\textbar{}aab} & \texttt{b} & \texttt{b} & \texttt{ab} & \texttt{01} & \texttt{2} & \texttt{+} & \texttt{S} & \texttt{none}\\
S41 & F2.7 & \texttt{TB4} & 3 & \texttt{aab\textbar{}aab} & \texttt{a} & \texttt{a} & \texttt{ab} & \texttt{01} & \texttt{01} & \texttt{+} & \texttt{S} & \texttt{none}\\
S42 & F2.7 & \texttt{TB4} & 3 & \texttt{abc\textbar{}cab} & \texttt{c} & \texttt{c} & \texttt{abc} & \texttt{01} & \texttt{2} & \texttt{+} & \texttt{S} & \texttt{none}\\
S43 & F2.7 & \texttt{TB4} & 3 & \texttt{abb\textbar{}bba} & \texttt{a} & \texttt{a} & \texttt{ab} & \texttt{01} & \texttt{0} & \texttt{+} & \texttt{S} & \texttt{none}\\
S44 & F2.7 & \texttt{TB4} & 3 & \texttt{abb\textbar{}bba} & \texttt{b} & \texttt{b} & \texttt{ab} & \texttt{01} & \texttt{12} & \texttt{+} & \texttt{S} & \texttt{none}\\
S45 & F2.7 & \texttt{TB4} & 3 & \texttt{aba\textbar{}bba} & \texttt{b} & \texttt{b} & \texttt{ab} & \texttt{01} & \texttt{1} & \texttt{+} & \texttt{S} & \texttt{none}\\
S46 & F2.7 & \texttt{TB4} & 3 & \texttt{aab\textbar{}bba} & \texttt{b} & \texttt{b} & \texttt{ab} & \texttt{01} & \texttt{2} & \texttt{+} & \texttt{S} & \texttt{none}\\
S47 & F2.7 & \texttt{TB4} & 3 & \texttt{aaa\textbar{}aaa} & \texttt{a} & \texttt{a} & \texttt{a} & \texttt{01} & \texttt{012} & \texttt{+} & \texttt{S} & \texttt{none}\\
S48 & F2.8 & \texttt{-{}B-{}} & 3 & \texttt{abc\textbar{}abc} & \texttt{a} & \texttt{a} & \texttt{abc} & \texttt{01} & \texttt{0} & \texttt{+} & \texttt{S} & \texttt{none}\\
S49 & F2.8 & \texttt{-{}B-{}} & 3 & \texttt{aab\textbar{}aab} & \texttt{b} & \texttt{b} & \texttt{ab} & \texttt{01} & \texttt{2} & \texttt{+} & \texttt{S} & \texttt{none}\\
S50 & F2.8 & \texttt{-{}B-{}} & 3 & \texttt{aab\textbar{}aab} & \texttt{a} & \texttt{a} & \texttt{ab} & \texttt{01} & \texttt{01} & \texttt{+} & \texttt{S} & \texttt{none}\\
S51 & F2.8 & \texttt{-{}B-{}} & 3 & \texttt{aaa\textbar{}aaa} & \texttt{a} & \texttt{a} & \texttt{a} & \texttt{01} & \texttt{012} & \texttt{+} & \texttt{S} & \texttt{none}\\
S52 & F2.10 & \texttt{TB4} & 3 & \texttt{abc\textbar{}def} & \texttt{ad} & \texttt{ad} & \texttt{abcdef} & \texttt{01} & \texttt{0} & \texttt{+} & \texttt{.} & \texttt{a@0}\\
S53 & F2.10 & \texttt{TB4} & 3 & \texttt{aab\textbar{}ccd} & \texttt{bd} & \texttt{bd} & \texttt{abcd} & \texttt{01} & \texttt{2} & \texttt{+} & \texttt{.} & \texttt{b@0}\\
S54 & F2.10 & \texttt{TB4} & 3 & \texttt{aab\textbar{}ccd} & \texttt{ac} & \texttt{ac} & \texttt{abcd} & \texttt{01} & \texttt{01} & \texttt{+} & \texttt{.} & \texttt{b@0}\\
S55 & F2.10 & \texttt{TB4} & 3 & \texttt{aaa\textbar{}bbb} & \texttt{ab} & \texttt{ab} & \texttt{ab} & \texttt{01} & \texttt{012} & \texttt{+} & \texttt{.} & \texttt{a@0}\\
S56 & F2.10 & \texttt{TB4} & 3 & \texttt{abc\textbar{}abc} & \texttt{a} & \texttt{a} & \texttt{abc} & \texttt{01} & \texttt{0} & \texttt{+} & \texttt{S} & \texttt{none}\\
S57 & F2.10 & \texttt{TB4} & 3 & \texttt{abc\textbar{}cba} & \texttt{b} & \texttt{b} & \texttt{abc} & \texttt{01} & \texttt{1} & \texttt{+} & \texttt{S} & \texttt{none}\\
S58 & F2.10 & \texttt{TB4} & 3 & \texttt{abc\textbar{}cba} & \texttt{ac} & \texttt{ac} & \texttt{abc} & \texttt{01} & \texttt{02} & \texttt{+} & \texttt{.} & \texttt{none}\\
S59 & F2.10 & \texttt{TB4} & 3 & \texttt{aab\textbar{}aab} & \texttt{b} & \texttt{b} & \texttt{ab} & \texttt{01} & \texttt{2} & \texttt{+} & \texttt{S} & \texttt{none}\\
S60 & F2.10 & \texttt{TB4} & 3 & \texttt{aab\textbar{}aab} & \texttt{a} & \texttt{a} & \texttt{ab} & \texttt{01} & \texttt{01} & \texttt{+} & \texttt{S} & \texttt{none}\\
S61 & F2.10 & \texttt{TB4} & 3 & \texttt{aab\textbar{}bba} & \texttt{ab} & \texttt{ab} & \texttt{ab} & \texttt{01} & \texttt{012} & \texttt{+} & \texttt{.} & \texttt{none}\\
S62 & F2.10 & \texttt{TB4} & 3 & \texttt{aaa\textbar{}aaa} & \texttt{a} & \texttt{a} & \texttt{a} & \texttt{01} & \texttt{012} & \texttt{+} & \texttt{S} & \texttt{none}\\
\multicolumn{13}{@{}l}{\rule{0pt}{2.8ex}\footnotesize\bfseries $n=2$ worlds, $m=4$ individuals: 85 models}\\*
S63 & F2.6 & \texttt{-{}B-{}} & 4 & \texttt{abcd\textbar{}abcd} & \texttt{a} & \texttt{a} & \texttt{abcd} & \texttt{01} & \texttt{0} & \texttt{+} & \texttt{S} & \texttt{none}\\
S64 & F2.6 & \texttt{-{}B-{}} & 4 & \texttt{abcd\textbar{}dbca} & \texttt{b} & \texttt{b} & \texttt{abcd} & \texttt{01} & \texttt{1} & \texttt{+} & \texttt{S} & \texttt{none}\\
S65 & F2.6 & \texttt{-{}B-{}} & 4 & \texttt{aabc\textbar{}aabc} & \texttt{b} & \texttt{b} & \texttt{abc} & \texttt{01} & \texttt{2} & \texttt{+} & \texttt{S} & \texttt{none}\\
S66 & F2.6 & \texttt{-{}B-{}} & 4 & \texttt{aabc\textbar{}aabc} & \texttt{a} & \texttt{a} & \texttt{abc} & \texttt{01} & \texttt{01} & \texttt{+} & \texttt{S} & \texttt{none}\\
S67 & F2.6 & \texttt{-{}B-{}} & 4 & \texttt{aabc\textbar{}ccba} & \texttt{b} & \texttt{b} & \texttt{abc} & \texttt{01} & \texttt{2} & \texttt{+} & \texttt{S} & \texttt{none}\\
S68 & F2.6 & \texttt{-{}B-{}} & 4 & \texttt{aaab\textbar{}aaab} & \texttt{b} & \texttt{b} & \texttt{ab} & \texttt{01} & \texttt{3} & \texttt{+} & \texttt{S} & \texttt{none}\\
S69 & F2.6 & \texttt{-{}B-{}} & 4 & \texttt{aaab\textbar{}aaab} & \texttt{a} & \texttt{a} & \texttt{ab} & \texttt{01} & \texttt{012} & \texttt{+} & \texttt{S} & \texttt{none}\\
S70 & F2.6 & \texttt{-{}B-{}} & 4 & \texttt{aabc\textbar{}aacb} & \texttt{a} & \texttt{a} & \texttt{abc} & \texttt{01} & \texttt{01} & \texttt{+} & \texttt{S} & \texttt{none}\\
S71 & F2.6 & \texttt{-{}B-{}} & 4 & \texttt{abba\textbar{}abba} & \texttt{b} & \texttt{b} & \texttt{ab} & \texttt{01} & \texttt{12} & \texttt{+} & \texttt{S} & \texttt{none}\\
S72 & F2.6 & \texttt{-{}B-{}} & 4 & \texttt{aaaa\textbar{}aaaa} & \texttt{a} & \texttt{a} & \texttt{a} & \texttt{01} & \texttt{0123} & \texttt{+} & \texttt{S} & \texttt{none}\\
S73 & F2.7 & \texttt{TB4} & 4 & \texttt{abcd\textbar{}abcd} & \texttt{a} & \texttt{a} & \texttt{abcd} & \texttt{01} & \texttt{0} & \texttt{+} & \texttt{S} & \texttt{none}\\
S74 & F2.7 & \texttt{TB4} & 4 & \texttt{abcd\textbar{}dbca} & \texttt{b} & \texttt{b} & \texttt{abcd} & \texttt{01} & \texttt{1} & \texttt{+} & \texttt{S} & \texttt{none}\\
S75 & F2.7 & \texttt{TB4} & 4 & \texttt{abcd\textbar{}dbca} & \texttt{d} & \texttt{d} & \texttt{abcd} & \texttt{01} & \texttt{3} & \texttt{+} & \texttt{S} & \texttt{none}\\
S76 & F2.7 & \texttt{TB4} & 4 & \texttt{aabc\textbar{}aabc} & \texttt{b} & \texttt{b} & \texttt{abc} & \texttt{01} & \texttt{2} & \texttt{+} & \texttt{S} & \texttt{none}\\
S77 & F2.7 & \texttt{TB4} & 4 & \texttt{aabc\textbar{}aabc} & \texttt{a} & \texttt{a} & \texttt{abc} & \texttt{01} & \texttt{01} & \texttt{+} & \texttt{S} & \texttt{none}\\
S78 & F2.7 & \texttt{TB4} & 4 & \texttt{abcd\textbar{}dacb} & \texttt{c} & \texttt{c} & \texttt{abcd} & \texttt{01} & \texttt{2} & \texttt{+} & \texttt{S} & \texttt{none}\\
S79 & F2.7 & \texttt{TB4} & 4 & \texttt{abcd\textbar{}dacb} & \texttt{d} & \texttt{d} & \texttt{abcd} & \texttt{01} & \texttt{3} & \texttt{+} & \texttt{S} & \texttt{none}\\
S80 & F2.7 & \texttt{TB4} & 4 & \texttt{abcb\textbar{}bbca} & \texttt{c} & \texttt{c} & \texttt{abc} & \texttt{01} & \texttt{2} & \texttt{+} & \texttt{S} & \texttt{none}\\
S81 & F2.7 & \texttt{TB4} & 4 & \texttt{abcb\textbar{}bbca} & \texttt{a} & \texttt{a} & \texttt{abc} & \texttt{01} & \texttt{0} & \texttt{+} & \texttt{S} & \texttt{none}\\
S82 & F2.7 & \texttt{TB4} & 4 & \texttt{abcb\textbar{}bbca} & \texttt{b} & \texttt{b} & \texttt{abc} & \texttt{01} & \texttt{13} & \texttt{+} & \texttt{S} & \texttt{none}\\
S83 & F2.7 & \texttt{TB4} & 4 & \texttt{abac\textbar{}bbac} & \texttt{c} & \texttt{c} & \texttt{abc} & \texttt{01} & \texttt{3} & \texttt{+} & \texttt{S} & \texttt{none}\\
S84 & F2.7 & \texttt{TB4} & 4 & \texttt{abac\textbar{}bbac} & \texttt{b} & \texttt{b} & \texttt{abc} & \texttt{01} & \texttt{1} & \texttt{+} & \texttt{S} & \texttt{none}\\
S85 & F2.7 & \texttt{TB4} & 4 & \texttt{aabc\textbar{}ccba} & \texttt{b} & \texttt{b} & \texttt{abc} & \texttt{01} & \texttt{2} & \texttt{+} & \texttt{S} & \texttt{none}\\
S86 & F2.7 & \texttt{TB4} & 4 & \texttt{aabc\textbar{}ccba} & \texttt{c} & \texttt{c} & \texttt{abc} & \texttt{01} & \texttt{3} & \texttt{+} & \texttt{S} & \texttt{none}\\
S87 & F2.7 & \texttt{TB4} & 4 & \texttt{aaab\textbar{}aaab} & \texttt{b} & \texttt{b} & \texttt{ab} & \texttt{01} & \texttt{3} & \texttt{+} & \texttt{S} & \texttt{none}\\
S88 & F2.7 & \texttt{TB4} & 4 & \texttt{aaab\textbar{}aaab} & \texttt{a} & \texttt{a} & \texttt{ab} & \texttt{01} & \texttt{012} & \texttt{+} & \texttt{S} & \texttt{none}\\
S89 & F2.7 & \texttt{TB4} & 4 & \texttt{abcd\textbar{}dcba} & \texttt{d} & \texttt{d} & \texttt{abcd} & \texttt{01} & \texttt{3} & \texttt{+} & \texttt{S} & \texttt{none}\\
S90 & F2.7 & \texttt{TB4} & 4 & \texttt{aabc\textbar{}aacb} & \texttt{c} & \texttt{c} & \texttt{abc} & \texttt{01} & \texttt{3} & \texttt{+} & \texttt{S} & \texttt{none}\\
S91 & F2.7 & \texttt{TB4} & 4 & \texttt{aabc\textbar{}aacb} & \texttt{a} & \texttt{a} & \texttt{abc} & \texttt{01} & \texttt{01} & \texttt{+} & \texttt{S} & \texttt{none}\\
S92 & F2.7 & \texttt{TB4} & 4 & \texttt{abcd\textbar{}dcab} & \texttt{d} & \texttt{d} & \texttt{abcd} & \texttt{01} & \texttt{3} & \texttt{+} & \texttt{S} & \texttt{none}\\
S93 & F2.7 & \texttt{TB4} & 4 & \texttt{abcb\textbar{}bbac} & \texttt{a} & \texttt{a} & \texttt{abc} & \texttt{01} & \texttt{0} & \texttt{+} & \texttt{S} & \texttt{none}\\
S94 & F2.7 & \texttt{TB4} & 4 & \texttt{abcb\textbar{}bbac} & \texttt{b} & \texttt{b} & \texttt{abc} & \texttt{01} & \texttt{13} & \texttt{+} & \texttt{S} & \texttt{none}\\
S95 & F2.7 & \texttt{TB4} & 4 & \texttt{abbc\textbar{}ccba} & \texttt{a} & \texttt{a} & \texttt{abc} & \texttt{01} & \texttt{0} & \texttt{+} & \texttt{S} & \texttt{none}\\
S96 & F2.7 & \texttt{TB4} & 4 & \texttt{abbc\textbar{}ccba} & \texttt{c} & \texttt{c} & \texttt{abc} & \texttt{01} & \texttt{3} & \texttt{+} & \texttt{S} & \texttt{none}\\
S97 & F2.7 & \texttt{TB4} & 4 & \texttt{abbc\textbar{}ccba} & \texttt{b} & \texttt{b} & \texttt{abc} & \texttt{01} & \texttt{12} & \texttt{+} & \texttt{S} & \texttt{none}\\
S98 & F2.7 & \texttt{TB4} & 4 & \texttt{aabc\textbar{}ccab} & \texttt{b} & \texttt{b} & \texttt{abc} & \texttt{01} & \texttt{2} & \texttt{+} & \texttt{S} & \texttt{none}\\
S99 & F2.7 & \texttt{TB4} & 4 & \texttt{aabc\textbar{}ccab} & \texttt{c} & \texttt{c} & \texttt{abc} & \texttt{01} & \texttt{3} & \texttt{+} & \texttt{S} & \texttt{none}\\
S100 & F2.7 & \texttt{TB4} & 4 & \texttt{abcc\textbar{}bbca} & \texttt{a} & \texttt{a} & \texttt{abc} & \texttt{01} & \texttt{0} & \texttt{+} & \texttt{S} & \texttt{none}\\
S101 & F2.7 & \texttt{TB4} & 4 & \texttt{abcc\textbar{}bbca} & \texttt{b} & \texttt{b} & \texttt{abc} & \texttt{01} & \texttt{1} & \texttt{+} & \texttt{S} & \texttt{none}\\
S102 & F2.7 & \texttt{TB4} & 4 & \texttt{abcc\textbar{}ccba} & \texttt{b} & \texttt{b} & \texttt{abc} & \texttt{01} & \texttt{1} & \texttt{+} & \texttt{S} & \texttt{none}\\
S103 & F2.7 & \texttt{TB4} & 4 & \texttt{abcc\textbar{}ccba} & \texttt{c} & \texttt{c} & \texttt{abc} & \texttt{01} & \texttt{23} & \texttt{+} & \texttt{S} & \texttt{none}\\
S104 & F2.7 & \texttt{TB4} & 4 & \texttt{abac\textbar{}ccba} & \texttt{b} & \texttt{b} & \texttt{abc} & \texttt{01} & \texttt{1} & \texttt{+} & \texttt{S} & \texttt{none}\\
S105 & F2.7 & \texttt{TB4} & 4 & \texttt{abac\textbar{}ccba} & \texttt{c} & \texttt{c} & \texttt{abc} & \texttt{01} & \texttt{3} & \texttt{+} & \texttt{S} & \texttt{none}\\
S106 & F2.7 & \texttt{TB4} & 4 & \texttt{abbb\textbar{}bbba} & \texttt{a} & \texttt{a} & \texttt{ab} & \texttt{01} & \texttt{0} & \texttt{+} & \texttt{S} & \texttt{none}\\
S107 & F2.7 & \texttt{TB4} & 4 & \texttt{abbb\textbar{}bbba} & \texttt{b} & \texttt{b} & \texttt{ab} & \texttt{01} & \texttt{123} & \texttt{+} & \texttt{S} & \texttt{none}\\
S108 & F2.7 & \texttt{TB4} & 4 & \texttt{abba\textbar{}bbba} & \texttt{a} & \texttt{a} & \texttt{ab} & \texttt{01} & \texttt{03} & \texttt{+} & \texttt{S} & \texttt{none}\\
S109 & F2.7 & \texttt{TB4} & 4 & \texttt{abba\textbar{}bbba} & \texttt{b} & \texttt{b} & \texttt{ab} & \texttt{01} & \texttt{12} & \texttt{+} & \texttt{S} & \texttt{none}\\
S110 & F2.7 & \texttt{TB4} & 4 & \texttt{aabb\textbar{}bbba} & \texttt{a} & \texttt{a} & \texttt{ab} & \texttt{01} & \texttt{01} & \texttt{+} & \texttt{S} & \texttt{none}\\
S111 & F2.7 & \texttt{TB4} & 4 & \texttt{aabb\textbar{}bbba} & \texttt{b} & \texttt{b} & \texttt{ab} & \texttt{01} & \texttt{23} & \texttt{+} & \texttt{S} & \texttt{none}\\
S112 & F2.7 & \texttt{TB4} & 4 & \texttt{aaba\textbar{}bbba} & \texttt{b} & \texttt{b} & \texttt{ab} & \texttt{01} & \texttt{2} & \texttt{+} & \texttt{S} & \texttt{none}\\
S113 & F2.7 & \texttt{TB4} & 4 & \texttt{aaab\textbar{}bbba} & \texttt{b} & \texttt{b} & \texttt{ab} & \texttt{01} & \texttt{3} & \texttt{+} & \texttt{S} & \texttt{none}\\
S114 & F2.7 & \texttt{TB4} & 4 & \texttt{abba\textbar{}abba} & \texttt{b} & \texttt{b} & \texttt{ab} & \texttt{01} & \texttt{12} & \texttt{+} & \texttt{S} & \texttt{none}\\
S115 & F2.7 & \texttt{TB4} & 4 & \texttt{aabb\textbar{}abba} & \texttt{b} & \texttt{b} & \texttt{ab} & \texttt{01} & \texttt{23} & \texttt{+} & \texttt{S} & \texttt{none}\\
S116 & F2.7 & \texttt{TB4} & 4 & \texttt{abba\textbar{}baab} & \texttt{a} & \texttt{a} & \texttt{ab} & \texttt{01} & \texttt{03} & \texttt{+} & \texttt{S} & \texttt{none}\\
S117 & F2.7 & \texttt{TB4} & 4 & \texttt{aaaa\textbar{}aaaa} & \texttt{a} & \texttt{a} & \texttt{a} & \texttt{01} & \texttt{0123} & \texttt{+} & \texttt{S} & \texttt{none}\\
S118 & F2.8 & \texttt{-{}B-{}} & 4 & \texttt{abcd\textbar{}abcd} & \texttt{a} & \texttt{a} & \texttt{abcd} & \texttt{01} & \texttt{0} & \texttt{+} & \texttt{S} & \texttt{none}\\
S119 & F2.8 & \texttt{-{}B-{}} & 4 & \texttt{aabc\textbar{}aabc} & \texttt{b} & \texttt{b} & \texttt{abc} & \texttt{01} & \texttt{2} & \texttt{+} & \texttt{S} & \texttt{none}\\
S120 & F2.8 & \texttt{-{}B-{}} & 4 & \texttt{aabc\textbar{}aabc} & \texttt{a} & \texttt{a} & \texttt{abc} & \texttt{01} & \texttt{01} & \texttt{+} & \texttt{S} & \texttt{none}\\
S121 & F2.8 & \texttt{-{}B-{}} & 4 & \texttt{aaab\textbar{}aaab} & \texttt{b} & \texttt{b} & \texttt{ab} & \texttt{01} & \texttt{3} & \texttt{+} & \texttt{S} & \texttt{none}\\
S122 & F2.8 & \texttt{-{}B-{}} & 4 & \texttt{aaab\textbar{}aaab} & \texttt{a} & \texttt{a} & \texttt{ab} & \texttt{01} & \texttt{012} & \texttt{+} & \texttt{S} & \texttt{none}\\
S123 & F2.8 & \texttt{-{}B-{}} & 4 & \texttt{abba\textbar{}abba} & \texttt{b} & \texttt{b} & \texttt{ab} & \texttt{01} & \texttt{12} & \texttt{+} & \texttt{S} & \texttt{none}\\
S124 & F2.8 & \texttt{-{}B-{}} & 4 & \texttt{aaaa\textbar{}aaaa} & \texttt{a} & \texttt{a} & \texttt{a} & \texttt{01} & \texttt{0123} & \texttt{+} & \texttt{S} & \texttt{none}\\
S125 & F2.10 & \texttt{TB4} & 4 & \texttt{abcd\textbar{}efgh} & \texttt{ae} & \texttt{ae} & \texttt{abcdefgh} & \texttt{01} & \texttt{0} & \texttt{+} & \texttt{.} & \texttt{a@0}\\
S126 & F2.10 & \texttt{TB4} & 4 & \texttt{aabc\textbar{}ddef} & \texttt{be} & \texttt{be} & \texttt{abcdef} & \texttt{01} & \texttt{2} & \texttt{+} & \texttt{.} & \texttt{b@0}\\
S127 & F2.10 & \texttt{TB4} & 4 & \texttt{aabc\textbar{}ddef} & \texttt{ad} & \texttt{ad} & \texttt{abcdef} & \texttt{01} & \texttt{01} & \texttt{+} & \texttt{.} & \texttt{b@0}\\
S128 & F2.10 & \texttt{TB4} & 4 & \texttt{aaab\textbar{}cccd} & \texttt{bd} & \texttt{bd} & \texttt{abcd} & \texttt{01} & \texttt{3} & \texttt{+} & \texttt{.} & \texttt{b@0}\\
S129 & F2.10 & \texttt{TB4} & 4 & \texttt{aaab\textbar{}cccd} & \texttt{ac} & \texttt{ac} & \texttt{abcd} & \texttt{01} & \texttt{012} & \texttt{+} & \texttt{.} & \texttt{b@0}\\
S130 & F2.10 & \texttt{TB4} & 4 & \texttt{abba\textbar{}cddc} & \texttt{bd} & \texttt{bd} & \texttt{abcd} & \texttt{01} & \texttt{12} & \texttt{+} & \texttt{.} & \texttt{b@0}\\
S131 & F2.10 & \texttt{TB4} & 4 & \texttt{aaaa\textbar{}bbbb} & \texttt{ab} & \texttt{ab} & \texttt{ab} & \texttt{01} & \texttt{0123} & \texttt{+} & \texttt{.} & \texttt{a@0}\\
S132 & F2.10 & \texttt{TB4} & 4 & \texttt{abcd\textbar{}abcd} & \texttt{a} & \texttt{a} & \texttt{abcd} & \texttt{01} & \texttt{0} & \texttt{+} & \texttt{S} & \texttt{none}\\
S133 & F2.10 & \texttt{TB4} & 4 & \texttt{abcd\textbar{}dbca} & \texttt{b} & \texttt{b} & \texttt{abcd} & \texttt{01} & \texttt{1} & \texttt{+} & \texttt{S} & \texttt{none}\\
S134 & F2.10 & \texttt{TB4} & 4 & \texttt{abcd\textbar{}dbca} & \texttt{ad} & \texttt{ad} & \texttt{abcd} & \texttt{01} & \texttt{03} & \texttt{+} & \texttt{.} & \texttt{none}\\
S135 & F2.10 & \texttt{TB4} & 4 & \texttt{aabc\textbar{}aabc} & \texttt{b} & \texttt{b} & \texttt{abc} & \texttt{01} & \texttt{2} & \texttt{+} & \texttt{S} & \texttt{none}\\
S136 & F2.10 & \texttt{TB4} & 4 & \texttt{aabc\textbar{}aabc} & \texttt{a} & \texttt{a} & \texttt{abc} & \texttt{01} & \texttt{01} & \texttt{+} & \texttt{S} & \texttt{none}\\
S137 & F2.10 & \texttt{TB4} & 4 & \texttt{aabc\textbar{}ccba} & \texttt{b} & \texttt{b} & \texttt{abc} & \texttt{01} & \texttt{2} & \texttt{+} & \texttt{S} & \texttt{none}\\
S138 & F2.10 & \texttt{TB4} & 4 & \texttt{aabc\textbar{}ccba} & \texttt{ac} & \texttt{ac} & \texttt{abc} & \texttt{01} & \texttt{013} & \texttt{+} & \texttt{.} & \texttt{none}\\
S139 & F2.10 & \texttt{TB4} & 4 & \texttt{aaab\textbar{}aaab} & \texttt{b} & \texttt{b} & \texttt{ab} & \texttt{01} & \texttt{3} & \texttt{+} & \texttt{S} & \texttt{none}\\
S140 & F2.10 & \texttt{TB4} & 4 & \texttt{aaab\textbar{}aaab} & \texttt{a} & \texttt{a} & \texttt{ab} & \texttt{01} & \texttt{012} & \texttt{+} & \texttt{S} & \texttt{none}\\
S141 & F2.10 & \texttt{TB4} & 4 & \texttt{abcd\textbar{}dcba} & \texttt{bc} & \texttt{bc} & \texttt{abcd} & \texttt{01} & \texttt{12} & \texttt{+} & \texttt{.} & \texttt{none}\\
S142 & F2.10 & \texttt{TB4} & 4 & \texttt{aabc\textbar{}aacb} & \texttt{bc} & \texttt{bc} & \texttt{abc} & \texttt{01} & \texttt{23} & \texttt{+} & \texttt{.} & \texttt{none}\\
S143 & F2.10 & \texttt{TB4} & 4 & \texttt{aabc\textbar{}aacb} & \texttt{a} & \texttt{a} & \texttt{abc} & \texttt{01} & \texttt{01} & \texttt{+} & \texttt{S} & \texttt{none}\\
S144 & F2.10 & \texttt{TB4} & 4 & \texttt{aaab\textbar{}bbba} & \texttt{ab} & \texttt{ab} & \texttt{ab} & \texttt{01} & \texttt{0123} & \texttt{+} & \texttt{.} & \texttt{none}\\
S145 & F2.10 & \texttt{TB4} & 4 & \texttt{abba\textbar{}abba} & \texttt{b} & \texttt{b} & \texttt{ab} & \texttt{01} & \texttt{12} & \texttt{+} & \texttt{S} & \texttt{none}\\
S146 & F2.10 & \texttt{TB4} & 4 & \texttt{abba\textbar{}baab} & \texttt{ab} & \texttt{ab} & \texttt{ab} & \texttt{01} & \texttt{0123} & \texttt{+} & \texttt{.} & \texttt{none}\\
S147 & F2.10 & \texttt{TB4} & 4 & \texttt{aaaa\textbar{}aaaa} & \texttt{a} & \texttt{a} & \texttt{a} & \texttt{01} & \texttt{0123} & \texttt{+} & \texttt{S} & \texttt{none}\\
\multicolumn{13}{@{}l}{\rule{0pt}{2.8ex}\footnotesize\bfseries $n=3$ worlds, $m=1$ individuals: 19 models}\\*
S148 & F3.26 & \texttt{-{}B-{}} & 1 & \texttt{a\textbar{}a\textbar{}a} & \texttt{a} & \texttt{a} & \texttt{a} & \texttt{012} & \texttt{0} & \texttt{+} & \texttt{S} & \texttt{none}\\
S149 & F3.28 & \texttt{TB4} & 1 & \texttt{a\textbar{}a\textbar{}a} & \texttt{a} & \texttt{a} & \texttt{a} & \texttt{012} & \texttt{0} & \texttt{+} & \texttt{S} & \texttt{none}\\
S150 & F3.39 & \texttt{-{}B-{}} & 1 & \texttt{a\textbar{}a\textbar{}a} & \texttt{a} & \texttt{a} & \texttt{a} & \texttt{012} & \texttt{0} & \texttt{+} & \texttt{S} & \texttt{none}\\
S151 & F3.41 & \texttt{-{}B-{}} & 1 & \texttt{a\textbar{}a\textbar{}a} & \texttt{a} & \texttt{a} & \texttt{a} & \texttt{012} & \texttt{0} & \texttt{+} & \texttt{S} & \texttt{none}\\
S152 & F3.56 & \texttt{-{}B-{}} & 1 & \texttt{a\textbar{}a\textbar{}a} & \texttt{a} & \texttt{a} & \texttt{a} & \texttt{012} & \texttt{0} & \texttt{+} & \texttt{S} & \texttt{none}\\
S153 & F3.60 & \texttt{-{}B-{}} & 1 & \texttt{a\textbar{}a\textbar{}a} & \texttt{a} & \texttt{a} & \texttt{a} & \texttt{012} & \texttt{0} & \texttt{+} & \texttt{S} & \texttt{none}\\
S154 & F3.74 & \texttt{TB4} & 1 & \texttt{a\textbar{}a\textbar{}a} & \texttt{a} & \texttt{a} & \texttt{a} & \texttt{012} & \texttt{0} & \texttt{+} & \texttt{S} & \texttt{none}\\
S155 & F3.78 & \texttt{-{}B-{}} & 1 & \texttt{a\textbar{}a\textbar{}a} & \texttt{a} & \texttt{a} & \texttt{a} & \texttt{012} & \texttt{0} & \texttt{+} & \texttt{S} & \texttt{none}\\
S156 & F3.88 & \texttt{-{}B-{}} & 1 & \texttt{a\textbar{}a\textbar{}a} & \texttt{a} & \texttt{a} & \texttt{a} & \texttt{012} & \texttt{0} & \texttt{+} & \texttt{S} & \texttt{none}\\
S157 & F3.89 & \texttt{-{}B-{}} & 1 & \texttt{a\textbar{}a\textbar{}a} & \texttt{a} & \texttt{a} & \texttt{a} & \texttt{012} & \texttt{0} & \texttt{+} & \texttt{S} & \texttt{none}\\
S158 & F3.96 & \texttt{-{}B-{}} & 1 & \texttt{a\textbar{}b\textbar{}b} & \texttt{ab} & \texttt{ab} & \texttt{ab} & \texttt{012} & \texttt{0} & \texttt{+} & \texttt{.} & \texttt{a@0}\\
S159 & F3.96 & \texttt{-{}B-{}} & 1 & \texttt{a\textbar{}a\textbar{}a} & \texttt{a} & \texttt{a} & \texttt{a} & \texttt{012} & \texttt{0} & \texttt{+} & \texttt{S} & \texttt{none}\\
S160 & F3.99 & \texttt{TB-{}} & 1 & \texttt{a\textbar{}b\textbar{}b} & \texttt{ab} & \texttt{ab} & \texttt{ab} & \texttt{012} & \texttt{0} & \texttt{+} & \texttt{.} & \texttt{a@0}\\
S161 & F3.99 & \texttt{TB-{}} & 1 & \texttt{a\textbar{}a\textbar{}a} & \texttt{a} & \texttt{a} & \texttt{a} & \texttt{012} & \texttt{0} & \texttt{+} & \texttt{S} & \texttt{none}\\
S162 & F3.102 & \texttt{-{}B-{}} & 1 & \texttt{a\textbar{}b\textbar{}a} & \texttt{ab} & \texttt{ab} & \texttt{ab} & \texttt{012} & \texttt{0} & \texttt{+} & \texttt{.} & \texttt{b@1}\\
S163 & F3.102 & \texttt{-{}B-{}} & 1 & \texttt{a\textbar{}a\textbar{}a} & \texttt{a} & \texttt{a} & \texttt{a} & \texttt{012} & \texttt{0} & \texttt{+} & \texttt{S} & \texttt{none}\\
S164 & F3.104 & \texttt{TB4} & 1 & \texttt{a\textbar{}b\textbar{}c} & \texttt{abc} & \texttt{abc} & \texttt{abc} & \texttt{012} & \texttt{0} & \texttt{+} & \texttt{.} & \texttt{a@0}\\
S165 & F3.104 & \texttt{TB4} & 1 & \texttt{a\textbar{}a\textbar{}b} & \texttt{ab} & \texttt{ab} & \texttt{ab} & \texttt{012} & \texttt{0} & \texttt{+} & \texttt{.} & \texttt{b@2}\\
S166 & F3.104 & \texttt{TB4} & 1 & \texttt{a\textbar{}a\textbar{}a} & \texttt{a} & \texttt{a} & \texttt{a} & \texttt{012} & \texttt{0} & \texttt{+} & \texttt{S} & \texttt{none}\\
\multicolumn{13}{@{}l}{\rule{0pt}{2.8ex}\footnotesize\bfseries $n=3$ worlds, $m=2$ individuals: 47 models}\\*
S167 & F3.26 & \texttt{-{}B-{}} & 2 & \texttt{ab\textbar{}ab\textbar{}ab} & \texttt{a} & \texttt{a} & \texttt{ab} & \texttt{012} & \texttt{0} & \texttt{+} & \texttt{S} & \texttt{none}\\
S168 & F3.26 & \texttt{-{}B-{}} & 2 & \texttt{ab\textbar{}ba\textbar{}ab} & \texttt{a} & \texttt{a} & \texttt{ab} & \texttt{012} & \texttt{0} & \texttt{+} & \texttt{S} & \texttt{none}\\
S169 & F3.26 & \texttt{-{}B-{}} & 2 & \texttt{aa\textbar{}aa\textbar{}aa} & \texttt{a} & \texttt{a} & \texttt{a} & \texttt{012} & \texttt{01} & \texttt{+} & \texttt{S} & \texttt{none}\\
S170 & F3.28 & \texttt{TB4} & 2 & \texttt{ab\textbar{}ab\textbar{}ab} & \texttt{a} & \texttt{a} & \texttt{ab} & \texttt{012} & \texttt{0} & \texttt{+} & \texttt{S} & \texttt{none}\\
S171 & F3.28 & \texttt{TB4} & 2 & \texttt{ab\textbar{}ab\textbar{}ba} & \texttt{b} & \texttt{b} & \texttt{ab} & \texttt{012} & \texttt{1} & \texttt{+} & \texttt{S} & \texttt{none}\\
S172 & F3.28 & \texttt{TB4} & 2 & \texttt{aa\textbar{}aa\textbar{}aa} & \texttt{a} & \texttt{a} & \texttt{a} & \texttt{012} & \texttt{01} & \texttt{+} & \texttt{S} & \texttt{none}\\
S173 & F3.39 & \texttt{-{}B-{}} & 2 & \texttt{ab\textbar{}ab\textbar{}ab} & \texttt{a} & \texttt{a} & \texttt{ab} & \texttt{012} & \texttt{0} & \texttt{+} & \texttt{S} & \texttt{none}\\
S174 & F3.39 & \texttt{-{}B-{}} & 2 & \texttt{aa\textbar{}aa\textbar{}aa} & \texttt{a} & \texttt{a} & \texttt{a} & \texttt{012} & \texttt{01} & \texttt{+} & \texttt{S} & \texttt{none}\\
S175 & F3.41 & \texttt{-{}B-{}} & 2 & \texttt{ab\textbar{}ab\textbar{}ab} & \texttt{a} & \texttt{a} & \texttt{ab} & \texttt{012} & \texttt{0} & \texttt{+} & \texttt{S} & \texttt{none}\\
S176 & F3.41 & \texttt{-{}B-{}} & 2 & \texttt{ab\textbar{}ba\textbar{}ab} & \texttt{a} & \texttt{a} & \texttt{ab} & \texttt{012} & \texttt{0} & \texttt{+} & \texttt{S} & \texttt{none}\\
S177 & F3.41 & \texttt{-{}B-{}} & 2 & \texttt{aa\textbar{}aa\textbar{}aa} & \texttt{a} & \texttt{a} & \texttt{a} & \texttt{012} & \texttt{01} & \texttt{+} & \texttt{S} & \texttt{none}\\
S178 & F3.56 & \texttt{-{}B-{}} & 2 & \texttt{ab\textbar{}ab\textbar{}ab} & \texttt{a} & \texttt{a} & \texttt{ab} & \texttt{012} & \texttt{0} & \texttt{+} & \texttt{S} & \texttt{none}\\
S179 & F3.56 & \texttt{-{}B-{}} & 2 & \texttt{aa\textbar{}aa\textbar{}aa} & \texttt{a} & \texttt{a} & \texttt{a} & \texttt{012} & \texttt{01} & \texttt{+} & \texttt{S} & \texttt{none}\\
S180 & F3.60 & \texttt{-{}B-{}} & 2 & \texttt{ab\textbar{}ab\textbar{}ab} & \texttt{a} & \texttt{a} & \texttt{ab} & \texttt{012} & \texttt{0} & \texttt{+} & \texttt{S} & \texttt{none}\\
S181 & F3.60 & \texttt{-{}B-{}} & 2 & \texttt{aa\textbar{}aa\textbar{}aa} & \texttt{a} & \texttt{a} & \texttt{a} & \texttt{012} & \texttt{01} & \texttt{+} & \texttt{S} & \texttt{none}\\
S182 & F3.74 & \texttt{TB4} & 2 & \texttt{ab\textbar{}ab\textbar{}ab} & \texttt{a} & \texttt{a} & \texttt{ab} & \texttt{012} & \texttt{0} & \texttt{+} & \texttt{S} & \texttt{none}\\
S183 & F3.74 & \texttt{TB4} & 2 & \texttt{ab\textbar{}ab\textbar{}ba} & \texttt{b} & \texttt{b} & \texttt{ab} & \texttt{012} & \texttt{1} & \texttt{+} & \texttt{S} & \texttt{none}\\
S184 & F3.74 & \texttt{TB4} & 2 & \texttt{aa\textbar{}aa\textbar{}aa} & \texttt{a} & \texttt{a} & \texttt{a} & \texttt{012} & \texttt{01} & \texttt{+} & \texttt{S} & \texttt{none}\\
S185 & F3.78 & \texttt{-{}B-{}} & 2 & \texttt{ab\textbar{}ab\textbar{}ab} & \texttt{a} & \texttt{a} & \texttt{ab} & \texttt{012} & \texttt{0} & \texttt{+} & \texttt{S} & \texttt{none}\\
S186 & F3.78 & \texttt{-{}B-{}} & 2 & \texttt{aa\textbar{}aa\textbar{}aa} & \texttt{a} & \texttt{a} & \texttt{a} & \texttt{012} & \texttt{01} & \texttt{+} & \texttt{S} & \texttt{none}\\
S187 & F3.88 & \texttt{-{}B-{}} & 2 & \texttt{ab\textbar{}ab\textbar{}ab} & \texttt{a} & \texttt{a} & \texttt{ab} & \texttt{012} & \texttt{0} & \texttt{+} & \texttt{S} & \texttt{none}\\
S188 & F3.88 & \texttt{-{}B-{}} & 2 & \texttt{aa\textbar{}aa\textbar{}aa} & \texttt{a} & \texttt{a} & \texttt{a} & \texttt{012} & \texttt{01} & \texttt{+} & \texttt{S} & \texttt{none}\\
S189 & F3.89 & \texttt{-{}B-{}} & 2 & \texttt{ab\textbar{}ab\textbar{}ab} & \texttt{a} & \texttt{a} & \texttt{ab} & \texttt{012} & \texttt{0} & \texttt{+} & \texttt{S} & \texttt{none}\\
S190 & F3.89 & \texttt{-{}B-{}} & 2 & \texttt{aa\textbar{}aa\textbar{}aa} & \texttt{a} & \texttt{a} & \texttt{a} & \texttt{012} & \texttt{01} & \texttt{+} & \texttt{S} & \texttt{none}\\
S191 & F3.96 & \texttt{-{}B-{}} & 2 & \texttt{ab\textbar{}cd\textbar{}cd} & \texttt{ac} & \texttt{ac} & \texttt{abcd} & \texttt{012} & \texttt{0} & \texttt{+} & \texttt{.} & \texttt{a@0}\\
S192 & F3.96 & \texttt{-{}B-{}} & 2 & \texttt{aa\textbar{}bb\textbar{}bb} & \texttt{ab} & \texttt{ab} & \texttt{ab} & \texttt{012} & \texttt{01} & \texttt{+} & \texttt{.} & \texttt{a@0}\\
S193 & F3.96 & \texttt{-{}B-{}} & 2 & \texttt{ab\textbar{}ab\textbar{}ab} & \texttt{a} & \texttt{a} & \texttt{ab} & \texttt{012} & \texttt{0} & \texttt{+} & \texttt{S} & \texttt{none}\\
S194 & F3.96 & \texttt{-{}B-{}} & 2 & \texttt{ab\textbar{}ba\textbar{}ba} & \texttt{ab} & \texttt{ab} & \texttt{ab} & \texttt{012} & \texttt{01} & \texttt{+} & \texttt{.} & \texttt{none}\\
S195 & F3.96 & \texttt{-{}B-{}} & 2 & \texttt{aa\textbar{}aa\textbar{}aa} & \texttt{a} & \texttt{a} & \texttt{a} & \texttt{012} & \texttt{01} & \texttt{+} & \texttt{S} & \texttt{none}\\
S196 & F3.99 & \texttt{TB-{}} & 2 & \texttt{ab\textbar{}cd\textbar{}cd} & \texttt{ac} & \texttt{ac} & \texttt{abcd} & \texttt{012} & \texttt{0} & \texttt{+} & \texttt{.} & \texttt{a@0}\\
S197 & F3.99 & \texttt{TB-{}} & 2 & \texttt{aa\textbar{}bb\textbar{}bb} & \texttt{ab} & \texttt{ab} & \texttt{ab} & \texttt{012} & \texttt{01} & \texttt{+} & \texttt{.} & \texttt{a@0}\\
S198 & F3.99 & \texttt{TB-{}} & 2 & \texttt{ab\textbar{}ab\textbar{}ab} & \texttt{a} & \texttt{a} & \texttt{ab} & \texttt{012} & \texttt{0} & \texttt{+} & \texttt{S} & \texttt{none}\\
S199 & F3.99 & \texttt{TB-{}} & 2 & \texttt{ab\textbar{}ba\textbar{}ba} & \texttt{ab} & \texttt{ab} & \texttt{ab} & \texttt{012} & \texttt{01} & \texttt{+} & \texttt{.} & \texttt{none}\\
S200 & F3.99 & \texttt{TB-{}} & 2 & \texttt{aa\textbar{}aa\textbar{}aa} & \texttt{a} & \texttt{a} & \texttt{a} & \texttt{012} & \texttt{01} & \texttt{+} & \texttt{S} & \texttt{none}\\
S201 & F3.102 & \texttt{-{}B-{}} & 2 & \texttt{ab\textbar{}cd\textbar{}ab} & \texttt{ac} & \texttt{ac} & \texttt{abcd} & \texttt{012} & \texttt{0} & \texttt{+} & \texttt{.} & \texttt{c@1}\\
S202 & F3.102 & \texttt{-{}B-{}} & 2 & \texttt{aa\textbar{}bb\textbar{}aa} & \texttt{ab} & \texttt{ab} & \texttt{ab} & \texttt{012} & \texttt{01} & \texttt{+} & \texttt{.} & \texttt{b@1}\\
S203 & F3.102 & \texttt{-{}B-{}} & 2 & \texttt{ab\textbar{}ab\textbar{}ab} & \texttt{a} & \texttt{a} & \texttt{ab} & \texttt{012} & \texttt{0} & \texttt{+} & \texttt{S} & \texttt{none}\\
S204 & F3.102 & \texttt{-{}B-{}} & 2 & \texttt{ab\textbar{}ba\textbar{}ab} & \texttt{ab} & \texttt{ab} & \texttt{ab} & \texttt{012} & \texttt{01} & \texttt{+} & \texttt{.} & \texttt{none}\\
S205 & F3.102 & \texttt{-{}B-{}} & 2 & \texttt{aa\textbar{}aa\textbar{}aa} & \texttt{a} & \texttt{a} & \texttt{a} & \texttt{012} & \texttt{01} & \texttt{+} & \texttt{S} & \texttt{none}\\
S206 & F3.104 & \texttt{TB4} & 2 & \texttt{ab\textbar{}cd\textbar{}ef} & \texttt{ace} & \texttt{ace} & \texttt{abcdef} & \texttt{012} & \texttt{0} & \texttt{+} & \texttt{.} & \texttt{a@0}\\
S207 & F3.104 & \texttt{TB4} & 2 & \texttt{ab\textbar{}ab\textbar{}cd} & \texttt{ac} & \texttt{ac} & \texttt{abcd} & \texttt{012} & \texttt{0} & \texttt{+} & \texttt{.} & \texttt{c@2}\\
S208 & F3.104 & \texttt{TB4} & 2 & \texttt{aa\textbar{}bb\textbar{}cc} & \texttt{abc} & \texttt{abc} & \texttt{abc} & \texttt{012} & \texttt{01} & \texttt{+} & \texttt{.} & \texttt{a@0}\\
S209 & F3.104 & \texttt{TB4} & 2 & \texttt{ab\textbar{}cc\textbar{}ba} & \texttt{abc} & \texttt{abc} & \texttt{abc} & \texttt{012} & \texttt{01} & \texttt{+} & \texttt{.} & \texttt{c@1}\\
S210 & F3.104 & \texttt{TB4} & 2 & \texttt{aa\textbar{}aa\textbar{}bb} & \texttt{ab} & \texttt{ab} & \texttt{ab} & \texttt{012} & \texttt{01} & \texttt{+} & \texttt{.} & \texttt{b@2}\\
S211 & F3.104 & \texttt{TB4} & 2 & \texttt{ab\textbar{}ab\textbar{}ab} & \texttt{a} & \texttt{a} & \texttt{ab} & \texttt{012} & \texttt{0} & \texttt{+} & \texttt{S} & \texttt{none}\\
S212 & F3.104 & \texttt{TB4} & 2 & \texttt{ab\textbar{}ab\textbar{}ba} & \texttt{ab} & \texttt{ab} & \texttt{ab} & \texttt{012} & \texttt{01} & \texttt{+} & \texttt{.} & \texttt{none}\\
S213 & F3.104 & \texttt{TB4} & 2 & \texttt{aa\textbar{}aa\textbar{}aa} & \texttt{a} & \texttt{a} & \texttt{a} & \texttt{012} & \texttt{01} & \texttt{+} & \texttt{S} & \texttt{none}\\
\multicolumn{13}{@{}l}{\rule{0pt}{2.8ex}\footnotesize\bfseries $n=3$ worlds, $m=3$ individuals: 143 models}\\*
S214 & F3.26 & \texttt{-{}B-{}} & 3 & \texttt{abc\textbar{}abc\textbar{}abc} & \texttt{a} & \texttt{a} & \texttt{abc} & \texttt{012} & \texttt{0} & \texttt{+} & \texttt{S} & \texttt{none}\\
S215 & F3.26 & \texttt{-{}B-{}} & 3 & \texttt{abc\textbar{}cba\textbar{}abc} & \texttt{b} & \texttt{b} & \texttt{abc} & \texttt{012} & \texttt{1} & \texttt{+} & \texttt{S} & \texttt{none}\\
S216 & F3.26 & \texttt{-{}B-{}} & 3 & \texttt{abc\textbar{}cba\textbar{}abc} & \texttt{a} & \texttt{a} & \texttt{abc} & \texttt{012} & \texttt{0} & \texttt{+} & \texttt{S} & \texttt{none}\\
S217 & F3.26 & \texttt{-{}B-{}} & 3 & \texttt{aab\textbar{}aab\textbar{}aab} & \texttt{b} & \texttt{b} & \texttt{ab} & \texttt{012} & \texttt{2} & \texttt{+} & \texttt{S} & \texttt{none}\\
S218 & F3.26 & \texttt{-{}B-{}} & 3 & \texttt{aab\textbar{}aab\textbar{}aab} & \texttt{a} & \texttt{a} & \texttt{ab} & \texttt{012} & \texttt{01} & \texttt{+} & \texttt{S} & \texttt{none}\\
S219 & F3.26 & \texttt{-{}B-{}} & 3 & \texttt{abc\textbar{}bca\textbar{}abc} & \texttt{a} & \texttt{a} & \texttt{abc} & \texttt{012} & \texttt{0} & \texttt{+} & \texttt{S} & \texttt{none}\\
S220 & F3.26 & \texttt{-{}B-{}} & 3 & \texttt{aab\textbar{}baa\textbar{}aab} & \texttt{b} & \texttt{b} & \texttt{ab} & \texttt{012} & \texttt{2} & \texttt{+} & \texttt{S} & \texttt{none}\\
S221 & F3.26 & \texttt{-{}B-{}} & 3 & \texttt{aab\textbar{}baa\textbar{}aab} & \texttt{a} & \texttt{a} & \texttt{ab} & \texttt{012} & \texttt{01} & \texttt{+} & \texttt{S} & \texttt{none}\\
S222 & F3.26 & \texttt{-{}B-{}} & 3 & \texttt{aab\textbar{}bab\textbar{}aab} & \texttt{b} & \texttt{b} & \texttt{ab} & \texttt{012} & \texttt{2} & \texttt{+} & \texttt{S} & \texttt{none}\\
S223 & F3.26 & \texttt{-{}B-{}} & 3 & \texttt{aab\textbar{}bab\textbar{}aab} & \texttt{a} & \texttt{a} & \texttt{ab} & \texttt{012} & \texttt{01} & \texttt{+} & \texttt{S} & \texttt{none}\\
S224 & F3.26 & \texttt{-{}B-{}} & 3 & \texttt{aab\textbar{}bba\textbar{}aab} & \texttt{b} & \texttt{b} & \texttt{ab} & \texttt{012} & \texttt{2} & \texttt{+} & \texttt{S} & \texttt{none}\\
S225 & F3.26 & \texttt{-{}B-{}} & 3 & \texttt{aab\textbar{}bba\textbar{}aab} & \texttt{a} & \texttt{a} & \texttt{ab} & \texttt{012} & \texttt{01} & \texttt{+} & \texttt{S} & \texttt{none}\\
S226 & F3.26 & \texttt{-{}B-{}} & 3 & \texttt{aaa\textbar{}aaa\textbar{}aaa} & \texttt{a} & \texttt{a} & \texttt{a} & \texttt{012} & \texttt{012} & \texttt{+} & \texttt{S} & \texttt{none}\\
S227 & F3.28 & \texttt{TB4} & 3 & \texttt{abc\textbar{}abc\textbar{}abc} & \texttt{a} & \texttt{a} & \texttt{abc} & \texttt{012} & \texttt{0} & \texttt{+} & \texttt{S} & \texttt{none}\\
S228 & F3.28 & \texttt{TB4} & 3 & \texttt{abc\textbar{}abc\textbar{}cba} & \texttt{b} & \texttt{b} & \texttt{abc} & \texttt{012} & \texttt{1} & \texttt{+} & \texttt{S} & \texttt{none}\\
S229 & F3.28 & \texttt{TB4} & 3 & \texttt{abc\textbar{}abc\textbar{}cba} & \texttt{c} & \texttt{c} & \texttt{abc} & \texttt{012} & \texttt{2} & \texttt{+} & \texttt{S} & \texttt{none}\\
S230 & F3.28 & \texttt{TB4} & 3 & \texttt{aab\textbar{}aab\textbar{}aab} & \texttt{b} & \texttt{b} & \texttt{ab} & \texttt{012} & \texttt{2} & \texttt{+} & \texttt{S} & \texttt{none}\\
S231 & F3.28 & \texttt{TB4} & 3 & \texttt{aab\textbar{}aab\textbar{}aab} & \texttt{a} & \texttt{a} & \texttt{ab} & \texttt{012} & \texttt{01} & \texttt{+} & \texttt{S} & \texttt{none}\\
S232 & F3.28 & \texttt{TB4} & 3 & \texttt{abc\textbar{}abc\textbar{}cab} & \texttt{c} & \texttt{c} & \texttt{abc} & \texttt{012} & \texttt{2} & \texttt{+} & \texttt{S} & \texttt{none}\\
S233 & F3.28 & \texttt{TB4} & 3 & \texttt{abc\textbar{}bac\textbar{}cab} & \texttt{c} & \texttt{c} & \texttt{abc} & \texttt{012} & \texttt{2} & \texttt{+} & \texttt{S} & \texttt{none}\\
S234 & F3.28 & \texttt{TB4} & 3 & \texttt{abc\textbar{}bac\textbar{}cab} & \texttt{b} & \texttt{b} & \texttt{abc} & \texttt{012} & \texttt{1} & \texttt{+} & \texttt{S} & \texttt{none}\\
S235 & F3.28 & \texttt{TB4} & 3 & \texttt{abb\textbar{}abb\textbar{}bba} & \texttt{a} & \texttt{a} & \texttt{ab} & \texttt{012} & \texttt{0} & \texttt{+} & \texttt{S} & \texttt{none}\\
S236 & F3.28 & \texttt{TB4} & 3 & \texttt{abb\textbar{}abb\textbar{}bba} & \texttt{b} & \texttt{b} & \texttt{ab} & \texttt{012} & \texttt{12} & \texttt{+} & \texttt{S} & \texttt{none}\\
S237 & F3.28 & \texttt{TB4} & 3 & \texttt{aba\textbar{}bba\textbar{}bba} & \texttt{a} & \texttt{a} & \texttt{ab} & \texttt{012} & \texttt{02} & \texttt{+} & \texttt{S} & \texttt{none}\\
S238 & F3.28 & \texttt{TB4} & 3 & \texttt{aba\textbar{}bba\textbar{}bba} & \texttt{b} & \texttt{b} & \texttt{ab} & \texttt{012} & \texttt{1} & \texttt{+} & \texttt{S} & \texttt{none}\\
S239 & F3.28 & \texttt{TB4} & 3 & \texttt{aab\textbar{}aab\textbar{}bba} & \texttt{b} & \texttt{b} & \texttt{ab} & \texttt{012} & \texttt{2} & \texttt{+} & \texttt{S} & \texttt{none}\\
S240 & F3.28 & \texttt{TB4} & 3 & \texttt{aab\textbar{}aab\textbar{}bba} & \texttt{a} & \texttt{a} & \texttt{ab} & \texttt{012} & \texttt{01} & \texttt{+} & \texttt{S} & \texttt{none}\\
S241 & F3.28 & \texttt{TB4} & 3 & \texttt{aba\textbar{}abb\textbar{}bba} & \texttt{a} & \texttt{a} & \texttt{ab} & \texttt{012} & \texttt{02} & \texttt{+} & \texttt{S} & \texttt{none}\\
S242 & F3.28 & \texttt{TB4} & 3 & \texttt{aba\textbar{}abb\textbar{}bba} & \texttt{b} & \texttt{b} & \texttt{ab} & \texttt{012} & \texttt{1} & \texttt{+} & \texttt{S} & \texttt{none}\\
S243 & F3.28 & \texttt{TB4} & 3 & \texttt{abc\textbar{}cab\textbar{}bca} & \texttt{b} & \texttt{b} & \texttt{abc} & \texttt{012} & \texttt{1} & \texttt{+} & \texttt{S} & \texttt{none}\\
S244 & F3.28 & \texttt{TB4} & 3 & \texttt{aba\textbar{}baa\textbar{}aab} & \texttt{b} & \texttt{b} & \texttt{ab} & \texttt{012} & \texttt{1} & \texttt{+} & \texttt{S} & \texttt{none}\\
S245 & F3.28 & \texttt{TB4} & 3 & \texttt{aba\textbar{}baa\textbar{}aab} & \texttt{a} & \texttt{a} & \texttt{ab} & \texttt{012} & \texttt{02} & \texttt{+} & \texttt{S} & \texttt{none}\\
S246 & F3.28 & \texttt{TB4} & 3 & \texttt{aab\textbar{}abb\textbar{}bba} & \texttt{a} & \texttt{a} & \texttt{ab} & \texttt{012} & \texttt{01} & \texttt{+} & \texttt{S} & \texttt{none}\\
S247 & F3.28 & \texttt{TB4} & 3 & \texttt{aab\textbar{}abb\textbar{}bba} & \texttt{b} & \texttt{b} & \texttt{ab} & \texttt{012} & \texttt{2} & \texttt{+} & \texttt{S} & \texttt{none}\\
S248 & F3.28 & \texttt{TB4} & 3 & \texttt{aaa\textbar{}aaa\textbar{}aaa} & \texttt{a} & \texttt{a} & \texttt{a} & \texttt{012} & \texttt{012} & \texttt{+} & \texttt{S} & \texttt{none}\\
S249 & F3.39 & \texttt{-{}B-{}} & 3 & \texttt{abc\textbar{}abc\textbar{}abc} & \texttt{a} & \texttt{a} & \texttt{abc} & \texttt{012} & \texttt{0} & \texttt{+} & \texttt{S} & \texttt{none}\\
S250 & F3.39 & \texttt{-{}B-{}} & 3 & \texttt{abc\textbar{}cba\textbar{}cba} & \texttt{b} & \texttt{b} & \texttt{abc} & \texttt{012} & \texttt{1} & \texttt{+} & \texttt{S} & \texttt{none}\\
S251 & F3.39 & \texttt{-{}B-{}} & 3 & \texttt{aab\textbar{}aab\textbar{}aab} & \texttt{b} & \texttt{b} & \texttt{ab} & \texttt{012} & \texttt{2} & \texttt{+} & \texttt{S} & \texttt{none}\\
S252 & F3.39 & \texttt{-{}B-{}} & 3 & \texttt{aab\textbar{}aab\textbar{}aab} & \texttt{a} & \texttt{a} & \texttt{ab} & \texttt{012} & \texttt{01} & \texttt{+} & \texttt{S} & \texttt{none}\\
S253 & F3.39 & \texttt{-{}B-{}} & 3 & \texttt{aaa\textbar{}aaa\textbar{}aaa} & \texttt{a} & \texttt{a} & \texttt{a} & \texttt{012} & \texttt{012} & \texttt{+} & \texttt{S} & \texttt{none}\\
S254 & F3.41 & \texttt{-{}B-{}} & 3 & \texttt{abc\textbar{}abc\textbar{}abc} & \texttt{a} & \texttt{a} & \texttt{abc} & \texttt{012} & \texttt{0} & \texttt{+} & \texttt{S} & \texttt{none}\\
S255 & F3.41 & \texttt{-{}B-{}} & 3 & \texttt{abc\textbar{}cba\textbar{}abc} & \texttt{b} & \texttt{b} & \texttt{abc} & \texttt{012} & \texttt{1} & \texttt{+} & \texttt{S} & \texttt{none}\\
S256 & F3.41 & \texttt{-{}B-{}} & 3 & \texttt{abc\textbar{}cba\textbar{}abc} & \texttt{a} & \texttt{a} & \texttt{abc} & \texttt{012} & \texttt{0} & \texttt{+} & \texttt{S} & \texttt{none}\\
S257 & F3.41 & \texttt{-{}B-{}} & 3 & \texttt{aab\textbar{}aab\textbar{}aab} & \texttt{b} & \texttt{b} & \texttt{ab} & \texttt{012} & \texttt{2} & \texttt{+} & \texttt{S} & \texttt{none}\\
S258 & F3.41 & \texttt{-{}B-{}} & 3 & \texttt{aab\textbar{}aab\textbar{}aab} & \texttt{a} & \texttt{a} & \texttt{ab} & \texttt{012} & \texttt{01} & \texttt{+} & \texttt{S} & \texttt{none}\\
S259 & F3.41 & \texttt{-{}B-{}} & 3 & \texttt{abc\textbar{}bca\textbar{}abc} & \texttt{a} & \texttt{a} & \texttt{abc} & \texttt{012} & \texttt{0} & \texttt{+} & \texttt{S} & \texttt{none}\\
S260 & F3.41 & \texttt{-{}B-{}} & 3 & \texttt{aab\textbar{}baa\textbar{}aab} & \texttt{b} & \texttt{b} & \texttt{ab} & \texttt{012} & \texttt{2} & \texttt{+} & \texttt{S} & \texttt{none}\\
S261 & F3.41 & \texttt{-{}B-{}} & 3 & \texttt{aab\textbar{}baa\textbar{}aab} & \texttt{a} & \texttt{a} & \texttt{ab} & \texttt{012} & \texttt{01} & \texttt{+} & \texttt{S} & \texttt{none}\\
S262 & F3.41 & \texttt{-{}B-{}} & 3 & \texttt{aab\textbar{}bab\textbar{}aab} & \texttt{b} & \texttt{b} & \texttt{ab} & \texttt{012} & \texttt{2} & \texttt{+} & \texttt{S} & \texttt{none}\\
S263 & F3.41 & \texttt{-{}B-{}} & 3 & \texttt{aab\textbar{}bab\textbar{}aab} & \texttt{a} & \texttt{a} & \texttt{ab} & \texttt{012} & \texttt{01} & \texttt{+} & \texttt{S} & \texttt{none}\\
S264 & F3.41 & \texttt{-{}B-{}} & 3 & \texttt{aab\textbar{}bba\textbar{}aab} & \texttt{b} & \texttt{b} & \texttt{ab} & \texttt{012} & \texttt{2} & \texttt{+} & \texttt{S} & \texttt{none}\\
S265 & F3.41 & \texttt{-{}B-{}} & 3 & \texttt{aab\textbar{}bba\textbar{}aab} & \texttt{a} & \texttt{a} & \texttt{ab} & \texttt{012} & \texttt{01} & \texttt{+} & \texttt{S} & \texttt{none}\\
S266 & F3.41 & \texttt{-{}B-{}} & 3 & \texttt{aaa\textbar{}aaa\textbar{}aaa} & \texttt{a} & \texttt{a} & \texttt{a} & \texttt{012} & \texttt{012} & \texttt{+} & \texttt{S} & \texttt{none}\\
S267 & F3.56 & \texttt{-{}B-{}} & 3 & \texttt{abc\textbar{}abc\textbar{}abc} & \texttt{a} & \texttt{a} & \texttt{abc} & \texttt{012} & \texttt{0} & \texttt{+} & \texttt{S} & \texttt{none}\\
S268 & F3.56 & \texttt{-{}B-{}} & 3 & \texttt{aab\textbar{}aab\textbar{}aab} & \texttt{b} & \texttt{b} & \texttt{ab} & \texttt{012} & \texttt{2} & \texttt{+} & \texttt{S} & \texttt{none}\\
S269 & F3.56 & \texttt{-{}B-{}} & 3 & \texttt{aab\textbar{}aab\textbar{}aab} & \texttt{a} & \texttt{a} & \texttt{ab} & \texttt{012} & \texttt{01} & \texttt{+} & \texttt{S} & \texttt{none}\\
S270 & F3.56 & \texttt{-{}B-{}} & 3 & \texttt{aaa\textbar{}aaa\textbar{}aaa} & \texttt{a} & \texttt{a} & \texttt{a} & \texttt{012} & \texttt{012} & \texttt{+} & \texttt{S} & \texttt{none}\\
S271 & F3.60 & \texttt{-{}B-{}} & 3 & \texttt{abc\textbar{}abc\textbar{}abc} & \texttt{a} & \texttt{a} & \texttt{abc} & \texttt{012} & \texttt{0} & \texttt{+} & \texttt{S} & \texttt{none}\\
S272 & F3.60 & \texttt{-{}B-{}} & 3 & \texttt{aab\textbar{}aab\textbar{}aab} & \texttt{b} & \texttt{b} & \texttt{ab} & \texttt{012} & \texttt{2} & \texttt{+} & \texttt{S} & \texttt{none}\\
S273 & F3.60 & \texttt{-{}B-{}} & 3 & \texttt{aab\textbar{}aab\textbar{}aab} & \texttt{a} & \texttt{a} & \texttt{ab} & \texttt{012} & \texttt{01} & \texttt{+} & \texttt{S} & \texttt{none}\\
S274 & F3.60 & \texttt{-{}B-{}} & 3 & \texttt{aaa\textbar{}aaa\textbar{}aaa} & \texttt{a} & \texttt{a} & \texttt{a} & \texttt{012} & \texttt{012} & \texttt{+} & \texttt{S} & \texttt{none}\\
S275 & F3.74 & \texttt{TB4} & 3 & \texttt{abc\textbar{}abc\textbar{}abc} & \texttt{a} & \texttt{a} & \texttt{abc} & \texttt{012} & \texttt{0} & \texttt{+} & \texttt{S} & \texttt{none}\\
S276 & F3.74 & \texttt{TB4} & 3 & \texttt{abc\textbar{}abc\textbar{}cba} & \texttt{b} & \texttt{b} & \texttt{abc} & \texttt{012} & \texttt{1} & \texttt{+} & \texttt{S} & \texttt{none}\\
S277 & F3.74 & \texttt{TB4} & 3 & \texttt{abc\textbar{}abc\textbar{}cba} & \texttt{c} & \texttt{c} & \texttt{abc} & \texttt{012} & \texttt{2} & \texttt{+} & \texttt{S} & \texttt{none}\\
S278 & F3.74 & \texttt{TB4} & 3 & \texttt{aab\textbar{}aab\textbar{}aab} & \texttt{b} & \texttt{b} & \texttt{ab} & \texttt{012} & \texttt{2} & \texttt{+} & \texttt{S} & \texttt{none}\\
S279 & F3.74 & \texttt{TB4} & 3 & \texttt{aab\textbar{}aab\textbar{}aab} & \texttt{a} & \texttt{a} & \texttt{ab} & \texttt{012} & \texttt{01} & \texttt{+} & \texttt{S} & \texttt{none}\\
S280 & F3.74 & \texttt{TB4} & 3 & \texttt{abc\textbar{}abc\textbar{}cab} & \texttt{c} & \texttt{c} & \texttt{abc} & \texttt{012} & \texttt{2} & \texttt{+} & \texttt{S} & \texttt{none}\\
S281 & F3.74 & \texttt{TB4} & 3 & \texttt{abb\textbar{}abb\textbar{}bba} & \texttt{a} & \texttt{a} & \texttt{ab} & \texttt{012} & \texttt{0} & \texttt{+} & \texttt{S} & \texttt{none}\\
S282 & F3.74 & \texttt{TB4} & 3 & \texttt{abb\textbar{}abb\textbar{}bba} & \texttt{b} & \texttt{b} & \texttt{ab} & \texttt{012} & \texttt{12} & \texttt{+} & \texttt{S} & \texttt{none}\\
S283 & F3.74 & \texttt{TB4} & 3 & \texttt{aba\textbar{}aba\textbar{}bba} & \texttt{a} & \texttt{a} & \texttt{ab} & \texttt{012} & \texttt{02} & \texttt{+} & \texttt{S} & \texttt{none}\\
S284 & F3.74 & \texttt{TB4} & 3 & \texttt{aba\textbar{}aba\textbar{}bba} & \texttt{b} & \texttt{b} & \texttt{ab} & \texttt{012} & \texttt{1} & \texttt{+} & \texttt{S} & \texttt{none}\\
S285 & F3.74 & \texttt{TB4} & 3 & \texttt{aab\textbar{}aab\textbar{}bba} & \texttt{a} & \texttt{a} & \texttt{ab} & \texttt{012} & \texttt{01} & \texttt{+} & \texttt{S} & \texttt{none}\\
S286 & F3.74 & \texttt{TB4} & 3 & \texttt{aab\textbar{}aab\textbar{}bba} & \texttt{b} & \texttt{b} & \texttt{ab} & \texttt{012} & \texttt{2} & \texttt{+} & \texttt{S} & \texttt{none}\\
S287 & F3.74 & \texttt{TB4} & 3 & \texttt{aaa\textbar{}aaa\textbar{}aaa} & \texttt{a} & \texttt{a} & \texttt{a} & \texttt{012} & \texttt{012} & \texttt{+} & \texttt{S} & \texttt{none}\\
S288 & F3.78 & \texttt{-{}B-{}} & 3 & \texttt{abc\textbar{}abc\textbar{}abc} & \texttt{a} & \texttt{a} & \texttt{abc} & \texttt{012} & \texttt{0} & \texttt{+} & \texttt{S} & \texttt{none}\\
S289 & F3.78 & \texttt{-{}B-{}} & 3 & \texttt{aab\textbar{}aab\textbar{}aab} & \texttt{b} & \texttt{b} & \texttt{ab} & \texttt{012} & \texttt{2} & \texttt{+} & \texttt{S} & \texttt{none}\\
S290 & F3.78 & \texttt{-{}B-{}} & 3 & \texttt{aab\textbar{}aab\textbar{}aab} & \texttt{a} & \texttt{a} & \texttt{ab} & \texttt{012} & \texttt{01} & \texttt{+} & \texttt{S} & \texttt{none}\\
S291 & F3.78 & \texttt{-{}B-{}} & 3 & \texttt{aaa\textbar{}aaa\textbar{}aaa} & \texttt{a} & \texttt{a} & \texttt{a} & \texttt{012} & \texttt{012} & \texttt{+} & \texttt{S} & \texttt{none}\\
S292 & F3.88 & \texttt{-{}B-{}} & 3 & \texttt{abc\textbar{}abc\textbar{}abc} & \texttt{a} & \texttt{a} & \texttt{abc} & \texttt{012} & \texttt{0} & \texttt{+} & \texttt{S} & \texttt{none}\\
S293 & F3.88 & \texttt{-{}B-{}} & 3 & \texttt{aab\textbar{}aab\textbar{}aab} & \texttt{b} & \texttt{b} & \texttt{ab} & \texttt{012} & \texttt{2} & \texttt{+} & \texttt{S} & \texttt{none}\\
S294 & F3.88 & \texttt{-{}B-{}} & 3 & \texttt{aab\textbar{}aab\textbar{}aab} & \texttt{a} & \texttt{a} & \texttt{ab} & \texttt{012} & \texttt{01} & \texttt{+} & \texttt{S} & \texttt{none}\\
S295 & F3.88 & \texttt{-{}B-{}} & 3 & \texttt{aaa\textbar{}aaa\textbar{}aaa} & \texttt{a} & \texttt{a} & \texttt{a} & \texttt{012} & \texttt{012} & \texttt{+} & \texttt{S} & \texttt{none}\\
S296 & F3.89 & \texttt{-{}B-{}} & 3 & \texttt{abc\textbar{}abc\textbar{}abc} & \texttt{a} & \texttt{a} & \texttt{abc} & \texttt{012} & \texttt{0} & \texttt{+} & \texttt{S} & \texttt{none}\\
S297 & F3.89 & \texttt{-{}B-{}} & 3 & \texttt{aab\textbar{}aab\textbar{}aab} & \texttt{b} & \texttt{b} & \texttt{ab} & \texttt{012} & \texttt{2} & \texttt{+} & \texttt{S} & \texttt{none}\\
S298 & F3.89 & \texttt{-{}B-{}} & 3 & \texttt{aab\textbar{}aab\textbar{}aab} & \texttt{a} & \texttt{a} & \texttt{ab} & \texttt{012} & \texttt{01} & \texttt{+} & \texttt{S} & \texttt{none}\\
S299 & F3.89 & \texttt{-{}B-{}} & 3 & \texttt{aaa\textbar{}aaa\textbar{}aaa} & \texttt{a} & \texttt{a} & \texttt{a} & \texttt{012} & \texttt{012} & \texttt{+} & \texttt{S} & \texttt{none}\\
S300 & F3.96 & \texttt{-{}B-{}} & 3 & \texttt{abc\textbar{}def\textbar{}def} & \texttt{ad} & \texttt{ad} & \texttt{abcdef} & \texttt{012} & \texttt{0} & \texttt{+} & \texttt{.} & \texttt{a@0}\\
S301 & F3.96 & \texttt{-{}B-{}} & 3 & \texttt{aba\textbar{}cde\textbar{}edc} & \texttt{bd} & \texttt{bd} & \texttt{abcde} & \texttt{012} & \texttt{1} & \texttt{+} & \texttt{.} & \texttt{b@0}\\
S302 & F3.96 & \texttt{-{}B-{}} & 3 & \texttt{aab\textbar{}ccd\textbar{}ccd} & \texttt{bd} & \texttt{bd} & \texttt{abcd} & \texttt{012} & \texttt{2} & \texttt{+} & \texttt{.} & \texttt{b@0}\\
S303 & F3.96 & \texttt{-{}B-{}} & 3 & \texttt{aab\textbar{}ccd\textbar{}ccd} & \texttt{ac} & \texttt{ac} & \texttt{abcd} & \texttt{012} & \texttt{01} & \texttt{+} & \texttt{.} & \texttt{b@0}\\
S304 & F3.96 & \texttt{-{}B-{}} & 3 & \texttt{aaa\textbar{}bbb\textbar{}bbb} & \texttt{ab} & \texttt{ab} & \texttt{ab} & \texttt{012} & \texttt{012} & \texttt{+} & \texttt{.} & \texttt{a@0}\\
S305 & F3.96 & \texttt{-{}B-{}} & 3 & \texttt{abc\textbar{}abc\textbar{}abc} & \texttt{a} & \texttt{a} & \texttt{abc} & \texttt{012} & \texttt{0} & \texttt{+} & \texttt{S} & \texttt{none}\\
S306 & F3.96 & \texttt{-{}B-{}} & 3 & \texttt{abc\textbar{}cba\textbar{}cba} & \texttt{b} & \texttt{b} & \texttt{abc} & \texttt{012} & \texttt{1} & \texttt{+} & \texttt{S} & \texttt{none}\\
S307 & F3.96 & \texttt{-{}B-{}} & 3 & \texttt{abc\textbar{}cba\textbar{}cba} & \texttt{ac} & \texttt{ac} & \texttt{abc} & \texttt{012} & \texttt{02} & \texttt{+} & \texttt{.} & \texttt{none}\\
S308 & F3.96 & \texttt{-{}B-{}} & 3 & \texttt{aab\textbar{}aab\textbar{}aab} & \texttt{b} & \texttt{b} & \texttt{ab} & \texttt{012} & \texttt{2} & \texttt{+} & \texttt{S} & \texttt{none}\\
S309 & F3.96 & \texttt{-{}B-{}} & 3 & \texttt{aab\textbar{}aab\textbar{}aab} & \texttt{a} & \texttt{a} & \texttt{ab} & \texttt{012} & \texttt{01} & \texttt{+} & \texttt{S} & \texttt{none}\\
S310 & F3.96 & \texttt{-{}B-{}} & 3 & \texttt{aab\textbar{}bba\textbar{}bba} & \texttt{ab} & \texttt{ab} & \texttt{ab} & \texttt{012} & \texttt{012} & \texttt{+} & \texttt{.} & \texttt{none}\\
S311 & F3.96 & \texttt{-{}B-{}} & 3 & \texttt{aaa\textbar{}aaa\textbar{}aaa} & \texttt{a} & \texttt{a} & \texttt{a} & \texttt{012} & \texttt{012} & \texttt{+} & \texttt{S} & \texttt{none}\\
S312 & F3.99 & \texttt{TB-{}} & 3 & \texttt{abc\textbar{}def\textbar{}def} & \texttt{ad} & \texttt{ad} & \texttt{abcdef} & \texttt{012} & \texttt{0} & \texttt{+} & \texttt{.} & \texttt{a@0}\\
S313 & F3.99 & \texttt{TB-{}} & 3 & \texttt{aba\textbar{}cde\textbar{}edc} & \texttt{bd} & \texttt{bd} & \texttt{abcde} & \texttt{012} & \texttt{1} & \texttt{+} & \texttt{.} & \texttt{b@0}\\
S314 & F3.99 & \texttt{TB-{}} & 3 & \texttt{aab\textbar{}ccd\textbar{}ccd} & \texttt{bd} & \texttt{bd} & \texttt{abcd} & \texttt{012} & \texttt{2} & \texttt{+} & \texttt{.} & \texttt{b@0}\\
S315 & F3.99 & \texttt{TB-{}} & 3 & \texttt{aab\textbar{}ccd\textbar{}ccd} & \texttt{ac} & \texttt{ac} & \texttt{abcd} & \texttt{012} & \texttt{01} & \texttt{+} & \texttt{.} & \texttt{b@0}\\
S316 & F3.99 & \texttt{TB-{}} & 3 & \texttt{aaa\textbar{}bbb\textbar{}bbb} & \texttt{ab} & \texttt{ab} & \texttt{ab} & \texttt{012} & \texttt{012} & \texttt{+} & \texttt{.} & \texttt{a@0}\\
S317 & F3.99 & \texttt{TB-{}} & 3 & \texttt{abc\textbar{}abc\textbar{}abc} & \texttt{a} & \texttt{a} & \texttt{abc} & \texttt{012} & \texttt{0} & \texttt{+} & \texttt{S} & \texttt{none}\\
S318 & F3.99 & \texttt{TB-{}} & 3 & \texttt{abc\textbar{}cba\textbar{}cba} & \texttt{b} & \texttt{b} & \texttt{abc} & \texttt{012} & \texttt{1} & \texttt{+} & \texttt{S} & \texttt{none}\\
S319 & F3.99 & \texttt{TB-{}} & 3 & \texttt{abc\textbar{}cba\textbar{}cba} & \texttt{ac} & \texttt{ac} & \texttt{abc} & \texttt{012} & \texttt{02} & \texttt{+} & \texttt{.} & \texttt{none}\\
S320 & F3.99 & \texttt{TB-{}} & 3 & \texttt{aab\textbar{}aab\textbar{}aab} & \texttt{b} & \texttt{b} & \texttt{ab} & \texttt{012} & \texttt{2} & \texttt{+} & \texttt{S} & \texttt{none}\\
S321 & F3.99 & \texttt{TB-{}} & 3 & \texttt{aab\textbar{}aab\textbar{}aab} & \texttt{a} & \texttt{a} & \texttt{ab} & \texttt{012} & \texttt{01} & \texttt{+} & \texttt{S} & \texttt{none}\\
S322 & F3.99 & \texttt{TB-{}} & 3 & \texttt{aab\textbar{}bba\textbar{}bba} & \texttt{ab} & \texttt{ab} & \texttt{ab} & \texttt{012} & \texttt{012} & \texttt{+} & \texttt{.} & \texttt{none}\\
S323 & F3.99 & \texttt{TB-{}} & 3 & \texttt{aaa\textbar{}aaa\textbar{}aaa} & \texttt{a} & \texttt{a} & \texttt{a} & \texttt{012} & \texttt{012} & \texttt{+} & \texttt{S} & \texttt{none}\\
S324 & F3.102 & \texttt{-{}B-{}} & 3 & \texttt{abc\textbar{}def\textbar{}abc} & \texttt{ad} & \texttt{ad} & \texttt{abcdef} & \texttt{012} & \texttt{0} & \texttt{+} & \texttt{.} & \texttt{d@1}\\
S325 & F3.102 & \texttt{-{}B-{}} & 3 & \texttt{aab\textbar{}ccd\textbar{}aab} & \texttt{bd} & \texttt{bd} & \texttt{abcd} & \texttt{012} & \texttt{2} & \texttt{+} & \texttt{.} & \texttt{d@1}\\
S326 & F3.102 & \texttt{-{}B-{}} & 3 & \texttt{aab\textbar{}ccd\textbar{}aab} & \texttt{ac} & \texttt{ac} & \texttt{abcd} & \texttt{012} & \texttt{01} & \texttt{+} & \texttt{.} & \texttt{d@1}\\
S327 & F3.102 & \texttt{-{}B-{}} & 3 & \texttt{aaa\textbar{}bbb\textbar{}aaa} & \texttt{ab} & \texttt{ab} & \texttt{ab} & \texttt{012} & \texttt{012} & \texttt{+} & \texttt{.} & \texttt{b@1}\\
S328 & F3.102 & \texttt{-{}B-{}} & 3 & \texttt{abc\textbar{}abc\textbar{}abc} & \texttt{a} & \texttt{a} & \texttt{abc} & \texttt{012} & \texttt{0} & \texttt{+} & \texttt{S} & \texttt{none}\\
S329 & F3.102 & \texttt{-{}B-{}} & 3 & \texttt{abc\textbar{}cba\textbar{}abc} & \texttt{b} & \texttt{b} & \texttt{abc} & \texttt{012} & \texttt{1} & \texttt{+} & \texttt{S} & \texttt{none}\\
S330 & F3.102 & \texttt{-{}B-{}} & 3 & \texttt{abc\textbar{}cba\textbar{}abc} & \texttt{ac} & \texttt{ac} & \texttt{abc} & \texttt{012} & \texttt{02} & \texttt{+} & \texttt{.} & \texttt{none}\\
S331 & F3.102 & \texttt{-{}B-{}} & 3 & \texttt{aab\textbar{}aab\textbar{}aab} & \texttt{b} & \texttt{b} & \texttt{ab} & \texttt{012} & \texttt{2} & \texttt{+} & \texttt{S} & \texttt{none}\\
S332 & F3.102 & \texttt{-{}B-{}} & 3 & \texttt{aab\textbar{}aab\textbar{}aab} & \texttt{a} & \texttt{a} & \texttt{ab} & \texttt{012} & \texttt{01} & \texttt{+} & \texttt{S} & \texttt{none}\\
S333 & F3.102 & \texttt{-{}B-{}} & 3 & \texttt{aab\textbar{}bba\textbar{}aab} & \texttt{ab} & \texttt{ab} & \texttt{ab} & \texttt{012} & \texttt{012} & \texttt{+} & \texttt{.} & \texttt{none}\\
S334 & F3.102 & \texttt{-{}B-{}} & 3 & \texttt{abb\textbar{}baa\textbar{}aab} & \texttt{ab} & \texttt{ab} & \texttt{ab} & \texttt{012} & \texttt{012} & \texttt{+} & \texttt{.} & \texttt{none}\\
S335 & F3.102 & \texttt{-{}B-{}} & 3 & \texttt{aaa\textbar{}aaa\textbar{}aaa} & \texttt{a} & \texttt{a} & \texttt{a} & \texttt{012} & \texttt{012} & \texttt{+} & \texttt{S} & \texttt{none}\\
S336 & F3.104 & \texttt{TB4} & 3 & \texttt{abc\textbar{}def\textbar{}ghi} & \texttt{adg} & \texttt{adg} & \texttt{abcdefghi} & \texttt{012} & \texttt{0} & \texttt{+} & \texttt{.} & \texttt{a@0}\\
S337 & F3.104 & \texttt{TB4} & 3 & \texttt{abc\textbar{}abc\textbar{}def} & \texttt{ad} & \texttt{ad} & \texttt{abcdef} & \texttt{012} & \texttt{0} & \texttt{+} & \texttt{.} & \texttt{d@2}\\
S338 & F3.104 & \texttt{TB4} & 3 & \texttt{aab\textbar{}ccd\textbar{}eef} & \texttt{bdf} & \texttt{bdf} & \texttt{abcdef} & \texttt{012} & \texttt{2} & \texttt{+} & \texttt{.} & \texttt{b@0}\\
S339 & F3.104 & \texttt{TB4} & 3 & \texttt{aab\textbar{}ccd\textbar{}eef} & \texttt{ace} & \texttt{ace} & \texttt{abcdef} & \texttt{012} & \texttt{01} & \texttt{+} & \texttt{.} & \texttt{b@0}\\
S340 & F3.104 & \texttt{TB4} & 3 & \texttt{abc\textbar{}dde\textbar{}bac} & \texttt{ce} & \texttt{ce} & \texttt{abcde} & \texttt{012} & \texttt{2} & \texttt{+} & \texttt{.} & \texttt{e@1}\\
S341 & F3.104 & \texttt{TB4} & 3 & \texttt{abc\textbar{}dde\textbar{}bac} & \texttt{abd} & \texttt{abd} & \texttt{abcde} & \texttt{012} & \texttt{01} & \texttt{+} & \texttt{.} & \texttt{e@1}\\
S342 & F3.104 & \texttt{TB4} & 3 & \texttt{aab\textbar{}aab\textbar{}ccd} & \texttt{bd} & \texttt{bd} & \texttt{abcd} & \texttt{012} & \texttt{2} & \texttt{+} & \texttt{.} & \texttt{d@2}\\
S343 & F3.104 & \texttt{TB4} & 3 & \texttt{aab\textbar{}aab\textbar{}ccd} & \texttt{ac} & \texttt{ac} & \texttt{abcd} & \texttt{012} & \texttt{01} & \texttt{+} & \texttt{.} & \texttt{d@2}\\
S344 & F3.104 & \texttt{TB4} & 3 & \texttt{aaa\textbar{}bbb\textbar{}ccc} & \texttt{abc} & \texttt{abc} & \texttt{abc} & \texttt{012} & \texttt{012} & \texttt{+} & \texttt{.} & \texttt{a@0}\\
S345 & F3.104 & \texttt{TB4} & 3 & \texttt{aab\textbar{}ccc\textbar{}bba} & \texttt{abc} & \texttt{abc} & \texttt{abc} & \texttt{012} & \texttt{012} & \texttt{+} & \texttt{.} & \texttt{c@1}\\
S346 & F3.104 & \texttt{TB4} & 3 & \texttt{aaa\textbar{}aaa\textbar{}bbb} & \texttt{ab} & \texttt{ab} & \texttt{ab} & \texttt{012} & \texttt{012} & \texttt{+} & \texttt{.} & \texttt{b@2}\\
S347 & F3.104 & \texttt{TB4} & 3 & \texttt{abc\textbar{}abc\textbar{}abc} & \texttt{a} & \texttt{a} & \texttt{abc} & \texttt{012} & \texttt{0} & \texttt{+} & \texttt{S} & \texttt{none}\\
S348 & F3.104 & \texttt{TB4} & 3 & \texttt{abc\textbar{}abc\textbar{}cba} & \texttt{b} & \texttt{b} & \texttt{abc} & \texttt{012} & \texttt{1} & \texttt{+} & \texttt{S} & \texttt{none}\\
S349 & F3.104 & \texttt{TB4} & 3 & \texttt{abc\textbar{}abc\textbar{}cba} & \texttt{ac} & \texttt{ac} & \texttt{abc} & \texttt{012} & \texttt{02} & \texttt{+} & \texttt{.} & \texttt{none}\\
S350 & F3.104 & \texttt{TB4} & 3 & \texttt{aab\textbar{}aab\textbar{}aab} & \texttt{b} & \texttt{b} & \texttt{ab} & \texttt{012} & \texttt{2} & \texttt{+} & \texttt{S} & \texttt{none}\\
S351 & F3.104 & \texttt{TB4} & 3 & \texttt{aab\textbar{}aab\textbar{}aab} & \texttt{a} & \texttt{a} & \texttt{ab} & \texttt{012} & \texttt{01} & \texttt{+} & \texttt{S} & \texttt{none}\\
S352 & F3.104 & \texttt{TB4} & 3 & \texttt{aab\textbar{}aab\textbar{}bba} & \texttt{ab} & \texttt{ab} & \texttt{ab} & \texttt{012} & \texttt{012} & \texttt{+} & \texttt{.} & \texttt{none}\\
S353 & F3.104 & \texttt{TB4} & 3 & \texttt{abc\textbar{}cab\textbar{}bca} & \texttt{abc} & \texttt{abc} & \texttt{abc} & \texttt{012} & \texttt{012} & \texttt{+} & \texttt{.} & \texttt{none}\\
S354 & F3.104 & \texttt{TB4} & 3 & \texttt{aba\textbar{}baa\textbar{}aab} & \texttt{ab} & \texttt{ab} & \texttt{ab} & \texttt{012} & \texttt{012} & \texttt{+} & \texttt{.} & \texttt{none}\\
S355 & F3.104 & \texttt{TB4} & 3 & \texttt{aab\textbar{}abb\textbar{}bba} & \texttt{ab} & \texttt{ab} & \texttt{ab} & \texttt{012} & \texttt{012} & \texttt{+} & \texttt{.} & \texttt{none}\\
S356 & F3.104 & \texttt{TB4} & 3 & \texttt{aaa\textbar{}aaa\textbar{}aaa} & \texttt{a} & \texttt{a} & \texttt{a} & \texttt{012} & \texttt{012} & \texttt{+} & \texttt{S} & \texttt{none}\\
\multicolumn{13}{@{}l}{\rule{0pt}{2.8ex}\footnotesize\bfseries $n=4$ worlds, $m=1$ individuals: 125 models}\\*
S357 & F4.133 & \texttt{-{}B-{}} & 1 & \texttt{a\textbar{}a\textbar{}a\textbar{}a} & \texttt{a} & \texttt{a} & \texttt{a} & \texttt{0123} & \texttt{0} & \texttt{+} & \texttt{S} & \texttt{none}\\
S358 & F4.137 & \texttt{-{}B-{}} & 1 & \texttt{a\textbar{}a\textbar{}a\textbar{}a} & \texttt{a} & \texttt{a} & \texttt{a} & \texttt{0123} & \texttt{0} & \texttt{+} & \texttt{S} & \texttt{none}\\
S359 & F4.139 & \texttt{TB4} & 1 & \texttt{a\textbar{}a\textbar{}a\textbar{}a} & \texttt{a} & \texttt{a} & \texttt{a} & \texttt{0123} & \texttt{0} & \texttt{+} & \texttt{S} & \texttt{none}\\
S360 & F4.276 & \texttt{-{}B-{}} & 1 & \texttt{a\textbar{}a\textbar{}a\textbar{}a} & \texttt{a} & \texttt{a} & \texttt{a} & \texttt{0123} & \texttt{0} & \texttt{+} & \texttt{S} & \texttt{none}\\
S361 & F4.306 & \texttt{-{}B-{}} & 1 & \texttt{a\textbar{}a\textbar{}a\textbar{}a} & \texttt{a} & \texttt{a} & \texttt{a} & \texttt{0123} & \texttt{0} & \texttt{+} & \texttt{S} & \texttt{none}\\
S362 & F4.321 & \texttt{-{}B-{}} & 1 & \texttt{a\textbar{}a\textbar{}a\textbar{}a} & \texttt{a} & \texttt{a} & \texttt{a} & \texttt{0123} & \texttt{0} & \texttt{+} & \texttt{S} & \texttt{none}\\
S363 & F4.505 & \texttt{-{}B-{}} & 1 & \texttt{a\textbar{}a\textbar{}a\textbar{}a} & \texttt{a} & \texttt{a} & \texttt{a} & \texttt{0123} & \texttt{0} & \texttt{+} & \texttt{S} & \texttt{none}\\
S364 & F4.507 & \texttt{-{}B-{}} & 1 & \texttt{a\textbar{}a\textbar{}a\textbar{}a} & \texttt{a} & \texttt{a} & \texttt{a} & \texttt{0123} & \texttt{0} & \texttt{+} & \texttt{S} & \texttt{none}\\
S365 & F4.559 & \texttt{-{}B-{}} & 1 & \texttt{a\textbar{}a\textbar{}a\textbar{}a} & \texttt{a} & \texttt{a} & \texttt{a} & \texttt{0123} & \texttt{0} & \texttt{+} & \texttt{S} & \texttt{none}\\
S366 & F4.563 & \texttt{-{}B-{}} & 1 & \texttt{a\textbar{}a\textbar{}a\textbar{}a} & \texttt{a} & \texttt{a} & \texttt{a} & \texttt{0123} & \texttt{0} & \texttt{+} & \texttt{S} & \texttt{none}\\
S367 & F4.616 & \texttt{-{}B-{}} & 1 & \texttt{a\textbar{}a\textbar{}a\textbar{}a} & \texttt{a} & \texttt{a} & \texttt{a} & \texttt{0123} & \texttt{0} & \texttt{+} & \texttt{S} & \texttt{none}\\
S368 & F4.679 & \texttt{-{}B-{}} & 1 & \texttt{a\textbar{}a\textbar{}a\textbar{}a} & \texttt{a} & \texttt{a} & \texttt{a} & \texttt{0123} & \texttt{0} & \texttt{+} & \texttt{S} & \texttt{none}\\
S369 & F4.712 & \texttt{TB4} & 1 & \texttt{a\textbar{}a\textbar{}a\textbar{}a} & \texttt{a} & \texttt{a} & \texttt{a} & \texttt{0123} & \texttt{0} & \texttt{+} & \texttt{S} & \texttt{none}\\
S370 & F4.857 & \texttt{-{}B-{}} & 1 & \texttt{a\textbar{}a\textbar{}a\textbar{}a} & \texttt{a} & \texttt{a} & \texttt{a} & \texttt{0123} & \texttt{0} & \texttt{+} & \texttt{S} & \texttt{none}\\
S371 & F4.861 & \texttt{-{}B-{}} & 1 & \texttt{a\textbar{}a\textbar{}a\textbar{}a} & \texttt{a} & \texttt{a} & \texttt{a} & \texttt{0123} & \texttt{0} & \texttt{+} & \texttt{S} & \texttt{none}\\
S372 & F4.898 & \texttt{-{}B-{}} & 1 & \texttt{a\textbar{}a\textbar{}a\textbar{}a} & \texttt{a} & \texttt{a} & \texttt{a} & \texttt{0123} & \texttt{0} & \texttt{+} & \texttt{S} & \texttt{none}\\
S373 & F4.937 & \texttt{-{}B-{}} & 1 & \texttt{a\textbar{}a\textbar{}a\textbar{}a} & \texttt{a} & \texttt{a} & \texttt{a} & \texttt{0123} & \texttt{0} & \texttt{+} & \texttt{S} & \texttt{none}\\
S374 & F4.978 & \texttt{-{}B-{}} & 1 & \texttt{a\textbar{}a\textbar{}a\textbar{}a} & \texttt{a} & \texttt{a} & \texttt{a} & \texttt{0123} & \texttt{0} & \texttt{+} & \texttt{S} & \texttt{none}\\
S375 & F4.1075 & \texttt{-{}B-{}} & 1 & \texttt{a\textbar{}a\textbar{}a\textbar{}a} & \texttt{a} & \texttt{a} & \texttt{a} & \texttt{0123} & \texttt{0} & \texttt{+} & \texttt{S} & \texttt{none}\\
S376 & F4.1195 & \texttt{-{}B-{}} & 1 & \texttt{a\textbar{}a\textbar{}a\textbar{}a} & \texttt{a} & \texttt{a} & \texttt{a} & \texttt{0123} & \texttt{0} & \texttt{+} & \texttt{S} & \texttt{none}\\
S377 & F4.1196 & \texttt{-{}B-{}} & 1 & \texttt{a\textbar{}a\textbar{}a\textbar{}a} & \texttt{a} & \texttt{a} & \texttt{a} & \texttt{0123} & \texttt{0} & \texttt{+} & \texttt{S} & \texttt{none}\\
S378 & F4.1321 & \texttt{-{}B-{}} & 1 & \texttt{a\textbar{}a\textbar{}a\textbar{}a} & \texttt{a} & \texttt{a} & \texttt{a} & \texttt{0123} & \texttt{0} & \texttt{+} & \texttt{S} & \texttt{none}\\
S379 & F4.1343 & \texttt{-{}B-{}} & 1 & \texttt{a\textbar{}a\textbar{}a\textbar{}a} & \texttt{a} & \texttt{a} & \texttt{a} & \texttt{0123} & \texttt{0} & \texttt{+} & \texttt{S} & \texttt{none}\\
S380 & F4.1362 & \texttt{-{}B-{}} & 1 & \texttt{a\textbar{}a\textbar{}a\textbar{}a} & \texttt{a} & \texttt{a} & \texttt{a} & \texttt{0123} & \texttt{0} & \texttt{+} & \texttt{S} & \texttt{none}\\
S381 & F4.1365 & \texttt{-{}B-{}} & 1 & \texttt{a\textbar{}a\textbar{}a\textbar{}a} & \texttt{a} & \texttt{a} & \texttt{a} & \texttt{0123} & \texttt{0} & \texttt{+} & \texttt{S} & \texttt{none}\\
S382 & F4.1437 & \texttt{-{}B-{}} & 1 & \texttt{a\textbar{}a\textbar{}a\textbar{}a} & \texttt{a} & \texttt{a} & \texttt{a} & \texttt{0123} & \texttt{0} & \texttt{+} & \texttt{S} & \texttt{none}\\
S383 & F4.1490 & \texttt{-{}B-{}} & 1 & \texttt{a\textbar{}a\textbar{}a\textbar{}a} & \texttt{a} & \texttt{a} & \texttt{a} & \texttt{0123} & \texttt{0} & \texttt{+} & \texttt{S} & \texttt{none}\\
S384 & F4.1675 & \texttt{-{}B-{}} & 1 & \texttt{a\textbar{}b\textbar{}b\textbar{}a} & \texttt{ab} & \texttt{ab} & \texttt{ab} & \texttt{0123} & \texttt{0} & \texttt{+} & \texttt{.} & \texttt{b@1}\\
S385 & F4.1675 & \texttt{-{}B-{}} & 1 & \texttt{a\textbar{}a\textbar{}a\textbar{}a} & \texttt{a} & \texttt{a} & \texttt{a} & \texttt{0123} & \texttt{0} & \texttt{+} & \texttt{S} & \texttt{none}\\
S386 & F4.1692 & \texttt{-{}B-{}} & 1 & \texttt{a\textbar{}a\textbar{}a\textbar{}a} & \texttt{a} & \texttt{a} & \texttt{a} & \texttt{0123} & \texttt{0} & \texttt{+} & \texttt{S} & \texttt{none}\\
S387 & F4.1732 & \texttt{-{}B-{}} & 1 & \texttt{a\textbar{}b\textbar{}b\textbar{}a} & \texttt{ab} & \texttt{ab} & \texttt{ab} & \texttt{0123} & \texttt{0} & \texttt{+} & \texttt{.} & \texttt{b@1}\\
S388 & F4.1732 & \texttt{-{}B-{}} & 1 & \texttt{a\textbar{}a\textbar{}a\textbar{}a} & \texttt{a} & \texttt{a} & \texttt{a} & \texttt{0123} & \texttt{0} & \texttt{+} & \texttt{S} & \texttt{none}\\
S389 & F4.1794 & \texttt{TB-{}} & 1 & \texttt{a\textbar{}a\textbar{}a\textbar{}a} & \texttt{a} & \texttt{a} & \texttt{a} & \texttt{0123} & \texttt{0} & \texttt{+} & \texttt{S} & \texttt{none}\\
S390 & F4.1813 & \texttt{TB4} & 1 & \texttt{a\textbar{}b\textbar{}a\textbar{}b} & \texttt{ab} & \texttt{ab} & \texttt{ab} & \texttt{0123} & \texttt{0} & \texttt{+} & \texttt{.} & \texttt{a@0}\\
S391 & F4.1813 & \texttt{TB4} & 1 & \texttt{a\textbar{}a\textbar{}a\textbar{}a} & \texttt{a} & \texttt{a} & \texttt{a} & \texttt{0123} & \texttt{0} & \texttt{+} & \texttt{S} & \texttt{none}\\
S392 & F4.1869 & \texttt{-{}B-{}} & 1 & \texttt{a\textbar{}a\textbar{}a\textbar{}a} & \texttt{a} & \texttt{a} & \texttt{a} & \texttt{0123} & \texttt{0} & \texttt{+} & \texttt{S} & \texttt{none}\\
S393 & F4.1874 & \texttt{-{}B-{}} & 1 & \texttt{a\textbar{}a\textbar{}a\textbar{}a} & \texttt{a} & \texttt{a} & \texttt{a} & \texttt{0123} & \texttt{0} & \texttt{+} & \texttt{S} & \texttt{none}\\
S394 & F4.1904 & \texttt{-{}B-{}} & 1 & \texttt{a\textbar{}a\textbar{}a\textbar{}a} & \texttt{a} & \texttt{a} & \texttt{a} & \texttt{0123} & \texttt{0} & \texttt{+} & \texttt{S} & \texttt{none}\\
S395 & F4.1933 & \texttt{-{}B-{}} & 1 & \texttt{a\textbar{}a\textbar{}a\textbar{}a} & \texttt{a} & \texttt{a} & \texttt{a} & \texttt{0123} & \texttt{0} & \texttt{+} & \texttt{S} & \texttt{none}\\
S396 & F4.2109 & \texttt{-{}B-{}} & 1 & \texttt{a\textbar{}b\textbar{}b\textbar{}a} & \texttt{ab} & \texttt{ab} & \texttt{ab} & \texttt{0123} & \texttt{0} & \texttt{+} & \texttt{.} & \texttt{b@1}\\
S397 & F4.2109 & \texttt{-{}B-{}} & 1 & \texttt{a\textbar{}a\textbar{}a\textbar{}a} & \texttt{a} & \texttt{a} & \texttt{a} & \texttt{0123} & \texttt{0} & \texttt{+} & \texttt{S} & \texttt{none}\\
S398 & F4.2120 & \texttt{-{}B-{}} & 1 & \texttt{a\textbar{}a\textbar{}a\textbar{}a} & \texttt{a} & \texttt{a} & \texttt{a} & \texttt{0123} & \texttt{0} & \texttt{+} & \texttt{S} & \texttt{none}\\
S399 & F4.2148 & \texttt{-{}B-{}} & 1 & \texttt{a\textbar{}a\textbar{}a\textbar{}a} & \texttt{a} & \texttt{a} & \texttt{a} & \texttt{0123} & \texttt{0} & \texttt{+} & \texttt{S} & \texttt{none}\\
S400 & F4.2151 & \texttt{-{}B-{}} & 1 & \texttt{a\textbar{}b\textbar{}b\textbar{}a} & \texttt{ab} & \texttt{ab} & \texttt{ab} & \texttt{0123} & \texttt{0} & \texttt{+} & \texttt{.} & \texttt{b@1}\\
S401 & F4.2151 & \texttt{-{}B-{}} & 1 & \texttt{a\textbar{}a\textbar{}a\textbar{}a} & \texttt{a} & \texttt{a} & \texttt{a} & \texttt{0123} & \texttt{0} & \texttt{+} & \texttt{S} & \texttt{none}\\
S402 & F4.2291 & \texttt{-{}B-{}} & 1 & \texttt{a\textbar{}a\textbar{}a\textbar{}a} & \texttt{a} & \texttt{a} & \texttt{a} & \texttt{0123} & \texttt{0} & \texttt{+} & \texttt{S} & \texttt{none}\\
S403 & F4.2346 & \texttt{-{}B-{}} & 1 & \texttt{a\textbar{}a\textbar{}a\textbar{}a} & \texttt{a} & \texttt{a} & \texttt{a} & \texttt{0123} & \texttt{0} & \texttt{+} & \texttt{S} & \texttt{none}\\
S404 & F4.2373 & \texttt{-{}B-{}} & 1 & \texttt{a\textbar{}a\textbar{}a\textbar{}a} & \texttt{a} & \texttt{a} & \texttt{a} & \texttt{0123} & \texttt{0} & \texttt{+} & \texttt{S} & \texttt{none}\\
S405 & F4.2467 & \texttt{-{}B-{}} & 1 & \texttt{a\textbar{}b\textbar{}b\textbar{}a} & \texttt{ab} & \texttt{ab} & \texttt{ab} & \texttt{0123} & \texttt{0} & \texttt{+} & \texttt{.} & \texttt{b@1}\\
S406 & F4.2467 & \texttt{-{}B-{}} & 1 & \texttt{a\textbar{}a\textbar{}a\textbar{}a} & \texttt{a} & \texttt{a} & \texttt{a} & \texttt{0123} & \texttt{0} & \texttt{+} & \texttt{S} & \texttt{none}\\
S407 & F4.2468 & \texttt{-{}B-{}} & 1 & \texttt{a\textbar{}b\textbar{}b\textbar{}a} & \texttt{ab} & \texttt{ab} & \texttt{ab} & \texttt{0123} & \texttt{0} & \texttt{+} & \texttt{.} & \texttt{b@1}\\
S408 & F4.2468 & \texttt{-{}B-{}} & 1 & \texttt{a\textbar{}a\textbar{}a\textbar{}a} & \texttt{a} & \texttt{a} & \texttt{a} & \texttt{0123} & \texttt{0} & \texttt{+} & \texttt{S} & \texttt{none}\\
S409 & F4.2480 & \texttt{-{}B-{}} & 1 & \texttt{a\textbar{}b\textbar{}b\textbar{}b} & \texttt{ab} & \texttt{ab} & \texttt{ab} & \texttt{0123} & \texttt{0} & \texttt{+} & \texttt{.} & \texttt{a@0}\\
S410 & F4.2480 & \texttt{-{}B-{}} & 1 & \texttt{a\textbar{}a\textbar{}a\textbar{}a} & \texttt{a} & \texttt{a} & \texttt{a} & \texttt{0123} & \texttt{0} & \texttt{+} & \texttt{S} & \texttt{none}\\
S411 & F4.2485 & \texttt{-{}B-{}} & 1 & \texttt{a\textbar{}b\textbar{}b\textbar{}a} & \texttt{ab} & \texttt{ab} & \texttt{ab} & \texttt{0123} & \texttt{0} & \texttt{+} & \texttt{.} & \texttt{b@1}\\
S412 & F4.2485 & \texttt{-{}B-{}} & 1 & \texttt{a\textbar{}a\textbar{}a\textbar{}a} & \texttt{a} & \texttt{a} & \texttt{a} & \texttt{0123} & \texttt{0} & \texttt{+} & \texttt{S} & \texttt{none}\\
S413 & F4.2524 & \texttt{-{}B-{}} & 1 & \texttt{a\textbar{}b\textbar{}a\textbar{}b} & \texttt{ab} & \texttt{ab} & \texttt{ab} & \texttt{0123} & \texttt{0} & \texttt{+} & \texttt{.} & \texttt{a@0}\\
S414 & F4.2524 & \texttt{-{}B-{}} & 1 & \texttt{a\textbar{}b\textbar{}b\textbar{}a} & \texttt{ab} & \texttt{ab} & \texttt{ab} & \texttt{0123} & \texttt{0} & \texttt{+} & \texttt{.} & \texttt{b@1}\\
S415 & F4.2524 & \texttt{-{}B-{}} & 1 & \texttt{a\textbar{}a\textbar{}a\textbar{}a} & \texttt{a} & \texttt{a} & \texttt{a} & \texttt{0123} & \texttt{0} & \texttt{+} & \texttt{S} & \texttt{none}\\
S416 & F4.2662 & \texttt{TB4} & 1 & \texttt{a\textbar{}a\textbar{}a\textbar{}a} & \texttt{a} & \texttt{a} & \texttt{a} & \texttt{0123} & \texttt{0} & \texttt{+} & \texttt{S} & \texttt{none}\\
S417 & F4.2672 & \texttt{TB-{}} & 1 & \texttt{a\textbar{}b\textbar{}b\textbar{}b} & \texttt{ab} & \texttt{ab} & \texttt{ab} & \texttt{0123} & \texttt{0} & \texttt{+} & \texttt{.} & \texttt{a@0}\\
S418 & F4.2672 & \texttt{TB-{}} & 1 & \texttt{a\textbar{}a\textbar{}a\textbar{}a} & \texttt{a} & \texttt{a} & \texttt{a} & \texttt{0123} & \texttt{0} & \texttt{+} & \texttt{S} & \texttt{none}\\
S419 & F4.2683 & \texttt{TB-{}} & 1 & \texttt{a\textbar{}a\textbar{}b\textbar{}b} & \texttt{ab} & \texttt{ab} & \texttt{ab} & \texttt{0123} & \texttt{0} & \texttt{+} & \texttt{.} & \texttt{a@0}\\
S420 & F4.2683 & \texttt{TB-{}} & 1 & \texttt{a\textbar{}b\textbar{}a\textbar{}b} & \texttt{ab} & \texttt{ab} & \texttt{ab} & \texttt{0123} & \texttt{0} & \texttt{+} & \texttt{.} & \texttt{a@0}\\
S421 & F4.2683 & \texttt{TB-{}} & 1 & \texttt{a\textbar{}a\textbar{}a\textbar{}a} & \texttt{a} & \texttt{a} & \texttt{a} & \texttt{0123} & \texttt{0} & \texttt{+} & \texttt{S} & \texttt{none}\\
S422 & F4.2709 & \texttt{-{}B-{}} & 1 & \texttt{a\textbar{}a\textbar{}a\textbar{}a} & \texttt{a} & \texttt{a} & \texttt{a} & \texttt{0123} & \texttt{0} & \texttt{+} & \texttt{S} & \texttt{none}\\
S423 & F4.2743 & \texttt{-{}B-{}} & 1 & \texttt{a\textbar{}b\textbar{}b\textbar{}b} & \texttt{ab} & \texttt{ab} & \texttt{ab} & \texttt{0123} & \texttt{0} & \texttt{+} & \texttt{.} & \texttt{a@0}\\
S424 & F4.2743 & \texttt{-{}B-{}} & 1 & \texttt{a\textbar{}b\textbar{}b\textbar{}a} & \texttt{ab} & \texttt{ab} & \texttt{ab} & \texttt{0123} & \texttt{0} & \texttt{+} & \texttt{.} & \texttt{b@1}\\
S425 & F4.2743 & \texttt{-{}B-{}} & 1 & \texttt{a\textbar{}a\textbar{}a\textbar{}a} & \texttt{a} & \texttt{a} & \texttt{a} & \texttt{0123} & \texttt{0} & \texttt{+} & \texttt{S} & \texttt{none}\\
S426 & F4.2747 & \texttt{-{}B-{}} & 1 & \texttt{a\textbar{}b\textbar{}b\textbar{}b} & \texttt{ab} & \texttt{ab} & \texttt{ab} & \texttt{0123} & \texttt{0} & \texttt{+} & \texttt{.} & \texttt{a@0}\\
S427 & F4.2747 & \texttt{-{}B-{}} & 1 & \texttt{a\textbar{}b\textbar{}b\textbar{}a} & \texttt{ab} & \texttt{ab} & \texttt{ab} & \texttt{0123} & \texttt{0} & \texttt{+} & \texttt{.} & \texttt{b@1}\\
S428 & F4.2747 & \texttt{-{}B-{}} & 1 & \texttt{a\textbar{}a\textbar{}a\textbar{}a} & \texttt{a} & \texttt{a} & \texttt{a} & \texttt{0123} & \texttt{0} & \texttt{+} & \texttt{S} & \texttt{none}\\
S429 & F4.2768 & \texttt{-{}B-{}} & 1 & \texttt{a\textbar{}b\textbar{}b\textbar{}a} & \texttt{ab} & \texttt{ab} & \texttt{ab} & \texttt{0123} & \texttt{0} & \texttt{+} & \texttt{.} & \texttt{b@1}\\
S430 & F4.2768 & \texttt{-{}B-{}} & 1 & \texttt{a\textbar{}a\textbar{}a\textbar{}a} & \texttt{a} & \texttt{a} & \texttt{a} & \texttt{0123} & \texttt{0} & \texttt{+} & \texttt{S} & \texttt{none}\\
S431 & F4.2769 & \texttt{-{}B-{}} & 1 & \texttt{a\textbar{}b\textbar{}b\textbar{}a} & \texttt{ab} & \texttt{ab} & \texttt{ab} & \texttt{0123} & \texttt{0} & \texttt{+} & \texttt{.} & \texttt{b@1}\\
S432 & F4.2769 & \texttt{-{}B-{}} & 1 & \texttt{a\textbar{}a\textbar{}a\textbar{}a} & \texttt{a} & \texttt{a} & \texttt{a} & \texttt{0123} & \texttt{0} & \texttt{+} & \texttt{S} & \texttt{none}\\
S433 & F4.2865 & \texttt{-{}B-{}} & 1 & \texttt{a\textbar{}b\textbar{}b\textbar{}a} & \texttt{ab} & \texttt{ab} & \texttt{ab} & \texttt{0123} & \texttt{0} & \texttt{+} & \texttt{.} & \texttt{b@1}\\
S434 & F4.2865 & \texttt{-{}B-{}} & 1 & \texttt{a\textbar{}a\textbar{}a\textbar{}a} & \texttt{a} & \texttt{a} & \texttt{a} & \texttt{0123} & \texttt{0} & \texttt{+} & \texttt{S} & \texttt{none}\\
S435 & F4.2914 & \texttt{-{}B-{}} & 1 & \texttt{a\textbar{}b\textbar{}a\textbar{}a} & \texttt{ab} & \texttt{ab} & \texttt{ab} & \texttt{0123} & \texttt{0} & \texttt{+} & \texttt{.} & \texttt{b@1}\\
S436 & F4.2914 & \texttt{-{}B-{}} & 1 & \texttt{a\textbar{}b\textbar{}b\textbar{}a} & \texttt{ab} & \texttt{ab} & \texttt{ab} & \texttt{0123} & \texttt{0} & \texttt{+} & \texttt{.} & \texttt{b@1}\\
S437 & F4.2914 & \texttt{-{}B-{}} & 1 & \texttt{a\textbar{}a\textbar{}a\textbar{}a} & \texttt{a} & \texttt{a} & \texttt{a} & \texttt{0123} & \texttt{0} & \texttt{+} & \texttt{S} & \texttt{none}\\
S438 & F4.2918 & \texttt{-{}B-{}} & 1 & \texttt{a\textbar{}b\textbar{}b\textbar{}b} & \texttt{ab} & \texttt{ab} & \texttt{ab} & \texttt{0123} & \texttt{0} & \texttt{+} & \texttt{.} & \texttt{a@0}\\
S439 & F4.2918 & \texttt{-{}B-{}} & 1 & \texttt{a\textbar{}b\textbar{}a\textbar{}b} & \texttt{ab} & \texttt{ab} & \texttt{ab} & \texttt{0123} & \texttt{0} & \texttt{+} & \texttt{.} & \texttt{a@0}\\
S440 & F4.2918 & \texttt{-{}B-{}} & 1 & \texttt{a\textbar{}b\textbar{}b\textbar{}a} & \texttt{ab} & \texttt{ab} & \texttt{ab} & \texttt{0123} & \texttt{0} & \texttt{+} & \texttt{.} & \texttt{b@1}\\
S441 & F4.2918 & \texttt{-{}B-{}} & 1 & \texttt{a\textbar{}a\textbar{}a\textbar{}a} & \texttt{a} & \texttt{a} & \texttt{a} & \texttt{0123} & \texttt{0} & \texttt{+} & \texttt{S} & \texttt{none}\\
S442 & F4.2968 & \texttt{-{}B-{}} & 1 & \texttt{a\textbar{}b\textbar{}b\textbar{}a} & \texttt{ab} & \texttt{ab} & \texttt{ab} & \texttt{0123} & \texttt{0} & \texttt{+} & \texttt{.} & \texttt{b@1}\\
S443 & F4.2968 & \texttt{-{}B-{}} & 1 & \texttt{a\textbar{}a\textbar{}a\textbar{}a} & \texttt{a} & \texttt{a} & \texttt{a} & \texttt{0123} & \texttt{0} & \texttt{+} & \texttt{S} & \texttt{none}\\
S444 & F4.2969 & \texttt{-{}B-{}} & 1 & \texttt{a\textbar{}a\textbar{}b\textbar{}b} & \texttt{ab} & \texttt{ab} & \texttt{ab} & \texttt{0123} & \texttt{0} & \texttt{+} & \texttt{.} & \texttt{a@0}\\
S445 & F4.2969 & \texttt{-{}B-{}} & 1 & \texttt{a\textbar{}b\textbar{}b\textbar{}a} & \texttt{ab} & \texttt{ab} & \texttt{ab} & \texttt{0123} & \texttt{0} & \texttt{+} & \texttt{.} & \texttt{b@1}\\
S446 & F4.2969 & \texttt{-{}B-{}} & 1 & \texttt{a\textbar{}a\textbar{}a\textbar{}a} & \texttt{a} & \texttt{a} & \texttt{a} & \texttt{0123} & \texttt{0} & \texttt{+} & \texttt{S} & \texttt{none}\\
S447 & F4.2978 & \texttt{TB-{}} & 1 & \texttt{a\textbar{}b\textbar{}b\textbar{}b} & \texttt{ab} & \texttt{ab} & \texttt{ab} & \texttt{0123} & \texttt{0} & \texttt{+} & \texttt{.} & \texttt{a@0}\\
S448 & F4.2978 & \texttt{TB-{}} & 1 & \texttt{a\textbar{}b\textbar{}a\textbar{}b} & \texttt{ab} & \texttt{ab} & \texttt{ab} & \texttt{0123} & \texttt{0} & \texttt{+} & \texttt{.} & \texttt{a@0}\\
S449 & F4.2978 & \texttt{TB-{}} & 1 & \texttt{a\textbar{}a\textbar{}a\textbar{}a} & \texttt{a} & \texttt{a} & \texttt{a} & \texttt{0123} & \texttt{0} & \texttt{+} & \texttt{S} & \texttt{none}\\
S450 & F4.2998 & \texttt{TB-{}} & 1 & \texttt{a\textbar{}a\textbar{}b\textbar{}b} & \texttt{ab} & \texttt{ab} & \texttt{ab} & \texttt{0123} & \texttt{0} & \texttt{+} & \texttt{.} & \texttt{a@0}\\
S451 & F4.2998 & \texttt{TB-{}} & 1 & \texttt{a\textbar{}b\textbar{}b\textbar{}a} & \texttt{ab} & \texttt{ab} & \texttt{ab} & \texttt{0123} & \texttt{0} & \texttt{+} & \texttt{.} & \texttt{b@1}\\
S452 & F4.2998 & \texttt{TB-{}} & 1 & \texttt{a\textbar{}a\textbar{}a\textbar{}a} & \texttt{a} & \texttt{a} & \texttt{a} & \texttt{0123} & \texttt{0} & \texttt{+} & \texttt{S} & \texttt{none}\\
S453 & F4.3002 & \texttt{-{}B-{}} & 1 & \texttt{a\textbar{}b\textbar{}a\textbar{}a} & \texttt{ab} & \texttt{ab} & \texttt{ab} & \texttt{0123} & \texttt{0} & \texttt{+} & \texttt{.} & \texttt{b@1}\\
S454 & F4.3002 & \texttt{-{}B-{}} & 1 & \texttt{a\textbar{}b\textbar{}b\textbar{}a} & \texttt{ab} & \texttt{ab} & \texttt{ab} & \texttt{0123} & \texttt{0} & \texttt{+} & \texttt{.} & \texttt{b@1}\\
S455 & F4.3002 & \texttt{-{}B-{}} & 1 & \texttt{a\textbar{}a\textbar{}a\textbar{}a} & \texttt{a} & \texttt{a} & \texttt{a} & \texttt{0123} & \texttt{0} & \texttt{+} & \texttt{S} & \texttt{none}\\
S456 & F4.3014 & \texttt{-{}B-{}} & 1 & \texttt{a\textbar{}b\textbar{}b\textbar{}b} & \texttt{ab} & \texttt{ab} & \texttt{ab} & \texttt{0123} & \texttt{0} & \texttt{+} & \texttt{.} & \texttt{a@0}\\
S457 & F4.3014 & \texttt{-{}B-{}} & 1 & \texttt{a\textbar{}b\textbar{}b\textbar{}a} & \texttt{ab} & \texttt{ab} & \texttt{ab} & \texttt{0123} & \texttt{0} & \texttt{+} & \texttt{.} & \texttt{b@1}\\
S458 & F4.3014 & \texttt{-{}B-{}} & 1 & \texttt{a\textbar{}a\textbar{}a\textbar{}a} & \texttt{a} & \texttt{a} & \texttt{a} & \texttt{0123} & \texttt{0} & \texttt{+} & \texttt{S} & \texttt{none}\\
S459 & F4.3017 & \texttt{-{}B-{}} & 1 & \texttt{a\textbar{}b\textbar{}b\textbar{}b} & \texttt{ab} & \texttt{ab} & \texttt{ab} & \texttt{0123} & \texttt{0} & \texttt{+} & \texttt{.} & \texttt{a@0}\\
S460 & F4.3017 & \texttt{-{}B-{}} & 1 & \texttt{a\textbar{}b\textbar{}a\textbar{}b} & \texttt{ab} & \texttt{ab} & \texttt{ab} & \texttt{0123} & \texttt{0} & \texttt{+} & \texttt{.} & \texttt{a@0}\\
S461 & F4.3017 & \texttt{-{}B-{}} & 1 & \texttt{a\textbar{}b\textbar{}b\textbar{}a} & \texttt{ab} & \texttt{ab} & \texttt{ab} & \texttt{0123} & \texttt{0} & \texttt{+} & \texttt{.} & \texttt{b@1}\\
S462 & F4.3017 & \texttt{-{}B-{}} & 1 & \texttt{a\textbar{}a\textbar{}a\textbar{}a} & \texttt{a} & \texttt{a} & \texttt{a} & \texttt{0123} & \texttt{0} & \texttt{+} & \texttt{S} & \texttt{none}\\
S463 & F4.3035 & \texttt{-{}B-{}} & 1 & \texttt{a\textbar{}b\textbar{}c\textbar{}c} & \texttt{abc} & \texttt{abc} & \texttt{abc} & \texttt{0123} & \texttt{0} & \texttt{+} & \texttt{.} & \texttt{a@0}\\
S464 & F4.3035 & \texttt{-{}B-{}} & 1 & \texttt{a\textbar{}b\textbar{}a\textbar{}a} & \texttt{ab} & \texttt{ab} & \texttt{ab} & \texttt{0123} & \texttt{0} & \texttt{+} & \texttt{.} & \texttt{b@1}\\
S465 & F4.3035 & \texttt{-{}B-{}} & 1 & \texttt{a\textbar{}a\textbar{}b\textbar{}b} & \texttt{ab} & \texttt{ab} & \texttt{ab} & \texttt{0123} & \texttt{0} & \texttt{+} & \texttt{.} & \texttt{a@0}\\
S466 & F4.3035 & \texttt{-{}B-{}} & 1 & \texttt{a\textbar{}b\textbar{}b\textbar{}a} & \texttt{ab} & \texttt{ab} & \texttt{ab} & \texttt{0123} & \texttt{0} & \texttt{+} & \texttt{.} & \texttt{b@1}\\
S467 & F4.3035 & \texttt{-{}B-{}} & 1 & \texttt{a\textbar{}a\textbar{}a\textbar{}a} & \texttt{a} & \texttt{a} & \texttt{a} & \texttt{0123} & \texttt{0} & \texttt{+} & \texttt{S} & \texttt{none}\\
S468 & F4.3038 & \texttt{TB-{}} & 1 & \texttt{a\textbar{}b\textbar{}c\textbar{}c} & \texttt{abc} & \texttt{abc} & \texttt{abc} & \texttt{0123} & \texttt{0} & \texttt{+} & \texttt{.} & \texttt{a@0}\\
S469 & F4.3038 & \texttt{TB-{}} & 1 & \texttt{a\textbar{}b\textbar{}a\textbar{}a} & \texttt{ab} & \texttt{ab} & \texttt{ab} & \texttt{0123} & \texttt{0} & \texttt{+} & \texttt{.} & \texttt{b@1}\\
S470 & F4.3038 & \texttt{TB-{}} & 1 & \texttt{a\textbar{}a\textbar{}b\textbar{}b} & \texttt{ab} & \texttt{ab} & \texttt{ab} & \texttt{0123} & \texttt{0} & \texttt{+} & \texttt{.} & \texttt{a@0}\\
S471 & F4.3038 & \texttt{TB-{}} & 1 & \texttt{a\textbar{}b\textbar{}b\textbar{}a} & \texttt{ab} & \texttt{ab} & \texttt{ab} & \texttt{0123} & \texttt{0} & \texttt{+} & \texttt{.} & \texttt{b@1}\\
S472 & F4.3038 & \texttt{TB-{}} & 1 & \texttt{a\textbar{}a\textbar{}a\textbar{}a} & \texttt{a} & \texttt{a} & \texttt{a} & \texttt{0123} & \texttt{0} & \texttt{+} & \texttt{S} & \texttt{none}\\
S473 & F4.3042 & \texttt{-{}B-{}} & 1 & \texttt{a\textbar{}b\textbar{}c\textbar{}a} & \texttt{abc} & \texttt{abc} & \texttt{abc} & \texttt{0123} & \texttt{0} & \texttt{+} & \texttt{.} & \texttt{b@1}\\
S474 & F4.3042 & \texttt{-{}B-{}} & 1 & \texttt{a\textbar{}a\textbar{}b\textbar{}a} & \texttt{ab} & \texttt{ab} & \texttt{ab} & \texttt{0123} & \texttt{0} & \texttt{+} & \texttt{.} & \texttt{b@2}\\
S475 & F4.3042 & \texttt{-{}B-{}} & 1 & \texttt{a\textbar{}b\textbar{}b\textbar{}a} & \texttt{ab} & \texttt{ab} & \texttt{ab} & \texttt{0123} & \texttt{0} & \texttt{+} & \texttt{.} & \texttt{b@1}\\
S476 & F4.3042 & \texttt{-{}B-{}} & 1 & \texttt{a\textbar{}a\textbar{}a\textbar{}a} & \texttt{a} & \texttt{a} & \texttt{a} & \texttt{0123} & \texttt{0} & \texttt{+} & \texttt{S} & \texttt{none}\\
S477 & F4.3044 & \texttt{TB4} & 1 & \texttt{a\textbar{}b\textbar{}c\textbar{}d} & \texttt{abcd} & \texttt{abcd} & \texttt{abcd} & \texttt{0123} & \texttt{0} & \texttt{+} & \texttt{.} & \texttt{a@0}\\
S478 & F4.3044 & \texttt{TB4} & 1 & \texttt{a\textbar{}a\textbar{}b\textbar{}c} & \texttt{abc} & \texttt{abc} & \texttt{abc} & \texttt{0123} & \texttt{0} & \texttt{+} & \texttt{.} & \texttt{b@2}\\
S479 & F4.3044 & \texttt{TB4} & 1 & \texttt{a\textbar{}a\textbar{}a\textbar{}b} & \texttt{ab} & \texttt{ab} & \texttt{ab} & \texttt{0123} & \texttt{0} & \texttt{+} & \texttt{.} & \texttt{b@3}\\
S480 & F4.3044 & \texttt{TB4} & 1 & \texttt{a\textbar{}b\textbar{}b\textbar{}a} & \texttt{ab} & \texttt{ab} & \texttt{ab} & \texttt{0123} & \texttt{0} & \texttt{+} & \texttt{.} & \texttt{b@1}\\
S481 & F4.3044 & \texttt{TB4} & 1 & \texttt{a\textbar{}a\textbar{}a\textbar{}a} & \texttt{a} & \texttt{a} & \texttt{a} & \texttt{0123} & \texttt{0} & \texttt{+} & \texttt{S} & \texttt{none}\\
\multicolumn{13}{@{}l}{\rule{0pt}{2.8ex}\footnotesize\bfseries $n=4$ worlds, $m=2$ individuals: 335 models}\\*
S482 & F4.133 & \texttt{-{}B-{}} & 2 & \texttt{ab\textbar{}ab\textbar{}ab\textbar{}ab} & \texttt{a} & \texttt{a} & \texttt{ab} & \texttt{0123} & \texttt{0} & \texttt{+} & \texttt{S} & \texttt{none}\\
S483 & F4.133 & \texttt{-{}B-{}} & 2 & \texttt{ab\textbar{}ba\textbar{}ba\textbar{}ab} & \texttt{a} & \texttt{a} & \texttt{ab} & \texttt{0123} & \texttt{0} & \texttt{+} & \texttt{S} & \texttt{none}\\
S484 & F4.133 & \texttt{-{}B-{}} & 2 & \texttt{aa\textbar{}aa\textbar{}aa\textbar{}aa} & \texttt{a} & \texttt{a} & \texttt{a} & \texttt{0123} & \texttt{01} & \texttt{+} & \texttt{S} & \texttt{none}\\
S485 & F4.137 & \texttt{-{}B-{}} & 2 & \texttt{ab\textbar{}ab\textbar{}ab\textbar{}ab} & \texttt{a} & \texttt{a} & \texttt{ab} & \texttt{0123} & \texttt{0} & \texttt{+} & \texttt{S} & \texttt{none}\\
S486 & F4.137 & \texttt{-{}B-{}} & 2 & \texttt{ab\textbar{}ab\textbar{}ba\textbar{}ab} & \texttt{a} & \texttt{a} & \texttt{ab} & \texttt{0123} & \texttt{0} & \texttt{+} & \texttt{S} & \texttt{none}\\
S487 & F4.137 & \texttt{-{}B-{}} & 2 & \texttt{ab\textbar{}ba\textbar{}ba\textbar{}ab} & \texttt{a} & \texttt{a} & \texttt{ab} & \texttt{0123} & \texttt{0} & \texttt{+} & \texttt{S} & \texttt{none}\\
S488 & F4.137 & \texttt{-{}B-{}} & 2 & \texttt{aa\textbar{}aa\textbar{}aa\textbar{}aa} & \texttt{a} & \texttt{a} & \texttt{a} & \texttt{0123} & \texttt{01} & \texttt{+} & \texttt{S} & \texttt{none}\\
S489 & F4.139 & \texttt{TB4} & 2 & \texttt{ab\textbar{}ab\textbar{}ab\textbar{}ab} & \texttt{a} & \texttt{a} & \texttt{ab} & \texttt{0123} & \texttt{0} & \texttt{+} & \texttt{S} & \texttt{none}\\
S490 & F4.139 & \texttt{TB4} & 2 & \texttt{ab\textbar{}ab\textbar{}ab\textbar{}ba} & \texttt{b} & \texttt{b} & \texttt{ab} & \texttt{0123} & \texttt{1} & \texttt{+} & \texttt{S} & \texttt{none}\\
S491 & F4.139 & \texttt{TB4} & 2 & \texttt{ab\textbar{}ba\textbar{}ba\textbar{}ab} & \texttt{a} & \texttt{a} & \texttt{ab} & \texttt{0123} & \texttt{0} & \texttt{+} & \texttt{S} & \texttt{none}\\
S492 & F4.139 & \texttt{TB4} & 2 & \texttt{aa\textbar{}aa\textbar{}aa\textbar{}aa} & \texttt{a} & \texttt{a} & \texttt{a} & \texttt{0123} & \texttt{01} & \texttt{+} & \texttt{S} & \texttt{none}\\
S493 & F4.276 & \texttt{-{}B-{}} & 2 & \texttt{ab\textbar{}ab\textbar{}ab\textbar{}ab} & \texttt{a} & \texttt{a} & \texttt{ab} & \texttt{0123} & \texttt{0} & \texttt{+} & \texttt{S} & \texttt{none}\\
S494 & F4.276 & \texttt{-{}B-{}} & 2 & \texttt{ab\textbar{}ba\textbar{}ab\textbar{}ab} & \texttt{a} & \texttt{a} & \texttt{ab} & \texttt{0123} & \texttt{0} & \texttt{+} & \texttt{S} & \texttt{none}\\
S495 & F4.276 & \texttt{-{}B-{}} & 2 & \texttt{aa\textbar{}aa\textbar{}aa\textbar{}aa} & \texttt{a} & \texttt{a} & \texttt{a} & \texttt{0123} & \texttt{01} & \texttt{+} & \texttt{S} & \texttt{none}\\
S496 & F4.306 & \texttt{-{}B-{}} & 2 & \texttt{ab\textbar{}ab\textbar{}ab\textbar{}ab} & \texttt{a} & \texttt{a} & \texttt{ab} & \texttt{0123} & \texttt{0} & \texttt{+} & \texttt{S} & \texttt{none}\\
S497 & F4.306 & \texttt{-{}B-{}} & 2 & \texttt{ab\textbar{}ba\textbar{}ba\textbar{}ab} & \texttt{a} & \texttt{a} & \texttt{ab} & \texttt{0123} & \texttt{0} & \texttt{+} & \texttt{S} & \texttt{none}\\
S498 & F4.306 & \texttt{-{}B-{}} & 2 & \texttt{aa\textbar{}aa\textbar{}aa\textbar{}aa} & \texttt{a} & \texttt{a} & \texttt{a} & \texttt{0123} & \texttt{01} & \texttt{+} & \texttt{S} & \texttt{none}\\
S499 & F4.321 & \texttt{-{}B-{}} & 2 & \texttt{ab\textbar{}ab\textbar{}ab\textbar{}ab} & \texttt{a} & \texttt{a} & \texttt{ab} & \texttt{0123} & \texttt{0} & \texttt{+} & \texttt{S} & \texttt{none}\\
S500 & F4.321 & \texttt{-{}B-{}} & 2 & \texttt{ab\textbar{}ab\textbar{}ba\textbar{}ab} & \texttt{a} & \texttt{a} & \texttt{ab} & \texttt{0123} & \texttt{0} & \texttt{+} & \texttt{S} & \texttt{none}\\
S501 & F4.321 & \texttt{-{}B-{}} & 2 & \texttt{ab\textbar{}ba\textbar{}ba\textbar{}ab} & \texttt{a} & \texttt{a} & \texttt{ab} & \texttt{0123} & \texttt{0} & \texttt{+} & \texttt{S} & \texttt{none}\\
S502 & F4.321 & \texttt{-{}B-{}} & 2 & \texttt{aa\textbar{}aa\textbar{}aa\textbar{}aa} & \texttt{a} & \texttt{a} & \texttt{a} & \texttt{0123} & \texttt{01} & \texttt{+} & \texttt{S} & \texttt{none}\\
S503 & F4.505 & \texttt{-{}B-{}} & 2 & \texttt{ab\textbar{}ab\textbar{}ab\textbar{}ab} & \texttt{a} & \texttt{a} & \texttt{ab} & \texttt{0123} & \texttt{0} & \texttt{+} & \texttt{S} & \texttt{none}\\
S504 & F4.505 & \texttt{-{}B-{}} & 2 & \texttt{aa\textbar{}aa\textbar{}aa\textbar{}aa} & \texttt{a} & \texttt{a} & \texttt{a} & \texttt{0123} & \texttt{01} & \texttt{+} & \texttt{S} & \texttt{none}\\
S505 & F4.507 & \texttt{-{}B-{}} & 2 & \texttt{ab\textbar{}ab\textbar{}ab\textbar{}ab} & \texttt{a} & \texttt{a} & \texttt{ab} & \texttt{0123} & \texttt{0} & \texttt{+} & \texttt{S} & \texttt{none}\\
S506 & F4.507 & \texttt{-{}B-{}} & 2 & \texttt{ab\textbar{}ba\textbar{}ab\textbar{}ab} & \texttt{a} & \texttt{a} & \texttt{ab} & \texttt{0123} & \texttt{0} & \texttt{+} & \texttt{S} & \texttt{none}\\
S507 & F4.507 & \texttt{-{}B-{}} & 2 & \texttt{aa\textbar{}aa\textbar{}aa\textbar{}aa} & \texttt{a} & \texttt{a} & \texttt{a} & \texttt{0123} & \texttt{01} & \texttt{+} & \texttt{S} & \texttt{none}\\
S508 & F4.559 & \texttt{-{}B-{}} & 2 & \texttt{ab\textbar{}ab\textbar{}ab\textbar{}ab} & \texttt{a} & \texttt{a} & \texttt{ab} & \texttt{0123} & \texttt{0} & \texttt{+} & \texttt{S} & \texttt{none}\\
S509 & F4.559 & \texttt{-{}B-{}} & 2 & \texttt{aa\textbar{}aa\textbar{}aa\textbar{}aa} & \texttt{a} & \texttt{a} & \texttt{a} & \texttt{0123} & \texttt{01} & \texttt{+} & \texttt{S} & \texttt{none}\\
S510 & F4.563 & \texttt{-{}B-{}} & 2 & \texttt{ab\textbar{}ab\textbar{}ab\textbar{}ab} & \texttt{a} & \texttt{a} & \texttt{ab} & \texttt{0123} & \texttt{0} & \texttt{+} & \texttt{S} & \texttt{none}\\
S511 & F4.563 & \texttt{-{}B-{}} & 2 & \texttt{ab\textbar{}ba\textbar{}ba\textbar{}ab} & \texttt{a} & \texttt{a} & \texttt{ab} & \texttt{0123} & \texttt{0} & \texttt{+} & \texttt{S} & \texttt{none}\\
S512 & F4.563 & \texttt{-{}B-{}} & 2 & \texttt{aa\textbar{}aa\textbar{}aa\textbar{}aa} & \texttt{a} & \texttt{a} & \texttt{a} & \texttt{0123} & \texttt{01} & \texttt{+} & \texttt{S} & \texttt{none}\\
S513 & F4.616 & \texttt{-{}B-{}} & 2 & \texttt{ab\textbar{}ab\textbar{}ab\textbar{}ab} & \texttt{a} & \texttt{a} & \texttt{ab} & \texttt{0123} & \texttt{0} & \texttt{+} & \texttt{S} & \texttt{none}\\
S514 & F4.616 & \texttt{-{}B-{}} & 2 & \texttt{ab\textbar{}ab\textbar{}ba\textbar{}ab} & \texttt{a} & \texttt{a} & \texttt{ab} & \texttt{0123} & \texttt{0} & \texttt{+} & \texttt{S} & \texttt{none}\\
S515 & F4.616 & \texttt{-{}B-{}} & 2 & \texttt{aa\textbar{}aa\textbar{}aa\textbar{}aa} & \texttt{a} & \texttt{a} & \texttt{a} & \texttt{0123} & \texttt{01} & \texttt{+} & \texttt{S} & \texttt{none}\\
S516 & F4.679 & \texttt{-{}B-{}} & 2 & \texttt{ab\textbar{}ab\textbar{}ab\textbar{}ab} & \texttt{a} & \texttt{a} & \texttt{ab} & \texttt{0123} & \texttt{0} & \texttt{+} & \texttt{S} & \texttt{none}\\
S517 & F4.679 & \texttt{-{}B-{}} & 2 & \texttt{ab\textbar{}ba\textbar{}ba\textbar{}ab} & \texttt{a} & \texttt{a} & \texttt{ab} & \texttt{0123} & \texttt{0} & \texttt{+} & \texttt{S} & \texttt{none}\\
S518 & F4.679 & \texttt{-{}B-{}} & 2 & \texttt{aa\textbar{}aa\textbar{}aa\textbar{}aa} & \texttt{a} & \texttt{a} & \texttt{a} & \texttt{0123} & \texttt{01} & \texttt{+} & \texttt{S} & \texttt{none}\\
S519 & F4.712 & \texttt{TB4} & 2 & \texttt{ab\textbar{}ab\textbar{}ab\textbar{}ab} & \texttt{a} & \texttt{a} & \texttt{ab} & \texttt{0123} & \texttt{0} & \texttt{+} & \texttt{S} & \texttt{none}\\
S520 & F4.712 & \texttt{TB4} & 2 & \texttt{ab\textbar{}ab\textbar{}ba\textbar{}ba} & \texttt{b} & \texttt{b} & \texttt{ab} & \texttt{0123} & \texttt{1} & \texttt{+} & \texttt{S} & \texttt{none}\\
S521 & F4.712 & \texttt{TB4} & 2 & \texttt{ab\textbar{}ab\textbar{}ab\textbar{}ba} & \texttt{b} & \texttt{b} & \texttt{ab} & \texttt{0123} & \texttt{1} & \texttt{+} & \texttt{S} & \texttt{none}\\
S522 & F4.712 & \texttt{TB4} & 2 & \texttt{aa\textbar{}aa\textbar{}aa\textbar{}aa} & \texttt{a} & \texttt{a} & \texttt{a} & \texttt{0123} & \texttt{01} & \texttt{+} & \texttt{S} & \texttt{none}\\
S523 & F4.857 & \texttt{-{}B-{}} & 2 & \texttt{ab\textbar{}ab\textbar{}ab\textbar{}ab} & \texttt{a} & \texttt{a} & \texttt{ab} & \texttt{0123} & \texttt{0} & \texttt{+} & \texttt{S} & \texttt{none}\\
S524 & F4.857 & \texttt{-{}B-{}} & 2 & \texttt{aa\textbar{}aa\textbar{}aa\textbar{}aa} & \texttt{a} & \texttt{a} & \texttt{a} & \texttt{0123} & \texttt{01} & \texttt{+} & \texttt{S} & \texttt{none}\\
S525 & F4.861 & \texttt{-{}B-{}} & 2 & \texttt{ab\textbar{}ab\textbar{}ab\textbar{}ab} & \texttt{a} & \texttt{a} & \texttt{ab} & \texttt{0123} & \texttt{0} & \texttt{+} & \texttt{S} & \texttt{none}\\
S526 & F4.861 & \texttt{-{}B-{}} & 2 & \texttt{aa\textbar{}aa\textbar{}aa\textbar{}aa} & \texttt{a} & \texttt{a} & \texttt{a} & \texttt{0123} & \texttt{01} & \texttt{+} & \texttt{S} & \texttt{none}\\
S527 & F4.898 & \texttt{-{}B-{}} & 2 & \texttt{ab\textbar{}ab\textbar{}ab\textbar{}ab} & \texttt{a} & \texttt{a} & \texttt{ab} & \texttt{0123} & \texttt{0} & \texttt{+} & \texttt{S} & \texttt{none}\\
S528 & F4.898 & \texttt{-{}B-{}} & 2 & \texttt{aa\textbar{}aa\textbar{}aa\textbar{}aa} & \texttt{a} & \texttt{a} & \texttt{a} & \texttt{0123} & \texttt{01} & \texttt{+} & \texttt{S} & \texttt{none}\\
S529 & F4.937 & \texttt{-{}B-{}} & 2 & \texttt{ab\textbar{}ab\textbar{}ab\textbar{}ab} & \texttt{a} & \texttt{a} & \texttt{ab} & \texttt{0123} & \texttt{0} & \texttt{+} & \texttt{S} & \texttt{none}\\
S530 & F4.937 & \texttt{-{}B-{}} & 2 & \texttt{aa\textbar{}aa\textbar{}aa\textbar{}aa} & \texttt{a} & \texttt{a} & \texttt{a} & \texttt{0123} & \texttt{01} & \texttt{+} & \texttt{S} & \texttt{none}\\
S531 & F4.978 & \texttt{-{}B-{}} & 2 & \texttt{ab\textbar{}ab\textbar{}ab\textbar{}ab} & \texttt{a} & \texttt{a} & \texttt{ab} & \texttt{0123} & \texttt{0} & \texttt{+} & \texttt{S} & \texttt{none}\\
S532 & F4.978 & \texttt{-{}B-{}} & 2 & \texttt{ab\textbar{}ab\textbar{}ba\textbar{}ab} & \texttt{a} & \texttt{a} & \texttt{ab} & \texttt{0123} & \texttt{0} & \texttt{+} & \texttt{S} & \texttt{none}\\
S533 & F4.978 & \texttt{-{}B-{}} & 2 & \texttt{aa\textbar{}aa\textbar{}aa\textbar{}aa} & \texttt{a} & \texttt{a} & \texttt{a} & \texttt{0123} & \texttt{01} & \texttt{+} & \texttt{S} & \texttt{none}\\
S534 & F4.1075 & \texttt{-{}B-{}} & 2 & \texttt{ab\textbar{}ab\textbar{}ab\textbar{}ab} & \texttt{a} & \texttt{a} & \texttt{ab} & \texttt{0123} & \texttt{0} & \texttt{+} & \texttt{S} & \texttt{none}\\
S535 & F4.1075 & \texttt{-{}B-{}} & 2 & \texttt{ab\textbar{}ba\textbar{}ba\textbar{}ab} & \texttt{a} & \texttt{a} & \texttt{ab} & \texttt{0123} & \texttt{0} & \texttt{+} & \texttt{S} & \texttt{none}\\
S536 & F4.1075 & \texttt{-{}B-{}} & 2 & \texttt{aa\textbar{}aa\textbar{}aa\textbar{}aa} & \texttt{a} & \texttt{a} & \texttt{a} & \texttt{0123} & \texttt{01} & \texttt{+} & \texttt{S} & \texttt{none}\\
S537 & F4.1195 & \texttt{-{}B-{}} & 2 & \texttt{ab\textbar{}ab\textbar{}ab\textbar{}ab} & \texttt{a} & \texttt{a} & \texttt{ab} & \texttt{0123} & \texttt{0} & \texttt{+} & \texttt{S} & \texttt{none}\\
S538 & F4.1195 & \texttt{-{}B-{}} & 2 & \texttt{ab\textbar{}ab\textbar{}ba\textbar{}ab} & \texttt{a} & \texttt{a} & \texttt{ab} & \texttt{0123} & \texttt{0} & \texttt{+} & \texttt{S} & \texttt{none}\\
S539 & F4.1195 & \texttt{-{}B-{}} & 2 & \texttt{aa\textbar{}aa\textbar{}aa\textbar{}aa} & \texttt{a} & \texttt{a} & \texttt{a} & \texttt{0123} & \texttt{01} & \texttt{+} & \texttt{S} & \texttt{none}\\
S540 & F4.1196 & \texttt{-{}B-{}} & 2 & \texttt{ab\textbar{}ab\textbar{}ab\textbar{}ab} & \texttt{a} & \texttt{a} & \texttt{ab} & \texttt{0123} & \texttt{0} & \texttt{+} & \texttt{S} & \texttt{none}\\
S541 & F4.1196 & \texttt{-{}B-{}} & 2 & \texttt{ab\textbar{}ab\textbar{}ba\textbar{}ab} & \texttt{a} & \texttt{a} & \texttt{ab} & \texttt{0123} & \texttt{0} & \texttt{+} & \texttt{S} & \texttt{none}\\
S542 & F4.1196 & \texttt{-{}B-{}} & 2 & \texttt{aa\textbar{}aa\textbar{}aa\textbar{}aa} & \texttt{a} & \texttt{a} & \texttt{a} & \texttt{0123} & \texttt{01} & \texttt{+} & \texttt{S} & \texttt{none}\\
S543 & F4.1321 & \texttt{-{}B-{}} & 2 & \texttt{ab\textbar{}ab\textbar{}ab\textbar{}ab} & \texttt{a} & \texttt{a} & \texttt{ab} & \texttt{0123} & \texttt{0} & \texttt{+} & \texttt{S} & \texttt{none}\\
S544 & F4.1321 & \texttt{-{}B-{}} & 2 & \texttt{aa\textbar{}aa\textbar{}aa\textbar{}aa} & \texttt{a} & \texttt{a} & \texttt{a} & \texttt{0123} & \texttt{01} & \texttt{+} & \texttt{S} & \texttt{none}\\
S545 & F4.1343 & \texttt{-{}B-{}} & 2 & \texttt{ab\textbar{}ab\textbar{}ab\textbar{}ab} & \texttt{a} & \texttt{a} & \texttt{ab} & \texttt{0123} & \texttt{0} & \texttt{+} & \texttt{S} & \texttt{none}\\
S546 & F4.1343 & \texttt{-{}B-{}} & 2 & \texttt{aa\textbar{}aa\textbar{}aa\textbar{}aa} & \texttt{a} & \texttt{a} & \texttt{a} & \texttt{0123} & \texttt{01} & \texttt{+} & \texttt{S} & \texttt{none}\\
S547 & F4.1362 & \texttt{-{}B-{}} & 2 & \texttt{ab\textbar{}ab\textbar{}ab\textbar{}ab} & \texttt{a} & \texttt{a} & \texttt{ab} & \texttt{0123} & \texttt{0} & \texttt{+} & \texttt{S} & \texttt{none}\\
S548 & F4.1362 & \texttt{-{}B-{}} & 2 & \texttt{aa\textbar{}aa\textbar{}aa\textbar{}aa} & \texttt{a} & \texttt{a} & \texttt{a} & \texttt{0123} & \texttt{01} & \texttt{+} & \texttt{S} & \texttt{none}\\
S549 & F4.1365 & \texttt{-{}B-{}} & 2 & \texttt{ab\textbar{}ab\textbar{}ab\textbar{}ab} & \texttt{a} & \texttt{a} & \texttt{ab} & \texttt{0123} & \texttt{0} & \texttt{+} & \texttt{S} & \texttt{none}\\
S550 & F4.1365 & \texttt{-{}B-{}} & 2 & \texttt{aa\textbar{}aa\textbar{}aa\textbar{}aa} & \texttt{a} & \texttt{a} & \texttt{a} & \texttt{0123} & \texttt{01} & \texttt{+} & \texttt{S} & \texttt{none}\\
S551 & F4.1437 & \texttt{-{}B-{}} & 2 & \texttt{ab\textbar{}ab\textbar{}ab\textbar{}ab} & \texttt{a} & \texttt{a} & \texttt{ab} & \texttt{0123} & \texttt{0} & \texttt{+} & \texttt{S} & \texttt{none}\\
S552 & F4.1437 & \texttt{-{}B-{}} & 2 & \texttt{aa\textbar{}aa\textbar{}aa\textbar{}aa} & \texttt{a} & \texttt{a} & \texttt{a} & \texttt{0123} & \texttt{01} & \texttt{+} & \texttt{S} & \texttt{none}\\
S553 & F4.1490 & \texttt{-{}B-{}} & 2 & \texttt{ab\textbar{}ab\textbar{}ab\textbar{}ab} & \texttt{a} & \texttt{a} & \texttt{ab} & \texttt{0123} & \texttt{0} & \texttt{+} & \texttt{S} & \texttt{none}\\
S554 & F4.1490 & \texttt{-{}B-{}} & 2 & \texttt{aa\textbar{}aa\textbar{}aa\textbar{}aa} & \texttt{a} & \texttt{a} & \texttt{a} & \texttt{0123} & \texttt{01} & \texttt{+} & \texttt{S} & \texttt{none}\\
S555 & F4.1675 & \texttt{-{}B-{}} & 2 & \texttt{ab\textbar{}cd\textbar{}cd\textbar{}ab} & \texttt{ac} & \texttt{ac} & \texttt{abcd} & \texttt{0123} & \texttt{0} & \texttt{+} & \texttt{.} & \texttt{c@1}\\
S556 & F4.1675 & \texttt{-{}B-{}} & 2 & \texttt{aa\textbar{}bb\textbar{}bb\textbar{}aa} & \texttt{ab} & \texttt{ab} & \texttt{ab} & \texttt{0123} & \texttt{01} & \texttt{+} & \texttt{.} & \texttt{b@1}\\
S557 & F4.1675 & \texttt{-{}B-{}} & 2 & \texttt{ab\textbar{}ab\textbar{}ab\textbar{}ab} & \texttt{a} & \texttt{a} & \texttt{ab} & \texttt{0123} & \texttt{0} & \texttt{+} & \texttt{S} & \texttt{none}\\
S558 & F4.1675 & \texttt{-{}B-{}} & 2 & \texttt{ab\textbar{}ba\textbar{}ba\textbar{}ab} & \texttt{ab} & \texttt{ab} & \texttt{ab} & \texttt{0123} & \texttt{01} & \texttt{+} & \texttt{.} & \texttt{none}\\
S559 & F4.1675 & \texttt{-{}B-{}} & 2 & \texttt{aa\textbar{}aa\textbar{}aa\textbar{}aa} & \texttt{a} & \texttt{a} & \texttt{a} & \texttt{0123} & \texttt{01} & \texttt{+} & \texttt{S} & \texttt{none}\\
S560 & F4.1692 & \texttt{-{}B-{}} & 2 & \texttt{ab\textbar{}ab\textbar{}ab\textbar{}ab} & \texttt{a} & \texttt{a} & \texttt{ab} & \texttt{0123} & \texttt{0} & \texttt{+} & \texttt{S} & \texttt{none}\\
S561 & F4.1692 & \texttt{-{}B-{}} & 2 & \texttt{ab\textbar{}ab\textbar{}ba\textbar{}ab} & \texttt{a} & \texttt{a} & \texttt{ab} & \texttt{0123} & \texttt{0} & \texttt{+} & \texttt{S} & \texttt{none}\\
S562 & F4.1692 & \texttt{-{}B-{}} & 2 & \texttt{aa\textbar{}aa\textbar{}aa\textbar{}aa} & \texttt{a} & \texttt{a} & \texttt{a} & \texttt{0123} & \texttt{01} & \texttt{+} & \texttt{S} & \texttt{none}\\
S563 & F4.1732 & \texttt{-{}B-{}} & 2 & \texttt{ab\textbar{}cd\textbar{}cd\textbar{}ab} & \texttt{ac} & \texttt{ac} & \texttt{abcd} & \texttt{0123} & \texttt{0} & \texttt{+} & \texttt{.} & \texttt{c@1}\\
S564 & F4.1732 & \texttt{-{}B-{}} & 2 & \texttt{aa\textbar{}bb\textbar{}bb\textbar{}aa} & \texttt{ab} & \texttt{ab} & \texttt{ab} & \texttt{0123} & \texttt{01} & \texttt{+} & \texttt{.} & \texttt{b@1}\\
S565 & F4.1732 & \texttt{-{}B-{}} & 2 & \texttt{ab\textbar{}ab\textbar{}ab\textbar{}ab} & \texttt{a} & \texttt{a} & \texttt{ab} & \texttt{0123} & \texttt{0} & \texttt{+} & \texttt{S} & \texttt{none}\\
S566 & F4.1732 & \texttt{-{}B-{}} & 2 & \texttt{ab\textbar{}ba\textbar{}ba\textbar{}ab} & \texttt{ab} & \texttt{ab} & \texttt{ab} & \texttt{0123} & \texttt{01} & \texttt{+} & \texttt{.} & \texttt{none}\\
S567 & F4.1732 & \texttt{-{}B-{}} & 2 & \texttt{aa\textbar{}aa\textbar{}aa\textbar{}aa} & \texttt{a} & \texttt{a} & \texttt{a} & \texttt{0123} & \texttt{01} & \texttt{+} & \texttt{S} & \texttt{none}\\
S568 & F4.1794 & \texttt{TB-{}} & 2 & \texttt{ab\textbar{}ab\textbar{}ab\textbar{}ab} & \texttt{a} & \texttt{a} & \texttt{ab} & \texttt{0123} & \texttt{0} & \texttt{+} & \texttt{S} & \texttt{none}\\
S569 & F4.1794 & \texttt{TB-{}} & 2 & \texttt{ab\textbar{}ab\textbar{}ab\textbar{}ba} & \texttt{b} & \texttt{b} & \texttt{ab} & \texttt{0123} & \texttt{1} & \texttt{+} & \texttt{S} & \texttt{none}\\
S570 & F4.1794 & \texttt{TB-{}} & 2 & \texttt{aa\textbar{}aa\textbar{}aa\textbar{}aa} & \texttt{a} & \texttt{a} & \texttt{a} & \texttt{0123} & \texttt{01} & \texttt{+} & \texttt{S} & \texttt{none}\\
S571 & F4.1813 & \texttt{TB4} & 2 & \texttt{ab\textbar{}cd\textbar{}ab\textbar{}cd} & \texttt{ac} & \texttt{ac} & \texttt{abcd} & \texttt{0123} & \texttt{0} & \texttt{+} & \texttt{.} & \texttt{a@0}\\
S572 & F4.1813 & \texttt{TB4} & 2 & \texttt{ab\textbar{}cd\textbar{}ba\textbar{}dc} & \texttt{ad} & \texttt{ad} & \texttt{abcd} & \texttt{0123} & \texttt{0} & \texttt{+} & \texttt{.} & \texttt{b@0}\\
S573 & F4.1813 & \texttt{TB4} & 2 & \texttt{aa\textbar{}bb\textbar{}aa\textbar{}bb} & \texttt{ab} & \texttt{ab} & \texttt{ab} & \texttt{0123} & \texttt{01} & \texttt{+} & \texttt{.} & \texttt{a@0}\\
S574 & F4.1813 & \texttt{TB4} & 2 & \texttt{ab\textbar{}ab\textbar{}ab\textbar{}ab} & \texttt{a} & \texttt{a} & \texttt{ab} & \texttt{0123} & \texttt{0} & \texttt{+} & \texttt{S} & \texttt{none}\\
S575 & F4.1813 & \texttt{TB4} & 2 & \texttt{ab\textbar{}ba\textbar{}ab\textbar{}ba} & \texttt{ab} & \texttt{ab} & \texttt{ab} & \texttt{0123} & \texttt{01} & \texttt{+} & \texttt{.} & \texttt{none}\\
S576 & F4.1813 & \texttt{TB4} & 2 & \texttt{ab\textbar{}ba\textbar{}ba\textbar{}ab} & \texttt{a} & \texttt{a} & \texttt{ab} & \texttt{0123} & \texttt{0} & \texttt{+} & \texttt{S} & \texttt{none}\\
S577 & F4.1813 & \texttt{TB4} & 2 & \texttt{aa\textbar{}aa\textbar{}aa\textbar{}aa} & \texttt{a} & \texttt{a} & \texttt{a} & \texttt{0123} & \texttt{01} & \texttt{+} & \texttt{S} & \texttt{none}\\
S578 & F4.1869 & \texttt{-{}B-{}} & 2 & \texttt{ab\textbar{}ab\textbar{}ab\textbar{}ab} & \texttt{a} & \texttt{a} & \texttt{ab} & \texttt{0123} & \texttt{0} & \texttt{+} & \texttt{S} & \texttt{none}\\
S579 & F4.1869 & \texttt{-{}B-{}} & 2 & \texttt{aa\textbar{}aa\textbar{}aa\textbar{}aa} & \texttt{a} & \texttt{a} & \texttt{a} & \texttt{0123} & \texttt{01} & \texttt{+} & \texttt{S} & \texttt{none}\\
S580 & F4.1874 & \texttt{-{}B-{}} & 2 & \texttt{ab\textbar{}ab\textbar{}ab\textbar{}ab} & \texttt{a} & \texttt{a} & \texttt{ab} & \texttt{0123} & \texttt{0} & \texttt{+} & \texttt{S} & \texttt{none}\\
S581 & F4.1874 & \texttt{-{}B-{}} & 2 & \texttt{aa\textbar{}aa\textbar{}aa\textbar{}aa} & \texttt{a} & \texttt{a} & \texttt{a} & \texttt{0123} & \texttt{01} & \texttt{+} & \texttt{S} & \texttt{none}\\
S582 & F4.1904 & \texttt{-{}B-{}} & 2 & \texttt{ab\textbar{}ab\textbar{}ab\textbar{}ab} & \texttt{a} & \texttt{a} & \texttt{ab} & \texttt{0123} & \texttt{0} & \texttt{+} & \texttt{S} & \texttt{none}\\
S583 & F4.1904 & \texttt{-{}B-{}} & 2 & \texttt{aa\textbar{}aa\textbar{}aa\textbar{}aa} & \texttt{a} & \texttt{a} & \texttt{a} & \texttt{0123} & \texttt{01} & \texttt{+} & \texttt{S} & \texttt{none}\\
S584 & F4.1933 & \texttt{-{}B-{}} & 2 & \texttt{ab\textbar{}ab\textbar{}ab\textbar{}ab} & \texttt{a} & \texttt{a} & \texttt{ab} & \texttt{0123} & \texttt{0} & \texttt{+} & \texttt{S} & \texttt{none}\\
S585 & F4.1933 & \texttt{-{}B-{}} & 2 & \texttt{aa\textbar{}aa\textbar{}aa\textbar{}aa} & \texttt{a} & \texttt{a} & \texttt{a} & \texttt{0123} & \texttt{01} & \texttt{+} & \texttt{S} & \texttt{none}\\
S586 & F4.2109 & \texttt{-{}B-{}} & 2 & \texttt{ab\textbar{}cd\textbar{}cd\textbar{}ab} & \texttt{ac} & \texttt{ac} & \texttt{abcd} & \texttt{0123} & \texttt{0} & \texttt{+} & \texttt{.} & \texttt{c@1}\\
S587 & F4.2109 & \texttt{-{}B-{}} & 2 & \texttt{aa\textbar{}bb\textbar{}bb\textbar{}aa} & \texttt{ab} & \texttt{ab} & \texttt{ab} & \texttt{0123} & \texttt{01} & \texttt{+} & \texttt{.} & \texttt{b@1}\\
S588 & F4.2109 & \texttt{-{}B-{}} & 2 & \texttt{ab\textbar{}ab\textbar{}ab\textbar{}ab} & \texttt{a} & \texttt{a} & \texttt{ab} & \texttt{0123} & \texttt{0} & \texttt{+} & \texttt{S} & \texttt{none}\\
S589 & F4.2109 & \texttt{-{}B-{}} & 2 & \texttt{ab\textbar{}ba\textbar{}ba\textbar{}ab} & \texttt{ab} & \texttt{ab} & \texttt{ab} & \texttt{0123} & \texttt{01} & \texttt{+} & \texttt{.} & \texttt{none}\\
S590 & F4.2109 & \texttt{-{}B-{}} & 2 & \texttt{aa\textbar{}aa\textbar{}aa\textbar{}aa} & \texttt{a} & \texttt{a} & \texttt{a} & \texttt{0123} & \texttt{01} & \texttt{+} & \texttt{S} & \texttt{none}\\
S591 & F4.2120 & \texttt{-{}B-{}} & 2 & \texttt{ab\textbar{}ab\textbar{}ab\textbar{}ab} & \texttt{a} & \texttt{a} & \texttt{ab} & \texttt{0123} & \texttt{0} & \texttt{+} & \texttt{S} & \texttt{none}\\
S592 & F4.2120 & \texttt{-{}B-{}} & 2 & \texttt{ab\textbar{}ab\textbar{}ba\textbar{}ab} & \texttt{a} & \texttt{a} & \texttt{ab} & \texttt{0123} & \texttt{0} & \texttt{+} & \texttt{S} & \texttt{none}\\
S593 & F4.2120 & \texttt{-{}B-{}} & 2 & \texttt{aa\textbar{}aa\textbar{}aa\textbar{}aa} & \texttt{a} & \texttt{a} & \texttt{a} & \texttt{0123} & \texttt{01} & \texttt{+} & \texttt{S} & \texttt{none}\\
S594 & F4.2148 & \texttt{-{}B-{}} & 2 & \texttt{ab\textbar{}ab\textbar{}ab\textbar{}ab} & \texttt{a} & \texttt{a} & \texttt{ab} & \texttt{0123} & \texttt{0} & \texttt{+} & \texttt{S} & \texttt{none}\\
S595 & F4.2148 & \texttt{-{}B-{}} & 2 & \texttt{aa\textbar{}aa\textbar{}aa\textbar{}aa} & \texttt{a} & \texttt{a} & \texttt{a} & \texttt{0123} & \texttt{01} & \texttt{+} & \texttt{S} & \texttt{none}\\
S596 & F4.2151 & \texttt{-{}B-{}} & 2 & \texttt{ab\textbar{}cd\textbar{}cd\textbar{}ab} & \texttt{ac} & \texttt{ac} & \texttt{abcd} & \texttt{0123} & \texttt{0} & \texttt{+} & \texttt{.} & \texttt{c@1}\\
S597 & F4.2151 & \texttt{-{}B-{}} & 2 & \texttt{aa\textbar{}bb\textbar{}bb\textbar{}aa} & \texttt{ab} & \texttt{ab} & \texttt{ab} & \texttt{0123} & \texttt{01} & \texttt{+} & \texttt{.} & \texttt{b@1}\\
S598 & F4.2151 & \texttt{-{}B-{}} & 2 & \texttt{ab\textbar{}ab\textbar{}ab\textbar{}ab} & \texttt{a} & \texttt{a} & \texttt{ab} & \texttt{0123} & \texttt{0} & \texttt{+} & \texttt{S} & \texttt{none}\\
S599 & F4.2151 & \texttt{-{}B-{}} & 2 & \texttt{ab\textbar{}ba\textbar{}ba\textbar{}ab} & \texttt{ab} & \texttt{ab} & \texttt{ab} & \texttt{0123} & \texttt{01} & \texttt{+} & \texttt{.} & \texttt{none}\\
S600 & F4.2151 & \texttt{-{}B-{}} & 2 & \texttt{aa\textbar{}aa\textbar{}aa\textbar{}aa} & \texttt{a} & \texttt{a} & \texttt{a} & \texttt{0123} & \texttt{01} & \texttt{+} & \texttt{S} & \texttt{none}\\
S601 & F4.2291 & \texttt{-{}B-{}} & 2 & \texttt{ab\textbar{}ab\textbar{}ab\textbar{}ab} & \texttt{a} & \texttt{a} & \texttt{ab} & \texttt{0123} & \texttt{0} & \texttt{+} & \texttt{S} & \texttt{none}\\
S602 & F4.2291 & \texttt{-{}B-{}} & 2 & \texttt{aa\textbar{}aa\textbar{}aa\textbar{}aa} & \texttt{a} & \texttt{a} & \texttt{a} & \texttt{0123} & \texttt{01} & \texttt{+} & \texttt{S} & \texttt{none}\\
S603 & F4.2346 & \texttt{-{}B-{}} & 2 & \texttt{ab\textbar{}ab\textbar{}ab\textbar{}ab} & \texttt{a} & \texttt{a} & \texttt{ab} & \texttt{0123} & \texttt{0} & \texttt{+} & \texttt{S} & \texttt{none}\\
S604 & F4.2346 & \texttt{-{}B-{}} & 2 & \texttt{aa\textbar{}aa\textbar{}aa\textbar{}aa} & \texttt{a} & \texttt{a} & \texttt{a} & \texttt{0123} & \texttt{01} & \texttt{+} & \texttt{S} & \texttt{none}\\
S605 & F4.2373 & \texttt{-{}B-{}} & 2 & \texttt{ab\textbar{}ab\textbar{}ab\textbar{}ab} & \texttt{a} & \texttt{a} & \texttt{ab} & \texttt{0123} & \texttt{0} & \texttt{+} & \texttt{S} & \texttt{none}\\
S606 & F4.2373 & \texttt{-{}B-{}} & 2 & \texttt{aa\textbar{}aa\textbar{}aa\textbar{}aa} & \texttt{a} & \texttt{a} & \texttt{a} & \texttt{0123} & \texttt{01} & \texttt{+} & \texttt{S} & \texttt{none}\\
S607 & F4.2467 & \texttt{-{}B-{}} & 2 & \texttt{ab\textbar{}cd\textbar{}cd\textbar{}ab} & \texttt{ac} & \texttt{ac} & \texttt{abcd} & \texttt{0123} & \texttt{0} & \texttt{+} & \texttt{.} & \texttt{c@1}\\
S608 & F4.2467 & \texttt{-{}B-{}} & 2 & \texttt{aa\textbar{}bb\textbar{}bb\textbar{}aa} & \texttt{ab} & \texttt{ab} & \texttt{ab} & \texttt{0123} & \texttt{01} & \texttt{+} & \texttt{.} & \texttt{b@1}\\
S609 & F4.2467 & \texttt{-{}B-{}} & 2 & \texttt{ab\textbar{}ab\textbar{}ab\textbar{}ab} & \texttt{a} & \texttt{a} & \texttt{ab} & \texttt{0123} & \texttt{0} & \texttt{+} & \texttt{S} & \texttt{none}\\
S610 & F4.2467 & \texttt{-{}B-{}} & 2 & \texttt{ab\textbar{}ba\textbar{}ba\textbar{}ab} & \texttt{ab} & \texttt{ab} & \texttt{ab} & \texttt{0123} & \texttt{01} & \texttt{+} & \texttt{.} & \texttt{none}\\
S611 & F4.2467 & \texttt{-{}B-{}} & 2 & \texttt{aa\textbar{}aa\textbar{}aa\textbar{}aa} & \texttt{a} & \texttt{a} & \texttt{a} & \texttt{0123} & \texttt{01} & \texttt{+} & \texttt{S} & \texttt{none}\\
S612 & F4.2468 & \texttt{-{}B-{}} & 2 & \texttt{ab\textbar{}cd\textbar{}cd\textbar{}ab} & \texttt{ac} & \texttt{ac} & \texttt{abcd} & \texttt{0123} & \texttt{0} & \texttt{+} & \texttt{.} & \texttt{c@1}\\
S613 & F4.2468 & \texttt{-{}B-{}} & 2 & \texttt{aa\textbar{}bb\textbar{}bb\textbar{}aa} & \texttt{ab} & \texttt{ab} & \texttt{ab} & \texttt{0123} & \texttt{01} & \texttt{+} & \texttt{.} & \texttt{b@1}\\
S614 & F4.2468 & \texttt{-{}B-{}} & 2 & \texttt{ab\textbar{}ab\textbar{}ab\textbar{}ab} & \texttt{a} & \texttt{a} & \texttt{ab} & \texttt{0123} & \texttt{0} & \texttt{+} & \texttt{S} & \texttt{none}\\
S615 & F4.2468 & \texttt{-{}B-{}} & 2 & \texttt{ab\textbar{}ba\textbar{}ba\textbar{}ab} & \texttt{ab} & \texttt{ab} & \texttt{ab} & \texttt{0123} & \texttt{01} & \texttt{+} & \texttt{.} & \texttt{none}\\
S616 & F4.2468 & \texttt{-{}B-{}} & 2 & \texttt{aa\textbar{}aa\textbar{}aa\textbar{}aa} & \texttt{a} & \texttt{a} & \texttt{a} & \texttt{0123} & \texttt{01} & \texttt{+} & \texttt{S} & \texttt{none}\\
S617 & F4.2480 & \texttt{-{}B-{}} & 2 & \texttt{ab\textbar{}cd\textbar{}cd\textbar{}cd} & \texttt{ac} & \texttt{ac} & \texttt{abcd} & \texttt{0123} & \texttt{0} & \texttt{+} & \texttt{.} & \texttt{a@0}\\
S618 & F4.2480 & \texttt{-{}B-{}} & 2 & \texttt{aa\textbar{}bb\textbar{}bb\textbar{}bb} & \texttt{ab} & \texttt{ab} & \texttt{ab} & \texttt{0123} & \texttt{01} & \texttt{+} & \texttt{.} & \texttt{a@0}\\
S619 & F4.2480 & \texttt{-{}B-{}} & 2 & \texttt{ab\textbar{}ab\textbar{}ab\textbar{}ab} & \texttt{a} & \texttt{a} & \texttt{ab} & \texttt{0123} & \texttt{0} & \texttt{+} & \texttt{S} & \texttt{none}\\
S620 & F4.2480 & \texttt{-{}B-{}} & 2 & \texttt{ab\textbar{}ba\textbar{}ba\textbar{}ba} & \texttt{ab} & \texttt{ab} & \texttt{ab} & \texttt{0123} & \texttt{01} & \texttt{+} & \texttt{.} & \texttt{none}\\
S621 & F4.2480 & \texttt{-{}B-{}} & 2 & \texttt{aa\textbar{}aa\textbar{}aa\textbar{}aa} & \texttt{a} & \texttt{a} & \texttt{a} & \texttt{0123} & \texttt{01} & \texttt{+} & \texttt{S} & \texttt{none}\\
S622 & F4.2485 & \texttt{-{}B-{}} & 2 & \texttt{ab\textbar{}cd\textbar{}cd\textbar{}ab} & \texttt{ac} & \texttt{ac} & \texttt{abcd} & \texttt{0123} & \texttt{0} & \texttt{+} & \texttt{.} & \texttt{c@1}\\
S623 & F4.2485 & \texttt{-{}B-{}} & 2 & \texttt{aa\textbar{}bb\textbar{}bb\textbar{}aa} & \texttt{ab} & \texttt{ab} & \texttt{ab} & \texttt{0123} & \texttt{01} & \texttt{+} & \texttt{.} & \texttt{b@1}\\
S624 & F4.2485 & \texttt{-{}B-{}} & 2 & \texttt{ab\textbar{}ab\textbar{}ab\textbar{}ab} & \texttt{a} & \texttt{a} & \texttt{ab} & \texttt{0123} & \texttt{0} & \texttt{+} & \texttt{S} & \texttt{none}\\
S625 & F4.2485 & \texttt{-{}B-{}} & 2 & \texttt{ab\textbar{}ba\textbar{}ba\textbar{}ab} & \texttt{ab} & \texttt{ab} & \texttt{ab} & \texttt{0123} & \texttt{01} & \texttt{+} & \texttt{.} & \texttt{none}\\
S626 & F4.2485 & \texttt{-{}B-{}} & 2 & \texttt{aa\textbar{}aa\textbar{}aa\textbar{}aa} & \texttt{a} & \texttt{a} & \texttt{a} & \texttt{0123} & \texttt{01} & \texttt{+} & \texttt{S} & \texttt{none}\\
S627 & F4.2524 & \texttt{-{}B-{}} & 2 & \texttt{ab\textbar{}cd\textbar{}ab\textbar{}cd} & \texttt{ac} & \texttt{ac} & \texttt{abcd} & \texttt{0123} & \texttt{0} & \texttt{+} & \texttt{.} & \texttt{a@0}\\
S628 & F4.2524 & \texttt{-{}B-{}} & 2 & \texttt{ab\textbar{}cd\textbar{}cd\textbar{}ab} & \texttt{ac} & \texttt{ac} & \texttt{abcd} & \texttt{0123} & \texttt{0} & \texttt{+} & \texttt{.} & \texttt{c@1}\\
S629 & F4.2524 & \texttt{-{}B-{}} & 2 & \texttt{aa\textbar{}bb\textbar{}aa\textbar{}bb} & \texttt{ab} & \texttt{ab} & \texttt{ab} & \texttt{0123} & \texttt{01} & \texttt{+} & \texttt{.} & \texttt{a@0}\\
S630 & F4.2524 & \texttt{-{}B-{}} & 2 & \texttt{aa\textbar{}bb\textbar{}bb\textbar{}aa} & \texttt{ab} & \texttt{ab} & \texttt{ab} & \texttt{0123} & \texttt{01} & \texttt{+} & \texttt{.} & \texttt{b@1}\\
S631 & F4.2524 & \texttt{-{}B-{}} & 2 & \texttt{ab\textbar{}ab\textbar{}ab\textbar{}ab} & \texttt{a} & \texttt{a} & \texttt{ab} & \texttt{0123} & \texttt{0} & \texttt{+} & \texttt{S} & \texttt{none}\\
S632 & F4.2524 & \texttt{-{}B-{}} & 2 & \texttt{ab\textbar{}ba\textbar{}ab\textbar{}ba} & \texttt{ab} & \texttt{ab} & \texttt{ab} & \texttt{0123} & \texttt{01} & \texttt{+} & \texttt{.} & \texttt{none}\\
S633 & F4.2524 & \texttt{-{}B-{}} & 2 & \texttt{ab\textbar{}ba\textbar{}ba\textbar{}ab} & \texttt{ab} & \texttt{ab} & \texttt{ab} & \texttt{0123} & \texttt{01} & \texttt{+} & \texttt{.} & \texttt{none}\\
S634 & F4.2524 & \texttt{-{}B-{}} & 2 & \texttt{aa\textbar{}aa\textbar{}aa\textbar{}aa} & \texttt{a} & \texttt{a} & \texttt{a} & \texttt{0123} & \texttt{01} & \texttt{+} & \texttt{S} & \texttt{none}\\
S635 & F4.2662 & \texttt{TB4} & 2 & \texttt{ab\textbar{}ab\textbar{}ab\textbar{}ab} & \texttt{a} & \texttt{a} & \texttt{ab} & \texttt{0123} & \texttt{0} & \texttt{+} & \texttt{S} & \texttt{none}\\
S636 & F4.2662 & \texttt{TB4} & 2 & \texttt{ab\textbar{}ab\textbar{}ab\textbar{}ba} & \texttt{b} & \texttt{b} & \texttt{ab} & \texttt{0123} & \texttt{1} & \texttt{+} & \texttt{S} & \texttt{none}\\
S637 & F4.2662 & \texttt{TB4} & 2 & \texttt{aa\textbar{}aa\textbar{}aa\textbar{}aa} & \texttt{a} & \texttt{a} & \texttt{a} & \texttt{0123} & \texttt{01} & \texttt{+} & \texttt{S} & \texttt{none}\\
S638 & F4.2672 & \texttt{TB-{}} & 2 & \texttt{ab\textbar{}cd\textbar{}cd\textbar{}cd} & \texttt{ac} & \texttt{ac} & \texttt{abcd} & \texttt{0123} & \texttt{0} & \texttt{+} & \texttt{.} & \texttt{a@0}\\
S639 & F4.2672 & \texttt{TB-{}} & 2 & \texttt{aa\textbar{}bb\textbar{}bb\textbar{}bb} & \texttt{ab} & \texttt{ab} & \texttt{ab} & \texttt{0123} & \texttt{01} & \texttt{+} & \texttt{.} & \texttt{a@0}\\
S640 & F4.2672 & \texttt{TB-{}} & 2 & \texttt{ab\textbar{}ab\textbar{}ab\textbar{}ab} & \texttt{a} & \texttt{a} & \texttt{ab} & \texttt{0123} & \texttt{0} & \texttt{+} & \texttt{S} & \texttt{none}\\
S641 & F4.2672 & \texttt{TB-{}} & 2 & \texttt{ab\textbar{}ba\textbar{}ba\textbar{}ba} & \texttt{ab} & \texttt{ab} & \texttt{ab} & \texttt{0123} & \texttt{01} & \texttt{+} & \texttt{.} & \texttt{none}\\
S642 & F4.2672 & \texttt{TB-{}} & 2 & \texttt{aa\textbar{}aa\textbar{}aa\textbar{}aa} & \texttt{a} & \texttt{a} & \texttt{a} & \texttt{0123} & \texttt{01} & \texttt{+} & \texttt{S} & \texttt{none}\\
S643 & F4.2683 & \texttt{TB-{}} & 2 & \texttt{ab\textbar{}ab\textbar{}cd\textbar{}cd} & \texttt{ac} & \texttt{ac} & \texttt{abcd} & \texttt{0123} & \texttt{0} & \texttt{+} & \texttt{.} & \texttt{a@0}\\
S644 & F4.2683 & \texttt{TB-{}} & 2 & \texttt{aa\textbar{}aa\textbar{}bb\textbar{}bb} & \texttt{ab} & \texttt{ab} & \texttt{ab} & \texttt{0123} & \texttt{01} & \texttt{+} & \texttt{.} & \texttt{a@0}\\
S645 & F4.2683 & \texttt{TB-{}} & 2 & \texttt{ab\textbar{}cd\textbar{}ab\textbar{}cd} & \texttt{ac} & \texttt{ac} & \texttt{abcd} & \texttt{0123} & \texttt{0} & \texttt{+} & \texttt{.} & \texttt{a@0}\\
S646 & F4.2683 & \texttt{TB-{}} & 2 & \texttt{aa\textbar{}bb\textbar{}aa\textbar{}bb} & \texttt{ab} & \texttt{ab} & \texttt{ab} & \texttt{0123} & \texttt{01} & \texttt{+} & \texttt{.} & \texttt{a@0}\\
S647 & F4.2683 & \texttt{TB-{}} & 2 & \texttt{ab\textbar{}ab\textbar{}ab\textbar{}ab} & \texttt{a} & \texttt{a} & \texttt{ab} & \texttt{0123} & \texttt{0} & \texttt{+} & \texttt{S} & \texttt{none}\\
S648 & F4.2683 & \texttt{TB-{}} & 2 & \texttt{ab\textbar{}ab\textbar{}ba\textbar{}ba} & \texttt{ab} & \texttt{ab} & \texttt{ab} & \texttt{0123} & \texttt{01} & \texttt{+} & \texttt{.} & \texttt{none}\\
S649 & F4.2683 & \texttt{TB-{}} & 2 & \texttt{ab\textbar{}ba\textbar{}ab\textbar{}ba} & \texttt{ab} & \texttt{ab} & \texttt{ab} & \texttt{0123} & \texttt{01} & \texttt{+} & \texttt{.} & \texttt{none}\\
S650 & F4.2683 & \texttt{TB-{}} & 2 & \texttt{aa\textbar{}aa\textbar{}aa\textbar{}aa} & \texttt{a} & \texttt{a} & \texttt{a} & \texttt{0123} & \texttt{01} & \texttt{+} & \texttt{S} & \texttt{none}\\
S651 & F4.2709 & \texttt{-{}B-{}} & 2 & \texttt{ab\textbar{}ab\textbar{}ab\textbar{}ab} & \texttt{a} & \texttt{a} & \texttt{ab} & \texttt{0123} & \texttt{0} & \texttt{+} & \texttt{S} & \texttt{none}\\
S652 & F4.2709 & \texttt{-{}B-{}} & 2 & \texttt{aa\textbar{}aa\textbar{}aa\textbar{}aa} & \texttt{a} & \texttt{a} & \texttt{a} & \texttt{0123} & \texttt{01} & \texttt{+} & \texttt{S} & \texttt{none}\\
S653 & F4.2743 & \texttt{-{}B-{}} & 2 & \texttt{ab\textbar{}cd\textbar{}cd\textbar{}cd} & \texttt{ac} & \texttt{ac} & \texttt{abcd} & \texttt{0123} & \texttt{0} & \texttt{+} & \texttt{.} & \texttt{a@0}\\
S654 & F4.2743 & \texttt{-{}B-{}} & 2 & \texttt{aa\textbar{}bb\textbar{}bb\textbar{}bb} & \texttt{ab} & \texttt{ab} & \texttt{ab} & \texttt{0123} & \texttt{01} & \texttt{+} & \texttt{.} & \texttt{a@0}\\
S655 & F4.2743 & \texttt{-{}B-{}} & 2 & \texttt{ab\textbar{}cd\textbar{}cd\textbar{}ab} & \texttt{ac} & \texttt{ac} & \texttt{abcd} & \texttt{0123} & \texttt{0} & \texttt{+} & \texttt{.} & \texttt{c@1}\\
S656 & F4.2743 & \texttt{-{}B-{}} & 2 & \texttt{aa\textbar{}bb\textbar{}bb\textbar{}aa} & \texttt{ab} & \texttt{ab} & \texttt{ab} & \texttt{0123} & \texttt{01} & \texttt{+} & \texttt{.} & \texttt{b@1}\\
S657 & F4.2743 & \texttt{-{}B-{}} & 2 & \texttt{ab\textbar{}ab\textbar{}ab\textbar{}ab} & \texttt{a} & \texttt{a} & \texttt{ab} & \texttt{0123} & \texttt{0} & \texttt{+} & \texttt{S} & \texttt{none}\\
S658 & F4.2743 & \texttt{-{}B-{}} & 2 & \texttt{ab\textbar{}ba\textbar{}ba\textbar{}ba} & \texttt{ab} & \texttt{ab} & \texttt{ab} & \texttt{0123} & \texttt{01} & \texttt{+} & \texttt{.} & \texttt{none}\\
S659 & F4.2743 & \texttt{-{}B-{}} & 2 & \texttt{ab\textbar{}ba\textbar{}ba\textbar{}ab} & \texttt{ab} & \texttt{ab} & \texttt{ab} & \texttt{0123} & \texttt{01} & \texttt{+} & \texttt{.} & \texttt{none}\\
S660 & F4.2743 & \texttt{-{}B-{}} & 2 & \texttt{aa\textbar{}aa\textbar{}aa\textbar{}aa} & \texttt{a} & \texttt{a} & \texttt{a} & \texttt{0123} & \texttt{01} & \texttt{+} & \texttt{S} & \texttt{none}\\
S661 & F4.2747 & \texttt{-{}B-{}} & 2 & \texttt{ab\textbar{}cd\textbar{}cd\textbar{}cd} & \texttt{ac} & \texttt{ac} & \texttt{abcd} & \texttt{0123} & \texttt{0} & \texttt{+} & \texttt{.} & \texttt{a@0}\\
S662 & F4.2747 & \texttt{-{}B-{}} & 2 & \texttt{aa\textbar{}bb\textbar{}bb\textbar{}bb} & \texttt{ab} & \texttt{ab} & \texttt{ab} & \texttt{0123} & \texttt{01} & \texttt{+} & \texttt{.} & \texttt{a@0}\\
S663 & F4.2747 & \texttt{-{}B-{}} & 2 & \texttt{ab\textbar{}cd\textbar{}cd\textbar{}ab} & \texttt{ac} & \texttt{ac} & \texttt{abcd} & \texttt{0123} & \texttt{0} & \texttt{+} & \texttt{.} & \texttt{c@1}\\
S664 & F4.2747 & \texttt{-{}B-{}} & 2 & \texttt{aa\textbar{}bb\textbar{}bb\textbar{}aa} & \texttt{ab} & \texttt{ab} & \texttt{ab} & \texttt{0123} & \texttt{01} & \texttt{+} & \texttt{.} & \texttt{b@1}\\
S665 & F4.2747 & \texttt{-{}B-{}} & 2 & \texttt{ab\textbar{}ab\textbar{}ab\textbar{}ab} & \texttt{a} & \texttt{a} & \texttt{ab} & \texttt{0123} & \texttt{0} & \texttt{+} & \texttt{S} & \texttt{none}\\
S666 & F4.2747 & \texttt{-{}B-{}} & 2 & \texttt{ab\textbar{}ba\textbar{}ba\textbar{}ba} & \texttt{ab} & \texttt{ab} & \texttt{ab} & \texttt{0123} & \texttt{01} & \texttt{+} & \texttt{.} & \texttt{none}\\
S667 & F4.2747 & \texttt{-{}B-{}} & 2 & \texttt{ab\textbar{}ba\textbar{}ba\textbar{}ab} & \texttt{ab} & \texttt{ab} & \texttt{ab} & \texttt{0123} & \texttt{01} & \texttt{+} & \texttt{.} & \texttt{none}\\
S668 & F4.2747 & \texttt{-{}B-{}} & 2 & \texttt{aa\textbar{}aa\textbar{}aa\textbar{}aa} & \texttt{a} & \texttt{a} & \texttt{a} & \texttt{0123} & \texttt{01} & \texttt{+} & \texttt{S} & \texttt{none}\\
S669 & F4.2768 & \texttt{-{}B-{}} & 2 & \texttt{ab\textbar{}cd\textbar{}cd\textbar{}ab} & \texttt{ac} & \texttt{ac} & \texttt{abcd} & \texttt{0123} & \texttt{0} & \texttt{+} & \texttt{.} & \texttt{c@1}\\
S670 & F4.2768 & \texttt{-{}B-{}} & 2 & \texttt{aa\textbar{}bb\textbar{}bb\textbar{}aa} & \texttt{ab} & \texttt{ab} & \texttt{ab} & \texttt{0123} & \texttt{01} & \texttt{+} & \texttt{.} & \texttt{b@1}\\
S671 & F4.2768 & \texttt{-{}B-{}} & 2 & \texttt{ab\textbar{}ab\textbar{}ab\textbar{}ab} & \texttt{a} & \texttt{a} & \texttt{ab} & \texttt{0123} & \texttt{0} & \texttt{+} & \texttt{S} & \texttt{none}\\
S672 & F4.2768 & \texttt{-{}B-{}} & 2 & \texttt{ab\textbar{}ba\textbar{}ba\textbar{}ab} & \texttt{ab} & \texttt{ab} & \texttt{ab} & \texttt{0123} & \texttt{01} & \texttt{+} & \texttt{.} & \texttt{none}\\
S673 & F4.2768 & \texttt{-{}B-{}} & 2 & \texttt{aa\textbar{}aa\textbar{}aa\textbar{}aa} & \texttt{a} & \texttt{a} & \texttt{a} & \texttt{0123} & \texttt{01} & \texttt{+} & \texttt{S} & \texttt{none}\\
S674 & F4.2769 & \texttt{-{}B-{}} & 2 & \texttt{ab\textbar{}cd\textbar{}cd\textbar{}ab} & \texttt{ac} & \texttt{ac} & \texttt{abcd} & \texttt{0123} & \texttt{0} & \texttt{+} & \texttt{.} & \texttt{c@1}\\
S675 & F4.2769 & \texttt{-{}B-{}} & 2 & \texttt{aa\textbar{}bb\textbar{}bb\textbar{}aa} & \texttt{ab} & \texttt{ab} & \texttt{ab} & \texttt{0123} & \texttt{01} & \texttt{+} & \texttt{.} & \texttt{b@1}\\
S676 & F4.2769 & \texttt{-{}B-{}} & 2 & \texttt{ab\textbar{}ab\textbar{}ab\textbar{}ab} & \texttt{a} & \texttt{a} & \texttt{ab} & \texttt{0123} & \texttt{0} & \texttt{+} & \texttt{S} & \texttt{none}\\
S677 & F4.2769 & \texttt{-{}B-{}} & 2 & \texttt{ab\textbar{}ba\textbar{}ba\textbar{}ab} & \texttt{ab} & \texttt{ab} & \texttt{ab} & \texttt{0123} & \texttt{01} & \texttt{+} & \texttt{.} & \texttt{none}\\
S678 & F4.2769 & \texttt{-{}B-{}} & 2 & \texttt{aa\textbar{}aa\textbar{}aa\textbar{}aa} & \texttt{a} & \texttt{a} & \texttt{a} & \texttt{0123} & \texttt{01} & \texttt{+} & \texttt{S} & \texttt{none}\\
S679 & F4.2865 & \texttt{-{}B-{}} & 2 & \texttt{ab\textbar{}cd\textbar{}cd\textbar{}ab} & \texttt{ac} & \texttt{ac} & \texttt{abcd} & \texttt{0123} & \texttt{0} & \texttt{+} & \texttt{.} & \texttt{c@1}\\
S680 & F4.2865 & \texttt{-{}B-{}} & 2 & \texttt{aa\textbar{}bb\textbar{}bb\textbar{}aa} & \texttt{ab} & \texttt{ab} & \texttt{ab} & \texttt{0123} & \texttt{01} & \texttt{+} & \texttt{.} & \texttt{b@1}\\
S681 & F4.2865 & \texttt{-{}B-{}} & 2 & \texttt{ab\textbar{}ab\textbar{}ab\textbar{}ab} & \texttt{a} & \texttt{a} & \texttt{ab} & \texttt{0123} & \texttt{0} & \texttt{+} & \texttt{S} & \texttt{none}\\
S682 & F4.2865 & \texttt{-{}B-{}} & 2 & \texttt{ab\textbar{}ba\textbar{}ba\textbar{}ab} & \texttt{ab} & \texttt{ab} & \texttt{ab} & \texttt{0123} & \texttt{01} & \texttt{+} & \texttt{.} & \texttt{none}\\
S683 & F4.2865 & \texttt{-{}B-{}} & 2 & \texttt{aa\textbar{}aa\textbar{}aa\textbar{}aa} & \texttt{a} & \texttt{a} & \texttt{a} & \texttt{0123} & \texttt{01} & \texttt{+} & \texttt{S} & \texttt{none}\\
S684 & F4.2914 & \texttt{-{}B-{}} & 2 & \texttt{ab\textbar{}cd\textbar{}ab\textbar{}ab} & \texttt{ac} & \texttt{ac} & \texttt{abcd} & \texttt{0123} & \texttt{0} & \texttt{+} & \texttt{.} & \texttt{c@1}\\
S685 & F4.2914 & \texttt{-{}B-{}} & 2 & \texttt{aa\textbar{}bb\textbar{}aa\textbar{}aa} & \texttt{ab} & \texttt{ab} & \texttt{ab} & \texttt{0123} & \texttt{01} & \texttt{+} & \texttt{.} & \texttt{b@1}\\
S686 & F4.2914 & \texttt{-{}B-{}} & 2 & \texttt{ab\textbar{}cd\textbar{}cd\textbar{}ab} & \texttt{ac} & \texttt{ac} & \texttt{abcd} & \texttt{0123} & \texttt{0} & \texttt{+} & \texttt{.} & \texttt{c@1}\\
S687 & F4.2914 & \texttt{-{}B-{}} & 2 & \texttt{aa\textbar{}bb\textbar{}bb\textbar{}aa} & \texttt{ab} & \texttt{ab} & \texttt{ab} & \texttt{0123} & \texttt{01} & \texttt{+} & \texttt{.} & \texttt{b@1}\\
S688 & F4.2914 & \texttt{-{}B-{}} & 2 & \texttt{ab\textbar{}ab\textbar{}ab\textbar{}ab} & \texttt{a} & \texttt{a} & \texttt{ab} & \texttt{0123} & \texttt{0} & \texttt{+} & \texttt{S} & \texttt{none}\\
S689 & F4.2914 & \texttt{-{}B-{}} & 2 & \texttt{ab\textbar{}ba\textbar{}ab\textbar{}ab} & \texttt{ab} & \texttt{ab} & \texttt{ab} & \texttt{0123} & \texttt{01} & \texttt{+} & \texttt{.} & \texttt{none}\\
S690 & F4.2914 & \texttt{-{}B-{}} & 2 & \texttt{ab\textbar{}ba\textbar{}ba\textbar{}ab} & \texttt{ab} & \texttt{ab} & \texttt{ab} & \texttt{0123} & \texttt{01} & \texttt{+} & \texttt{.} & \texttt{none}\\
S691 & F4.2914 & \texttt{-{}B-{}} & 2 & \texttt{aa\textbar{}aa\textbar{}aa\textbar{}aa} & \texttt{a} & \texttt{a} & \texttt{a} & \texttt{0123} & \texttt{01} & \texttt{+} & \texttt{S} & \texttt{none}\\
S692 & F4.2918 & \texttt{-{}B-{}} & 2 & \texttt{ab\textbar{}cd\textbar{}cd\textbar{}cd} & \texttt{ac} & \texttt{ac} & \texttt{abcd} & \texttt{0123} & \texttt{0} & \texttt{+} & \texttt{.} & \texttt{a@0}\\
S693 & F4.2918 & \texttt{-{}B-{}} & 2 & \texttt{aa\textbar{}bb\textbar{}bb\textbar{}bb} & \texttt{ab} & \texttt{ab} & \texttt{ab} & \texttt{0123} & \texttt{01} & \texttt{+} & \texttt{.} & \texttt{a@0}\\
S694 & F4.2918 & \texttt{-{}B-{}} & 2 & \texttt{ab\textbar{}cd\textbar{}ab\textbar{}cd} & \texttt{ac} & \texttt{ac} & \texttt{abcd} & \texttt{0123} & \texttt{0} & \texttt{+} & \texttt{.} & \texttt{a@0}\\
S695 & F4.2918 & \texttt{-{}B-{}} & 2 & \texttt{ab\textbar{}cd\textbar{}cd\textbar{}ab} & \texttt{ac} & \texttt{ac} & \texttt{abcd} & \texttt{0123} & \texttt{0} & \texttt{+} & \texttt{.} & \texttt{c@1}\\
S696 & F4.2918 & \texttt{-{}B-{}} & 2 & \texttt{aa\textbar{}bb\textbar{}aa\textbar{}bb} & \texttt{ab} & \texttt{ab} & \texttt{ab} & \texttt{0123} & \texttt{01} & \texttt{+} & \texttt{.} & \texttt{a@0}\\
S697 & F4.2918 & \texttt{-{}B-{}} & 2 & \texttt{aa\textbar{}bb\textbar{}bb\textbar{}aa} & \texttt{ab} & \texttt{ab} & \texttt{ab} & \texttt{0123} & \texttt{01} & \texttt{+} & \texttt{.} & \texttt{b@1}\\
S698 & F4.2918 & \texttt{-{}B-{}} & 2 & \texttt{ab\textbar{}ab\textbar{}ab\textbar{}ab} & \texttt{a} & \texttt{a} & \texttt{ab} & \texttt{0123} & \texttt{0} & \texttt{+} & \texttt{S} & \texttt{none}\\
S699 & F4.2918 & \texttt{-{}B-{}} & 2 & \texttt{ab\textbar{}ba\textbar{}ba\textbar{}ba} & \texttt{ab} & \texttt{ab} & \texttt{ab} & \texttt{0123} & \texttt{01} & \texttt{+} & \texttt{.} & \texttt{none}\\
S700 & F4.2918 & \texttt{-{}B-{}} & 2 & \texttt{ab\textbar{}ba\textbar{}ab\textbar{}ba} & \texttt{ab} & \texttt{ab} & \texttt{ab} & \texttt{0123} & \texttt{01} & \texttt{+} & \texttt{.} & \texttt{none}\\
S701 & F4.2918 & \texttt{-{}B-{}} & 2 & \texttt{ab\textbar{}ba\textbar{}ba\textbar{}ab} & \texttt{ab} & \texttt{ab} & \texttt{ab} & \texttt{0123} & \texttt{01} & \texttt{+} & \texttt{.} & \texttt{none}\\
S702 & F4.2918 & \texttt{-{}B-{}} & 2 & \texttt{aa\textbar{}aa\textbar{}aa\textbar{}aa} & \texttt{a} & \texttt{a} & \texttt{a} & \texttt{0123} & \texttt{01} & \texttt{+} & \texttt{S} & \texttt{none}\\
S703 & F4.2968 & \texttt{-{}B-{}} & 2 & \texttt{ab\textbar{}cd\textbar{}cd\textbar{}ab} & \texttt{ac} & \texttt{ac} & \texttt{abcd} & \texttt{0123} & \texttt{0} & \texttt{+} & \texttt{.} & \texttt{c@1}\\
S704 & F4.2968 & \texttt{-{}B-{}} & 2 & \texttt{aa\textbar{}bb\textbar{}bb\textbar{}aa} & \texttt{ab} & \texttt{ab} & \texttt{ab} & \texttt{0123} & \texttt{01} & \texttt{+} & \texttt{.} & \texttt{b@1}\\
S705 & F4.2968 & \texttt{-{}B-{}} & 2 & \texttt{ab\textbar{}ab\textbar{}ab\textbar{}ab} & \texttt{a} & \texttt{a} & \texttt{ab} & \texttt{0123} & \texttt{0} & \texttt{+} & \texttt{S} & \texttt{none}\\
S706 & F4.2968 & \texttt{-{}B-{}} & 2 & \texttt{ab\textbar{}ba\textbar{}ba\textbar{}ab} & \texttt{ab} & \texttt{ab} & \texttt{ab} & \texttt{0123} & \texttt{01} & \texttt{+} & \texttt{.} & \texttt{none}\\
S707 & F4.2968 & \texttt{-{}B-{}} & 2 & \texttt{aa\textbar{}aa\textbar{}aa\textbar{}aa} & \texttt{a} & \texttt{a} & \texttt{a} & \texttt{0123} & \texttt{01} & \texttt{+} & \texttt{S} & \texttt{none}\\
S708 & F4.2969 & \texttt{-{}B-{}} & 2 & \texttt{ab\textbar{}ab\textbar{}cd\textbar{}cd} & \texttt{ac} & \texttt{ac} & \texttt{abcd} & \texttt{0123} & \texttt{0} & \texttt{+} & \texttt{.} & \texttt{a@0}\\
S709 & F4.2969 & \texttt{-{}B-{}} & 2 & \texttt{aa\textbar{}aa\textbar{}bb\textbar{}bb} & \texttt{ab} & \texttt{ab} & \texttt{ab} & \texttt{0123} & \texttt{01} & \texttt{+} & \texttt{.} & \texttt{a@0}\\
S710 & F4.2969 & \texttt{-{}B-{}} & 2 & \texttt{ab\textbar{}cd\textbar{}cd\textbar{}ab} & \texttt{ac} & \texttt{ac} & \texttt{abcd} & \texttt{0123} & \texttt{0} & \texttt{+} & \texttt{.} & \texttt{c@1}\\
S711 & F4.2969 & \texttt{-{}B-{}} & 2 & \texttt{aa\textbar{}bb\textbar{}bb\textbar{}aa} & \texttt{ab} & \texttt{ab} & \texttt{ab} & \texttt{0123} & \texttt{01} & \texttt{+} & \texttt{.} & \texttt{b@1}\\
S712 & F4.2969 & \texttt{-{}B-{}} & 2 & \texttt{ab\textbar{}ab\textbar{}ab\textbar{}ab} & \texttt{a} & \texttt{a} & \texttt{ab} & \texttt{0123} & \texttt{0} & \texttt{+} & \texttt{S} & \texttt{none}\\
S713 & F4.2969 & \texttt{-{}B-{}} & 2 & \texttt{ab\textbar{}ab\textbar{}ba\textbar{}ba} & \texttt{ab} & \texttt{ab} & \texttt{ab} & \texttt{0123} & \texttt{01} & \texttt{+} & \texttt{.} & \texttt{none}\\
S714 & F4.2969 & \texttt{-{}B-{}} & 2 & \texttt{ab\textbar{}ba\textbar{}ba\textbar{}ab} & \texttt{ab} & \texttt{ab} & \texttt{ab} & \texttt{0123} & \texttt{01} & \texttt{+} & \texttt{.} & \texttt{none}\\
S715 & F4.2969 & \texttt{-{}B-{}} & 2 & \texttt{aa\textbar{}aa\textbar{}aa\textbar{}aa} & \texttt{a} & \texttt{a} & \texttt{a} & \texttt{0123} & \texttt{01} & \texttt{+} & \texttt{S} & \texttt{none}\\
S716 & F4.2978 & \texttt{TB-{}} & 2 & \texttt{ab\textbar{}cd\textbar{}cd\textbar{}cd} & \texttt{ac} & \texttt{ac} & \texttt{abcd} & \texttt{0123} & \texttt{0} & \texttt{+} & \texttt{.} & \texttt{a@0}\\
S717 & F4.2978 & \texttt{TB-{}} & 2 & \texttt{aa\textbar{}bb\textbar{}bb\textbar{}bb} & \texttt{ab} & \texttt{ab} & \texttt{ab} & \texttt{0123} & \texttt{01} & \texttt{+} & \texttt{.} & \texttt{a@0}\\
S718 & F4.2978 & \texttt{TB-{}} & 2 & \texttt{ab\textbar{}cd\textbar{}ab\textbar{}cd} & \texttt{ac} & \texttt{ac} & \texttt{abcd} & \texttt{0123} & \texttt{0} & \texttt{+} & \texttt{.} & \texttt{a@0}\\
S719 & F4.2978 & \texttt{TB-{}} & 2 & \texttt{aa\textbar{}bb\textbar{}aa\textbar{}bb} & \texttt{ab} & \texttt{ab} & \texttt{ab} & \texttt{0123} & \texttt{01} & \texttt{+} & \texttt{.} & \texttt{a@0}\\
S720 & F4.2978 & \texttt{TB-{}} & 2 & \texttt{ab\textbar{}ab\textbar{}ab\textbar{}ab} & \texttt{a} & \texttt{a} & \texttt{ab} & \texttt{0123} & \texttt{0} & \texttt{+} & \texttt{S} & \texttt{none}\\
S721 & F4.2978 & \texttt{TB-{}} & 2 & \texttt{ab\textbar{}ba\textbar{}ba\textbar{}ba} & \texttt{ab} & \texttt{ab} & \texttt{ab} & \texttt{0123} & \texttt{01} & \texttt{+} & \texttt{.} & \texttt{none}\\
S722 & F4.2978 & \texttt{TB-{}} & 2 & \texttt{ab\textbar{}ba\textbar{}ab\textbar{}ba} & \texttt{ab} & \texttt{ab} & \texttt{ab} & \texttt{0123} & \texttt{01} & \texttt{+} & \texttt{.} & \texttt{none}\\
S723 & F4.2978 & \texttt{TB-{}} & 2 & \texttt{aa\textbar{}aa\textbar{}aa\textbar{}aa} & \texttt{a} & \texttt{a} & \texttt{a} & \texttt{0123} & \texttt{01} & \texttt{+} & \texttt{S} & \texttt{none}\\
S724 & F4.2998 & \texttt{TB-{}} & 2 & \texttt{ab\textbar{}ab\textbar{}cd\textbar{}cd} & \texttt{ac} & \texttt{ac} & \texttt{abcd} & \texttt{0123} & \texttt{0} & \texttt{+} & \texttt{.} & \texttt{a@0}\\
S725 & F4.2998 & \texttt{TB-{}} & 2 & \texttt{aa\textbar{}aa\textbar{}bb\textbar{}bb} & \texttt{ab} & \texttt{ab} & \texttt{ab} & \texttt{0123} & \texttt{01} & \texttt{+} & \texttt{.} & \texttt{a@0}\\
S726 & F4.2998 & \texttt{TB-{}} & 2 & \texttt{ab\textbar{}cd\textbar{}cd\textbar{}ab} & \texttt{ac} & \texttt{ac} & \texttt{abcd} & \texttt{0123} & \texttt{0} & \texttt{+} & \texttt{.} & \texttt{c@1}\\
S727 & F4.2998 & \texttt{TB-{}} & 2 & \texttt{aa\textbar{}bb\textbar{}bb\textbar{}aa} & \texttt{ab} & \texttt{ab} & \texttt{ab} & \texttt{0123} & \texttt{01} & \texttt{+} & \texttt{.} & \texttt{b@1}\\
S728 & F4.2998 & \texttt{TB-{}} & 2 & \texttt{ab\textbar{}ab\textbar{}ab\textbar{}ab} & \texttt{a} & \texttt{a} & \texttt{ab} & \texttt{0123} & \texttt{0} & \texttt{+} & \texttt{S} & \texttt{none}\\
S729 & F4.2998 & \texttt{TB-{}} & 2 & \texttt{ab\textbar{}ab\textbar{}ba\textbar{}ba} & \texttt{ab} & \texttt{ab} & \texttt{ab} & \texttt{0123} & \texttt{01} & \texttt{+} & \texttt{.} & \texttt{none}\\
S730 & F4.2998 & \texttt{TB-{}} & 2 & \texttt{ab\textbar{}ba\textbar{}ba\textbar{}ab} & \texttt{ab} & \texttt{ab} & \texttt{ab} & \texttt{0123} & \texttt{01} & \texttt{+} & \texttt{.} & \texttt{none}\\
S731 & F4.2998 & \texttt{TB-{}} & 2 & \texttt{aa\textbar{}aa\textbar{}aa\textbar{}aa} & \texttt{a} & \texttt{a} & \texttt{a} & \texttt{0123} & \texttt{01} & \texttt{+} & \texttt{S} & \texttt{none}\\
S732 & F4.3002 & \texttt{-{}B-{}} & 2 & \texttt{ab\textbar{}cd\textbar{}ab\textbar{}ab} & \texttt{ac} & \texttt{ac} & \texttt{abcd} & \texttt{0123} & \texttt{0} & \texttt{+} & \texttt{.} & \texttt{c@1}\\
S733 & F4.3002 & \texttt{-{}B-{}} & 2 & \texttt{aa\textbar{}bb\textbar{}aa\textbar{}aa} & \texttt{ab} & \texttt{ab} & \texttt{ab} & \texttt{0123} & \texttt{01} & \texttt{+} & \texttt{.} & \texttt{b@1}\\
S734 & F4.3002 & \texttt{-{}B-{}} & 2 & \texttt{ab\textbar{}cd\textbar{}cd\textbar{}ab} & \texttt{ac} & \texttt{ac} & \texttt{abcd} & \texttt{0123} & \texttt{0} & \texttt{+} & \texttt{.} & \texttt{c@1}\\
S735 & F4.3002 & \texttt{-{}B-{}} & 2 & \texttt{aa\textbar{}bb\textbar{}bb\textbar{}aa} & \texttt{ab} & \texttt{ab} & \texttt{ab} & \texttt{0123} & \texttt{01} & \texttt{+} & \texttt{.} & \texttt{b@1}\\
S736 & F4.3002 & \texttt{-{}B-{}} & 2 & \texttt{ab\textbar{}ab\textbar{}ab\textbar{}ab} & \texttt{a} & \texttt{a} & \texttt{ab} & \texttt{0123} & \texttt{0} & \texttt{+} & \texttt{S} & \texttt{none}\\
S737 & F4.3002 & \texttt{-{}B-{}} & 2 & \texttt{ab\textbar{}ba\textbar{}ab\textbar{}ab} & \texttt{ab} & \texttt{ab} & \texttt{ab} & \texttt{0123} & \texttt{01} & \texttt{+} & \texttt{.} & \texttt{none}\\
S738 & F4.3002 & \texttt{-{}B-{}} & 2 & \texttt{ab\textbar{}ba\textbar{}ba\textbar{}ab} & \texttt{ab} & \texttt{ab} & \texttt{ab} & \texttt{0123} & \texttt{01} & \texttt{+} & \texttt{.} & \texttt{none}\\
S739 & F4.3002 & \texttt{-{}B-{}} & 2 & \texttt{aa\textbar{}aa\textbar{}aa\textbar{}aa} & \texttt{a} & \texttt{a} & \texttt{a} & \texttt{0123} & \texttt{01} & \texttt{+} & \texttt{S} & \texttt{none}\\
S740 & F4.3014 & \texttt{-{}B-{}} & 2 & \texttt{ab\textbar{}cd\textbar{}cd\textbar{}cd} & \texttt{ac} & \texttt{ac} & \texttt{abcd} & \texttt{0123} & \texttt{0} & \texttt{+} & \texttt{.} & \texttt{a@0}\\
S741 & F4.3014 & \texttt{-{}B-{}} & 2 & \texttt{aa\textbar{}bb\textbar{}bb\textbar{}bb} & \texttt{ab} & \texttt{ab} & \texttt{ab} & \texttt{0123} & \texttt{01} & \texttt{+} & \texttt{.} & \texttt{a@0}\\
S742 & F4.3014 & \texttt{-{}B-{}} & 2 & \texttt{ab\textbar{}cd\textbar{}cd\textbar{}ab} & \texttt{ac} & \texttt{ac} & \texttt{abcd} & \texttt{0123} & \texttt{0} & \texttt{+} & \texttt{.} & \texttt{c@1}\\
S743 & F4.3014 & \texttt{-{}B-{}} & 2 & \texttt{aa\textbar{}bb\textbar{}bb\textbar{}aa} & \texttt{ab} & \texttt{ab} & \texttt{ab} & \texttt{0123} & \texttt{01} & \texttt{+} & \texttt{.} & \texttt{b@1}\\
S744 & F4.3014 & \texttt{-{}B-{}} & 2 & \texttt{ab\textbar{}ab\textbar{}ab\textbar{}ab} & \texttt{a} & \texttt{a} & \texttt{ab} & \texttt{0123} & \texttt{0} & \texttt{+} & \texttt{S} & \texttt{none}\\
S745 & F4.3014 & \texttt{-{}B-{}} & 2 & \texttt{ab\textbar{}ba\textbar{}ba\textbar{}ba} & \texttt{ab} & \texttt{ab} & \texttt{ab} & \texttt{0123} & \texttt{01} & \texttt{+} & \texttt{.} & \texttt{none}\\
S746 & F4.3014 & \texttt{-{}B-{}} & 2 & \texttt{ab\textbar{}ba\textbar{}ba\textbar{}ab} & \texttt{ab} & \texttt{ab} & \texttt{ab} & \texttt{0123} & \texttt{01} & \texttt{+} & \texttt{.} & \texttt{none}\\
S747 & F4.3014 & \texttt{-{}B-{}} & 2 & \texttt{aa\textbar{}aa\textbar{}aa\textbar{}aa} & \texttt{a} & \texttt{a} & \texttt{a} & \texttt{0123} & \texttt{01} & \texttt{+} & \texttt{S} & \texttt{none}\\
S748 & F4.3017 & \texttt{-{}B-{}} & 2 & \texttt{ab\textbar{}cd\textbar{}cd\textbar{}cd} & \texttt{ac} & \texttt{ac} & \texttt{abcd} & \texttt{0123} & \texttt{0} & \texttt{+} & \texttt{.} & \texttt{a@0}\\
S749 & F4.3017 & \texttt{-{}B-{}} & 2 & \texttt{aa\textbar{}bb\textbar{}bb\textbar{}bb} & \texttt{ab} & \texttt{ab} & \texttt{ab} & \texttt{0123} & \texttt{01} & \texttt{+} & \texttt{.} & \texttt{a@0}\\
S750 & F4.3017 & \texttt{-{}B-{}} & 2 & \texttt{ab\textbar{}cd\textbar{}ab\textbar{}cd} & \texttt{ac} & \texttt{ac} & \texttt{abcd} & \texttt{0123} & \texttt{0} & \texttt{+} & \texttt{.} & \texttt{a@0}\\
S751 & F4.3017 & \texttt{-{}B-{}} & 2 & \texttt{aa\textbar{}bb\textbar{}aa\textbar{}bb} & \texttt{ab} & \texttt{ab} & \texttt{ab} & \texttt{0123} & \texttt{01} & \texttt{+} & \texttt{.} & \texttt{a@0}\\
S752 & F4.3017 & \texttt{-{}B-{}} & 2 & \texttt{ab\textbar{}cd\textbar{}cd\textbar{}ab} & \texttt{ac} & \texttt{ac} & \texttt{abcd} & \texttt{0123} & \texttt{0} & \texttt{+} & \texttt{.} & \texttt{c@1}\\
S753 & F4.3017 & \texttt{-{}B-{}} & 2 & \texttt{aa\textbar{}bb\textbar{}bb\textbar{}aa} & \texttt{ab} & \texttt{ab} & \texttt{ab} & \texttt{0123} & \texttt{01} & \texttt{+} & \texttt{.} & \texttt{b@1}\\
S754 & F4.3017 & \texttt{-{}B-{}} & 2 & \texttt{ab\textbar{}ab\textbar{}ab\textbar{}ab} & \texttt{a} & \texttt{a} & \texttt{ab} & \texttt{0123} & \texttt{0} & \texttt{+} & \texttt{S} & \texttt{none}\\
S755 & F4.3017 & \texttt{-{}B-{}} & 2 & \texttt{ab\textbar{}ba\textbar{}ba\textbar{}ba} & \texttt{ab} & \texttt{ab} & \texttt{ab} & \texttt{0123} & \texttt{01} & \texttt{+} & \texttt{.} & \texttt{none}\\
S756 & F4.3017 & \texttt{-{}B-{}} & 2 & \texttt{ab\textbar{}ba\textbar{}ab\textbar{}ba} & \texttt{ab} & \texttt{ab} & \texttt{ab} & \texttt{0123} & \texttt{01} & \texttt{+} & \texttt{.} & \texttt{none}\\
S757 & F4.3017 & \texttt{-{}B-{}} & 2 & \texttt{ab\textbar{}ba\textbar{}ba\textbar{}ab} & \texttt{ab} & \texttt{ab} & \texttt{ab} & \texttt{0123} & \texttt{01} & \texttt{+} & \texttt{.} & \texttt{none}\\
S758 & F4.3017 & \texttt{-{}B-{}} & 2 & \texttt{aa\textbar{}aa\textbar{}aa\textbar{}aa} & \texttt{a} & \texttt{a} & \texttt{a} & \texttt{0123} & \texttt{01} & \texttt{+} & \texttt{S} & \texttt{none}\\
S759 & F4.3035 & \texttt{-{}B-{}} & 2 & \texttt{ab\textbar{}cd\textbar{}ef\textbar{}ef} & \texttt{ace} & \texttt{ace} & \texttt{abcdef} & \texttt{0123} & \texttt{0} & \texttt{+} & \texttt{.} & \texttt{a@0}\\
S760 & F4.3035 & \texttt{-{}B-{}} & 2 & \texttt{ab\textbar{}cd\textbar{}ab\textbar{}ab} & \texttt{ac} & \texttt{ac} & \texttt{abcd} & \texttt{0123} & \texttt{0} & \texttt{+} & \texttt{.} & \texttt{c@1}\\
S761 & F4.3035 & \texttt{-{}B-{}} & 2 & \texttt{aa\textbar{}bb\textbar{}cc\textbar{}cc} & \texttt{abc} & \texttt{abc} & \texttt{abc} & \texttt{0123} & \texttt{01} & \texttt{+} & \texttt{.} & \texttt{a@0}\\
S762 & F4.3035 & \texttt{-{}B-{}} & 2 & \texttt{ab\textbar{}cc\textbar{}ba\textbar{}ba} & \texttt{abc} & \texttt{abc} & \texttt{abc} & \texttt{0123} & \texttt{01} & \texttt{+} & \texttt{.} & \texttt{c@1}\\
S763 & F4.3035 & \texttt{-{}B-{}} & 2 & \texttt{aa\textbar{}bb\textbar{}aa\textbar{}aa} & \texttt{ab} & \texttt{ab} & \texttt{ab} & \texttt{0123} & \texttt{01} & \texttt{+} & \texttt{.} & \texttt{b@1}\\
S764 & F4.3035 & \texttt{-{}B-{}} & 2 & \texttt{ab\textbar{}ab\textbar{}cd\textbar{}cd} & \texttt{ac} & \texttt{ac} & \texttt{abcd} & \texttt{0123} & \texttt{0} & \texttt{+} & \texttt{.} & \texttt{a@0}\\
S765 & F4.3035 & \texttt{-{}B-{}} & 2 & \texttt{ab\textbar{}ba\textbar{}cc\textbar{}cc} & \texttt{abc} & \texttt{abc} & \texttt{abc} & \texttt{0123} & \texttt{01} & \texttt{+} & \texttt{.} & \texttt{b@0}\\
S766 & F4.3035 & \texttt{-{}B-{}} & 2 & \texttt{aa\textbar{}aa\textbar{}bb\textbar{}bb} & \texttt{ab} & \texttt{ab} & \texttt{ab} & \texttt{0123} & \texttt{01} & \texttt{+} & \texttt{.} & \texttt{a@0}\\
S767 & F4.3035 & \texttt{-{}B-{}} & 2 & \texttt{ab\textbar{}cd\textbar{}cd\textbar{}ab} & \texttt{ac} & \texttt{ac} & \texttt{abcd} & \texttt{0123} & \texttt{0} & \texttt{+} & \texttt{.} & \texttt{c@1}\\
S768 & F4.3035 & \texttt{-{}B-{}} & 2 & \texttt{aa\textbar{}bb\textbar{}bb\textbar{}aa} & \texttt{ab} & \texttt{ab} & \texttt{ab} & \texttt{0123} & \texttt{01} & \texttt{+} & \texttt{.} & \texttt{b@1}\\
S769 & F4.3035 & \texttt{-{}B-{}} & 2 & \texttt{ab\textbar{}ab\textbar{}ab\textbar{}ab} & \texttt{a} & \texttt{a} & \texttt{ab} & \texttt{0123} & \texttt{0} & \texttt{+} & \texttt{S} & \texttt{none}\\
S770 & F4.3035 & \texttt{-{}B-{}} & 2 & \texttt{ab\textbar{}ba\textbar{}ab\textbar{}ab} & \texttt{ab} & \texttt{ab} & \texttt{ab} & \texttt{0123} & \texttt{01} & \texttt{+} & \texttt{.} & \texttt{none}\\
S771 & F4.3035 & \texttt{-{}B-{}} & 2 & \texttt{ab\textbar{}ab\textbar{}ba\textbar{}ba} & \texttt{ab} & \texttt{ab} & \texttt{ab} & \texttt{0123} & \texttt{01} & \texttt{+} & \texttt{.} & \texttt{none}\\
S772 & F4.3035 & \texttt{-{}B-{}} & 2 & \texttt{ab\textbar{}ba\textbar{}ba\textbar{}ab} & \texttt{ab} & \texttt{ab} & \texttt{ab} & \texttt{0123} & \texttt{01} & \texttt{+} & \texttt{.} & \texttt{none}\\
S773 & F4.3035 & \texttt{-{}B-{}} & 2 & \texttt{aa\textbar{}aa\textbar{}aa\textbar{}aa} & \texttt{a} & \texttt{a} & \texttt{a} & \texttt{0123} & \texttt{01} & \texttt{+} & \texttt{S} & \texttt{none}\\
S774 & F4.3038 & \texttt{TB-{}} & 2 & \texttt{ab\textbar{}cd\textbar{}ef\textbar{}ef} & \texttt{ace} & \texttt{ace} & \texttt{abcdef} & \texttt{0123} & \texttt{0} & \texttt{+} & \texttt{.} & \texttt{a@0}\\
S775 & F4.3038 & \texttt{TB-{}} & 2 & \texttt{ab\textbar{}cd\textbar{}ab\textbar{}ab} & \texttt{ac} & \texttt{ac} & \texttt{abcd} & \texttt{0123} & \texttt{0} & \texttt{+} & \texttt{.} & \texttt{c@1}\\
S776 & F4.3038 & \texttt{TB-{}} & 2 & \texttt{aa\textbar{}bb\textbar{}cc\textbar{}cc} & \texttt{abc} & \texttt{abc} & \texttt{abc} & \texttt{0123} & \texttt{01} & \texttt{+} & \texttt{.} & \texttt{a@0}\\
S777 & F4.3038 & \texttt{TB-{}} & 2 & \texttt{ab\textbar{}cc\textbar{}ba\textbar{}ba} & \texttt{abc} & \texttt{abc} & \texttt{abc} & \texttt{0123} & \texttt{01} & \texttt{+} & \texttt{.} & \texttt{c@1}\\
S778 & F4.3038 & \texttt{TB-{}} & 2 & \texttt{aa\textbar{}bb\textbar{}aa\textbar{}aa} & \texttt{ab} & \texttt{ab} & \texttt{ab} & \texttt{0123} & \texttt{01} & \texttt{+} & \texttt{.} & \texttt{b@1}\\
S779 & F4.3038 & \texttt{TB-{}} & 2 & \texttt{ab\textbar{}ab\textbar{}cd\textbar{}cd} & \texttt{ac} & \texttt{ac} & \texttt{abcd} & \texttt{0123} & \texttt{0} & \texttt{+} & \texttt{.} & \texttt{a@0}\\
S780 & F4.3038 & \texttt{TB-{}} & 2 & \texttt{ab\textbar{}ba\textbar{}cc\textbar{}cc} & \texttt{abc} & \texttt{abc} & \texttt{abc} & \texttt{0123} & \texttt{01} & \texttt{+} & \texttt{.} & \texttt{b@0}\\
S781 & F4.3038 & \texttt{TB-{}} & 2 & \texttt{aa\textbar{}aa\textbar{}bb\textbar{}bb} & \texttt{ab} & \texttt{ab} & \texttt{ab} & \texttt{0123} & \texttt{01} & \texttt{+} & \texttt{.} & \texttt{a@0}\\
S782 & F4.3038 & \texttt{TB-{}} & 2 & \texttt{ab\textbar{}cd\textbar{}cd\textbar{}ab} & \texttt{ac} & \texttt{ac} & \texttt{abcd} & \texttt{0123} & \texttt{0} & \texttt{+} & \texttt{.} & \texttt{c@1}\\
S783 & F4.3038 & \texttt{TB-{}} & 2 & \texttt{aa\textbar{}bb\textbar{}bb\textbar{}aa} & \texttt{ab} & \texttt{ab} & \texttt{ab} & \texttt{0123} & \texttt{01} & \texttt{+} & \texttt{.} & \texttt{b@1}\\
S784 & F4.3038 & \texttt{TB-{}} & 2 & \texttt{ab\textbar{}ab\textbar{}ab\textbar{}ab} & \texttt{a} & \texttt{a} & \texttt{ab} & \texttt{0123} & \texttt{0} & \texttt{+} & \texttt{S} & \texttt{none}\\
S785 & F4.3038 & \texttt{TB-{}} & 2 & \texttt{ab\textbar{}ba\textbar{}ab\textbar{}ab} & \texttt{ab} & \texttt{ab} & \texttt{ab} & \texttt{0123} & \texttt{01} & \texttt{+} & \texttt{.} & \texttt{none}\\
S786 & F4.3038 & \texttt{TB-{}} & 2 & \texttt{ab\textbar{}ab\textbar{}ba\textbar{}ba} & \texttt{ab} & \texttt{ab} & \texttt{ab} & \texttt{0123} & \texttt{01} & \texttt{+} & \texttt{.} & \texttt{none}\\
S787 & F4.3038 & \texttt{TB-{}} & 2 & \texttt{ab\textbar{}ba\textbar{}ba\textbar{}ab} & \texttt{ab} & \texttt{ab} & \texttt{ab} & \texttt{0123} & \texttt{01} & \texttt{+} & \texttt{.} & \texttt{none}\\
S788 & F4.3038 & \texttt{TB-{}} & 2 & \texttt{aa\textbar{}aa\textbar{}aa\textbar{}aa} & \texttt{a} & \texttt{a} & \texttt{a} & \texttt{0123} & \texttt{01} & \texttt{+} & \texttt{S} & \texttt{none}\\
S789 & F4.3042 & \texttt{-{}B-{}} & 2 & \texttt{ab\textbar{}cd\textbar{}ef\textbar{}ab} & \texttt{ace} & \texttt{ace} & \texttt{abcdef} & \texttt{0123} & \texttt{0} & \texttt{+} & \texttt{.} & \texttt{c@1}\\
S790 & F4.3042 & \texttt{-{}B-{}} & 2 & \texttt{ab\textbar{}ab\textbar{}cd\textbar{}ab} & \texttt{ac} & \texttt{ac} & \texttt{abcd} & \texttt{0123} & \texttt{0} & \texttt{+} & \texttt{.} & \texttt{c@2}\\
S791 & F4.3042 & \texttt{-{}B-{}} & 2 & \texttt{aa\textbar{}bb\textbar{}cc\textbar{}aa} & \texttt{abc} & \texttt{abc} & \texttt{abc} & \texttt{0123} & \texttt{01} & \texttt{+} & \texttt{.} & \texttt{b@1}\\
S792 & F4.3042 & \texttt{-{}B-{}} & 2 & \texttt{ab\textbar{}ba\textbar{}cc\textbar{}ab} & \texttt{abc} & \texttt{abc} & \texttt{abc} & \texttt{0123} & \texttt{01} & \texttt{+} & \texttt{.} & \texttt{c@2}\\
S793 & F4.3042 & \texttt{-{}B-{}} & 2 & \texttt{aa\textbar{}aa\textbar{}bb\textbar{}aa} & \texttt{ab} & \texttt{ab} & \texttt{ab} & \texttt{0123} & \texttt{01} & \texttt{+} & \texttt{.} & \texttt{b@2}\\
S794 & F4.3042 & \texttt{-{}B-{}} & 2 & \texttt{ab\textbar{}cd\textbar{}cd\textbar{}ab} & \texttt{ac} & \texttt{ac} & \texttt{abcd} & \texttt{0123} & \texttt{0} & \texttt{+} & \texttt{.} & \texttt{c@1}\\
S795 & F4.3042 & \texttt{-{}B-{}} & 2 & \texttt{aa\textbar{}bc\textbar{}cb\textbar{}aa} & \texttt{abc} & \texttt{abc} & \texttt{abc} & \texttt{0123} & \texttt{01} & \texttt{+} & \texttt{.} & \texttt{c@1}\\
S796 & F4.3042 & \texttt{-{}B-{}} & 2 & \texttt{aa\textbar{}bb\textbar{}bb\textbar{}aa} & \texttt{ab} & \texttt{ab} & \texttt{ab} & \texttt{0123} & \texttt{01} & \texttt{+} & \texttt{.} & \texttt{b@1}\\
S797 & F4.3042 & \texttt{-{}B-{}} & 2 & \texttt{ab\textbar{}ab\textbar{}ab\textbar{}ab} & \texttt{a} & \texttt{a} & \texttt{ab} & \texttt{0123} & \texttt{0} & \texttt{+} & \texttt{S} & \texttt{none}\\
S798 & F4.3042 & \texttt{-{}B-{}} & 2 & \texttt{ab\textbar{}ab\textbar{}ba\textbar{}ab} & \texttt{ab} & \texttt{ab} & \texttt{ab} & \texttt{0123} & \texttt{01} & \texttt{+} & \texttt{.} & \texttt{none}\\
S799 & F4.3042 & \texttt{-{}B-{}} & 2 & \texttt{ab\textbar{}ba\textbar{}ba\textbar{}ab} & \texttt{ab} & \texttt{ab} & \texttt{ab} & \texttt{0123} & \texttt{01} & \texttt{+} & \texttt{.} & \texttt{none}\\
S800 & F4.3042 & \texttt{-{}B-{}} & 2 & \texttt{aa\textbar{}aa\textbar{}aa\textbar{}aa} & \texttt{a} & \texttt{a} & \texttt{a} & \texttt{0123} & \texttt{01} & \texttt{+} & \texttt{S} & \texttt{none}\\
S801 & F4.3044 & \texttt{TB4} & 2 & \texttt{ab\textbar{}cd\textbar{}ef\textbar{}gh} & \texttt{aceg} & \texttt{aceg} & \texttt{abcdefgh} & \texttt{0123} & \texttt{0} & \texttt{+} & \texttt{.} & \texttt{a@0}\\
S802 & F4.3044 & \texttt{TB4} & 2 & \texttt{ab\textbar{}ab\textbar{}cd\textbar{}ef} & \texttt{ace} & \texttt{ace} & \texttt{abcdef} & \texttt{0123} & \texttt{0} & \texttt{+} & \texttt{.} & \texttt{c@2}\\
S803 & F4.3044 & \texttt{TB4} & 2 & \texttt{ab\textbar{}ab\textbar{}ab\textbar{}cd} & \texttt{ac} & \texttt{ac} & \texttt{abcd} & \texttt{0123} & \texttt{0} & \texttt{+} & \texttt{.} & \texttt{c@3}\\
S804 & F4.3044 & \texttt{TB4} & 2 & \texttt{aa\textbar{}bb\textbar{}cc\textbar{}dd} & \texttt{abcd} & \texttt{abcd} & \texttt{abcd} & \texttt{0123} & \texttt{01} & \texttt{+} & \texttt{.} & \texttt{a@0}\\
S805 & F4.3044 & \texttt{TB4} & 2 & \texttt{ab\textbar{}cc\textbar{}dd\textbar{}ba} & \texttt{abcd} & \texttt{abcd} & \texttt{abcd} & \texttt{0123} & \texttt{01} & \texttt{+} & \texttt{.} & \texttt{c@1}\\
S806 & F4.3044 & \texttt{TB4} & 2 & \texttt{aa\textbar{}aa\textbar{}bb\textbar{}cc} & \texttt{abc} & \texttt{abc} & \texttt{abc} & \texttt{0123} & \texttt{01} & \texttt{+} & \texttt{.} & \texttt{b@2}\\
S807 & F4.3044 & \texttt{TB4} & 2 & \texttt{ab\textbar{}ab\textbar{}cc\textbar{}ba} & \texttt{abc} & \texttt{abc} & \texttt{abc} & \texttt{0123} & \texttt{01} & \texttt{+} & \texttt{.} & \texttt{c@2}\\
S808 & F4.3044 & \texttt{TB4} & 2 & \texttt{aa\textbar{}aa\textbar{}aa\textbar{}bb} & \texttt{ab} & \texttt{ab} & \texttt{ab} & \texttt{0123} & \texttt{01} & \texttt{+} & \texttt{.} & \texttt{b@3}\\
S809 & F4.3044 & \texttt{TB4} & 2 & \texttt{ab\textbar{}cd\textbar{}cd\textbar{}ab} & \texttt{ac} & \texttt{ac} & \texttt{abcd} & \texttt{0123} & \texttt{0} & \texttt{+} & \texttt{.} & \texttt{c@1}\\
S810 & F4.3044 & \texttt{TB4} & 2 & \texttt{ab\textbar{}cd\textbar{}dc\textbar{}ba} & \texttt{abcd} & \texttt{abcd} & \texttt{abcd} & \texttt{0123} & \texttt{01} & \texttt{+} & \texttt{.} & \texttt{d@1}\\
S811 & F4.3044 & \texttt{TB4} & 2 & \texttt{ab\textbar{}cc\textbar{}cc\textbar{}ba} & \texttt{abc} & \texttt{abc} & \texttt{abc} & \texttt{0123} & \texttt{01} & \texttt{+} & \texttt{.} & \texttt{b@0}\\
S812 & F4.3044 & \texttt{TB4} & 2 & \texttt{aa\textbar{}bb\textbar{}bb\textbar{}aa} & \texttt{ab} & \texttt{ab} & \texttt{ab} & \texttt{0123} & \texttt{01} & \texttt{+} & \texttt{.} & \texttt{b@1}\\
S813 & F4.3044 & \texttt{TB4} & 2 & \texttt{ab\textbar{}ab\textbar{}ab\textbar{}ab} & \texttt{a} & \texttt{a} & \texttt{ab} & \texttt{0123} & \texttt{0} & \texttt{+} & \texttt{S} & \texttt{none}\\
S814 & F4.3044 & \texttt{TB4} & 2 & \texttt{ab\textbar{}ab\textbar{}ab\textbar{}ba} & \texttt{ab} & \texttt{ab} & \texttt{ab} & \texttt{0123} & \texttt{01} & \texttt{+} & \texttt{.} & \texttt{none}\\
S815 & F4.3044 & \texttt{TB4} & 2 & \texttt{ab\textbar{}ba\textbar{}ba\textbar{}ab} & \texttt{ab} & \texttt{ab} & \texttt{ab} & \texttt{0123} & \texttt{01} & \texttt{+} & \texttt{.} & \texttt{none}\\
S816 & F4.3044 & \texttt{TB4} & 2 & \texttt{aa\textbar{}aa\textbar{}aa\textbar{}aa} & \texttt{a} & \texttt{a} & \texttt{a} & \texttt{0123} & \texttt{01} & \texttt{+} & \texttt{S} & \texttt{none}\\
\bottomrule\end{longtable}}


\subsection{Models on non-symmetric frames}

% generated by gen/generate.py -- do not edit
{\scriptsize\setlength{\tabcolsep}{2.6pt}
\begin{longtable}{@{}lllrlllllllll@{}}\toprule
id & frame & TB4 & $m$ & $\Pi$ (atoms by point) & core & $G$ & $\mathit{NE}$ & $\exists G$ at & gods@0 & $\nec\exists G$ & Sc & collapse wit.\\
\midrule\endfirsthead
\toprule
id & frame & TB4 & $m$ & $\Pi$ (atoms by point) & core & $G$ & $\mathit{NE}$ & $\exists G$ at & gods@0 & $\nec\exists G$ & Sc & collapse wit.\\
\midrule\endhead
\multicolumn{13}{@{}l}{\rule{0pt}{2.8ex}\footnotesize\bfseries $n=2$ worlds, $m=1$ individuals: 4 models}\\*
N1 & F2.5 & \texttt{-{}-{}4} & 1 & \texttt{a\textbar{}b} & \texttt{a} & \texttt{ab} & \texttt{ab} & \texttt{01} & \texttt{0} & \texttt{+} & \texttt{S} & \texttt{b@1}\\
N2 & F2.5 & \texttt{-{}-{}4} & 1 & \texttt{a\textbar{}a} & \texttt{a} & \texttt{a} & \texttt{a} & \texttt{01} & \texttt{0} & \texttt{+} & \texttt{S} & \texttt{none}\\
N3 & F2.9 & \texttt{T-{}4} & 1 & \texttt{a\textbar{}b} & \texttt{b} & \texttt{b} & \texttt{ab} & \texttt{1} & \texttt{-{}} & \texttt{-{}0} & \texttt{s} & \texttt{a@0}\\
N4 & F2.9 & \texttt{T-{}4} & 1 & \texttt{a\textbar{}a} & \texttt{a} & \texttt{a} & \texttt{a} & \texttt{01} & \texttt{0} & \texttt{+} & \texttt{S} & \texttt{none}\\
\multicolumn{13}{@{}l}{\rule{0pt}{2.8ex}\footnotesize\bfseries $n=2$ worlds, $m=2$ individuals: 10 models}\\*
N5 & F2.5 & \texttt{-{}-{}4} & 2 & \texttt{ab\textbar{}cd} & \texttt{a} & \texttt{ac} & \texttt{abcd} & \texttt{01} & \texttt{0} & \texttt{+} & \texttt{S} & \texttt{c@1}\\
N6 & F2.5 & \texttt{-{}-{}4} & 2 & \texttt{aa\textbar{}bc} & \texttt{a} & \texttt{abc} & \texttt{abc} & \texttt{01} & \texttt{01} & \texttt{+} & \texttt{S} & \texttt{b@1}\\
N7 & F2.5 & \texttt{-{}-{}4} & 2 & \texttt{aa\textbar{}bb} & \texttt{a} & \texttt{ab} & \texttt{ab} & \texttt{01} & \texttt{01} & \texttt{+} & \texttt{S} & \texttt{b@1}\\
N8 & F2.5 & \texttt{-{}-{}4} & 2 & \texttt{ab\textbar{}ab} & \texttt{a} & \texttt{a} & \texttt{ab} & \texttt{01} & \texttt{0} & \texttt{+} & \texttt{S} & \texttt{none}\\
N9 & F2.5 & \texttt{-{}-{}4} & 2 & \texttt{aa\textbar{}aa} & \texttt{a} & \texttt{a} & \texttt{a} & \texttt{01} & \texttt{01} & \texttt{+} & \texttt{S} & \texttt{none}\\
N10 & F2.9 & \texttt{T-{}4} & 2 & \texttt{ab\textbar{}cd} & \texttt{c} & \texttt{c} & \texttt{abcd} & \texttt{1} & \texttt{-{}} & \texttt{-{}0} & \texttt{s} & \texttt{a@0}\\
N11 & F2.9 & \texttt{T-{}4} & 2 & \texttt{ab\textbar{}cc} & \texttt{c} & \texttt{c} & \texttt{abc} & \texttt{1} & \texttt{-{}} & \texttt{-{}0} & \texttt{s} & \texttt{a@0}\\
N12 & F2.9 & \texttt{T-{}4} & 2 & \texttt{aa\textbar{}bb} & \texttt{b} & \texttt{b} & \texttt{ab} & \texttt{1} & \texttt{-{}} & \texttt{-{}0} & \texttt{s} & \texttt{a@0}\\
N13 & F2.9 & \texttt{T-{}4} & 2 & \texttt{ab\textbar{}ab} & \texttt{a} & \texttt{a} & \texttt{ab} & \texttt{01} & \texttt{0} & \texttt{+} & \texttt{S} & \texttt{none}\\
N14 & F2.9 & \texttt{T-{}4} & 2 & \texttt{aa\textbar{}aa} & \texttt{a} & \texttt{a} & \texttt{a} & \texttt{01} & \texttt{01} & \texttt{+} & \texttt{S} & \texttt{none}\\
\multicolumn{13}{@{}l}{\rule{0pt}{2.8ex}\footnotesize\bfseries $n=2$ worlds, $m=3$ individuals: 24 models}\\*
N15 & F2.5 & \texttt{-{}-{}4} & 3 & \texttt{abc\textbar{}def} & \texttt{a} & \texttt{ad} & \texttt{abcdef} & \texttt{01} & \texttt{0} & \texttt{+} & \texttt{S} & \texttt{d@1}\\
N16 & F2.5 & \texttt{-{}-{}4} & 3 & \texttt{aab\textbar{}cde} & \texttt{b} & \texttt{be} & \texttt{abcde} & \texttt{01} & \texttt{2} & \texttt{+} & \texttt{S} & \texttt{c@1}\\
N17 & F2.5 & \texttt{-{}-{}4} & 3 & \texttt{aab\textbar{}cde} & \texttt{a} & \texttt{acd} & \texttt{abcde} & \texttt{01} & \texttt{01} & \texttt{+} & \texttt{S} & \texttt{c@1}\\
N18 & F2.5 & \texttt{-{}-{}4} & 3 & \texttt{aab\textbar{}ccd} & \texttt{b} & \texttt{bd} & \texttt{abcd} & \texttt{01} & \texttt{2} & \texttt{+} & \texttt{S} & \texttt{d@1}\\
N19 & F2.5 & \texttt{-{}-{}4} & 3 & \texttt{aab\textbar{}ccd} & \texttt{a} & \texttt{ac} & \texttt{abcd} & \texttt{01} & \texttt{01} & \texttt{+} & \texttt{S} & \texttt{d@1}\\
N20 & F2.5 & \texttt{-{}-{}4} & 3 & \texttt{aaa\textbar{}bcd} & \texttt{a} & \texttt{abcd} & \texttt{abcd} & \texttt{01} & \texttt{012} & \texttt{+} & \texttt{S} & \texttt{b@1}\\
N21 & F2.5 & \texttt{-{}-{}4} & 3 & \texttt{aaa\textbar{}bbc} & \texttt{a} & \texttt{abc} & \texttt{abc} & \texttt{01} & \texttt{012} & \texttt{+} & \texttt{S} & \texttt{c@1}\\
N22 & F2.5 & \texttt{-{}-{}4} & 3 & \texttt{aaa\textbar{}bbb} & \texttt{a} & \texttt{ab} & \texttt{ab} & \texttt{01} & \texttt{012} & \texttt{+} & \texttt{S} & \texttt{b@1}\\
N23 & F2.5 & \texttt{-{}-{}4} & 3 & \texttt{abc\textbar{}abc} & \texttt{a} & \texttt{a} & \texttt{abc} & \texttt{01} & \texttt{0} & \texttt{+} & \texttt{S} & \texttt{none}\\
N24 & F2.5 & \texttt{-{}-{}4} & 3 & \texttt{aab\textbar{}aab} & \texttt{b} & \texttt{b} & \texttt{ab} & \texttt{01} & \texttt{2} & \texttt{+} & \texttt{S} & \texttt{none}\\
N25 & F2.5 & \texttt{-{}-{}4} & 3 & \texttt{aab\textbar{}aab} & \texttt{a} & \texttt{a} & \texttt{ab} & \texttt{01} & \texttt{01} & \texttt{+} & \texttt{S} & \texttt{none}\\
N26 & F2.5 & \texttt{-{}-{}4} & 3 & \texttt{aaa\textbar{}aaa} & \texttt{a} & \texttt{a} & \texttt{a} & \texttt{01} & \texttt{012} & \texttt{+} & \texttt{S} & \texttt{none}\\
N27 & F2.9 & \texttt{T-{}4} & 3 & \texttt{abc\textbar{}def} & \texttt{d} & \texttt{d} & \texttt{abcdef} & \texttt{1} & \texttt{-{}} & \texttt{-{}0} & \texttt{s} & \texttt{a@0}\\
N28 & F2.9 & \texttt{T-{}4} & 3 & \texttt{abc\textbar{}dde} & \texttt{e} & \texttt{e} & \texttt{abcde} & \texttt{1} & \texttt{-{}} & \texttt{-{}0} & \texttt{s} & \texttt{a@0}\\
N29 & F2.9 & \texttt{T-{}4} & 3 & \texttt{abc\textbar{}dde} & \texttt{d} & \texttt{d} & \texttt{abcde} & \texttt{1} & \texttt{-{}} & \texttt{-{}0} & \texttt{s} & \texttt{a@0}\\
N30 & F2.9 & \texttt{T-{}4} & 3 & \texttt{abc\textbar{}ddd} & \texttt{d} & \texttt{d} & \texttt{abcd} & \texttt{1} & \texttt{-{}} & \texttt{-{}0} & \texttt{s} & \texttt{a@0}\\
N31 & F2.9 & \texttt{T-{}4} & 3 & \texttt{aab\textbar{}ccd} & \texttt{d} & \texttt{d} & \texttt{abcd} & \texttt{1} & \texttt{-{}} & \texttt{-{}0} & \texttt{s} & \texttt{b@0}\\
N32 & F2.9 & \texttt{T-{}4} & 3 & \texttt{aab\textbar{}ccd} & \texttt{c} & \texttt{c} & \texttt{abcd} & \texttt{1} & \texttt{-{}} & \texttt{-{}0} & \texttt{s} & \texttt{b@0}\\
N33 & F2.9 & \texttt{T-{}4} & 3 & \texttt{aab\textbar{}ccc} & \texttt{c} & \texttt{c} & \texttt{abc} & \texttt{1} & \texttt{-{}} & \texttt{-{}0} & \texttt{s} & \texttt{b@0}\\
N34 & F2.9 & \texttt{T-{}4} & 3 & \texttt{aaa\textbar{}bbb} & \texttt{b} & \texttt{b} & \texttt{ab} & \texttt{1} & \texttt{-{}} & \texttt{-{}0} & \texttt{s} & \texttt{a@0}\\
N35 & F2.9 & \texttt{T-{}4} & 3 & \texttt{abc\textbar{}abc} & \texttt{a} & \texttt{a} & \texttt{abc} & \texttt{01} & \texttt{0} & \texttt{+} & \texttt{S} & \texttt{none}\\
N36 & F2.9 & \texttt{T-{}4} & 3 & \texttt{aab\textbar{}aab} & \texttt{b} & \texttt{b} & \texttt{ab} & \texttt{01} & \texttt{2} & \texttt{+} & \texttt{S} & \texttt{none}\\
N37 & F2.9 & \texttt{T-{}4} & 3 & \texttt{aab\textbar{}aab} & \texttt{a} & \texttt{a} & \texttt{ab} & \texttt{01} & \texttt{01} & \texttt{+} & \texttt{S} & \texttt{none}\\
N38 & F2.9 & \texttt{T-{}4} & 3 & \texttt{aaa\textbar{}aaa} & \texttt{a} & \texttt{a} & \texttt{a} & \texttt{01} & \texttt{012} & \texttt{+} & \texttt{S} & \texttt{none}\\
\multicolumn{13}{@{}l}{\rule{0pt}{2.8ex}\footnotesize\bfseries $n=3$ worlds, $m=1$ individuals: 116 models}\\*
N39 & F3.22 & \texttt{-{}-{}4} & 1 & \texttt{a\textbar{}b\textbar{}c} & \texttt{a} & \texttt{abc} & \texttt{abc} & \texttt{012} & \texttt{0} & \texttt{+} & \texttt{S} & \texttt{b@1}\\
N40 & F3.22 & \texttt{-{}-{}4} & 1 & \texttt{a\textbar{}b\textbar{}b} & \texttt{a} & \texttt{ab} & \texttt{ab} & \texttt{012} & \texttt{0} & \texttt{+} & \texttt{S} & \texttt{b@1}\\
N41 & F3.22 & \texttt{-{}-{}4} & 1 & \texttt{a\textbar{}a\textbar{}b} & \texttt{a} & \texttt{ab} & \texttt{ab} & \texttt{012} & \texttt{0} & \texttt{+} & \texttt{S} & \texttt{b@2}\\
N42 & F3.22 & \texttt{-{}-{}4} & 1 & \texttt{a\textbar{}a\textbar{}a} & \texttt{a} & \texttt{a} & \texttt{a} & \texttt{012} & \texttt{0} & \texttt{+} & \texttt{S} & \texttt{none}\\
N43 & F3.23 & \texttt{-{}-{}-{}} & 1 & \texttt{a\textbar{}a\textbar{}b} & \texttt{a} & \texttt{ab} & \texttt{ab} & \texttt{012} & \texttt{0} & \texttt{+} & \texttt{S} & \texttt{b@2}\\
N44 & F3.23 & \texttt{-{}-{}-{}} & 1 & \texttt{a\textbar{}a\textbar{}a} & \texttt{a} & \texttt{a} & \texttt{a} & \texttt{012} & \texttt{0} & \texttt{+} & \texttt{S} & \texttt{none}\\
N45 & F3.24 & \texttt{-{}-{}4} & 1 & \texttt{a\textbar{}a\textbar{}b} & \texttt{a} & \texttt{ab} & \texttt{ab} & \texttt{012} & \texttt{0} & \texttt{+} & \texttt{S} & \texttt{b@2}\\
N46 & F3.24 & \texttt{-{}-{}4} & 1 & \texttt{a\textbar{}a\textbar{}a} & \texttt{a} & \texttt{a} & \texttt{a} & \texttt{012} & \texttt{0} & \texttt{+} & \texttt{S} & \texttt{none}\\
N47 & F3.25 & \texttt{-{}-{}-{}} & 1 & \texttt{a\textbar{}a\textbar{}b} & \texttt{a} & \texttt{ab} & \texttt{ab} & \texttt{012} & \texttt{0} & \texttt{+} & \texttt{S} & \texttt{b@2}\\
N48 & F3.25 & \texttt{-{}-{}-{}} & 1 & \texttt{a\textbar{}a\textbar{}a} & \texttt{a} & \texttt{a} & \texttt{a} & \texttt{012} & \texttt{0} & \texttt{+} & \texttt{S} & \texttt{none}\\
N49 & F3.27 & \texttt{-{}-{}-{}} & 1 & \texttt{a\textbar{}a\textbar{}a} & \texttt{a} & \texttt{a} & \texttt{a} & \texttt{012} & \texttt{0} & \texttt{+} & \texttt{S} & \texttt{none}\\
N50 & F3.38 & \texttt{-{}-{}-{}} & 1 & \texttt{a\textbar{}a\textbar{}b} & \texttt{a} & \texttt{ab} & \texttt{ab} & \texttt{012} & \texttt{0} & \texttt{+} & \texttt{S} & \texttt{b@2}\\
N51 & F3.38 & \texttt{-{}-{}-{}} & 1 & \texttt{a\textbar{}a\textbar{}a} & \texttt{a} & \texttt{a} & \texttt{a} & \texttt{012} & \texttt{0} & \texttt{+} & \texttt{S} & \texttt{none}\\
N52 & F3.40 & \texttt{-{}-{}-{}} & 1 & \texttt{a\textbar{}a\textbar{}b} & \texttt{a} & \texttt{ab} & \texttt{ab} & \texttt{012} & \texttt{0} & \texttt{+} & \texttt{S} & \texttt{b@2}\\
N53 & F3.40 & \texttt{-{}-{}-{}} & 1 & \texttt{a\textbar{}a\textbar{}a} & \texttt{a} & \texttt{a} & \texttt{a} & \texttt{012} & \texttt{0} & \texttt{+} & \texttt{S} & \texttt{none}\\
N54 & F3.42 & \texttt{-{}-{}-{}} & 1 & \texttt{a\textbar{}a\textbar{}a} & \texttt{a} & \texttt{a} & \texttt{a} & \texttt{012} & \texttt{0} & \texttt{+} & \texttt{S} & \texttt{none}\\
N55 & F3.43 & \texttt{-{}-{}4} & 1 & \texttt{a\textbar{}b\textbar{}c} & \texttt{a} & \texttt{ac} & \texttt{abc} & \texttt{02} & \texttt{0} & \texttt{-{}1} & \texttt{s} & \texttt{b@1}\\
N56 & F3.43 & \texttt{-{}-{}4} & 1 & \texttt{a\textbar{}b\textbar{}a} & \texttt{a} & \texttt{a} & \texttt{ab} & \texttt{02} & \texttt{0} & \texttt{-{}1} & \texttt{s} & \texttt{b@1}\\
N57 & F3.43 & \texttt{-{}-{}4} & 1 & \texttt{a\textbar{}a\textbar{}b} & \texttt{a} & \texttt{ab} & \texttt{ab} & \texttt{012} & \texttt{0} & \texttt{+} & \texttt{S} & \texttt{b@2}\\
N58 & F3.43 & \texttt{-{}-{}4} & 1 & \texttt{a\textbar{}a\textbar{}a} & \texttt{a} & \texttt{a} & \texttt{a} & \texttt{012} & \texttt{0} & \texttt{+} & \texttt{S} & \texttt{none}\\
N59 & F3.44 & \texttt{-{}-{}-{}} & 1 & \texttt{a\textbar{}a\textbar{}b} & \texttt{a} & \texttt{ab} & \texttt{ab} & \texttt{012} & \texttt{0} & \texttt{+} & \texttt{S} & \texttt{b@2}\\
N60 & F3.44 & \texttt{-{}-{}-{}} & 1 & \texttt{a\textbar{}a\textbar{}a} & \texttt{a} & \texttt{a} & \texttt{a} & \texttt{012} & \texttt{0} & \texttt{+} & \texttt{S} & \texttt{none}\\
N61 & F3.45 & \texttt{-{}-{}-{}} & 1 & \texttt{a\textbar{}b\textbar{}a} & \texttt{a} & \texttt{a} & \texttt{ab} & \texttt{02} & \texttt{0} & \texttt{-{}1} & \texttt{s} & \texttt{b@1}\\
N62 & F3.45 & \texttt{-{}-{}-{}} & 1 & \texttt{a\textbar{}a\textbar{}a} & \texttt{a} & \texttt{a} & \texttt{a} & \texttt{012} & \texttt{0} & \texttt{+} & \texttt{S} & \texttt{none}\\
N63 & F3.46 & \texttt{-{}-{}-{}} & 1 & \texttt{a\textbar{}a\textbar{}a} & \texttt{a} & \texttt{a} & \texttt{a} & \texttt{012} & \texttt{0} & \texttt{+} & \texttt{S} & \texttt{none}\\
N64 & F3.47 & \texttt{-{}-{}-{}} & 1 & \texttt{a\textbar{}a\textbar{}a} & \texttt{a} & \texttt{a} & \texttt{a} & \texttt{012} & \texttt{0} & \texttt{+} & \texttt{S} & \texttt{none}\\
N65 & F3.48 & \texttt{-{}-{}4} & 1 & \texttt{a\textbar{}b\textbar{}c} & \texttt{a} & \texttt{ac} & \texttt{abc} & \texttt{02} & \texttt{0} & \texttt{+} & \texttt{s} & \texttt{b@1}\\
N66 & F3.48 & \texttt{-{}-{}4} & 1 & \texttt{a\textbar{}b\textbar{}a} & \texttt{a} & \texttt{ab} & \texttt{ab} & \texttt{012} & \texttt{0} & \texttt{+} & \texttt{S} & \texttt{b@1}\\
N67 & F3.48 & \texttt{-{}-{}4} & 1 & \texttt{a\textbar{}a\textbar{}a} & \texttt{a} & \texttt{a} & \texttt{a} & \texttt{012} & \texttt{0} & \texttt{+} & \texttt{S} & \texttt{none}\\
N68 & F3.49 & \texttt{-{}-{}-{}} & 1 & \texttt{a\textbar{}b\textbar{}a} & \texttt{a} & \texttt{ab} & \texttt{ab} & \texttt{012} & \texttt{0} & \texttt{+} & \texttt{S} & \texttt{b@1}\\
N69 & F3.49 & \texttt{-{}-{}-{}} & 1 & \texttt{a\textbar{}a\textbar{}a} & \texttt{a} & \texttt{a} & \texttt{a} & \texttt{012} & \texttt{0} & \texttt{+} & \texttt{S} & \texttt{none}\\
N70 & F3.50 & \texttt{-{}-{}-{}} & 1 & \texttt{a\textbar{}b\textbar{}a} & \texttt{a} & \texttt{a} & \texttt{ab} & \texttt{02} & \texttt{0} & \texttt{-{}1} & \texttt{s} & \texttt{b@1}\\
N71 & F3.50 & \texttt{-{}-{}-{}} & 1 & \texttt{a\textbar{}a\textbar{}a} & \texttt{a} & \texttt{a} & \texttt{a} & \texttt{012} & \texttt{0} & \texttt{+} & \texttt{S} & \texttt{none}\\
N72 & F3.51 & \texttt{-{}-{}4} & 1 & \texttt{a\textbar{}a\textbar{}b} & \texttt{a} & \texttt{ab} & \texttt{ab} & \texttt{012} & \texttt{0} & \texttt{+} & \texttt{S} & \texttt{b@2}\\
N73 & F3.51 & \texttt{-{}-{}4} & 1 & \texttt{a\textbar{}a\textbar{}a} & \texttt{a} & \texttt{a} & \texttt{a} & \texttt{012} & \texttt{0} & \texttt{+} & \texttt{S} & \texttt{none}\\
N74 & F3.52 & \texttt{T-{}4} & 1 & \texttt{a\textbar{}b\textbar{}b} & \texttt{b} & \texttt{b} & \texttt{ab} & \texttt{12} & \texttt{-{}} & \texttt{-{}0} & \texttt{s} & \texttt{a@0}\\
N75 & F3.52 & \texttt{T-{}4} & 1 & \texttt{a\textbar{}a\textbar{}a} & \texttt{a} & \texttt{a} & \texttt{a} & \texttt{012} & \texttt{0} & \texttt{+} & \texttt{S} & \texttt{none}\\
N76 & F3.57 & \texttt{-{}-{}-{}} & 1 & \texttt{a\textbar{}a\textbar{}a} & \texttt{a} & \texttt{a} & \texttt{a} & \texttt{012} & \texttt{0} & \texttt{+} & \texttt{S} & \texttt{none}\\
N77 & F3.58 & \texttt{-{}-{}-{}} & 1 & \texttt{a\textbar{}a\textbar{}b} & \texttt{a} & \texttt{ab} & \texttt{ab} & \texttt{012} & \texttt{0} & \texttt{+} & \texttt{S} & \texttt{b@2}\\
N78 & F3.58 & \texttt{-{}-{}-{}} & 1 & \texttt{a\textbar{}a\textbar{}a} & \texttt{a} & \texttt{a} & \texttt{a} & \texttt{012} & \texttt{0} & \texttt{+} & \texttt{S} & \texttt{none}\\
N79 & F3.59 & \texttt{-{}-{}-{}} & 1 & \texttt{a\textbar{}b\textbar{}a} & \texttt{a} & \texttt{a} & \texttt{ab} & \texttt{02} & \texttt{0} & \texttt{-{}1} & \texttt{s} & \texttt{b@1}\\
N80 & F3.59 & \texttt{-{}-{}-{}} & 1 & \texttt{a\textbar{}a\textbar{}a} & \texttt{a} & \texttt{a} & \texttt{a} & \texttt{012} & \texttt{0} & \texttt{+} & \texttt{S} & \texttt{none}\\
N81 & F3.61 & \texttt{-{}-{}-{}} & 1 & \texttt{a\textbar{}a\textbar{}a} & \texttt{a} & \texttt{a} & \texttt{a} & \texttt{012} & \texttt{0} & \texttt{+} & \texttt{S} & \texttt{none}\\
N82 & F3.62 & \texttt{-{}-{}-{}} & 1 & \texttt{a\textbar{}a\textbar{}a} & \texttt{a} & \texttt{a} & \texttt{a} & \texttt{012} & \texttt{0} & \texttt{+} & \texttt{S} & \texttt{none}\\
N83 & F3.63 & \texttt{-{}-{}-{}} & 1 & \texttt{a\textbar{}b\textbar{}a} & \texttt{a} & \texttt{ab} & \texttt{ab} & \texttt{012} & \texttt{0} & \texttt{+} & \texttt{S} & \texttt{b@1}\\
N84 & F3.63 & \texttt{-{}-{}-{}} & 1 & \texttt{a\textbar{}a\textbar{}a} & \texttt{a} & \texttt{a} & \texttt{a} & \texttt{012} & \texttt{0} & \texttt{+} & \texttt{S} & \texttt{none}\\
N85 & F3.64 & \texttt{-{}-{}-{}} & 1 & \texttt{a\textbar{}a\textbar{}a} & \texttt{a} & \texttt{a} & \texttt{a} & \texttt{012} & \texttt{0} & \texttt{+} & \texttt{S} & \texttt{none}\\
N86 & F3.65 & \texttt{-{}-{}-{}} & 1 & \texttt{a\textbar{}a\textbar{}a} & \texttt{a} & \texttt{a} & \texttt{a} & \texttt{012} & \texttt{0} & \texttt{+} & \texttt{S} & \texttt{none}\\
N87 & F3.66 & \texttt{-{}-{}-{}} & 1 & \texttt{a\textbar{}b\textbar{}a} & \texttt{a} & \texttt{a} & \texttt{ab} & \texttt{02} & \texttt{0} & \texttt{-{}1} & \texttt{s} & \texttt{b@1}\\
N88 & F3.66 & \texttt{-{}-{}-{}} & 1 & \texttt{a\textbar{}a\textbar{}a} & \texttt{a} & \texttt{a} & \texttt{a} & \texttt{012} & \texttt{0} & \texttt{+} & \texttt{S} & \texttt{none}\\
N89 & F3.67 & \texttt{-{}-{}-{}} & 1 & \texttt{a\textbar{}a\textbar{}a} & \texttt{a} & \texttt{a} & \texttt{a} & \texttt{012} & \texttt{0} & \texttt{+} & \texttt{S} & \texttt{none}\\
N90 & F3.68 & \texttt{-{}-{}4} & 1 & \texttt{a\textbar{}b\textbar{}c} & \texttt{a} & \texttt{ac} & \texttt{abc} & \texttt{02} & \texttt{0} & \texttt{-{}1} & \texttt{s} & \texttt{b@1}\\
N91 & F3.68 & \texttt{-{}-{}4} & 1 & \texttt{a\textbar{}b\textbar{}a} & \texttt{a} & \texttt{a} & \texttt{ab} & \texttt{02} & \texttt{0} & \texttt{-{}1} & \texttt{s} & \texttt{b@1}\\
N92 & F3.68 & \texttt{-{}-{}4} & 1 & \texttt{a\textbar{}a\textbar{}a} & \texttt{a} & \texttt{a} & \texttt{a} & \texttt{012} & \texttt{0} & \texttt{+} & \texttt{S} & \texttt{none}\\
N93 & F3.69 & \texttt{-{}-{}-{}} & 1 & \texttt{a\textbar{}b\textbar{}a} & \texttt{a} & \texttt{a} & \texttt{ab} & \texttt{02} & \texttt{0} & \texttt{-{}1} & \texttt{s} & \texttt{b@1}\\
N94 & F3.69 & \texttt{-{}-{}-{}} & 1 & \texttt{a\textbar{}a\textbar{}a} & \texttt{a} & \texttt{a} & \texttt{a} & \texttt{012} & \texttt{0} & \texttt{+} & \texttt{S} & \texttt{none}\\
N95 & F3.70 & \texttt{-{}-{}4} & 1 & \texttt{a\textbar{}b\textbar{}c} & \texttt{b} & \texttt{b} & \texttt{abc} & \texttt{1} & \texttt{-{}} & \texttt{-{}0} & \texttt{s} & \texttt{a@0}\\
N96 & F3.70 & \texttt{-{}-{}4} & 1 & \texttt{a\textbar{}b\textbar{}a} & \texttt{b} & \texttt{b} & \texttt{ab} & \texttt{1} & \texttt{-{}} & \texttt{-{}0} & \texttt{s} & \texttt{a@0}\\
N97 & F3.70 & \texttt{-{}-{}4} & 1 & \texttt{a\textbar{}a\textbar{}b} & \texttt{a} & \texttt{ab} & \texttt{ab} & \texttt{012} & \texttt{0} & \texttt{+} & \texttt{S} & \texttt{b@2}\\
N98 & F3.70 & \texttt{-{}-{}4} & 1 & \texttt{a\textbar{}a\textbar{}a} & \texttt{a} & \texttt{a} & \texttt{a} & \texttt{012} & \texttt{0} & \texttt{+} & \texttt{S} & \texttt{none}\\
N99 & F3.71 & \texttt{-{}-{}-{}} & 1 & \texttt{a\textbar{}a\textbar{}a} & \texttt{a} & \texttt{a} & \texttt{a} & \texttt{012} & \texttt{0} & \texttt{+} & \texttt{S} & \texttt{none}\\
N100 & F3.72 & \texttt{-{}-{}-{}} & 1 & \texttt{a\textbar{}b\textbar{}c} & \texttt{b} & \texttt{b} & \texttt{abc} & \texttt{1} & \texttt{-{}} & \texttt{-{}0} & \texttt{s} & \texttt{a@0}\\
N101 & F3.72 & \texttt{-{}-{}-{}} & 1 & \texttt{a\textbar{}b\textbar{}a} & \texttt{b} & \texttt{b} & \texttt{ab} & \texttt{1} & \texttt{-{}} & \texttt{-{}0} & \texttt{s} & \texttt{a@0}\\
N102 & F3.72 & \texttt{-{}-{}-{}} & 1 & \texttt{a\textbar{}a\textbar{}a} & \texttt{a} & \texttt{a} & \texttt{a} & \texttt{012} & \texttt{0} & \texttt{+} & \texttt{S} & \texttt{none}\\
N103 & F3.73 & \texttt{T-{}4} & 1 & \texttt{a\textbar{}b\textbar{}b} & \texttt{b} & \texttt{b} & \texttt{ab} & \texttt{12} & \texttt{-{}} & \texttt{-{}0} & \texttt{s} & \texttt{a@0}\\
N104 & F3.73 & \texttt{T-{}4} & 1 & \texttt{a\textbar{}a\textbar{}a} & \texttt{a} & \texttt{a} & \texttt{a} & \texttt{012} & \texttt{0} & \texttt{+} & \texttt{S} & \texttt{none}\\
N105 & F3.75 & \texttt{T-{}-{}} & 1 & \texttt{a\textbar{}b\textbar{}a} & \texttt{a} & \texttt{a} & \texttt{ab} & \texttt{02} & \texttt{0} & \texttt{-{}1} & \texttt{s} & \texttt{b@1}\\
N106 & F3.75 & \texttt{T-{}-{}} & 1 & \texttt{a\textbar{}a\textbar{}a} & \texttt{a} & \texttt{a} & \texttt{a} & \texttt{012} & \texttt{0} & \texttt{+} & \texttt{S} & \texttt{none}\\
N107 & F3.76 & \texttt{T-{}4} & 1 & \texttt{a\textbar{}b\textbar{}c} & \texttt{c} & \texttt{c} & \texttt{abc} & \texttt{2} & \texttt{-{}} & \texttt{-{}0} & \texttt{s} & \texttt{a@0}\\
N108 & F3.76 & \texttt{T-{}4} & 1 & \texttt{a\textbar{}b\textbar{}a} & \texttt{a} & \texttt{a} & \texttt{ab} & \texttt{02} & \texttt{0} & \texttt{-{}1} & \texttt{s} & \texttt{b@1}\\
N109 & F3.76 & \texttt{T-{}4} & 1 & \texttt{a\textbar{}a\textbar{}b} & \texttt{b} & \texttt{b} & \texttt{ab} & \texttt{2} & \texttt{-{}} & \texttt{-{}0} & \texttt{s} & \texttt{a@0}\\
N110 & F3.76 & \texttt{T-{}4} & 1 & \texttt{a\textbar{}a\textbar{}a} & \texttt{a} & \texttt{a} & \texttt{a} & \texttt{012} & \texttt{0} & \texttt{+} & \texttt{S} & \texttt{none}\\
N111 & F3.79 & \texttt{-{}-{}-{}} & 1 & \texttt{a\textbar{}a\textbar{}a} & \texttt{a} & \texttt{a} & \texttt{a} & \texttt{012} & \texttt{0} & \texttt{+} & \texttt{S} & \texttt{none}\\
N112 & F3.80 & \texttt{-{}-{}-{}} & 1 & \texttt{a\textbar{}a\textbar{}a} & \texttt{a} & \texttt{a} & \texttt{a} & \texttt{012} & \texttt{0} & \texttt{+} & \texttt{S} & \texttt{none}\\
N113 & F3.81 & \texttt{-{}-{}-{}} & 1 & \texttt{a\textbar{}a\textbar{}a} & \texttt{a} & \texttt{a} & \texttt{a} & \texttt{012} & \texttt{0} & \texttt{+} & \texttt{S} & \texttt{none}\\
N114 & F3.82 & \texttt{-{}-{}-{}} & 1 & \texttt{a\textbar{}b\textbar{}a} & \texttt{a} & \texttt{a} & \texttt{ab} & \texttt{02} & \texttt{0} & \texttt{-{}1} & \texttt{s} & \texttt{b@1}\\
N115 & F3.82 & \texttt{-{}-{}-{}} & 1 & \texttt{a\textbar{}a\textbar{}a} & \texttt{a} & \texttt{a} & \texttt{a} & \texttt{012} & \texttt{0} & \texttt{+} & \texttt{S} & \texttt{none}\\
N116 & F3.83 & \texttt{-{}-{}-{}} & 1 & \texttt{a\textbar{}a\textbar{}a} & \texttt{a} & \texttt{a} & \texttt{a} & \texttt{012} & \texttt{0} & \texttt{+} & \texttt{S} & \texttt{none}\\
N117 & F3.84 & \texttt{-{}-{}-{}} & 1 & \texttt{a\textbar{}b\textbar{}c} & \texttt{b} & \texttt{b} & \texttt{abc} & \texttt{1} & \texttt{-{}} & \texttt{-{}0} & \texttt{s} & \texttt{a@0}\\
N118 & F3.84 & \texttt{-{}-{}-{}} & 1 & \texttt{a\textbar{}b\textbar{}a} & \texttt{b} & \texttt{b} & \texttt{ab} & \texttt{1} & \texttt{-{}} & \texttt{-{}0} & \texttt{s} & \texttt{a@0}\\
N119 & F3.84 & \texttt{-{}-{}-{}} & 1 & \texttt{a\textbar{}a\textbar{}a} & \texttt{a} & \texttt{a} & \texttt{a} & \texttt{012} & \texttt{0} & \texttt{+} & \texttt{S} & \texttt{none}\\
N120 & F3.85 & \texttt{-{}-{}4} & 1 & \texttt{a\textbar{}b\textbar{}c} & \texttt{ab} & \texttt{abc} & \texttt{abc} & \texttt{012} & \texttt{0} & \texttt{+} & \texttt{.} & \texttt{a@0}\\
N121 & F3.85 & \texttt{-{}-{}4} & 1 & \texttt{a\textbar{}b\textbar{}a} & \texttt{ab} & \texttt{ab} & \texttt{ab} & \texttt{012} & \texttt{0} & \texttt{+} & \texttt{.} & \texttt{b@1}\\
N122 & F3.85 & \texttt{-{}-{}4} & 1 & \texttt{a\textbar{}a\textbar{}b} & \texttt{a} & \texttt{ab} & \texttt{ab} & \texttt{012} & \texttt{0} & \texttt{+} & \texttt{S} & \texttt{b@2}\\
N123 & F3.85 & \texttt{-{}-{}4} & 1 & \texttt{a\textbar{}a\textbar{}a} & \texttt{a} & \texttt{a} & \texttt{a} & \texttt{012} & \texttt{0} & \texttt{+} & \texttt{S} & \texttt{none}\\
N124 & F3.86 & \texttt{-{}-{}-{}} & 1 & \texttt{a\textbar{}b\textbar{}b} & \texttt{ab} & \texttt{ab} & \texttt{ab} & \texttt{012} & \texttt{0} & \texttt{+} & \texttt{.} & \texttt{a@0}\\
N125 & F3.86 & \texttt{-{}-{}-{}} & 1 & \texttt{a\textbar{}a\textbar{}a} & \texttt{a} & \texttt{a} & \texttt{a} & \texttt{012} & \texttt{0} & \texttt{+} & \texttt{S} & \texttt{none}\\
N126 & F3.87 & \texttt{-{}-{}-{}} & 1 & \texttt{a\textbar{}b\textbar{}a} & \texttt{ab} & \texttt{ab} & \texttt{ab} & \texttt{012} & \texttt{0} & \texttt{+} & \texttt{.} & \texttt{b@1}\\
N127 & F3.87 & \texttt{-{}-{}-{}} & 1 & \texttt{a\textbar{}a\textbar{}a} & \texttt{a} & \texttt{a} & \texttt{a} & \texttt{012} & \texttt{0} & \texttt{+} & \texttt{S} & \texttt{none}\\
N128 & F3.90 & \texttt{T-{}-{}} & 1 & \texttt{a\textbar{}a\textbar{}a} & \texttt{a} & \texttt{a} & \texttt{a} & \texttt{012} & \texttt{0} & \texttt{+} & \texttt{S} & \texttt{none}\\
N129 & F3.91 & \texttt{T-{}4} & 1 & \texttt{a\textbar{}b\textbar{}c} & \texttt{c} & \texttt{c} & \texttt{abc} & \texttt{2} & \texttt{-{}} & \texttt{-{}0} & \texttt{s} & \texttt{a@0}\\
N130 & F3.91 & \texttt{T-{}4} & 1 & \texttt{a\textbar{}b\textbar{}b} & \texttt{b} & \texttt{b} & \texttt{ab} & \texttt{12} & \texttt{-{}} & \texttt{-{}0} & \texttt{s} & \texttt{a@0}\\
N131 & F3.91 & \texttt{T-{}4} & 1 & \texttt{a\textbar{}a\textbar{}b} & \texttt{b} & \texttt{b} & \texttt{ab} & \texttt{2} & \texttt{-{}} & \texttt{-{}0} & \texttt{s} & \texttt{a@0}\\
N132 & F3.91 & \texttt{T-{}4} & 1 & \texttt{a\textbar{}a\textbar{}a} & \texttt{a} & \texttt{a} & \texttt{a} & \texttt{012} & \texttt{0} & \texttt{+} & \texttt{S} & \texttt{none}\\
N133 & F3.92 & \texttt{T-{}-{}} & 1 & \texttt{a\textbar{}b\textbar{}b} & \texttt{ab} & \texttt{ab} & \texttt{ab} & \texttt{012} & \texttt{0} & \texttt{+} & \texttt{.} & \texttt{a@0}\\
N134 & F3.92 & \texttt{T-{}-{}} & 1 & \texttt{a\textbar{}a\textbar{}b} & \texttt{a} & \texttt{a} & \texttt{ab} & \texttt{01} & \texttt{0} & \texttt{-{}2} & \texttt{s} & \texttt{b@2}\\
N135 & F3.92 & \texttt{T-{}-{}} & 1 & \texttt{a\textbar{}a\textbar{}a} & \texttt{a} & \texttt{a} & \texttt{a} & \texttt{012} & \texttt{0} & \texttt{+} & \texttt{S} & \texttt{none}\\
N136 & F3.93 & \texttt{T-{}-{}} & 1 & \texttt{a\textbar{}a\textbar{}a} & \texttt{a} & \texttt{a} & \texttt{a} & \texttt{012} & \texttt{0} & \texttt{+} & \texttt{S} & \texttt{none}\\
N137 & F3.94 & \texttt{-{}-{}-{}} & 1 & \texttt{a\textbar{}a\textbar{}a} & \texttt{a} & \texttt{a} & \texttt{a} & \texttt{012} & \texttt{0} & \texttt{+} & \texttt{S} & \texttt{none}\\
N138 & F3.95 & \texttt{-{}-{}-{}} & 1 & \texttt{a\textbar{}b\textbar{}b} & \texttt{ab} & \texttt{ab} & \texttt{ab} & \texttt{012} & \texttt{0} & \texttt{+} & \texttt{.} & \texttt{a@0}\\
N139 & F3.95 & \texttt{-{}-{}-{}} & 1 & \texttt{a\textbar{}b\textbar{}a} & \texttt{ab} & \texttt{ab} & \texttt{ab} & \texttt{012} & \texttt{0} & \texttt{+} & \texttt{.} & \texttt{b@1}\\
N140 & F3.95 & \texttt{-{}-{}-{}} & 1 & \texttt{a\textbar{}a\textbar{}a} & \texttt{a} & \texttt{a} & \texttt{a} & \texttt{012} & \texttt{0} & \texttt{+} & \texttt{S} & \texttt{none}\\
N141 & F3.97 & \texttt{-{}-{}-{}} & 1 & \texttt{a\textbar{}b\textbar{}a} & \texttt{ab} & \texttt{ab} & \texttt{ab} & \texttt{012} & \texttt{0} & \texttt{+} & \texttt{.} & \texttt{b@1}\\
N142 & F3.97 & \texttt{-{}-{}-{}} & 1 & \texttt{a\textbar{}a\textbar{}a} & \texttt{a} & \texttt{a} & \texttt{a} & \texttt{012} & \texttt{0} & \texttt{+} & \texttt{S} & \texttt{none}\\
N143 & F3.98 & \texttt{T-{}4} & 1 & \texttt{a\textbar{}b\textbar{}c} & \texttt{c} & \texttt{c} & \texttt{abc} & \texttt{2} & \texttt{-{}} & \texttt{-{}0} & \texttt{s} & \texttt{a@0}\\
N144 & F3.98 & \texttt{T-{}4} & 1 & \texttt{a\textbar{}a\textbar{}b} & \texttt{b} & \texttt{b} & \texttt{ab} & \texttt{2} & \texttt{-{}} & \texttt{-{}0} & \texttt{s} & \texttt{a@0}\\
N145 & F3.98 & \texttt{T-{}4} & 1 & \texttt{a\textbar{}a\textbar{}a} & \texttt{a} & \texttt{a} & \texttt{a} & \texttt{012} & \texttt{0} & \texttt{+} & \texttt{S} & \texttt{none}\\
N146 & F3.100 & \texttt{T-{}-{}} & 1 & \texttt{a\textbar{}a\textbar{}b} & \texttt{ab} & \texttt{ab} & \texttt{ab} & \texttt{012} & \texttt{0} & \texttt{+} & \texttt{.} & \texttt{b@2}\\
N147 & F3.100 & \texttt{T-{}-{}} & 1 & \texttt{a\textbar{}a\textbar{}a} & \texttt{a} & \texttt{a} & \texttt{a} & \texttt{012} & \texttt{0} & \texttt{+} & \texttt{S} & \texttt{none}\\
N148 & F3.101 & \texttt{T-{}4} & 1 & \texttt{a\textbar{}b\textbar{}c} & \texttt{ac} & \texttt{ac} & \texttt{abc} & \texttt{02} & \texttt{0} & \texttt{-{}1} & \texttt{.} & \texttt{a@0}\\
N149 & F3.101 & \texttt{T-{}4} & 1 & \texttt{a\textbar{}a\textbar{}b} & \texttt{ab} & \texttt{ab} & \texttt{ab} & \texttt{012} & \texttt{0} & \texttt{+} & \texttt{.} & \texttt{b@2}\\
N150 & F3.101 & \texttt{T-{}4} & 1 & \texttt{a\textbar{}b\textbar{}a} & \texttt{a} & \texttt{a} & \texttt{ab} & \texttt{02} & \texttt{0} & \texttt{-{}1} & \texttt{s} & \texttt{b@1}\\
N151 & F3.101 & \texttt{T-{}4} & 1 & \texttt{a\textbar{}a\textbar{}a} & \texttt{a} & \texttt{a} & \texttt{a} & \texttt{012} & \texttt{0} & \texttt{+} & \texttt{S} & \texttt{none}\\
N152 & F3.103 & \texttt{T-{}-{}} & 1 & \texttt{a\textbar{}b\textbar{}b} & \texttt{ab} & \texttt{ab} & \texttt{ab} & \texttt{012} & \texttt{0} & \texttt{+} & \texttt{.} & \texttt{a@0}\\
N153 & F3.103 & \texttt{T-{}-{}} & 1 & \texttt{a\textbar{}a\textbar{}b} & \texttt{ab} & \texttt{ab} & \texttt{ab} & \texttt{012} & \texttt{0} & \texttt{+} & \texttt{.} & \texttt{b@2}\\
N154 & F3.103 & \texttt{T-{}-{}} & 1 & \texttt{a\textbar{}a\textbar{}a} & \texttt{a} & \texttt{a} & \texttt{a} & \texttt{012} & \texttt{0} & \texttt{+} & \texttt{S} & \texttt{none}\\
\multicolumn{13}{@{}l}{\rule{0pt}{2.8ex}\footnotesize\bfseries $n=3$ worlds, $m=2$ individuals: 318 models}\\*
N155 & F3.22 & \texttt{-{}-{}4} & 2 & \texttt{ab\textbar{}cd\textbar{}ef} & \texttt{a} & \texttt{ace} & \texttt{abcdef} & \texttt{012} & \texttt{0} & \texttt{+} & \texttt{S} & \texttt{c@1}\\
N156 & F3.22 & \texttt{-{}-{}4} & 2 & \texttt{ab\textbar{}cd\textbar{}cd} & \texttt{a} & \texttt{ac} & \texttt{abcd} & \texttt{012} & \texttt{0} & \texttt{+} & \texttt{S} & \texttt{c@1}\\
N157 & F3.22 & \texttt{-{}-{}4} & 2 & \texttt{aa\textbar{}bc\textbar{}de} & \texttt{a} & \texttt{abcde} & \texttt{abcde} & \texttt{012} & \texttt{01} & \texttt{+} & \texttt{S} & \texttt{b@1}\\
N158 & F3.22 & \texttt{-{}-{}4} & 2 & \texttt{aa\textbar{}bb\textbar{}cd} & \texttt{a} & \texttt{abcd} & \texttt{abcd} & \texttt{012} & \texttt{01} & \texttt{+} & \texttt{S} & \texttt{c@2}\\
N159 & F3.22 & \texttt{-{}-{}4} & 2 & \texttt{aa\textbar{}bb\textbar{}cc} & \texttt{a} & \texttt{abc} & \texttt{abc} & \texttt{012} & \texttt{01} & \texttt{+} & \texttt{S} & \texttt{b@1}\\
N160 & F3.22 & \texttt{-{}-{}4} & 2 & \texttt{aa\textbar{}bc\textbar{}bc} & \texttt{a} & \texttt{abc} & \texttt{abc} & \texttt{012} & \texttt{01} & \texttt{+} & \texttt{S} & \texttt{b@1}\\
N161 & F3.22 & \texttt{-{}-{}4} & 2 & \texttt{aa\textbar{}bc\textbar{}cb} & \texttt{a} & \texttt{abc} & \texttt{abc} & \texttt{012} & \texttt{01} & \texttt{+} & \texttt{S} & \texttt{c@1}\\
N162 & F3.22 & \texttt{-{}-{}4} & 2 & \texttt{aa\textbar{}bb\textbar{}bb} & \texttt{a} & \texttt{ab} & \texttt{ab} & \texttt{012} & \texttt{01} & \texttt{+} & \texttt{S} & \texttt{b@1}\\
N163 & F3.22 & \texttt{-{}-{}4} & 2 & \texttt{ab\textbar{}ab\textbar{}cd} & \texttt{a} & \texttt{ac} & \texttt{abcd} & \texttt{012} & \texttt{0} & \texttt{+} & \texttt{S} & \texttt{c@2}\\
N164 & F3.22 & \texttt{-{}-{}4} & 2 & \texttt{aa\textbar{}aa\textbar{}bc} & \texttt{a} & \texttt{abc} & \texttt{abc} & \texttt{012} & \texttt{01} & \texttt{+} & \texttt{S} & \texttt{b@2}\\
N165 & F3.22 & \texttt{-{}-{}4} & 2 & \texttt{aa\textbar{}aa\textbar{}bb} & \texttt{a} & \texttt{ab} & \texttt{ab} & \texttt{012} & \texttt{01} & \texttt{+} & \texttt{S} & \texttt{b@2}\\
N166 & F3.22 & \texttt{-{}-{}4} & 2 & \texttt{ab\textbar{}ab\textbar{}ab} & \texttt{a} & \texttt{a} & \texttt{ab} & \texttt{012} & \texttt{0} & \texttt{+} & \texttt{S} & \texttt{none}\\
N167 & F3.22 & \texttt{-{}-{}4} & 2 & \texttt{aa\textbar{}aa\textbar{}aa} & \texttt{a} & \texttt{a} & \texttt{a} & \texttt{012} & \texttt{01} & \texttt{+} & \texttt{S} & \texttt{none}\\
N168 & F3.23 & \texttt{-{}-{}-{}} & 2 & \texttt{ab\textbar{}ab\textbar{}cd} & \texttt{a} & \texttt{ac} & \texttt{abcd} & \texttt{012} & \texttt{0} & \texttt{+} & \texttt{S} & \texttt{c@2}\\
N169 & F3.23 & \texttt{-{}-{}-{}} & 2 & \texttt{aa\textbar{}aa\textbar{}bc} & \texttt{a} & \texttt{abc} & \texttt{abc} & \texttt{012} & \texttt{01} & \texttt{+} & \texttt{S} & \texttt{b@2}\\
N170 & F3.23 & \texttt{-{}-{}-{}} & 2 & \texttt{aa\textbar{}aa\textbar{}bb} & \texttt{a} & \texttt{ab} & \texttt{ab} & \texttt{012} & \texttt{01} & \texttt{+} & \texttt{S} & \texttt{b@2}\\
N171 & F3.23 & \texttt{-{}-{}-{}} & 2 & \texttt{ab\textbar{}ab\textbar{}ab} & \texttt{a} & \texttt{a} & \texttt{ab} & \texttt{012} & \texttt{0} & \texttt{+} & \texttt{S} & \texttt{none}\\
N172 & F3.23 & \texttt{-{}-{}-{}} & 2 & \texttt{aa\textbar{}aa\textbar{}aa} & \texttt{a} & \texttt{a} & \texttt{a} & \texttt{012} & \texttt{01} & \texttt{+} & \texttt{S} & \texttt{none}\\
N173 & F3.24 & \texttt{-{}-{}4} & 2 & \texttt{ab\textbar{}ab\textbar{}cd} & \texttt{a} & \texttt{ac} & \texttt{abcd} & \texttt{012} & \texttt{0} & \texttt{+} & \texttt{S} & \texttt{c@2}\\
N174 & F3.24 & \texttt{-{}-{}4} & 2 & \texttt{ab\textbar{}ba\textbar{}cd} & \texttt{b} & \texttt{bd} & \texttt{abcd} & \texttt{012} & \texttt{1} & \texttt{+} & \texttt{S} & \texttt{c@2}\\
N175 & F3.24 & \texttt{-{}-{}4} & 2 & \texttt{aa\textbar{}aa\textbar{}bc} & \texttt{a} & \texttt{abc} & \texttt{abc} & \texttt{012} & \texttt{01} & \texttt{+} & \texttt{S} & \texttt{b@2}\\
N176 & F3.24 & \texttt{-{}-{}4} & 2 & \texttt{aa\textbar{}aa\textbar{}bb} & \texttt{a} & \texttt{ab} & \texttt{ab} & \texttt{012} & \texttt{01} & \texttt{+} & \texttt{S} & \texttt{b@2}\\
N177 & F3.24 & \texttt{-{}-{}4} & 2 & \texttt{ab\textbar{}ab\textbar{}ab} & \texttt{a} & \texttt{a} & \texttt{ab} & \texttt{012} & \texttt{0} & \texttt{+} & \texttt{S} & \texttt{none}\\
N178 & F3.24 & \texttt{-{}-{}4} & 2 & \texttt{ab\textbar{}ba\textbar{}ab} & \texttt{a} & \texttt{a} & \texttt{ab} & \texttt{012} & \texttt{0} & \texttt{+} & \texttt{S} & \texttt{none}\\
N179 & F3.24 & \texttt{-{}-{}4} & 2 & \texttt{aa\textbar{}aa\textbar{}aa} & \texttt{a} & \texttt{a} & \texttt{a} & \texttt{012} & \texttt{01} & \texttt{+} & \texttt{S} & \texttt{none}\\
N180 & F3.25 & \texttt{-{}-{}-{}} & 2 & \texttt{ab\textbar{}ab\textbar{}cd} & \texttt{a} & \texttt{ac} & \texttt{abcd} & \texttt{012} & \texttt{0} & \texttt{+} & \texttt{S} & \texttt{c@2}\\
N181 & F3.25 & \texttt{-{}-{}-{}} & 2 & \texttt{aa\textbar{}aa\textbar{}bc} & \texttt{a} & \texttt{abc} & \texttt{abc} & \texttt{012} & \texttt{01} & \texttt{+} & \texttt{S} & \texttt{b@2}\\
N182 & F3.25 & \texttt{-{}-{}-{}} & 2 & \texttt{aa\textbar{}aa\textbar{}bb} & \texttt{a} & \texttt{ab} & \texttt{ab} & \texttt{012} & \texttt{01} & \texttt{+} & \texttt{S} & \texttt{b@2}\\
N183 & F3.25 & \texttt{-{}-{}-{}} & 2 & \texttt{ab\textbar{}ab\textbar{}ab} & \texttt{a} & \texttt{a} & \texttt{ab} & \texttt{012} & \texttt{0} & \texttt{+} & \texttt{S} & \texttt{none}\\
N184 & F3.25 & \texttt{-{}-{}-{}} & 2 & \texttt{aa\textbar{}aa\textbar{}aa} & \texttt{a} & \texttt{a} & \texttt{a} & \texttt{012} & \texttt{01} & \texttt{+} & \texttt{S} & \texttt{none}\\
N185 & F3.27 & \texttt{-{}-{}-{}} & 2 & \texttt{ab\textbar{}ab\textbar{}ab} & \texttt{a} & \texttt{a} & \texttt{ab} & \texttt{012} & \texttt{0} & \texttt{+} & \texttt{S} & \texttt{none}\\
N186 & F3.27 & \texttt{-{}-{}-{}} & 2 & \texttt{aa\textbar{}aa\textbar{}aa} & \texttt{a} & \texttt{a} & \texttt{a} & \texttt{012} & \texttt{01} & \texttt{+} & \texttt{S} & \texttt{none}\\
N187 & F3.38 & \texttt{-{}-{}-{}} & 2 & \texttt{ab\textbar{}ab\textbar{}cd} & \texttt{a} & \texttt{ac} & \texttt{abcd} & \texttt{012} & \texttt{0} & \texttt{+} & \texttt{S} & \texttt{c@2}\\
N188 & F3.38 & \texttt{-{}-{}-{}} & 2 & \texttt{aa\textbar{}aa\textbar{}bc} & \texttt{a} & \texttt{abc} & \texttt{abc} & \texttt{012} & \texttt{01} & \texttt{+} & \texttt{S} & \texttt{b@2}\\
N189 & F3.38 & \texttt{-{}-{}-{}} & 2 & \texttt{aa\textbar{}aa\textbar{}bb} & \texttt{a} & \texttt{ab} & \texttt{ab} & \texttt{012} & \texttt{01} & \texttt{+} & \texttt{S} & \texttt{b@2}\\
N190 & F3.38 & \texttt{-{}-{}-{}} & 2 & \texttt{ab\textbar{}ab\textbar{}ab} & \texttt{a} & \texttt{a} & \texttt{ab} & \texttt{012} & \texttt{0} & \texttt{+} & \texttt{S} & \texttt{none}\\
N191 & F3.38 & \texttt{-{}-{}-{}} & 2 & \texttt{aa\textbar{}aa\textbar{}aa} & \texttt{a} & \texttt{a} & \texttt{a} & \texttt{012} & \texttt{01} & \texttt{+} & \texttt{S} & \texttt{none}\\
N192 & F3.40 & \texttt{-{}-{}-{}} & 2 & \texttt{ab\textbar{}ab\textbar{}cd} & \texttt{a} & \texttt{ac} & \texttt{abcd} & \texttt{012} & \texttt{0} & \texttt{+} & \texttt{S} & \texttt{c@2}\\
N193 & F3.40 & \texttt{-{}-{}-{}} & 2 & \texttt{aa\textbar{}aa\textbar{}bc} & \texttt{a} & \texttt{abc} & \texttt{abc} & \texttt{012} & \texttt{01} & \texttt{+} & \texttt{S} & \texttt{b@2}\\
N194 & F3.40 & \texttt{-{}-{}-{}} & 2 & \texttt{aa\textbar{}aa\textbar{}bb} & \texttt{a} & \texttt{ab} & \texttt{ab} & \texttt{012} & \texttt{01} & \texttt{+} & \texttt{S} & \texttt{b@2}\\
N195 & F3.40 & \texttt{-{}-{}-{}} & 2 & \texttt{ab\textbar{}ab\textbar{}ab} & \texttt{a} & \texttt{a} & \texttt{ab} & \texttt{012} & \texttt{0} & \texttt{+} & \texttt{S} & \texttt{none}\\
N196 & F3.40 & \texttt{-{}-{}-{}} & 2 & \texttt{aa\textbar{}aa\textbar{}aa} & \texttt{a} & \texttt{a} & \texttt{a} & \texttt{012} & \texttt{01} & \texttt{+} & \texttt{S} & \texttt{none}\\
N197 & F3.42 & \texttt{-{}-{}-{}} & 2 & \texttt{ab\textbar{}ab\textbar{}ab} & \texttt{a} & \texttt{a} & \texttt{ab} & \texttt{012} & \texttt{0} & \texttt{+} & \texttt{S} & \texttt{none}\\
N198 & F3.42 & \texttt{-{}-{}-{}} & 2 & \texttt{aa\textbar{}aa\textbar{}aa} & \texttt{a} & \texttt{a} & \texttt{a} & \texttt{012} & \texttt{01} & \texttt{+} & \texttt{S} & \texttt{none}\\
N199 & F3.43 & \texttt{-{}-{}4} & 2 & \texttt{ab\textbar{}cd\textbar{}ef} & \texttt{a} & \texttt{ae} & \texttt{abcdef} & \texttt{02} & \texttt{0} & \texttt{-{}1} & \texttt{s} & \texttt{c@1}\\
N200 & F3.43 & \texttt{-{}-{}4} & 2 & \texttt{aa\textbar{}bc\textbar{}de} & \texttt{a} & \texttt{ade} & \texttt{abcde} & \texttt{02} & \texttt{01} & \texttt{-{}1} & \texttt{s} & \texttt{b@1}\\
N201 & F3.43 & \texttt{-{}-{}4} & 2 & \texttt{aa\textbar{}bc\textbar{}dd} & \texttt{a} & \texttt{ad} & \texttt{abcd} & \texttt{02} & \texttt{01} & \texttt{-{}1} & \texttt{s} & \texttt{b@1}\\
N202 & F3.43 & \texttt{-{}-{}4} & 2 & \texttt{aa\textbar{}bb\textbar{}cd} & \texttt{a} & \texttt{acd} & \texttt{abcd} & \texttt{02} & \texttt{01} & \texttt{-{}1} & \texttt{s} & \texttt{c@2}\\
N203 & F3.43 & \texttt{-{}-{}4} & 2 & \texttt{aa\textbar{}bb\textbar{}cc} & \texttt{a} & \texttt{ac} & \texttt{abc} & \texttt{02} & \texttt{01} & \texttt{-{}1} & \texttt{s} & \texttt{b@1}\\
N204 & F3.43 & \texttt{-{}-{}4} & 2 & \texttt{ab\textbar{}cd\textbar{}ab} & \texttt{a} & \texttt{a} & \texttt{abcd} & \texttt{02} & \texttt{0} & \texttt{-{}1} & \texttt{s} & \texttt{c@1}\\
N205 & F3.43 & \texttt{-{}-{}4} & 2 & \texttt{aa\textbar{}bc\textbar{}aa} & \texttt{a} & \texttt{a} & \texttt{abc} & \texttt{02} & \texttt{01} & \texttt{-{}1} & \texttt{s} & \texttt{b@1}\\
N206 & F3.43 & \texttt{-{}-{}4} & 2 & \texttt{ab\textbar{}ab\textbar{}cd} & \texttt{a} & \texttt{ac} & \texttt{abcd} & \texttt{012} & \texttt{0} & \texttt{+} & \texttt{S} & \texttt{c@2}\\
N207 & F3.43 & \texttt{-{}-{}4} & 2 & \texttt{aa\textbar{}bb\textbar{}aa} & \texttt{a} & \texttt{a} & \texttt{ab} & \texttt{02} & \texttt{01} & \texttt{-{}1} & \texttt{s} & \texttt{b@1}\\
N208 & F3.43 & \texttt{-{}-{}4} & 2 & \texttt{aa\textbar{}aa\textbar{}bc} & \texttt{a} & \texttt{abc} & \texttt{abc} & \texttt{012} & \texttt{01} & \texttt{+} & \texttt{S} & \texttt{b@2}\\
N209 & F3.43 & \texttt{-{}-{}4} & 2 & \texttt{aa\textbar{}aa\textbar{}bb} & \texttt{a} & \texttt{ab} & \texttt{ab} & \texttt{012} & \texttt{01} & \texttt{+} & \texttt{S} & \texttt{b@2}\\
N210 & F3.43 & \texttt{-{}-{}4} & 2 & \texttt{ab\textbar{}ab\textbar{}ab} & \texttt{a} & \texttt{a} & \texttt{ab} & \texttt{012} & \texttt{0} & \texttt{+} & \texttt{S} & \texttt{none}\\
N211 & F3.43 & \texttt{-{}-{}4} & 2 & \texttt{aa\textbar{}aa\textbar{}aa} & \texttt{a} & \texttt{a} & \texttt{a} & \texttt{012} & \texttt{01} & \texttt{+} & \texttt{S} & \texttt{none}\\
N212 & F3.44 & \texttt{-{}-{}-{}} & 2 & \texttt{ab\textbar{}ab\textbar{}cd} & \texttt{a} & \texttt{ac} & \texttt{abcd} & \texttt{012} & \texttt{0} & \texttt{+} & \texttt{S} & \texttt{c@2}\\
N213 & F3.44 & \texttt{-{}-{}-{}} & 2 & \texttt{aa\textbar{}aa\textbar{}bc} & \texttt{a} & \texttt{abc} & \texttt{abc} & \texttt{012} & \texttt{01} & \texttt{+} & \texttt{S} & \texttt{b@2}\\
N214 & F3.44 & \texttt{-{}-{}-{}} & 2 & \texttt{aa\textbar{}aa\textbar{}bb} & \texttt{a} & \texttt{ab} & \texttt{ab} & \texttt{012} & \texttt{01} & \texttt{+} & \texttt{S} & \texttt{b@2}\\
N215 & F3.44 & \texttt{-{}-{}-{}} & 2 & \texttt{ab\textbar{}ab\textbar{}ab} & \texttt{a} & \texttt{a} & \texttt{ab} & \texttt{012} & \texttt{0} & \texttt{+} & \texttt{S} & \texttt{none}\\
N216 & F3.44 & \texttt{-{}-{}-{}} & 2 & \texttt{aa\textbar{}aa\textbar{}aa} & \texttt{a} & \texttt{a} & \texttt{a} & \texttt{012} & \texttt{01} & \texttt{+} & \texttt{S} & \texttt{none}\\
N217 & F3.45 & \texttt{-{}-{}-{}} & 2 & \texttt{ab\textbar{}cd\textbar{}ab} & \texttt{a} & \texttt{a} & \texttt{abcd} & \texttt{02} & \texttt{0} & \texttt{-{}1} & \texttt{s} & \texttt{c@1}\\
N218 & F3.45 & \texttt{-{}-{}-{}} & 2 & \texttt{aa\textbar{}bc\textbar{}aa} & \texttt{a} & \texttt{a} & \texttt{abc} & \texttt{02} & \texttt{01} & \texttt{-{}1} & \texttt{s} & \texttt{b@1}\\
N219 & F3.45 & \texttt{-{}-{}-{}} & 2 & \texttt{aa\textbar{}bb\textbar{}aa} & \texttt{a} & \texttt{a} & \texttt{ab} & \texttt{02} & \texttt{01} & \texttt{-{}1} & \texttt{s} & \texttt{b@1}\\
N220 & F3.45 & \texttt{-{}-{}-{}} & 2 & \texttt{ab\textbar{}ab\textbar{}ab} & \texttt{a} & \texttt{a} & \texttt{ab} & \texttt{012} & \texttt{0} & \texttt{+} & \texttt{S} & \texttt{none}\\
N221 & F3.45 & \texttt{-{}-{}-{}} & 2 & \texttt{aa\textbar{}aa\textbar{}aa} & \texttt{a} & \texttt{a} & \texttt{a} & \texttt{012} & \texttt{01} & \texttt{+} & \texttt{S} & \texttt{none}\\
N222 & F3.46 & \texttt{-{}-{}-{}} & 2 & \texttt{ab\textbar{}ab\textbar{}ab} & \texttt{a} & \texttt{a} & \texttt{ab} & \texttt{012} & \texttt{0} & \texttt{+} & \texttt{S} & \texttt{none}\\
N223 & F3.46 & \texttt{-{}-{}-{}} & 2 & \texttt{aa\textbar{}aa\textbar{}aa} & \texttt{a} & \texttt{a} & \texttt{a} & \texttt{012} & \texttt{01} & \texttt{+} & \texttt{S} & \texttt{none}\\
N224 & F3.47 & \texttt{-{}-{}-{}} & 2 & \texttt{ab\textbar{}ab\textbar{}ab} & \texttt{a} & \texttt{a} & \texttt{ab} & \texttt{012} & \texttt{0} & \texttt{+} & \texttt{S} & \texttt{none}\\
N225 & F3.47 & \texttt{-{}-{}-{}} & 2 & \texttt{aa\textbar{}aa\textbar{}aa} & \texttt{a} & \texttt{a} & \texttt{a} & \texttt{012} & \texttt{01} & \texttt{+} & \texttt{S} & \texttt{none}\\
N226 & F3.48 & \texttt{-{}-{}4} & 2 & \texttt{ab\textbar{}cd\textbar{}ef} & \texttt{a} & \texttt{ae} & \texttt{abcdef} & \texttt{02} & \texttt{0} & \texttt{+} & \texttt{s} & \texttt{c@1}\\
N227 & F3.48 & \texttt{-{}-{}4} & 2 & \texttt{aa\textbar{}bc\textbar{}de} & \texttt{a} & \texttt{ade} & \texttt{abcde} & \texttt{02} & \texttt{01} & \texttt{+} & \texttt{s} & \texttt{b@1}\\
N228 & F3.48 & \texttt{-{}-{}4} & 2 & \texttt{aa\textbar{}bc\textbar{}dd} & \texttt{a} & \texttt{ad} & \texttt{abcd} & \texttt{02} & \texttt{01} & \texttt{+} & \texttt{s} & \texttt{b@1}\\
N229 & F3.48 & \texttt{-{}-{}4} & 2 & \texttt{aa\textbar{}bb\textbar{}cc} & \texttt{a} & \texttt{ac} & \texttt{abc} & \texttt{02} & \texttt{01} & \texttt{+} & \texttt{s} & \texttt{b@1}\\
N230 & F3.48 & \texttt{-{}-{}4} & 2 & \texttt{ab\textbar{}cd\textbar{}ab} & \texttt{a} & \texttt{ac} & \texttt{abcd} & \texttt{012} & \texttt{0} & \texttt{+} & \texttt{S} & \texttt{c@1}\\
N231 & F3.48 & \texttt{-{}-{}4} & 2 & \texttt{aa\textbar{}bc\textbar{}aa} & \texttt{a} & \texttt{abc} & \texttt{abc} & \texttt{012} & \texttt{01} & \texttt{+} & \texttt{S} & \texttt{b@1}\\
N232 & F3.48 & \texttt{-{}-{}4} & 2 & \texttt{aa\textbar{}bb\textbar{}aa} & \texttt{a} & \texttt{ab} & \texttt{ab} & \texttt{012} & \texttt{01} & \texttt{+} & \texttt{S} & \texttt{b@1}\\
N233 & F3.48 & \texttt{-{}-{}4} & 2 & \texttt{ab\textbar{}ab\textbar{}ab} & \texttt{a} & \texttt{a} & \texttt{ab} & \texttt{012} & \texttt{0} & \texttt{+} & \texttt{S} & \texttt{none}\\
N234 & F3.48 & \texttt{-{}-{}4} & 2 & \texttt{aa\textbar{}aa\textbar{}aa} & \texttt{a} & \texttt{a} & \texttt{a} & \texttt{012} & \texttt{01} & \texttt{+} & \texttt{S} & \texttt{none}\\
N235 & F3.49 & \texttt{-{}-{}-{}} & 2 & \texttt{ab\textbar{}cd\textbar{}ab} & \texttt{a} & \texttt{ac} & \texttt{abcd} & \texttt{012} & \texttt{0} & \texttt{+} & \texttt{S} & \texttt{c@1}\\
N236 & F3.49 & \texttt{-{}-{}-{}} & 2 & \texttt{aa\textbar{}bc\textbar{}aa} & \texttt{a} & \texttt{abc} & \texttt{abc} & \texttt{012} & \texttt{01} & \texttt{+} & \texttt{S} & \texttt{b@1}\\
N237 & F3.49 & \texttt{-{}-{}-{}} & 2 & \texttt{aa\textbar{}bb\textbar{}aa} & \texttt{a} & \texttt{ab} & \texttt{ab} & \texttt{012} & \texttt{01} & \texttt{+} & \texttt{S} & \texttt{b@1}\\
N238 & F3.49 & \texttt{-{}-{}-{}} & 2 & \texttt{ab\textbar{}ab\textbar{}ab} & \texttt{a} & \texttt{a} & \texttt{ab} & \texttt{012} & \texttt{0} & \texttt{+} & \texttt{S} & \texttt{none}\\
N239 & F3.49 & \texttt{-{}-{}-{}} & 2 & \texttt{aa\textbar{}aa\textbar{}aa} & \texttt{a} & \texttt{a} & \texttt{a} & \texttt{012} & \texttt{01} & \texttt{+} & \texttt{S} & \texttt{none}\\
N240 & F3.50 & \texttt{-{}-{}-{}} & 2 & \texttt{ab\textbar{}cd\textbar{}ab} & \texttt{a} & \texttt{a} & \texttt{abcd} & \texttt{02} & \texttt{0} & \texttt{-{}1} & \texttt{s} & \texttt{c@1}\\
N241 & F3.50 & \texttt{-{}-{}-{}} & 2 & \texttt{aa\textbar{}bc\textbar{}aa} & \texttt{a} & \texttt{a} & \texttt{abc} & \texttt{02} & \texttt{01} & \texttt{-{}1} & \texttt{s} & \texttt{b@1}\\
N242 & F3.50 & \texttt{-{}-{}-{}} & 2 & \texttt{aa\textbar{}bb\textbar{}aa} & \texttt{a} & \texttt{a} & \texttt{ab} & \texttt{02} & \texttt{01} & \texttt{-{}1} & \texttt{s} & \texttt{b@1}\\
N243 & F3.50 & \texttt{-{}-{}-{}} & 2 & \texttt{ab\textbar{}ab\textbar{}ab} & \texttt{a} & \texttt{a} & \texttt{ab} & \texttt{012} & \texttt{0} & \texttt{+} & \texttt{S} & \texttt{none}\\
N244 & F3.50 & \texttt{-{}-{}-{}} & 2 & \texttt{aa\textbar{}aa\textbar{}aa} & \texttt{a} & \texttt{a} & \texttt{a} & \texttt{012} & \texttt{01} & \texttt{+} & \texttt{S} & \texttt{none}\\
N245 & F3.51 & \texttt{-{}-{}4} & 2 & \texttt{ab\textbar{}ab\textbar{}cd} & \texttt{a} & \texttt{ac} & \texttt{abcd} & \texttt{012} & \texttt{0} & \texttt{+} & \texttt{S} & \texttt{c@2}\\
N246 & F3.51 & \texttt{-{}-{}4} & 2 & \texttt{ab\textbar{}ba\textbar{}cd} & \texttt{b} & \texttt{b} & \texttt{abcd} & \texttt{01} & \texttt{1} & \texttt{+} & \texttt{S} & \texttt{c@2}\\
N247 & F3.51 & \texttt{-{}-{}4} & 2 & \texttt{ab\textbar{}ba\textbar{}cc} & \texttt{b} & \texttt{b} & \texttt{abc} & \texttt{01} & \texttt{1} & \texttt{+} & \texttt{S} & \texttt{c@2}\\
N248 & F3.51 & \texttt{-{}-{}4} & 2 & \texttt{aa\textbar{}aa\textbar{}bc} & \texttt{a} & \texttt{abc} & \texttt{abc} & \texttt{012} & \texttt{01} & \texttt{+} & \texttt{S} & \texttt{b@2}\\
N249 & F3.51 & \texttt{-{}-{}4} & 2 & \texttt{aa\textbar{}aa\textbar{}bb} & \texttt{a} & \texttt{ab} & \texttt{ab} & \texttt{012} & \texttt{01} & \texttt{+} & \texttt{S} & \texttt{b@2}\\
N250 & F3.51 & \texttt{-{}-{}4} & 2 & \texttt{ab\textbar{}ab\textbar{}ab} & \texttt{a} & \texttt{a} & \texttt{ab} & \texttt{012} & \texttt{0} & \texttt{+} & \texttt{S} & \texttt{none}\\
N251 & F3.51 & \texttt{-{}-{}4} & 2 & \texttt{aa\textbar{}aa\textbar{}aa} & \texttt{a} & \texttt{a} & \texttt{a} & \texttt{012} & \texttt{01} & \texttt{+} & \texttt{S} & \texttt{none}\\
N252 & F3.52 & \texttt{T-{}4} & 2 & \texttt{ab\textbar{}cd\textbar{}cd} & \texttt{c} & \texttt{c} & \texttt{abcd} & \texttt{12} & \texttt{-{}} & \texttt{-{}0} & \texttt{s} & \texttt{a@0}\\
N253 & F3.52 & \texttt{T-{}4} & 2 & \texttt{ab\textbar{}cd\textbar{}dc} & \texttt{d} & \texttt{d} & \texttt{abcd} & \texttt{12} & \texttt{-{}} & \texttt{-{}0} & \texttt{s} & \texttt{a@0}\\
N254 & F3.52 & \texttt{T-{}4} & 2 & \texttt{ab\textbar{}cc\textbar{}cc} & \texttt{c} & \texttt{c} & \texttt{abc} & \texttt{12} & \texttt{-{}} & \texttt{-{}0} & \texttt{s} & \texttt{a@0}\\
N255 & F3.52 & \texttt{T-{}4} & 2 & \texttt{aa\textbar{}bb\textbar{}bb} & \texttt{b} & \texttt{b} & \texttt{ab} & \texttt{12} & \texttt{-{}} & \texttt{-{}0} & \texttt{s} & \texttt{a@0}\\
N256 & F3.52 & \texttt{T-{}4} & 2 & \texttt{ab\textbar{}ab\textbar{}ab} & \texttt{a} & \texttt{a} & \texttt{ab} & \texttt{012} & \texttt{0} & \texttt{+} & \texttt{S} & \texttt{none}\\
N257 & F3.52 & \texttt{T-{}4} & 2 & \texttt{ab\textbar{}ab\textbar{}ba} & \texttt{b} & \texttt{b} & \texttt{ab} & \texttt{012} & \texttt{1} & \texttt{+} & \texttt{S} & \texttt{none}\\
N258 & F3.52 & \texttt{T-{}4} & 2 & \texttt{aa\textbar{}aa\textbar{}aa} & \texttt{a} & \texttt{a} & \texttt{a} & \texttt{012} & \texttt{01} & \texttt{+} & \texttt{S} & \texttt{none}\\
N259 & F3.57 & \texttt{-{}-{}-{}} & 2 & \texttt{ab\textbar{}ab\textbar{}ab} & \texttt{a} & \texttt{a} & \texttt{ab} & \texttt{012} & \texttt{0} & \texttt{+} & \texttt{S} & \texttt{none}\\
N260 & F3.57 & \texttt{-{}-{}-{}} & 2 & \texttt{aa\textbar{}aa\textbar{}aa} & \texttt{a} & \texttt{a} & \texttt{a} & \texttt{012} & \texttt{01} & \texttt{+} & \texttt{S} & \texttt{none}\\
N261 & F3.58 & \texttt{-{}-{}-{}} & 2 & \texttt{ab\textbar{}ab\textbar{}cd} & \texttt{a} & \texttt{ac} & \texttt{abcd} & \texttt{012} & \texttt{0} & \texttt{+} & \texttt{S} & \texttt{c@2}\\
N262 & F3.58 & \texttt{-{}-{}-{}} & 2 & \texttt{ab\textbar{}ba\textbar{}cd} & \texttt{ab} & \texttt{ab} & \texttt{abcd} & \texttt{01} & \texttt{01} & \texttt{+} & \texttt{.} & \texttt{c@2}\\
N263 & F3.58 & \texttt{-{}-{}-{}} & 2 & \texttt{aa\textbar{}aa\textbar{}bc} & \texttt{a} & \texttt{abc} & \texttt{abc} & \texttt{012} & \texttt{01} & \texttt{+} & \texttt{S} & \texttt{b@2}\\
N264 & F3.58 & \texttt{-{}-{}-{}} & 2 & \texttt{aa\textbar{}aa\textbar{}bb} & \texttt{a} & \texttt{ab} & \texttt{ab} & \texttt{012} & \texttt{01} & \texttt{+} & \texttt{S} & \texttt{b@2}\\
N265 & F3.58 & \texttt{-{}-{}-{}} & 2 & \texttt{ab\textbar{}ab\textbar{}ab} & \texttt{a} & \texttt{a} & \texttt{ab} & \texttt{012} & \texttt{0} & \texttt{+} & \texttt{S} & \texttt{none}\\
N266 & F3.58 & \texttt{-{}-{}-{}} & 2 & \texttt{aa\textbar{}aa\textbar{}aa} & \texttt{a} & \texttt{a} & \texttt{a} & \texttt{012} & \texttt{01} & \texttt{+} & \texttt{S} & \texttt{none}\\
N267 & F3.59 & \texttt{-{}-{}-{}} & 2 & \texttt{ab\textbar{}cd\textbar{}ab} & \texttt{a} & \texttt{a} & \texttt{abcd} & \texttt{02} & \texttt{0} & \texttt{-{}1} & \texttt{s} & \texttt{c@1}\\
N268 & F3.59 & \texttt{-{}-{}-{}} & 2 & \texttt{aa\textbar{}bc\textbar{}aa} & \texttt{a} & \texttt{a} & \texttt{abc} & \texttt{02} & \texttt{01} & \texttt{-{}1} & \texttt{s} & \texttt{b@1}\\
N269 & F3.59 & \texttt{-{}-{}-{}} & 2 & \texttt{aa\textbar{}bb\textbar{}aa} & \texttt{a} & \texttt{a} & \texttt{ab} & \texttt{02} & \texttt{01} & \texttt{-{}1} & \texttt{s} & \texttt{b@1}\\
N270 & F3.59 & \texttt{-{}-{}-{}} & 2 & \texttt{ab\textbar{}ab\textbar{}ab} & \texttt{a} & \texttt{a} & \texttt{ab} & \texttt{012} & \texttt{0} & \texttt{+} & \texttt{S} & \texttt{none}\\
N271 & F3.59 & \texttt{-{}-{}-{}} & 2 & \texttt{aa\textbar{}aa\textbar{}aa} & \texttt{a} & \texttt{a} & \texttt{a} & \texttt{012} & \texttt{01} & \texttt{+} & \texttt{S} & \texttt{none}\\
N272 & F3.61 & \texttt{-{}-{}-{}} & 2 & \texttt{ab\textbar{}ab\textbar{}ab} & \texttt{a} & \texttt{a} & \texttt{ab} & \texttt{012} & \texttt{0} & \texttt{+} & \texttt{S} & \texttt{none}\\
N273 & F3.61 & \texttt{-{}-{}-{}} & 2 & \texttt{aa\textbar{}aa\textbar{}aa} & \texttt{a} & \texttt{a} & \texttt{a} & \texttt{012} & \texttt{01} & \texttt{+} & \texttt{S} & \texttt{none}\\
N274 & F3.62 & \texttt{-{}-{}-{}} & 2 & \texttt{ab\textbar{}ab\textbar{}ab} & \texttt{a} & \texttt{a} & \texttt{ab} & \texttt{012} & \texttt{0} & \texttt{+} & \texttt{S} & \texttt{none}\\
N275 & F3.62 & \texttt{-{}-{}-{}} & 2 & \texttt{aa\textbar{}aa\textbar{}aa} & \texttt{a} & \texttt{a} & \texttt{a} & \texttt{012} & \texttt{01} & \texttt{+} & \texttt{S} & \texttt{none}\\
N276 & F3.63 & \texttt{-{}-{}-{}} & 2 & \texttt{ab\textbar{}cd\textbar{}ab} & \texttt{a} & \texttt{ac} & \texttt{abcd} & \texttt{012} & \texttt{0} & \texttt{+} & \texttt{S} & \texttt{c@1}\\
N277 & F3.63 & \texttt{-{}-{}-{}} & 2 & \texttt{aa\textbar{}bc\textbar{}aa} & \texttt{a} & \texttt{abc} & \texttt{abc} & \texttt{012} & \texttt{01} & \texttt{+} & \texttt{S} & \texttt{b@1}\\
N278 & F3.63 & \texttt{-{}-{}-{}} & 2 & \texttt{aa\textbar{}bb\textbar{}aa} & \texttt{a} & \texttt{ab} & \texttt{ab} & \texttt{012} & \texttt{01} & \texttt{+} & \texttt{S} & \texttt{b@1}\\
N279 & F3.63 & \texttt{-{}-{}-{}} & 2 & \texttt{ab\textbar{}ab\textbar{}ab} & \texttt{a} & \texttt{a} & \texttt{ab} & \texttt{012} & \texttt{0} & \texttt{+} & \texttt{S} & \texttt{none}\\
N280 & F3.63 & \texttt{-{}-{}-{}} & 2 & \texttt{aa\textbar{}aa\textbar{}aa} & \texttt{a} & \texttt{a} & \texttt{a} & \texttt{012} & \texttt{01} & \texttt{+} & \texttt{S} & \texttt{none}\\
N281 & F3.64 & \texttt{-{}-{}-{}} & 2 & \texttt{ab\textbar{}ab\textbar{}ab} & \texttt{a} & \texttt{a} & \texttt{ab} & \texttt{012} & \texttt{0} & \texttt{+} & \texttt{S} & \texttt{none}\\
N282 & F3.64 & \texttt{-{}-{}-{}} & 2 & \texttt{aa\textbar{}aa\textbar{}aa} & \texttt{a} & \texttt{a} & \texttt{a} & \texttt{012} & \texttt{01} & \texttt{+} & \texttt{S} & \texttt{none}\\
N283 & F3.65 & \texttt{-{}-{}-{}} & 2 & \texttt{ab\textbar{}ab\textbar{}ab} & \texttt{a} & \texttt{a} & \texttt{ab} & \texttt{012} & \texttt{0} & \texttt{+} & \texttt{S} & \texttt{none}\\
N284 & F3.65 & \texttt{-{}-{}-{}} & 2 & \texttt{aa\textbar{}aa\textbar{}aa} & \texttt{a} & \texttt{a} & \texttt{a} & \texttt{012} & \texttt{01} & \texttt{+} & \texttt{S} & \texttt{none}\\
N285 & F3.66 & \texttt{-{}-{}-{}} & 2 & \texttt{ab\textbar{}cd\textbar{}ab} & \texttt{a} & \texttt{a} & \texttt{abcd} & \texttt{02} & \texttt{0} & \texttt{-{}1} & \texttt{s} & \texttt{c@1}\\
N286 & F3.66 & \texttt{-{}-{}-{}} & 2 & \texttt{aa\textbar{}bc\textbar{}aa} & \texttt{a} & \texttt{a} & \texttt{abc} & \texttt{02} & \texttt{01} & \texttt{-{}1} & \texttt{s} & \texttt{b@1}\\
N287 & F3.66 & \texttt{-{}-{}-{}} & 2 & \texttt{aa\textbar{}bb\textbar{}aa} & \texttt{a} & \texttt{a} & \texttt{ab} & \texttt{02} & \texttt{01} & \texttt{-{}1} & \texttt{s} & \texttt{b@1}\\
N288 & F3.66 & \texttt{-{}-{}-{}} & 2 & \texttt{ab\textbar{}ab\textbar{}ab} & \texttt{a} & \texttt{a} & \texttt{ab} & \texttt{012} & \texttt{0} & \texttt{+} & \texttt{S} & \texttt{none}\\
N289 & F3.66 & \texttt{-{}-{}-{}} & 2 & \texttt{aa\textbar{}aa\textbar{}aa} & \texttt{a} & \texttt{a} & \texttt{a} & \texttt{012} & \texttt{01} & \texttt{+} & \texttt{S} & \texttt{none}\\
N290 & F3.67 & \texttt{-{}-{}-{}} & 2 & \texttt{ab\textbar{}ab\textbar{}ab} & \texttt{a} & \texttt{a} & \texttt{ab} & \texttt{012} & \texttt{0} & \texttt{+} & \texttt{S} & \texttt{none}\\
N291 & F3.67 & \texttt{-{}-{}-{}} & 2 & \texttt{aa\textbar{}aa\textbar{}aa} & \texttt{a} & \texttt{a} & \texttt{a} & \texttt{012} & \texttt{01} & \texttt{+} & \texttt{S} & \texttt{none}\\
N292 & F3.68 & \texttt{-{}-{}4} & 2 & \texttt{ab\textbar{}cd\textbar{}ef} & \texttt{a} & \texttt{ae} & \texttt{abcdef} & \texttt{02} & \texttt{0} & \texttt{-{}1} & \texttt{s} & \texttt{c@1}\\
N293 & F3.68 & \texttt{-{}-{}4} & 2 & \texttt{aa\textbar{}bc\textbar{}de} & \texttt{a} & \texttt{ade} & \texttt{abcde} & \texttt{02} & \texttt{01} & \texttt{-{}1} & \texttt{s} & \texttt{b@1}\\
N294 & F3.68 & \texttt{-{}-{}4} & 2 & \texttt{aa\textbar{}bc\textbar{}dd} & \texttt{a} & \texttt{ad} & \texttt{abcd} & \texttt{02} & \texttt{01} & \texttt{-{}1} & \texttt{s} & \texttt{b@1}\\
N295 & F3.68 & \texttt{-{}-{}4} & 2 & \texttt{aa\textbar{}bb\textbar{}cc} & \texttt{a} & \texttt{ac} & \texttt{abc} & \texttt{02} & \texttt{01} & \texttt{-{}1} & \texttt{s} & \texttt{b@1}\\
N296 & F3.68 & \texttt{-{}-{}4} & 2 & \texttt{ab\textbar{}cd\textbar{}ab} & \texttt{a} & \texttt{a} & \texttt{abcd} & \texttt{02} & \texttt{0} & \texttt{-{}1} & \texttt{s} & \texttt{c@1}\\
N297 & F3.68 & \texttt{-{}-{}4} & 2 & \texttt{aa\textbar{}bc\textbar{}aa} & \texttt{a} & \texttt{a} & \texttt{abc} & \texttt{02} & \texttt{01} & \texttt{-{}1} & \texttt{s} & \texttt{b@1}\\
N298 & F3.68 & \texttt{-{}-{}4} & 2 & \texttt{aa\textbar{}bb\textbar{}aa} & \texttt{a} & \texttt{a} & \texttt{ab} & \texttt{02} & \texttt{01} & \texttt{-{}1} & \texttt{s} & \texttt{b@1}\\
N299 & F3.68 & \texttt{-{}-{}4} & 2 & \texttt{ab\textbar{}ab\textbar{}ab} & \texttt{a} & \texttt{a} & \texttt{ab} & \texttt{012} & \texttt{0} & \texttt{+} & \texttt{S} & \texttt{none}\\
N300 & F3.68 & \texttt{-{}-{}4} & 2 & \texttt{aa\textbar{}aa\textbar{}aa} & \texttt{a} & \texttt{a} & \texttt{a} & \texttt{012} & \texttt{01} & \texttt{+} & \texttt{S} & \texttt{none}\\
N301 & F3.69 & \texttt{-{}-{}-{}} & 2 & \texttt{ab\textbar{}cd\textbar{}ab} & \texttt{a} & \texttt{a} & \texttt{abcd} & \texttt{02} & \texttt{0} & \texttt{-{}1} & \texttt{s} & \texttt{c@1}\\
N302 & F3.69 & \texttt{-{}-{}-{}} & 2 & \texttt{aa\textbar{}bc\textbar{}aa} & \texttt{a} & \texttt{a} & \texttt{abc} & \texttt{02} & \texttt{01} & \texttt{-{}1} & \texttt{s} & \texttt{b@1}\\
N303 & F3.69 & \texttt{-{}-{}-{}} & 2 & \texttt{aa\textbar{}bb\textbar{}aa} & \texttt{a} & \texttt{a} & \texttt{ab} & \texttt{02} & \texttt{01} & \texttt{-{}1} & \texttt{s} & \texttt{b@1}\\
N304 & F3.69 & \texttt{-{}-{}-{}} & 2 & \texttt{ab\textbar{}ab\textbar{}ab} & \texttt{a} & \texttt{a} & \texttt{ab} & \texttt{012} & \texttt{0} & \texttt{+} & \texttt{S} & \texttt{none}\\
N305 & F3.69 & \texttt{-{}-{}-{}} & 2 & \texttt{aa\textbar{}aa\textbar{}aa} & \texttt{a} & \texttt{a} & \texttt{a} & \texttt{012} & \texttt{01} & \texttt{+} & \texttt{S} & \texttt{none}\\
N306 & F3.70 & \texttt{-{}-{}4} & 2 & \texttt{ab\textbar{}cd\textbar{}ef} & \texttt{c} & \texttt{c} & \texttt{abcdef} & \texttt{1} & \texttt{-{}} & \texttt{-{}0} & \texttt{s} & \texttt{a@0}\\
N307 & F3.70 & \texttt{-{}-{}4} & 2 & \texttt{ab\textbar{}cc\textbar{}de} & \texttt{c} & \texttt{c} & \texttt{abcde} & \texttt{1} & \texttt{-{}} & \texttt{-{}0} & \texttt{s} & \texttt{a@0}\\
N308 & F3.70 & \texttt{-{}-{}4} & 2 & \texttt{aa\textbar{}bb\textbar{}cd} & \texttt{b} & \texttt{b} & \texttt{abcd} & \texttt{1} & \texttt{-{}} & \texttt{-{}0} & \texttt{s} & \texttt{c@2}\\
N309 & F3.70 & \texttt{-{}-{}4} & 2 & \texttt{aa\textbar{}bb\textbar{}cc} & \texttt{b} & \texttt{b} & \texttt{abc} & \texttt{1} & \texttt{-{}} & \texttt{-{}0} & \texttt{s} & \texttt{a@0}\\
N310 & F3.70 & \texttt{-{}-{}4} & 2 & \texttt{ab\textbar{}cd\textbar{}ab} & \texttt{c} & \texttt{c} & \texttt{abcd} & \texttt{1} & \texttt{-{}} & \texttt{-{}0} & \texttt{s} & \texttt{a@0}\\
N311 & F3.70 & \texttt{-{}-{}4} & 2 & \texttt{ab\textbar{}ab\textbar{}cd} & \texttt{a} & \texttt{ac} & \texttt{abcd} & \texttt{012} & \texttt{0} & \texttt{+} & \texttt{S} & \texttt{c@2}\\
N312 & F3.70 & \texttt{-{}-{}4} & 2 & \texttt{ab\textbar{}cc\textbar{}ab} & \texttt{c} & \texttt{c} & \texttt{abc} & \texttt{1} & \texttt{-{}} & \texttt{-{}0} & \texttt{s} & \texttt{a@0}\\
N313 & F3.70 & \texttt{-{}-{}4} & 2 & \texttt{aa\textbar{}bb\textbar{}aa} & \texttt{b} & \texttt{b} & \texttt{ab} & \texttt{1} & \texttt{-{}} & \texttt{-{}0} & \texttt{s} & \texttt{a@0}\\
N314 & F3.70 & \texttt{-{}-{}4} & 2 & \texttt{aa\textbar{}aa\textbar{}bc} & \texttt{a} & \texttt{abc} & \texttt{abc} & \texttt{012} & \texttt{01} & \texttt{+} & \texttt{S} & \texttt{b@2}\\
N315 & F3.70 & \texttt{-{}-{}4} & 2 & \texttt{aa\textbar{}aa\textbar{}bb} & \texttt{a} & \texttt{ab} & \texttt{ab} & \texttt{012} & \texttt{01} & \texttt{+} & \texttt{S} & \texttt{b@2}\\
N316 & F3.70 & \texttt{-{}-{}4} & 2 & \texttt{ab\textbar{}ab\textbar{}ab} & \texttt{a} & \texttt{a} & \texttt{ab} & \texttt{012} & \texttt{0} & \texttt{+} & \texttt{S} & \texttt{none}\\
N317 & F3.70 & \texttt{-{}-{}4} & 2 & \texttt{aa\textbar{}aa\textbar{}aa} & \texttt{a} & \texttt{a} & \texttt{a} & \texttt{012} & \texttt{01} & \texttt{+} & \texttt{S} & \texttt{none}\\
N318 & F3.71 & \texttt{-{}-{}-{}} & 2 & \texttt{ab\textbar{}ab\textbar{}ab} & \texttt{a} & \texttt{a} & \texttt{ab} & \texttt{012} & \texttt{0} & \texttt{+} & \texttt{S} & \texttt{none}\\
N319 & F3.71 & \texttt{-{}-{}-{}} & 2 & \texttt{aa\textbar{}aa\textbar{}aa} & \texttt{a} & \texttt{a} & \texttt{a} & \texttt{012} & \texttt{01} & \texttt{+} & \texttt{S} & \texttt{none}\\
N320 & F3.72 & \texttt{-{}-{}-{}} & 2 & \texttt{ab\textbar{}cd\textbar{}ef} & \texttt{c} & \texttt{c} & \texttt{abcdef} & \texttt{1} & \texttt{-{}} & \texttt{-{}0} & \texttt{s} & \texttt{a@0}\\
N321 & F3.72 & \texttt{-{}-{}-{}} & 2 & \texttt{ab\textbar{}cc\textbar{}de} & \texttt{c} & \texttt{c} & \texttt{abcde} & \texttt{1} & \texttt{-{}} & \texttt{-{}0} & \texttt{s} & \texttt{a@0}\\
N322 & F3.72 & \texttt{-{}-{}-{}} & 2 & \texttt{aa\textbar{}bb\textbar{}cc} & \texttt{b} & \texttt{b} & \texttt{abc} & \texttt{1} & \texttt{-{}} & \texttt{-{}0} & \texttt{s} & \texttt{a@0}\\
N323 & F3.72 & \texttt{-{}-{}-{}} & 2 & \texttt{ab\textbar{}cd\textbar{}ab} & \texttt{c} & \texttt{c} & \texttt{abcd} & \texttt{1} & \texttt{-{}} & \texttt{-{}0} & \texttt{s} & \texttt{a@0}\\
N324 & F3.72 & \texttt{-{}-{}-{}} & 2 & \texttt{ab\textbar{}cc\textbar{}ab} & \texttt{c} & \texttt{c} & \texttt{abc} & \texttt{1} & \texttt{-{}} & \texttt{-{}0} & \texttt{s} & \texttt{a@0}\\
N325 & F3.72 & \texttt{-{}-{}-{}} & 2 & \texttt{ab\textbar{}cc\textbar{}ba} & \texttt{c} & \texttt{c} & \texttt{abc} & \texttt{1} & \texttt{-{}} & \texttt{-{}0} & \texttt{s} & \texttt{b@0}\\
N326 & F3.72 & \texttt{-{}-{}-{}} & 2 & \texttt{aa\textbar{}bb\textbar{}aa} & \texttt{b} & \texttt{b} & \texttt{ab} & \texttt{1} & \texttt{-{}} & \texttt{-{}0} & \texttt{s} & \texttt{a@0}\\
N327 & F3.72 & \texttt{-{}-{}-{}} & 2 & \texttt{ab\textbar{}ab\textbar{}ab} & \texttt{a} & \texttt{a} & \texttt{ab} & \texttt{012} & \texttt{0} & \texttt{+} & \texttt{S} & \texttt{none}\\
N328 & F3.72 & \texttt{-{}-{}-{}} & 2 & \texttt{aa\textbar{}aa\textbar{}aa} & \texttt{a} & \texttt{a} & \texttt{a} & \texttt{012} & \texttt{01} & \texttt{+} & \texttt{S} & \texttt{none}\\
N329 & F3.73 & \texttt{T-{}4} & 2 & \texttt{ab\textbar{}cd\textbar{}cd} & \texttt{c} & \texttt{c} & \texttt{abcd} & \texttt{12} & \texttt{-{}} & \texttt{-{}0} & \texttt{s} & \texttt{a@0}\\
N330 & F3.73 & \texttt{T-{}4} & 2 & \texttt{ab\textbar{}cd\textbar{}dc} & \texttt{d} & \texttt{d} & \texttt{abcd} & \texttt{12} & \texttt{-{}} & \texttt{-{}0} & \texttt{s} & \texttt{a@0}\\
N331 & F3.73 & \texttt{T-{}4} & 2 & \texttt{ab\textbar{}cc\textbar{}cc} & \texttt{c} & \texttt{c} & \texttt{abc} & \texttt{12} & \texttt{-{}} & \texttt{-{}0} & \texttt{s} & \texttt{a@0}\\
N332 & F3.73 & \texttt{T-{}4} & 2 & \texttt{aa\textbar{}bc\textbar{}cb} & \texttt{c} & \texttt{c} & \texttt{abc} & \texttt{12} & \texttt{-{}} & \texttt{-{}0} & \texttt{s} & \texttt{a@0}\\
N333 & F3.73 & \texttt{T-{}4} & 2 & \texttt{aa\textbar{}bb\textbar{}bb} & \texttt{b} & \texttt{b} & \texttt{ab} & \texttt{12} & \texttt{-{}} & \texttt{-{}0} & \texttt{s} & \texttt{a@0}\\
N334 & F3.73 & \texttt{T-{}4} & 2 & \texttt{ab\textbar{}ab\textbar{}ab} & \texttt{a} & \texttt{a} & \texttt{ab} & \texttt{012} & \texttt{0} & \texttt{+} & \texttt{S} & \texttt{none}\\
N335 & F3.73 & \texttt{T-{}4} & 2 & \texttt{aa\textbar{}aa\textbar{}aa} & \texttt{a} & \texttt{a} & \texttt{a} & \texttt{012} & \texttt{01} & \texttt{+} & \texttt{S} & \texttt{none}\\
N336 & F3.75 & \texttt{T-{}-{}} & 2 & \texttt{ab\textbar{}cd\textbar{}ab} & \texttt{a} & \texttt{a} & \texttt{abcd} & \texttt{02} & \texttt{0} & \texttt{-{}1} & \texttt{s} & \texttt{c@1}\\
N337 & F3.75 & \texttt{T-{}-{}} & 2 & \texttt{aa\textbar{}bc\textbar{}aa} & \texttt{a} & \texttt{a} & \texttt{abc} & \texttt{02} & \texttt{01} & \texttt{-{}1} & \texttt{s} & \texttt{b@1}\\
N338 & F3.75 & \texttt{T-{}-{}} & 2 & \texttt{aa\textbar{}bb\textbar{}aa} & \texttt{a} & \texttt{a} & \texttt{ab} & \texttt{02} & \texttt{01} & \texttt{-{}1} & \texttt{s} & \texttt{b@1}\\
N339 & F3.75 & \texttt{T-{}-{}} & 2 & \texttt{ab\textbar{}ab\textbar{}ab} & \texttt{a} & \texttt{a} & \texttt{ab} & \texttt{012} & \texttt{0} & \texttt{+} & \texttt{S} & \texttt{none}\\
N340 & F3.75 & \texttt{T-{}-{}} & 2 & \texttt{aa\textbar{}aa\textbar{}aa} & \texttt{a} & \texttt{a} & \texttt{a} & \texttt{012} & \texttt{01} & \texttt{+} & \texttt{S} & \texttt{none}\\
N341 & F3.76 & \texttt{T-{}4} & 2 & \texttt{ab\textbar{}cd\textbar{}ef} & \texttt{e} & \texttt{e} & \texttt{abcdef} & \texttt{2} & \texttt{-{}} & \texttt{-{}0} & \texttt{s} & \texttt{a@0}\\
N342 & F3.76 & \texttt{T-{}4} & 2 & \texttt{ab\textbar{}cd\textbar{}ee} & \texttt{e} & \texttt{e} & \texttt{abcde} & \texttt{2} & \texttt{-{}} & \texttt{-{}0} & \texttt{s} & \texttt{a@0}\\
N343 & F3.76 & \texttt{T-{}4} & 2 & \texttt{aa\textbar{}bc\textbar{}dd} & \texttt{d} & \texttt{d} & \texttt{abcd} & \texttt{2} & \texttt{-{}} & \texttt{-{}0} & \texttt{s} & \texttt{b@1}\\
N344 & F3.76 & \texttt{T-{}4} & 2 & \texttt{ab\textbar{}cd\textbar{}ab} & \texttt{a} & \texttt{a} & \texttt{abcd} & \texttt{02} & \texttt{0} & \texttt{-{}1} & \texttt{s} & \texttt{c@1}\\
N345 & F3.76 & \texttt{T-{}4} & 2 & \texttt{aa\textbar{}bc\textbar{}aa} & \texttt{a} & \texttt{a} & \texttt{abc} & \texttt{02} & \texttt{01} & \texttt{-{}1} & \texttt{s} & \texttt{b@1}\\
N346 & F3.76 & \texttt{T-{}4} & 2 & \texttt{aa\textbar{}bb\textbar{}cc} & \texttt{c} & \texttt{c} & \texttt{abc} & \texttt{2} & \texttt{-{}} & \texttt{-{}0} & \texttt{s} & \texttt{a@0}\\
N347 & F3.76 & \texttt{T-{}4} & 2 & \texttt{aa\textbar{}bb\textbar{}aa} & \texttt{a} & \texttt{a} & \texttt{ab} & \texttt{02} & \texttt{01} & \texttt{-{}1} & \texttt{s} & \texttt{b@1}\\
N348 & F3.76 & \texttt{T-{}4} & 2 & \texttt{ab\textbar{}ab\textbar{}cd} & \texttt{c} & \texttt{c} & \texttt{abcd} & \texttt{2} & \texttt{-{}} & \texttt{-{}0} & \texttt{s} & \texttt{a@0}\\
N349 & F3.76 & \texttt{T-{}4} & 2 & \texttt{ab\textbar{}ab\textbar{}cc} & \texttt{c} & \texttt{c} & \texttt{abc} & \texttt{2} & \texttt{-{}} & \texttt{-{}0} & \texttt{s} & \texttt{a@0}\\
N350 & F3.76 & \texttt{T-{}4} & 2 & \texttt{ab\textbar{}ba\textbar{}cc} & \texttt{c} & \texttt{c} & \texttt{abc} & \texttt{2} & \texttt{-{}} & \texttt{-{}0} & \texttt{s} & \texttt{b@0}\\
N351 & F3.76 & \texttt{T-{}4} & 2 & \texttt{aa\textbar{}aa\textbar{}bb} & \texttt{b} & \texttt{b} & \texttt{ab} & \texttt{2} & \texttt{-{}} & \texttt{-{}0} & \texttt{s} & \texttt{a@0}\\
N352 & F3.76 & \texttt{T-{}4} & 2 & \texttt{ab\textbar{}ab\textbar{}ab} & \texttt{a} & \texttt{a} & \texttt{ab} & \texttt{012} & \texttt{0} & \texttt{+} & \texttt{S} & \texttt{none}\\
N353 & F3.76 & \texttt{T-{}4} & 2 & \texttt{aa\textbar{}aa\textbar{}aa} & \texttt{a} & \texttt{a} & \texttt{a} & \texttt{012} & \texttt{01} & \texttt{+} & \texttt{S} & \texttt{none}\\
N354 & F3.79 & \texttt{-{}-{}-{}} & 2 & \texttt{ab\textbar{}ab\textbar{}ab} & \texttt{a} & \texttt{a} & \texttt{ab} & \texttt{012} & \texttt{0} & \texttt{+} & \texttt{S} & \texttt{none}\\
N355 & F3.79 & \texttt{-{}-{}-{}} & 2 & \texttt{aa\textbar{}aa\textbar{}aa} & \texttt{a} & \texttt{a} & \texttt{a} & \texttt{012} & \texttt{01} & \texttt{+} & \texttt{S} & \texttt{none}\\
N356 & F3.80 & \texttt{-{}-{}-{}} & 2 & \texttt{ab\textbar{}ab\textbar{}ab} & \texttt{a} & \texttt{a} & \texttt{ab} & \texttt{012} & \texttt{0} & \texttt{+} & \texttt{S} & \texttt{none}\\
N357 & F3.80 & \texttt{-{}-{}-{}} & 2 & \texttt{aa\textbar{}aa\textbar{}aa} & \texttt{a} & \texttt{a} & \texttt{a} & \texttt{012} & \texttt{01} & \texttt{+} & \texttt{S} & \texttt{none}\\
N358 & F3.81 & \texttt{-{}-{}-{}} & 2 & \texttt{ab\textbar{}ab\textbar{}ab} & \texttt{a} & \texttt{a} & \texttt{ab} & \texttt{012} & \texttt{0} & \texttt{+} & \texttt{S} & \texttt{none}\\
N359 & F3.81 & \texttt{-{}-{}-{}} & 2 & \texttt{aa\textbar{}aa\textbar{}aa} & \texttt{a} & \texttt{a} & \texttt{a} & \texttt{012} & \texttt{01} & \texttt{+} & \texttt{S} & \texttt{none}\\
N360 & F3.82 & \texttt{-{}-{}-{}} & 2 & \texttt{ab\textbar{}cd\textbar{}ab} & \texttt{a} & \texttt{a} & \texttt{abcd} & \texttt{02} & \texttt{0} & \texttt{-{}1} & \texttt{s} & \texttt{c@1}\\
N361 & F3.82 & \texttt{-{}-{}-{}} & 2 & \texttt{aa\textbar{}bc\textbar{}aa} & \texttt{a} & \texttt{a} & \texttt{abc} & \texttt{02} & \texttt{01} & \texttt{-{}1} & \texttt{s} & \texttt{b@1}\\
N362 & F3.82 & \texttt{-{}-{}-{}} & 2 & \texttt{aa\textbar{}bb\textbar{}aa} & \texttt{a} & \texttt{a} & \texttt{ab} & \texttt{02} & \texttt{01} & \texttt{-{}1} & \texttt{s} & \texttt{b@1}\\
N363 & F3.82 & \texttt{-{}-{}-{}} & 2 & \texttt{ab\textbar{}ab\textbar{}ab} & \texttt{a} & \texttt{a} & \texttt{ab} & \texttt{012} & \texttt{0} & \texttt{+} & \texttt{S} & \texttt{none}\\
N364 & F3.82 & \texttt{-{}-{}-{}} & 2 & \texttt{aa\textbar{}aa\textbar{}aa} & \texttt{a} & \texttt{a} & \texttt{a} & \texttt{012} & \texttt{01} & \texttt{+} & \texttt{S} & \texttt{none}\\
N365 & F3.83 & \texttt{-{}-{}-{}} & 2 & \texttt{ab\textbar{}ab\textbar{}ab} & \texttt{a} & \texttt{a} & \texttt{ab} & \texttt{012} & \texttt{0} & \texttt{+} & \texttt{S} & \texttt{none}\\
N366 & F3.83 & \texttt{-{}-{}-{}} & 2 & \texttt{aa\textbar{}aa\textbar{}aa} & \texttt{a} & \texttt{a} & \texttt{a} & \texttt{012} & \texttt{01} & \texttt{+} & \texttt{S} & \texttt{none}\\
N367 & F3.84 & \texttt{-{}-{}-{}} & 2 & \texttt{ab\textbar{}cd\textbar{}ef} & \texttt{c} & \texttt{c} & \texttt{abcdef} & \texttt{1} & \texttt{-{}} & \texttt{-{}0} & \texttt{s} & \texttt{a@0}\\
N368 & F3.84 & \texttt{-{}-{}-{}} & 2 & \texttt{ab\textbar{}cc\textbar{}de} & \texttt{c} & \texttt{c} & \texttt{abcde} & \texttt{1} & \texttt{-{}} & \texttt{-{}0} & \texttt{s} & \texttt{a@0}\\
N369 & F3.84 & \texttt{-{}-{}-{}} & 2 & \texttt{aa\textbar{}bb\textbar{}cc} & \texttt{b} & \texttt{b} & \texttt{abc} & \texttt{1} & \texttt{-{}} & \texttt{-{}0} & \texttt{s} & \texttt{a@0}\\
N370 & F3.84 & \texttt{-{}-{}-{}} & 2 & \texttt{ab\textbar{}cd\textbar{}ab} & \texttt{c} & \texttt{c} & \texttt{abcd} & \texttt{1} & \texttt{-{}} & \texttt{-{}0} & \texttt{s} & \texttt{a@0}\\
N371 & F3.84 & \texttt{-{}-{}-{}} & 2 & \texttt{ab\textbar{}cc\textbar{}ab} & \texttt{c} & \texttt{c} & \texttt{abc} & \texttt{1} & \texttt{-{}} & \texttt{-{}0} & \texttt{s} & \texttt{a@0}\\
N372 & F3.84 & \texttt{-{}-{}-{}} & 2 & \texttt{aa\textbar{}bb\textbar{}aa} & \texttt{b} & \texttt{b} & \texttt{ab} & \texttt{1} & \texttt{-{}} & \texttt{-{}0} & \texttt{s} & \texttt{a@0}\\
N373 & F3.84 & \texttt{-{}-{}-{}} & 2 & \texttt{ab\textbar{}ab\textbar{}ab} & \texttt{a} & \texttt{a} & \texttt{ab} & \texttt{012} & \texttt{0} & \texttt{+} & \texttt{S} & \texttt{none}\\
N374 & F3.84 & \texttt{-{}-{}-{}} & 2 & \texttt{aa\textbar{}aa\textbar{}aa} & \texttt{a} & \texttt{a} & \texttt{a} & \texttt{012} & \texttt{01} & \texttt{+} & \texttt{S} & \texttt{none}\\
N375 & F3.85 & \texttt{-{}-{}4} & 2 & \texttt{ab\textbar{}cd\textbar{}ef} & \texttt{ac} & \texttt{ace} & \texttt{abcdef} & \texttt{012} & \texttt{0} & \texttt{+} & \texttt{.} & \texttt{a@0}\\
N376 & F3.85 & \texttt{-{}-{}4} & 2 & \texttt{ab\textbar{}cd\textbar{}ab} & \texttt{ac} & \texttt{ac} & \texttt{abcd} & \texttt{012} & \texttt{0} & \texttt{+} & \texttt{.} & \texttt{c@1}\\
N377 & F3.85 & \texttt{-{}-{}4} & 2 & \texttt{aa\textbar{}bb\textbar{}cd} & \texttt{ab} & \texttt{abcd} & \texttt{abcd} & \texttt{012} & \texttt{01} & \texttt{+} & \texttt{.} & \texttt{c@2}\\
N378 & F3.85 & \texttt{-{}-{}4} & 2 & \texttt{aa\textbar{}bb\textbar{}cc} & \texttt{ab} & \texttt{abc} & \texttt{abc} & \texttt{012} & \texttt{01} & \texttt{+} & \texttt{.} & \texttt{a@0}\\
N379 & F3.85 & \texttt{-{}-{}4} & 2 & \texttt{aa\textbar{}bb\textbar{}aa} & \texttt{ab} & \texttt{ab} & \texttt{ab} & \texttt{012} & \texttt{01} & \texttt{+} & \texttt{.} & \texttt{b@1}\\
N380 & F3.85 & \texttt{-{}-{}4} & 2 & \texttt{ab\textbar{}ab\textbar{}cd} & \texttt{a} & \texttt{ac} & \texttt{abcd} & \texttt{012} & \texttt{0} & \texttt{+} & \texttt{S} & \texttt{c@2}\\
N381 & F3.85 & \texttt{-{}-{}4} & 2 & \texttt{ab\textbar{}ba\textbar{}cd} & \texttt{ab} & \texttt{abcd} & \texttt{abcd} & \texttt{012} & \texttt{01} & \texttt{+} & \texttt{.} & \texttt{c@2}\\
N382 & F3.85 & \texttt{-{}-{}4} & 2 & \texttt{ab\textbar{}ba\textbar{}cc} & \texttt{ab} & \texttt{abc} & \texttt{abc} & \texttt{012} & \texttt{01} & \texttt{+} & \texttt{.} & \texttt{c@2}\\
N383 & F3.85 & \texttt{-{}-{}4} & 2 & \texttt{aa\textbar{}aa\textbar{}bc} & \texttt{a} & \texttt{abc} & \texttt{abc} & \texttt{012} & \texttt{01} & \texttt{+} & \texttt{S} & \texttt{b@2}\\
N384 & F3.85 & \texttt{-{}-{}4} & 2 & \texttt{aa\textbar{}aa\textbar{}bb} & \texttt{a} & \texttt{ab} & \texttt{ab} & \texttt{012} & \texttt{01} & \texttt{+} & \texttt{S} & \texttt{b@2}\\
N385 & F3.85 & \texttt{-{}-{}4} & 2 & \texttt{ab\textbar{}ab\textbar{}ab} & \texttt{a} & \texttt{a} & \texttt{ab} & \texttt{012} & \texttt{0} & \texttt{+} & \texttt{S} & \texttt{none}\\
N386 & F3.85 & \texttt{-{}-{}4} & 2 & \texttt{ab\textbar{}ba\textbar{}ab} & \texttt{ab} & \texttt{ab} & \texttt{ab} & \texttt{012} & \texttt{01} & \texttt{+} & \texttt{.} & \texttt{none}\\
N387 & F3.85 & \texttt{-{}-{}4} & 2 & \texttt{aa\textbar{}aa\textbar{}aa} & \texttt{a} & \texttt{a} & \texttt{a} & \texttt{012} & \texttt{01} & \texttt{+} & \texttt{S} & \texttt{none}\\
N388 & F3.86 & \texttt{-{}-{}-{}} & 2 & \texttt{ab\textbar{}cd\textbar{}cd} & \texttt{ac} & \texttt{ac} & \texttt{abcd} & \texttt{012} & \texttt{0} & \texttt{+} & \texttt{.} & \texttt{a@0}\\
N389 & F3.86 & \texttt{-{}-{}-{}} & 2 & \texttt{aa\textbar{}bb\textbar{}bb} & \texttt{ab} & \texttt{ab} & \texttt{ab} & \texttt{012} & \texttt{01} & \texttt{+} & \texttt{.} & \texttt{a@0}\\
N390 & F3.86 & \texttt{-{}-{}-{}} & 2 & \texttt{ab\textbar{}ab\textbar{}ab} & \texttt{a} & \texttt{a} & \texttt{ab} & \texttt{012} & \texttt{0} & \texttt{+} & \texttt{S} & \texttt{none}\\
N391 & F3.86 & \texttt{-{}-{}-{}} & 2 & \texttt{ab\textbar{}ba\textbar{}ba} & \texttt{ab} & \texttt{ab} & \texttt{ab} & \texttt{012} & \texttt{01} & \texttt{+} & \texttt{.} & \texttt{none}\\
N392 & F3.86 & \texttt{-{}-{}-{}} & 2 & \texttt{aa\textbar{}aa\textbar{}aa} & \texttt{a} & \texttt{a} & \texttt{a} & \texttt{012} & \texttt{01} & \texttt{+} & \texttt{S} & \texttt{none}\\
N393 & F3.87 & \texttt{-{}-{}-{}} & 2 & \texttt{ab\textbar{}cd\textbar{}ab} & \texttt{ac} & \texttt{ac} & \texttt{abcd} & \texttt{012} & \texttt{0} & \texttt{+} & \texttt{.} & \texttt{c@1}\\
N394 & F3.87 & \texttt{-{}-{}-{}} & 2 & \texttt{aa\textbar{}bb\textbar{}aa} & \texttt{ab} & \texttt{ab} & \texttt{ab} & \texttt{012} & \texttt{01} & \texttt{+} & \texttt{.} & \texttt{b@1}\\
N395 & F3.87 & \texttt{-{}-{}-{}} & 2 & \texttt{ab\textbar{}ab\textbar{}ab} & \texttt{a} & \texttt{a} & \texttt{ab} & \texttt{012} & \texttt{0} & \texttt{+} & \texttt{S} & \texttt{none}\\
N396 & F3.87 & \texttt{-{}-{}-{}} & 2 & \texttt{ab\textbar{}ba\textbar{}ab} & \texttt{ab} & \texttt{ab} & \texttt{ab} & \texttt{012} & \texttt{01} & \texttt{+} & \texttt{.} & \texttt{none}\\
N397 & F3.87 & \texttt{-{}-{}-{}} & 2 & \texttt{aa\textbar{}aa\textbar{}aa} & \texttt{a} & \texttt{a} & \texttt{a} & \texttt{012} & \texttt{01} & \texttt{+} & \texttt{S} & \texttt{none}\\
N398 & F3.90 & \texttt{T-{}-{}} & 2 & \texttt{ab\textbar{}ab\textbar{}ab} & \texttt{a} & \texttt{a} & \texttt{ab} & \texttt{012} & \texttt{0} & \texttt{+} & \texttt{S} & \texttt{none}\\
N399 & F3.90 & \texttt{T-{}-{}} & 2 & \texttt{aa\textbar{}aa\textbar{}aa} & \texttt{a} & \texttt{a} & \texttt{a} & \texttt{012} & \texttt{01} & \texttt{+} & \texttt{S} & \texttt{none}\\
N400 & F3.91 & \texttt{T-{}4} & 2 & \texttt{ab\textbar{}cd\textbar{}ef} & \texttt{e} & \texttt{e} & \texttt{abcdef} & \texttt{2} & \texttt{-{}} & \texttt{-{}0} & \texttt{s} & \texttt{a@0}\\
N401 & F3.91 & \texttt{T-{}4} & 2 & \texttt{ab\textbar{}cd\textbar{}ee} & \texttt{e} & \texttt{e} & \texttt{abcde} & \texttt{2} & \texttt{-{}} & \texttt{-{}0} & \texttt{s} & \texttt{a@0}\\
N402 & F3.91 & \texttt{T-{}4} & 2 & \texttt{ab\textbar{}cc\textbar{}dd} & \texttt{d} & \texttt{d} & \texttt{abcd} & \texttt{2} & \texttt{-{}} & \texttt{-{}0} & \texttt{s} & \texttt{a@0}\\
N403 & F3.91 & \texttt{T-{}4} & 2 & \texttt{ab\textbar{}cd\textbar{}cd} & \texttt{c} & \texttt{c} & \texttt{abcd} & \texttt{12} & \texttt{-{}} & \texttt{-{}0} & \texttt{s} & \texttt{a@0}\\
N404 & F3.91 & \texttt{T-{}4} & 2 & \texttt{ab\textbar{}cc\textbar{}cc} & \texttt{c} & \texttt{c} & \texttt{abc} & \texttt{12} & \texttt{-{}} & \texttt{-{}0} & \texttt{s} & \texttt{a@0}\\
N405 & F3.91 & \texttt{T-{}4} & 2 & \texttt{aa\textbar{}bb\textbar{}cc} & \texttt{c} & \texttt{c} & \texttt{abc} & \texttt{2} & \texttt{-{}} & \texttt{-{}0} & \texttt{s} & \texttt{a@0}\\
N406 & F3.91 & \texttt{T-{}4} & 2 & \texttt{aa\textbar{}bb\textbar{}bb} & \texttt{b} & \texttt{b} & \texttt{ab} & \texttt{12} & \texttt{-{}} & \texttt{-{}0} & \texttt{s} & \texttt{a@0}\\
N407 & F3.91 & \texttt{T-{}4} & 2 & \texttt{ab\textbar{}ab\textbar{}cd} & \texttt{c} & \texttt{c} & \texttt{abcd} & \texttt{2} & \texttt{-{}} & \texttt{-{}0} & \texttt{s} & \texttt{a@0}\\
N408 & F3.91 & \texttt{T-{}4} & 2 & \texttt{ab\textbar{}ab\textbar{}cc} & \texttt{c} & \texttt{c} & \texttt{abc} & \texttt{2} & \texttt{-{}} & \texttt{-{}0} & \texttt{s} & \texttt{a@0}\\
N409 & F3.91 & \texttt{T-{}4} & 2 & \texttt{aa\textbar{}aa\textbar{}bb} & \texttt{b} & \texttt{b} & \texttt{ab} & \texttt{2} & \texttt{-{}} & \texttt{-{}0} & \texttt{s} & \texttt{a@0}\\
N410 & F3.91 & \texttt{T-{}4} & 2 & \texttt{ab\textbar{}ab\textbar{}ab} & \texttt{a} & \texttt{a} & \texttt{ab} & \texttt{012} & \texttt{0} & \texttt{+} & \texttt{S} & \texttt{none}\\
N411 & F3.91 & \texttt{T-{}4} & 2 & \texttt{aa\textbar{}aa\textbar{}aa} & \texttt{a} & \texttt{a} & \texttt{a} & \texttt{012} & \texttt{01} & \texttt{+} & \texttt{S} & \texttt{none}\\
N412 & F3.92 & \texttt{T-{}-{}} & 2 & \texttt{ab\textbar{}cd\textbar{}cd} & \texttt{ac} & \texttt{ac} & \texttt{abcd} & \texttt{012} & \texttt{0} & \texttt{+} & \texttt{.} & \texttt{a@0}\\
N413 & F3.92 & \texttt{T-{}-{}} & 2 & \texttt{aa\textbar{}bb\textbar{}bb} & \texttt{ab} & \texttt{ab} & \texttt{ab} & \texttt{012} & \texttt{01} & \texttt{+} & \texttt{.} & \texttt{a@0}\\
N414 & F3.92 & \texttt{T-{}-{}} & 2 & \texttt{ab\textbar{}ab\textbar{}cd} & \texttt{a} & \texttt{a} & \texttt{abcd} & \texttt{01} & \texttt{0} & \texttt{-{}2} & \texttt{s} & \texttt{c@2}\\
N415 & F3.92 & \texttt{T-{}-{}} & 2 & \texttt{ab\textbar{}ba\textbar{}cd} & \texttt{ab} & \texttt{ab} & \texttt{abcd} & \texttt{01} & \texttt{01} & \texttt{-{}2} & \texttt{.} & \texttt{c@2}\\
N416 & F3.92 & \texttt{T-{}-{}} & 2 & \texttt{aa\textbar{}aa\textbar{}bc} & \texttt{a} & \texttt{a} & \texttt{abc} & \texttt{01} & \texttt{01} & \texttt{-{}2} & \texttt{s} & \texttt{b@2}\\
N417 & F3.92 & \texttt{T-{}-{}} & 2 & \texttt{aa\textbar{}aa\textbar{}bb} & \texttt{a} & \texttt{a} & \texttt{ab} & \texttt{01} & \texttt{01} & \texttt{-{}2} & \texttt{s} & \texttt{b@2}\\
N418 & F3.92 & \texttt{T-{}-{}} & 2 & \texttt{ab\textbar{}ab\textbar{}ab} & \texttt{a} & \texttt{a} & \texttt{ab} & \texttt{012} & \texttt{0} & \texttt{+} & \texttt{S} & \texttt{none}\\
N419 & F3.92 & \texttt{T-{}-{}} & 2 & \texttt{ab\textbar{}ba\textbar{}ba} & \texttt{ab} & \texttt{ab} & \texttt{ab} & \texttt{012} & \texttt{01} & \texttt{+} & \texttt{.} & \texttt{none}\\
N420 & F3.92 & \texttt{T-{}-{}} & 2 & \texttt{aa\textbar{}aa\textbar{}aa} & \texttt{a} & \texttt{a} & \texttt{a} & \texttt{012} & \texttt{01} & \texttt{+} & \texttt{S} & \texttt{none}\\
N421 & F3.93 & \texttt{T-{}-{}} & 2 & \texttt{ab\textbar{}ab\textbar{}ab} & \texttt{a} & \texttt{a} & \texttt{ab} & \texttt{012} & \texttt{0} & \texttt{+} & \texttt{S} & \texttt{none}\\
N422 & F3.93 & \texttt{T-{}-{}} & 2 & \texttt{aa\textbar{}aa\textbar{}aa} & \texttt{a} & \texttt{a} & \texttt{a} & \texttt{012} & \texttt{01} & \texttt{+} & \texttt{S} & \texttt{none}\\
N423 & F3.94 & \texttt{-{}-{}-{}} & 2 & \texttt{ab\textbar{}ab\textbar{}ab} & \texttt{a} & \texttt{a} & \texttt{ab} & \texttt{012} & \texttt{0} & \texttt{+} & \texttt{S} & \texttt{none}\\
N424 & F3.94 & \texttt{-{}-{}-{}} & 2 & \texttt{aa\textbar{}aa\textbar{}aa} & \texttt{a} & \texttt{a} & \texttt{a} & \texttt{012} & \texttt{01} & \texttt{+} & \texttt{S} & \texttt{none}\\
N425 & F3.95 & \texttt{-{}-{}-{}} & 2 & \texttt{ab\textbar{}cd\textbar{}cd} & \texttt{ac} & \texttt{ac} & \texttt{abcd} & \texttt{012} & \texttt{0} & \texttt{+} & \texttt{.} & \texttt{a@0}\\
N426 & F3.95 & \texttt{-{}-{}-{}} & 2 & \texttt{aa\textbar{}bb\textbar{}bb} & \texttt{ab} & \texttt{ab} & \texttt{ab} & \texttt{012} & \texttt{01} & \texttt{+} & \texttt{.} & \texttt{a@0}\\
N427 & F3.95 & \texttt{-{}-{}-{}} & 2 & \texttt{ab\textbar{}cd\textbar{}ab} & \texttt{ac} & \texttt{ac} & \texttt{abcd} & \texttt{012} & \texttt{0} & \texttt{+} & \texttt{.} & \texttt{c@1}\\
N428 & F3.95 & \texttt{-{}-{}-{}} & 2 & \texttt{aa\textbar{}bb\textbar{}aa} & \texttt{ab} & \texttt{ab} & \texttt{ab} & \texttt{012} & \texttt{01} & \texttt{+} & \texttt{.} & \texttt{b@1}\\
N429 & F3.95 & \texttt{-{}-{}-{}} & 2 & \texttt{ab\textbar{}ab\textbar{}ab} & \texttt{a} & \texttt{a} & \texttt{ab} & \texttt{012} & \texttt{0} & \texttt{+} & \texttt{S} & \texttt{none}\\
N430 & F3.95 & \texttt{-{}-{}-{}} & 2 & \texttt{ab\textbar{}ba\textbar{}ba} & \texttt{ab} & \texttt{ab} & \texttt{ab} & \texttt{012} & \texttt{01} & \texttt{+} & \texttt{.} & \texttt{none}\\
N431 & F3.95 & \texttt{-{}-{}-{}} & 2 & \texttt{ab\textbar{}ba\textbar{}ab} & \texttt{ab} & \texttt{ab} & \texttt{ab} & \texttt{012} & \texttt{01} & \texttt{+} & \texttt{.} & \texttt{none}\\
N432 & F3.95 & \texttt{-{}-{}-{}} & 2 & \texttt{aa\textbar{}aa\textbar{}aa} & \texttt{a} & \texttt{a} & \texttt{a} & \texttt{012} & \texttt{01} & \texttt{+} & \texttt{S} & \texttt{none}\\
N433 & F3.97 & \texttt{-{}-{}-{}} & 2 & \texttt{ab\textbar{}cd\textbar{}ab} & \texttt{ac} & \texttt{ac} & \texttt{abcd} & \texttt{012} & \texttt{0} & \texttt{+} & \texttt{.} & \texttt{c@1}\\
N434 & F3.97 & \texttt{-{}-{}-{}} & 2 & \texttt{aa\textbar{}bb\textbar{}aa} & \texttt{ab} & \texttt{ab} & \texttt{ab} & \texttt{012} & \texttt{01} & \texttt{+} & \texttt{.} & \texttt{b@1}\\
N435 & F3.97 & \texttt{-{}-{}-{}} & 2 & \texttt{ab\textbar{}ab\textbar{}ab} & \texttt{a} & \texttt{a} & \texttt{ab} & \texttt{012} & \texttt{0} & \texttt{+} & \texttt{S} & \texttt{none}\\
N436 & F3.97 & \texttt{-{}-{}-{}} & 2 & \texttt{ab\textbar{}ba\textbar{}ab} & \texttt{ab} & \texttt{ab} & \texttt{ab} & \texttt{012} & \texttt{01} & \texttt{+} & \texttt{.} & \texttt{none}\\
N437 & F3.97 & \texttt{-{}-{}-{}} & 2 & \texttt{aa\textbar{}aa\textbar{}aa} & \texttt{a} & \texttt{a} & \texttt{a} & \texttt{012} & \texttt{01} & \texttt{+} & \texttt{S} & \texttt{none}\\
N438 & F3.98 & \texttt{T-{}4} & 2 & \texttt{ab\textbar{}cd\textbar{}ef} & \texttt{e} & \texttt{e} & \texttt{abcdef} & \texttt{2} & \texttt{-{}} & \texttt{-{}0} & \texttt{s} & \texttt{a@0}\\
N439 & F3.98 & \texttt{T-{}4} & 2 & \texttt{ab\textbar{}cd\textbar{}ee} & \texttt{e} & \texttt{e} & \texttt{abcde} & \texttt{2} & \texttt{-{}} & \texttt{-{}0} & \texttt{s} & \texttt{a@0}\\
N440 & F3.98 & \texttt{T-{}4} & 2 & \texttt{aa\textbar{}bb\textbar{}cc} & \texttt{c} & \texttt{c} & \texttt{abc} & \texttt{2} & \texttt{-{}} & \texttt{-{}0} & \texttt{s} & \texttt{a@0}\\
N441 & F3.98 & \texttt{T-{}4} & 2 & \texttt{ab\textbar{}ab\textbar{}cd} & \texttt{c} & \texttt{c} & \texttt{abcd} & \texttt{2} & \texttt{-{}} & \texttt{-{}0} & \texttt{s} & \texttt{a@0}\\
N442 & F3.98 & \texttt{T-{}4} & 2 & \texttt{ab\textbar{}ab\textbar{}cc} & \texttt{c} & \texttt{c} & \texttt{abc} & \texttt{2} & \texttt{-{}} & \texttt{-{}0} & \texttt{s} & \texttt{a@0}\\
N443 & F3.98 & \texttt{T-{}4} & 2 & \texttt{ab\textbar{}ba\textbar{}cc} & \texttt{c} & \texttt{c} & \texttt{abc} & \texttt{2} & \texttt{-{}} & \texttt{-{}0} & \texttt{s} & \texttt{b@0}\\
N444 & F3.98 & \texttt{T-{}4} & 2 & \texttt{aa\textbar{}aa\textbar{}bb} & \texttt{b} & \texttt{b} & \texttt{ab} & \texttt{2} & \texttt{-{}} & \texttt{-{}0} & \texttt{s} & \texttt{a@0}\\
N445 & F3.98 & \texttt{T-{}4} & 2 & \texttt{ab\textbar{}ab\textbar{}ab} & \texttt{a} & \texttt{a} & \texttt{ab} & \texttt{012} & \texttt{0} & \texttt{+} & \texttt{S} & \texttt{none}\\
N446 & F3.98 & \texttt{T-{}4} & 2 & \texttt{aa\textbar{}aa\textbar{}aa} & \texttt{a} & \texttt{a} & \texttt{a} & \texttt{012} & \texttt{01} & \texttt{+} & \texttt{S} & \texttt{none}\\
N447 & F3.100 & \texttt{T-{}-{}} & 2 & \texttt{ab\textbar{}ab\textbar{}cd} & \texttt{ac} & \texttt{ac} & \texttt{abcd} & \texttt{012} & \texttt{0} & \texttt{+} & \texttt{.} & \texttt{c@2}\\
N448 & F3.100 & \texttt{T-{}-{}} & 2 & \texttt{aa\textbar{}aa\textbar{}bb} & \texttt{ab} & \texttt{ab} & \texttt{ab} & \texttt{012} & \texttt{01} & \texttt{+} & \texttt{.} & \texttt{b@2}\\
N449 & F3.100 & \texttt{T-{}-{}} & 2 & \texttt{ab\textbar{}ab\textbar{}ab} & \texttt{a} & \texttt{a} & \texttt{ab} & \texttt{012} & \texttt{0} & \texttt{+} & \texttt{S} & \texttt{none}\\
N450 & F3.100 & \texttt{T-{}-{}} & 2 & \texttt{ab\textbar{}ab\textbar{}ba} & \texttt{ab} & \texttt{ab} & \texttt{ab} & \texttt{012} & \texttt{01} & \texttt{+} & \texttt{.} & \texttt{none}\\
N451 & F3.100 & \texttt{T-{}-{}} & 2 & \texttt{aa\textbar{}aa\textbar{}aa} & \texttt{a} & \texttt{a} & \texttt{a} & \texttt{012} & \texttt{01} & \texttt{+} & \texttt{S} & \texttt{none}\\
N452 & F3.101 & \texttt{T-{}4} & 2 & \texttt{ab\textbar{}cd\textbar{}ef} & \texttt{ae} & \texttt{ae} & \texttt{abcdef} & \texttt{02} & \texttt{0} & \texttt{-{}1} & \texttt{.} & \texttt{a@0}\\
N453 & F3.101 & \texttt{T-{}4} & 2 & \texttt{ab\textbar{}ab\textbar{}cd} & \texttt{ac} & \texttt{ac} & \texttt{abcd} & \texttt{012} & \texttt{0} & \texttt{+} & \texttt{.} & \texttt{c@2}\\
N454 & F3.101 & \texttt{T-{}4} & 2 & \texttt{aa\textbar{}bc\textbar{}dd} & \texttt{ad} & \texttt{ad} & \texttt{abcd} & \texttt{02} & \texttt{01} & \texttt{-{}1} & \texttt{.} & \texttt{b@1}\\
N455 & F3.101 & \texttt{T-{}4} & 2 & \texttt{aa\textbar{}bb\textbar{}cc} & \texttt{ac} & \texttt{ac} & \texttt{abc} & \texttt{02} & \texttt{01} & \texttt{-{}1} & \texttt{.} & \texttt{a@0}\\
N456 & F3.101 & \texttt{T-{}4} & 2 & \texttt{aa\textbar{}aa\textbar{}bb} & \texttt{ab} & \texttt{ab} & \texttt{ab} & \texttt{012} & \texttt{01} & \texttt{+} & \texttt{.} & \texttt{b@2}\\
N457 & F3.101 & \texttt{T-{}4} & 2 & \texttt{ab\textbar{}cd\textbar{}ab} & \texttt{a} & \texttt{a} & \texttt{abcd} & \texttt{02} & \texttt{0} & \texttt{-{}1} & \texttt{s} & \texttt{c@1}\\
N458 & F3.101 & \texttt{T-{}4} & 2 & \texttt{ab\textbar{}cd\textbar{}ba} & \texttt{ab} & \texttt{ab} & \texttt{abcd} & \texttt{02} & \texttt{01} & \texttt{-{}1} & \texttt{.} & \texttt{c@1}\\
N459 & F3.101 & \texttt{T-{}4} & 2 & \texttt{aa\textbar{}bc\textbar{}aa} & \texttt{a} & \texttt{a} & \texttt{abc} & \texttt{02} & \texttt{01} & \texttt{-{}1} & \texttt{s} & \texttt{b@1}\\
N460 & F3.101 & \texttt{T-{}4} & 2 & \texttt{ab\textbar{}cc\textbar{}ba} & \texttt{ab} & \texttt{ab} & \texttt{abc} & \texttt{02} & \texttt{01} & \texttt{-{}1} & \texttt{.} & \texttt{c@1}\\
N461 & F3.101 & \texttt{T-{}4} & 2 & \texttt{aa\textbar{}bb\textbar{}aa} & \texttt{a} & \texttt{a} & \texttt{ab} & \texttt{02} & \texttt{01} & \texttt{-{}1} & \texttt{s} & \texttt{b@1}\\
N462 & F3.101 & \texttt{T-{}4} & 2 & \texttt{ab\textbar{}ab\textbar{}ab} & \texttt{a} & \texttt{a} & \texttt{ab} & \texttt{012} & \texttt{0} & \texttt{+} & \texttt{S} & \texttt{none}\\
N463 & F3.101 & \texttt{T-{}4} & 2 & \texttt{ab\textbar{}ab\textbar{}ba} & \texttt{ab} & \texttt{ab} & \texttt{ab} & \texttt{012} & \texttt{01} & \texttt{+} & \texttt{.} & \texttt{none}\\
N464 & F3.101 & \texttt{T-{}4} & 2 & \texttt{aa\textbar{}aa\textbar{}aa} & \texttt{a} & \texttt{a} & \texttt{a} & \texttt{012} & \texttt{01} & \texttt{+} & \texttt{S} & \texttt{none}\\
N465 & F3.103 & \texttt{T-{}-{}} & 2 & \texttt{ab\textbar{}cd\textbar{}cd} & \texttt{ac} & \texttt{ac} & \texttt{abcd} & \texttt{012} & \texttt{0} & \texttt{+} & \texttt{.} & \texttt{a@0}\\
N466 & F3.103 & \texttt{T-{}-{}} & 2 & \texttt{aa\textbar{}bb\textbar{}bb} & \texttt{ab} & \texttt{ab} & \texttt{ab} & \texttt{012} & \texttt{01} & \texttt{+} & \texttt{.} & \texttt{a@0}\\
N467 & F3.103 & \texttt{T-{}-{}} & 2 & \texttt{ab\textbar{}ab\textbar{}cd} & \texttt{ac} & \texttt{ac} & \texttt{abcd} & \texttt{012} & \texttt{0} & \texttt{+} & \texttt{.} & \texttt{c@2}\\
N468 & F3.103 & \texttt{T-{}-{}} & 2 & \texttt{aa\textbar{}aa\textbar{}bb} & \texttt{ab} & \texttt{ab} & \texttt{ab} & \texttt{012} & \texttt{01} & \texttt{+} & \texttt{.} & \texttt{b@2}\\
N469 & F3.103 & \texttt{T-{}-{}} & 2 & \texttt{ab\textbar{}ab\textbar{}ab} & \texttt{a} & \texttt{a} & \texttt{ab} & \texttt{012} & \texttt{0} & \texttt{+} & \texttt{S} & \texttt{none}\\
N470 & F3.103 & \texttt{T-{}-{}} & 2 & \texttt{ab\textbar{}ba\textbar{}ba} & \texttt{ab} & \texttt{ab} & \texttt{ab} & \texttt{012} & \texttt{01} & \texttt{+} & \texttt{.} & \texttt{none}\\
N471 & F3.103 & \texttt{T-{}-{}} & 2 & \texttt{ab\textbar{}ab\textbar{}ba} & \texttt{ab} & \texttt{ab} & \texttt{ab} & \texttt{012} & \texttt{01} & \texttt{+} & \texttt{.} & \texttt{none}\\
N472 & F3.103 & \texttt{T-{}-{}} & 2 & \texttt{aa\textbar{}aa\textbar{}aa} & \texttt{a} & \texttt{a} & \texttt{a} & \texttt{012} & \texttt{01} & \texttt{+} & \texttt{S} & \texttt{none}\\
\multicolumn{13}{@{}l}{\rule{0pt}{2.8ex}\footnotesize\bfseries $n=4$ worlds, $m=1$ individuals: 737 models}\\*
N473 & F4.342 & \texttt{T-{}4} & 1 & \texttt{a\textbar{}b\textbar{}b\textbar{}b} & \texttt{b} & \texttt{b} & \texttt{ab} & \texttt{123} & \texttt{-{}} & \texttt{-{}0} & \texttt{s} & \texttt{a@0}\\
N474 & F4.342 & \texttt{T-{}4} & 1 & \texttt{a\textbar{}a\textbar{}a\textbar{}a} & \texttt{a} & \texttt{a} & \texttt{a} & \texttt{0123} & \texttt{0} & \texttt{+} & \texttt{S} & \texttt{none}\\
N475 & F4.711 & \texttt{T-{}4} & 1 & \texttt{a\textbar{}b\textbar{}b\textbar{}b} & \texttt{b} & \texttt{b} & \texttt{ab} & \texttt{123} & \texttt{-{}} & \texttt{-{}0} & \texttt{s} & \texttt{a@0}\\
N476 & F4.711 & \texttt{T-{}4} & 1 & \texttt{a\textbar{}a\textbar{}a\textbar{}a} & \texttt{a} & \texttt{a} & \texttt{a} & \texttt{0123} & \texttt{0} & \texttt{+} & \texttt{S} & \texttt{none}\\
N477 & F4.713 & \texttt{T-{}-{}} & 1 & \texttt{a\textbar{}b\textbar{}a\textbar{}a} & \texttt{a} & \texttt{a} & \texttt{ab} & \texttt{023} & \texttt{0} & \texttt{-{}1} & \texttt{s} & \texttt{b@1}\\
N478 & F4.713 & \texttt{T-{}-{}} & 1 & \texttt{a\textbar{}a\textbar{}a\textbar{}a} & \texttt{a} & \texttt{a} & \texttt{a} & \texttt{0123} & \texttt{0} & \texttt{+} & \texttt{S} & \texttt{none}\\
N479 & F4.714 & \texttt{T-{}4} & 1 & \texttt{a\textbar{}b\textbar{}c\textbar{}c} & \texttt{c} & \texttt{c} & \texttt{abc} & \texttt{23} & \texttt{-{}} & \texttt{-{}0} & \texttt{s} & \texttt{a@0}\\
N480 & F4.714 & \texttt{T-{}4} & 1 & \texttt{a\textbar{}b\textbar{}a\textbar{}a} & \texttt{a} & \texttt{a} & \texttt{ab} & \texttt{023} & \texttt{0} & \texttt{-{}1} & \texttt{s} & \texttt{b@1}\\
N481 & F4.714 & \texttt{T-{}4} & 1 & \texttt{a\textbar{}a\textbar{}b\textbar{}b} & \texttt{b} & \texttt{b} & \texttt{ab} & \texttt{23} & \texttt{-{}} & \texttt{-{}0} & \texttt{s} & \texttt{a@0}\\
N482 & F4.714 & \texttt{T-{}4} & 1 & \texttt{a\textbar{}a\textbar{}a\textbar{}a} & \texttt{a} & \texttt{a} & \texttt{a} & \texttt{0123} & \texttt{0} & \texttt{+} & \texttt{S} & \texttt{none}\\
N483 & F4.715 & \texttt{T-{}4} & 1 & \texttt{a\textbar{}b\textbar{}c\textbar{}c} & \texttt{c} & \texttt{c} & \texttt{abc} & \texttt{23} & \texttt{-{}} & \texttt{-{}0} & \texttt{s} & \texttt{a@0}\\
N484 & F4.715 & \texttt{T-{}4} & 1 & \texttt{a\textbar{}b\textbar{}a\textbar{}a} & \texttt{a} & \texttt{a} & \texttt{ab} & \texttt{023} & \texttt{0} & \texttt{-{}1} & \texttt{s} & \texttt{b@1}\\
N485 & F4.715 & \texttt{T-{}4} & 1 & \texttt{a\textbar{}a\textbar{}b\textbar{}b} & \texttt{b} & \texttt{b} & \texttt{ab} & \texttt{23} & \texttt{-{}} & \texttt{-{}0} & \texttt{s} & \texttt{a@0}\\
N486 & F4.715 & \texttt{T-{}4} & 1 & \texttt{a\textbar{}a\textbar{}a\textbar{}a} & \texttt{a} & \texttt{a} & \texttt{a} & \texttt{0123} & \texttt{0} & \texttt{+} & \texttt{S} & \texttt{none}\\
N487 & F4.1218 & \texttt{T-{}4} & 1 & \texttt{a\textbar{}b\textbar{}b\textbar{}b} & \texttt{b} & \texttt{b} & \texttt{ab} & \texttt{123} & \texttt{-{}} & \texttt{-{}0} & \texttt{s} & \texttt{a@0}\\
N488 & F4.1218 & \texttt{T-{}4} & 1 & \texttt{a\textbar{}a\textbar{}a\textbar{}a} & \texttt{a} & \texttt{a} & \texttt{a} & \texttt{0123} & \texttt{0} & \texttt{+} & \texttt{S} & \texttt{none}\\
N489 & F4.1219 & \texttt{T-{}-{}} & 1 & \texttt{a\textbar{}a\textbar{}a\textbar{}a} & \texttt{a} & \texttt{a} & \texttt{a} & \texttt{0123} & \texttt{0} & \texttt{+} & \texttt{S} & \texttt{none}\\
N490 & F4.1220 & \texttt{T-{}-{}} & 1 & \texttt{a\textbar{}b\textbar{}a\textbar{}a} & \texttt{a} & \texttt{a} & \texttt{ab} & \texttt{023} & \texttt{0} & \texttt{-{}1} & \texttt{s} & \texttt{b@1}\\
N491 & F4.1220 & \texttt{T-{}-{}} & 1 & \texttt{a\textbar{}a\textbar{}a\textbar{}a} & \texttt{a} & \texttt{a} & \texttt{a} & \texttt{0123} & \texttt{0} & \texttt{+} & \texttt{S} & \texttt{none}\\
N492 & F4.1221 & \texttt{T-{}4} & 1 & \texttt{a\textbar{}b\textbar{}c\textbar{}c} & \texttt{c} & \texttt{c} & \texttt{abc} & \texttt{23} & \texttt{-{}} & \texttt{-{}0} & \texttt{s} & \texttt{a@0}\\
N493 & F4.1221 & \texttt{T-{}4} & 1 & \texttt{a\textbar{}b\textbar{}b\textbar{}b} & \texttt{b} & \texttt{b} & \texttt{ab} & \texttt{123} & \texttt{-{}} & \texttt{-{}0} & \texttt{s} & \texttt{a@0}\\
N494 & F4.1221 & \texttt{T-{}4} & 1 & \texttt{a\textbar{}a\textbar{}b\textbar{}b} & \texttt{b} & \texttt{b} & \texttt{ab} & \texttt{23} & \texttt{-{}} & \texttt{-{}0} & \texttt{s} & \texttt{a@0}\\
N495 & F4.1221 & \texttt{T-{}4} & 1 & \texttt{a\textbar{}a\textbar{}a\textbar{}a} & \texttt{a} & \texttt{a} & \texttt{a} & \texttt{0123} & \texttt{0} & \texttt{+} & \texttt{S} & \texttt{none}\\
N496 & F4.1222 & \texttt{T-{}-{}} & 1 & \texttt{a\textbar{}b\textbar{}c\textbar{}c} & \texttt{c} & \texttt{c} & \texttt{abc} & \texttt{23} & \texttt{-{}} & \texttt{-{}0} & \texttt{s} & \texttt{a@0}\\
N497 & F4.1222 & \texttt{T-{}-{}} & 1 & \texttt{a\textbar{}b\textbar{}b\textbar{}b} & \texttt{b} & \texttt{b} & \texttt{ab} & \texttt{123} & \texttt{-{}} & \texttt{-{}0} & \texttt{s} & \texttt{a@0}\\
N498 & F4.1222 & \texttt{T-{}-{}} & 1 & \texttt{a\textbar{}a\textbar{}b\textbar{}b} & \texttt{b} & \texttt{b} & \texttt{ab} & \texttt{23} & \texttt{-{}} & \texttt{-{}0} & \texttt{s} & \texttt{a@0}\\
N499 & F4.1222 & \texttt{T-{}-{}} & 1 & \texttt{a\textbar{}a\textbar{}a\textbar{}a} & \texttt{a} & \texttt{a} & \texttt{a} & \texttt{0123} & \texttt{0} & \texttt{+} & \texttt{S} & \texttt{none}\\
N500 & F4.1223 & \texttt{T-{}4} & 1 & \texttt{a\textbar{}b\textbar{}c\textbar{}c} & \texttt{c} & \texttt{c} & \texttt{abc} & \texttt{23} & \texttt{-{}} & \texttt{-{}0} & \texttt{s} & \texttt{a@0}\\
N501 & F4.1223 & \texttt{T-{}4} & 1 & \texttt{a\textbar{}b\textbar{}b\textbar{}b} & \texttt{b} & \texttt{b} & \texttt{ab} & \texttt{123} & \texttt{-{}} & \texttt{-{}0} & \texttt{s} & \texttt{a@0}\\
N502 & F4.1223 & \texttt{T-{}4} & 1 & \texttt{a\textbar{}b\textbar{}a\textbar{}a} & \texttt{a} & \texttt{a} & \texttt{ab} & \texttt{023} & \texttt{0} & \texttt{-{}1} & \texttt{s} & \texttt{b@1}\\
N503 & F4.1223 & \texttt{T-{}4} & 1 & \texttt{a\textbar{}a\textbar{}b\textbar{}b} & \texttt{b} & \texttt{b} & \texttt{ab} & \texttt{23} & \texttt{-{}} & \texttt{-{}0} & \texttt{s} & \texttt{a@0}\\
N504 & F4.1223 & \texttt{T-{}4} & 1 & \texttt{a\textbar{}a\textbar{}a\textbar{}a} & \texttt{a} & \texttt{a} & \texttt{a} & \texttt{0123} & \texttt{0} & \texttt{+} & \texttt{S} & \texttt{none}\\
N505 & F4.1224 & \texttt{T-{}-{}} & 1 & \texttt{a\textbar{}a\textbar{}b\textbar{}a} & \texttt{a} & \texttt{a} & \texttt{ab} & \texttt{013} & \texttt{0} & \texttt{-{}2} & \texttt{s} & \texttt{b@2}\\
N506 & F4.1224 & \texttt{T-{}-{}} & 1 & \texttt{a\textbar{}a\textbar{}a\textbar{}a} & \texttt{a} & \texttt{a} & \texttt{a} & \texttt{0123} & \texttt{0} & \texttt{+} & \texttt{S} & \texttt{none}\\
N507 & F4.1225 & \texttt{T-{}-{}} & 1 & \texttt{a\textbar{}b\textbar{}c\textbar{}a} & \texttt{a} & \texttt{a} & \texttt{abc} & \texttt{03} & \texttt{0} & \texttt{-{}1} & \texttt{s} & \texttt{b@1}\\
N508 & F4.1225 & \texttt{T-{}-{}} & 1 & \texttt{a\textbar{}a\textbar{}b\textbar{}a} & \texttt{a} & \texttt{a} & \texttt{ab} & \texttt{013} & \texttt{0} & \texttt{-{}2} & \texttt{s} & \texttt{b@2}\\
N509 & F4.1225 & \texttt{T-{}-{}} & 1 & \texttt{a\textbar{}b\textbar{}b\textbar{}a} & \texttt{a} & \texttt{a} & \texttt{ab} & \texttt{03} & \texttt{0} & \texttt{-{}1} & \texttt{s} & \texttt{b@1}\\
N510 & F4.1225 & \texttt{T-{}-{}} & 1 & \texttt{a\textbar{}a\textbar{}a\textbar{}a} & \texttt{a} & \texttt{a} & \texttt{a} & \texttt{0123} & \texttt{0} & \texttt{+} & \texttt{S} & \texttt{none}\\
N511 & F4.1226 & \texttt{T-{}-{}} & 1 & \texttt{a\textbar{}a\textbar{}a\textbar{}a} & \texttt{a} & \texttt{a} & \texttt{a} & \texttt{0123} & \texttt{0} & \texttt{+} & \texttt{S} & \texttt{none}\\
N512 & F4.1227 & \texttt{T-{}-{}} & 1 & \texttt{a\textbar{}b\textbar{}a\textbar{}a} & \texttt{a} & \texttt{a} & \texttt{ab} & \texttt{023} & \texttt{0} & \texttt{-{}1} & \texttt{s} & \texttt{b@1}\\
N513 & F4.1227 & \texttt{T-{}-{}} & 1 & \texttt{a\textbar{}a\textbar{}a\textbar{}a} & \texttt{a} & \texttt{a} & \texttt{a} & \texttt{0123} & \texttt{0} & \texttt{+} & \texttt{S} & \texttt{none}\\
N514 & F4.1228 & \texttt{T-{}4} & 1 & \texttt{a\textbar{}b\textbar{}a\textbar{}a} & \texttt{a} & \texttt{a} & \texttt{ab} & \texttt{023} & \texttt{0} & \texttt{-{}1} & \texttt{s} & \texttt{b@1}\\
N515 & F4.1228 & \texttt{T-{}4} & 1 & \texttt{a\textbar{}a\textbar{}a\textbar{}a} & \texttt{a} & \texttt{a} & \texttt{a} & \texttt{0123} & \texttt{0} & \texttt{+} & \texttt{S} & \texttt{none}\\
N516 & F4.1229 & \texttt{T-{}-{}} & 1 & \texttt{a\textbar{}b\textbar{}c\textbar{}a} & \texttt{a} & \texttt{a} & \texttt{abc} & \texttt{03} & \texttt{0} & \texttt{-{}1} & \texttt{s} & \texttt{b@1}\\
N517 & F4.1229 & \texttt{T-{}-{}} & 1 & \texttt{a\textbar{}b\textbar{}a\textbar{}a} & \texttt{a} & \texttt{a} & \texttt{ab} & \texttt{023} & \texttt{0} & \texttt{-{}1} & \texttt{s} & \texttt{b@1}\\
N518 & F4.1229 & \texttt{T-{}-{}} & 1 & \texttt{a\textbar{}a\textbar{}b\textbar{}a} & \texttt{a} & \texttt{a} & \texttt{ab} & \texttt{013} & \texttt{0} & \texttt{-{}2} & \texttt{s} & \texttt{b@2}\\
N519 & F4.1229 & \texttt{T-{}-{}} & 1 & \texttt{a\textbar{}b\textbar{}b\textbar{}a} & \texttt{a} & \texttt{a} & \texttt{ab} & \texttt{03} & \texttt{0} & \texttt{-{}1} & \texttt{s} & \texttt{b@1}\\
N520 & F4.1229 & \texttt{T-{}-{}} & 1 & \texttt{a\textbar{}a\textbar{}a\textbar{}a} & \texttt{a} & \texttt{a} & \texttt{a} & \texttt{0123} & \texttt{0} & \texttt{+} & \texttt{S} & \texttt{none}\\
N521 & F4.1230 & \texttt{T-{}4} & 1 & \texttt{a\textbar{}b\textbar{}c\textbar{}d} & \texttt{d} & \texttt{d} & \texttt{abcd} & \texttt{3} & \texttt{-{}} & \texttt{-{}0} & \texttt{s} & \texttt{a@0}\\
N522 & F4.1230 & \texttt{T-{}4} & 1 & \texttt{a\textbar{}b\textbar{}c\textbar{}a} & \texttt{a} & \texttt{a} & \texttt{abc} & \texttt{03} & \texttt{0} & \texttt{-{}1} & \texttt{s} & \texttt{b@1}\\
N523 & F4.1230 & \texttt{T-{}4} & 1 & \texttt{a\textbar{}a\textbar{}b\textbar{}c} & \texttt{c} & \texttt{c} & \texttt{abc} & \texttt{3} & \texttt{-{}} & \texttt{-{}0} & \texttt{s} & \texttt{b@2}\\
N524 & F4.1230 & \texttt{T-{}4} & 1 & \texttt{a\textbar{}a\textbar{}b\textbar{}a} & \texttt{a} & \texttt{a} & \texttt{ab} & \texttt{013} & \texttt{0} & \texttt{-{}2} & \texttt{s} & \texttt{b@2}\\
N525 & F4.1230 & \texttt{T-{}4} & 1 & \texttt{a\textbar{}b\textbar{}b\textbar{}a} & \texttt{a} & \texttt{a} & \texttt{ab} & \texttt{03} & \texttt{0} & \texttt{-{}1} & \texttt{s} & \texttt{b@1}\\
N526 & F4.1230 & \texttt{T-{}4} & 1 & \texttt{a\textbar{}a\textbar{}a\textbar{}b} & \texttt{b} & \texttt{b} & \texttt{ab} & \texttt{3} & \texttt{-{}} & \texttt{-{}0} & \texttt{s} & \texttt{a@0}\\
N527 & F4.1230 & \texttt{T-{}4} & 1 & \texttt{a\textbar{}a\textbar{}a\textbar{}a} & \texttt{a} & \texttt{a} & \texttt{a} & \texttt{0123} & \texttt{0} & \texttt{+} & \texttt{S} & \texttt{none}\\
N528 & F4.1788 & \texttt{T-{}-{}} & 1 & \texttt{a\textbar{}a\textbar{}a\textbar{}a} & \texttt{a} & \texttt{a} & \texttt{a} & \texttt{0123} & \texttt{0} & \texttt{+} & \texttt{S} & \texttt{none}\\
N529 & F4.1789 & \texttt{T-{}4} & 1 & \texttt{a\textbar{}b\textbar{}c\textbar{}c} & \texttt{c} & \texttt{c} & \texttt{abc} & \texttt{23} & \texttt{-{}} & \texttt{-{}0} & \texttt{s} & \texttt{a@0}\\
N530 & F4.1789 & \texttt{T-{}4} & 1 & \texttt{a\textbar{}b\textbar{}b\textbar{}b} & \texttt{b} & \texttt{b} & \texttt{ab} & \texttt{123} & \texttt{-{}} & \texttt{-{}0} & \texttt{s} & \texttt{a@0}\\
N531 & F4.1789 & \texttt{T-{}4} & 1 & \texttt{a\textbar{}a\textbar{}b\textbar{}b} & \texttt{b} & \texttt{b} & \texttt{ab} & \texttt{23} & \texttt{-{}} & \texttt{-{}0} & \texttt{s} & \texttt{a@0}\\
N532 & F4.1789 & \texttt{T-{}4} & 1 & \texttt{a\textbar{}a\textbar{}a\textbar{}a} & \texttt{a} & \texttt{a} & \texttt{a} & \texttt{0123} & \texttt{0} & \texttt{+} & \texttt{S} & \texttt{none}\\
N533 & F4.1790 & \texttt{T-{}4} & 1 & \texttt{a\textbar{}b\textbar{}c\textbar{}c} & \texttt{c} & \texttt{c} & \texttt{abc} & \texttt{23} & \texttt{-{}} & \texttt{-{}0} & \texttt{s} & \texttt{a@0}\\
N534 & F4.1790 & \texttt{T-{}4} & 1 & \texttt{a\textbar{}a\textbar{}b\textbar{}b} & \texttt{b} & \texttt{b} & \texttt{ab} & \texttt{23} & \texttt{-{}} & \texttt{-{}0} & \texttt{s} & \texttt{a@0}\\
N535 & F4.1790 & \texttt{T-{}4} & 1 & \texttt{a\textbar{}a\textbar{}a\textbar{}a} & \texttt{a} & \texttt{a} & \texttt{a} & \texttt{0123} & \texttt{0} & \texttt{+} & \texttt{S} & \texttt{none}\\
N536 & F4.1791 & \texttt{T-{}-{}} & 1 & \texttt{a\textbar{}b\textbar{}c\textbar{}c} & \texttt{c} & \texttt{c} & \texttt{abc} & \texttt{23} & \texttt{-{}} & \texttt{-{}0} & \texttt{s} & \texttt{a@0}\\
N537 & F4.1791 & \texttt{T-{}-{}} & 1 & \texttt{a\textbar{}a\textbar{}b\textbar{}b} & \texttt{b} & \texttt{b} & \texttt{ab} & \texttt{23} & \texttt{-{}} & \texttt{-{}0} & \texttt{s} & \texttt{a@0}\\
N538 & F4.1791 & \texttt{T-{}-{}} & 1 & \texttt{a\textbar{}a\textbar{}a\textbar{}a} & \texttt{a} & \texttt{a} & \texttt{a} & \texttt{0123} & \texttt{0} & \texttt{+} & \texttt{S} & \texttt{none}\\
N539 & F4.1792 & \texttt{T-{}-{}} & 1 & \texttt{a\textbar{}b\textbar{}c\textbar{}c} & \texttt{c} & \texttt{c} & \texttt{abc} & \texttt{23} & \texttt{-{}} & \texttt{-{}0} & \texttt{s} & \texttt{a@0}\\
N540 & F4.1792 & \texttt{T-{}-{}} & 1 & \texttt{a\textbar{}b\textbar{}a\textbar{}a} & \texttt{a} & \texttt{a} & \texttt{ab} & \texttt{023} & \texttt{0} & \texttt{-{}1} & \texttt{s} & \texttt{b@1}\\
N541 & F4.1792 & \texttt{T-{}-{}} & 1 & \texttt{a\textbar{}a\textbar{}b\textbar{}b} & \texttt{b} & \texttt{b} & \texttt{ab} & \texttt{23} & \texttt{-{}} & \texttt{-{}0} & \texttt{s} & \texttt{a@0}\\
N542 & F4.1792 & \texttt{T-{}-{}} & 1 & \texttt{a\textbar{}a\textbar{}a\textbar{}a} & \texttt{a} & \texttt{a} & \texttt{a} & \texttt{0123} & \texttt{0} & \texttt{+} & \texttt{S} & \texttt{none}\\
N543 & F4.1793 & \texttt{T-{}4} & 1 & \texttt{a\textbar{}b\textbar{}c\textbar{}c} & \texttt{c} & \texttt{c} & \texttt{abc} & \texttt{23} & \texttt{-{}} & \texttt{-{}0} & \texttt{s} & \texttt{a@0}\\
N544 & F4.1793 & \texttt{T-{}4} & 1 & \texttt{a\textbar{}b\textbar{}a\textbar{}a} & \texttt{a} & \texttt{a} & \texttt{ab} & \texttt{023} & \texttt{0} & \texttt{-{}1} & \texttt{s} & \texttt{b@1}\\
N545 & F4.1793 & \texttt{T-{}4} & 1 & \texttt{a\textbar{}a\textbar{}b\textbar{}b} & \texttt{b} & \texttt{b} & \texttt{ab} & \texttt{23} & \texttt{-{}} & \texttt{-{}0} & \texttt{s} & \texttt{a@0}\\
N546 & F4.1793 & \texttt{T-{}4} & 1 & \texttt{a\textbar{}a\textbar{}a\textbar{}a} & \texttt{a} & \texttt{a} & \texttt{a} & \texttt{0123} & \texttt{0} & \texttt{+} & \texttt{S} & \texttt{none}\\
N547 & F4.1795 & \texttt{T-{}-{}} & 1 & \texttt{a\textbar{}a\textbar{}b\textbar{}a} & \texttt{a} & \texttt{a} & \texttt{ab} & \texttt{013} & \texttt{0} & \texttt{-{}2} & \texttt{s} & \texttt{b@2}\\
N548 & F4.1795 & \texttt{T-{}-{}} & 1 & \texttt{a\textbar{}a\textbar{}a\textbar{}a} & \texttt{a} & \texttt{a} & \texttt{a} & \texttt{0123} & \texttt{0} & \texttt{+} & \texttt{S} & \texttt{none}\\
N549 & F4.1796 & \texttt{T-{}-{}} & 1 & \texttt{a\textbar{}a\textbar{}a\textbar{}a} & \texttt{a} & \texttt{a} & \texttt{a} & \texttt{0123} & \texttt{0} & \texttt{+} & \texttt{S} & \texttt{none}\\
N550 & F4.1797 & \texttt{T-{}-{}} & 1 & \texttt{a\textbar{}a\textbar{}a\textbar{}a} & \texttt{a} & \texttt{a} & \texttt{a} & \texttt{0123} & \texttt{0} & \texttt{+} & \texttt{S} & \texttt{none}\\
N551 & F4.1798 & \texttt{T-{}-{}} & 1 & \texttt{a\textbar{}b\textbar{}a\textbar{}a} & \texttt{a} & \texttt{a} & \texttt{ab} & \texttt{023} & \texttt{0} & \texttt{-{}1} & \texttt{s} & \texttt{b@1}\\
N552 & F4.1798 & \texttt{T-{}-{}} & 1 & \texttt{a\textbar{}a\textbar{}a\textbar{}a} & \texttt{a} & \texttt{a} & \texttt{a} & \texttt{0123} & \texttt{0} & \texttt{+} & \texttt{S} & \texttt{none}\\
N553 & F4.1799 & \texttt{T-{}4} & 1 & \texttt{a\textbar{}b\textbar{}a\textbar{}a} & \texttt{a} & \texttt{a} & \texttt{ab} & \texttt{023} & \texttt{0} & \texttt{-{}1} & \texttt{s} & \texttt{b@1}\\
N554 & F4.1799 & \texttt{T-{}4} & 1 & \texttt{a\textbar{}a\textbar{}a\textbar{}a} & \texttt{a} & \texttt{a} & \texttt{a} & \texttt{0123} & \texttt{0} & \texttt{+} & \texttt{S} & \texttt{none}\\
N555 & F4.1800 & \texttt{T-{}-{}} & 1 & \texttt{a\textbar{}b\textbar{}c\textbar{}a} & \texttt{a} & \texttt{a} & \texttt{abc} & \texttt{03} & \texttt{0} & \texttt{-{}1} & \texttt{s} & \texttt{b@1}\\
N556 & F4.1800 & \texttt{T-{}-{}} & 1 & \texttt{a\textbar{}b\textbar{}a\textbar{}a} & \texttt{a} & \texttt{a} & \texttt{ab} & \texttt{023} & \texttt{0} & \texttt{-{}1} & \texttt{s} & \texttt{b@1}\\
N557 & F4.1800 & \texttt{T-{}-{}} & 1 & \texttt{a\textbar{}b\textbar{}b\textbar{}a} & \texttt{a} & \texttt{a} & \texttt{ab} & \texttt{03} & \texttt{0} & \texttt{-{}1} & \texttt{s} & \texttt{b@1}\\
N558 & F4.1800 & \texttt{T-{}-{}} & 1 & \texttt{a\textbar{}a\textbar{}a\textbar{}a} & \texttt{a} & \texttt{a} & \texttt{a} & \texttt{0123} & \texttt{0} & \texttt{+} & \texttt{S} & \texttt{none}\\
N559 & F4.1801 & \texttt{T-{}-{}} & 1 & \texttt{a\textbar{}a\textbar{}a\textbar{}a} & \texttt{a} & \texttt{a} & \texttt{a} & \texttt{0123} & \texttt{0} & \texttt{+} & \texttt{S} & \texttt{none}\\
N560 & F4.1802 & \texttt{T-{}-{}} & 1 & \texttt{a\textbar{}a\textbar{}b\textbar{}a} & \texttt{a} & \texttt{a} & \texttt{ab} & \texttt{013} & \texttt{0} & \texttt{-{}2} & \texttt{s} & \texttt{b@2}\\
N561 & F4.1802 & \texttt{T-{}-{}} & 1 & \texttt{a\textbar{}a\textbar{}a\textbar{}a} & \texttt{a} & \texttt{a} & \texttt{a} & \texttt{0123} & \texttt{0} & \texttt{+} & \texttt{S} & \texttt{none}\\
N562 & F4.1803 & \texttt{T-{}-{}} & 1 & \texttt{a\textbar{}b\textbar{}a\textbar{}a} & \texttt{a} & \texttt{a} & \texttt{ab} & \texttt{023} & \texttt{0} & \texttt{-{}1} & \texttt{s} & \texttt{b@1}\\
N563 & F4.1803 & \texttt{T-{}-{}} & 1 & \texttt{a\textbar{}a\textbar{}a\textbar{}a} & \texttt{a} & \texttt{a} & \texttt{a} & \texttt{0123} & \texttt{0} & \texttt{+} & \texttt{S} & \texttt{none}\\
N564 & F4.1804 & \texttt{T-{}-{}} & 1 & \texttt{a\textbar{}b\textbar{}a\textbar{}a} & \texttt{a} & \texttt{a} & \texttt{ab} & \texttt{023} & \texttt{0} & \texttt{-{}1} & \texttt{s} & \texttt{b@1}\\
N565 & F4.1804 & \texttt{T-{}-{}} & 1 & \texttt{a\textbar{}a\textbar{}a\textbar{}a} & \texttt{a} & \texttt{a} & \texttt{a} & \texttt{0123} & \texttt{0} & \texttt{+} & \texttt{S} & \texttt{none}\\
N566 & F4.1805 & \texttt{T-{}-{}} & 1 & \texttt{a\textbar{}b\textbar{}c\textbar{}a} & \texttt{a} & \texttt{a} & \texttt{abc} & \texttt{03} & \texttt{0} & \texttt{-{}1} & \texttt{s} & \texttt{b@1}\\
N567 & F4.1805 & \texttt{T-{}-{}} & 1 & \texttt{a\textbar{}b\textbar{}a\textbar{}a} & \texttt{a} & \texttt{a} & \texttt{ab} & \texttt{023} & \texttt{0} & \texttt{-{}1} & \texttt{s} & \texttt{b@1}\\
N568 & F4.1805 & \texttt{T-{}-{}} & 1 & \texttt{a\textbar{}a\textbar{}b\textbar{}a} & \texttt{a} & \texttt{a} & \texttt{ab} & \texttt{013} & \texttt{0} & \texttt{-{}2} & \texttt{s} & \texttt{b@2}\\
N569 & F4.1805 & \texttt{T-{}-{}} & 1 & \texttt{a\textbar{}b\textbar{}b\textbar{}a} & \texttt{a} & \texttt{a} & \texttt{ab} & \texttt{03} & \texttt{0} & \texttt{-{}1} & \texttt{s} & \texttt{b@1}\\
N570 & F4.1805 & \texttt{T-{}-{}} & 1 & \texttt{a\textbar{}a\textbar{}a\textbar{}a} & \texttt{a} & \texttt{a} & \texttt{a} & \texttt{0123} & \texttt{0} & \texttt{+} & \texttt{S} & \texttt{none}\\
N571 & F4.1806 & \texttt{T-{}-{}} & 1 & \texttt{a\textbar{}b\textbar{}c\textbar{}a} & \texttt{a} & \texttt{a} & \texttt{abc} & \texttt{03} & \texttt{0} & \texttt{-{}1} & \texttt{s} & \texttt{b@1}\\
N572 & F4.1806 & \texttt{T-{}-{}} & 1 & \texttt{a\textbar{}b\textbar{}a\textbar{}a} & \texttt{a} & \texttt{a} & \texttt{ab} & \texttt{023} & \texttt{0} & \texttt{-{}1} & \texttt{s} & \texttt{b@1}\\
N573 & F4.1806 & \texttt{T-{}-{}} & 1 & \texttt{a\textbar{}b\textbar{}b\textbar{}a} & \texttt{a} & \texttt{a} & \texttt{ab} & \texttt{03} & \texttt{0} & \texttt{-{}1} & \texttt{s} & \texttt{b@1}\\
N574 & F4.1806 & \texttt{T-{}-{}} & 1 & \texttt{a\textbar{}a\textbar{}a\textbar{}a} & \texttt{a} & \texttt{a} & \texttt{a} & \texttt{0123} & \texttt{0} & \texttt{+} & \texttt{S} & \texttt{none}\\
N575 & F4.1807 & \texttt{T-{}-{}} & 1 & \texttt{a\textbar{}b\textbar{}c\textbar{}a} & \texttt{a} & \texttt{a} & \texttt{abc} & \texttt{03} & \texttt{0} & \texttt{-{}1} & \texttt{s} & \texttt{b@1}\\
N576 & F4.1807 & \texttt{T-{}-{}} & 1 & \texttt{a\textbar{}a\textbar{}b\textbar{}a} & \texttt{a} & \texttt{a} & \texttt{ab} & \texttt{013} & \texttt{0} & \texttt{-{}2} & \texttt{s} & \texttt{b@2}\\
N577 & F4.1807 & \texttt{T-{}-{}} & 1 & \texttt{a\textbar{}b\textbar{}b\textbar{}a} & \texttt{a} & \texttt{a} & \texttt{ab} & \texttt{03} & \texttt{0} & \texttt{-{}1} & \texttt{s} & \texttt{b@1}\\
N578 & F4.1807 & \texttt{T-{}-{}} & 1 & \texttt{a\textbar{}a\textbar{}a\textbar{}a} & \texttt{a} & \texttt{a} & \texttt{a} & \texttt{0123} & \texttt{0} & \texttt{+} & \texttt{S} & \texttt{none}\\
N579 & F4.1808 & \texttt{T-{}4} & 1 & \texttt{a\textbar{}b\textbar{}c\textbar{}d} & \texttt{d} & \texttt{d} & \texttt{abcd} & \texttt{3} & \texttt{-{}} & \texttt{-{}0} & \texttt{s} & \texttt{a@0}\\
N580 & F4.1808 & \texttt{T-{}4} & 1 & \texttt{a\textbar{}b\textbar{}c\textbar{}c} & \texttt{c} & \texttt{c} & \texttt{abc} & \texttt{23} & \texttt{-{}} & \texttt{-{}0} & \texttt{s} & \texttt{a@0}\\
N581 & F4.1808 & \texttt{T-{}4} & 1 & \texttt{a\textbar{}b\textbar{}c\textbar{}b} & \texttt{b} & \texttt{b} & \texttt{abc} & \texttt{13} & \texttt{-{}} & \texttt{-{}0} & \texttt{s} & \texttt{a@0}\\
N582 & F4.1808 & \texttt{T-{}4} & 1 & \texttt{a\textbar{}b\textbar{}b\textbar{}c} & \texttt{c} & \texttt{c} & \texttt{abc} & \texttt{3} & \texttt{-{}} & \texttt{-{}0} & \texttt{s} & \texttt{a@0}\\
N583 & F4.1808 & \texttt{T-{}4} & 1 & \texttt{a\textbar{}b\textbar{}b\textbar{}b} & \texttt{b} & \texttt{b} & \texttt{ab} & \texttt{123} & \texttt{-{}} & \texttt{-{}0} & \texttt{s} & \texttt{a@0}\\
N584 & F4.1808 & \texttt{T-{}4} & 1 & \texttt{a\textbar{}a\textbar{}b\textbar{}c} & \texttt{c} & \texttt{c} & \texttt{abc} & \texttt{3} & \texttt{-{}} & \texttt{-{}0} & \texttt{s} & \texttt{b@2}\\
N585 & F4.1808 & \texttt{T-{}4} & 1 & \texttt{a\textbar{}a\textbar{}b\textbar{}b} & \texttt{b} & \texttt{b} & \texttt{ab} & \texttt{23} & \texttt{-{}} & \texttt{-{}0} & \texttt{s} & \texttt{a@0}\\
N586 & F4.1808 & \texttt{T-{}4} & 1 & \texttt{a\textbar{}a\textbar{}b\textbar{}a} & \texttt{a} & \texttt{a} & \texttt{ab} & \texttt{013} & \texttt{0} & \texttt{-{}2} & \texttt{s} & \texttt{b@2}\\
N587 & F4.1808 & \texttt{T-{}4} & 1 & \texttt{a\textbar{}b\textbar{}a\textbar{}b} & \texttt{b} & \texttt{b} & \texttt{ab} & \texttt{13} & \texttt{-{}} & \texttt{-{}0} & \texttt{s} & \texttt{a@0}\\
N588 & F4.1808 & \texttt{T-{}4} & 1 & \texttt{a\textbar{}a\textbar{}a\textbar{}b} & \texttt{b} & \texttt{b} & \texttt{ab} & \texttt{3} & \texttt{-{}} & \texttt{-{}0} & \texttt{s} & \texttt{a@0}\\
N589 & F4.1808 & \texttt{T-{}4} & 1 & \texttt{a\textbar{}a\textbar{}a\textbar{}a} & \texttt{a} & \texttt{a} & \texttt{a} & \texttt{0123} & \texttt{0} & \texttt{+} & \texttt{S} & \texttt{none}\\
N590 & F4.1809 & \texttt{T-{}-{}} & 1 & \texttt{a\textbar{}b\textbar{}b\textbar{}b} & \texttt{ab} & \texttt{ab} & \texttt{ab} & \texttt{0123} & \texttt{0} & \texttt{+} & \texttt{.} & \texttt{a@0}\\
N591 & F4.1809 & \texttt{T-{}-{}} & 1 & \texttt{a\textbar{}a\textbar{}b\textbar{}c} & \texttt{a} & \texttt{a} & \texttt{abc} & \texttt{01} & \texttt{0} & \texttt{-{}2} & \texttt{s} & \texttt{b@2}\\
N592 & F4.1809 & \texttt{T-{}-{}} & 1 & \texttt{a\textbar{}a\textbar{}b\textbar{}b} & \texttt{a} & \texttt{a} & \texttt{ab} & \texttt{01} & \texttt{0} & \texttt{-{}2} & \texttt{s} & \texttt{b@2}\\
N593 & F4.1809 & \texttt{T-{}-{}} & 1 & \texttt{a\textbar{}a\textbar{}a\textbar{}b} & \texttt{a} & \texttt{a} & \texttt{ab} & \texttt{012} & \texttt{0} & \texttt{-{}3} & \texttt{s} & \texttt{b@3}\\
N594 & F4.1809 & \texttt{T-{}-{}} & 1 & \texttt{a\textbar{}a\textbar{}a\textbar{}a} & \texttt{a} & \texttt{a} & \texttt{a} & \texttt{0123} & \texttt{0} & \texttt{+} & \texttt{S} & \texttt{none}\\
N595 & F4.1810 & \texttt{T-{}-{}} & 1 & \texttt{a\textbar{}a\textbar{}a\textbar{}b} & \texttt{a} & \texttt{a} & \texttt{ab} & \texttt{012} & \texttt{0} & \texttt{-{}3} & \texttt{s} & \texttt{b@3}\\
N596 & F4.1810 & \texttt{T-{}-{}} & 1 & \texttt{a\textbar{}a\textbar{}a\textbar{}a} & \texttt{a} & \texttt{a} & \texttt{a} & \texttt{0123} & \texttt{0} & \texttt{+} & \texttt{S} & \texttt{none}\\
N597 & F4.1811 & \texttt{T-{}-{}} & 1 & \texttt{a\textbar{}b\textbar{}a\textbar{}c} & \texttt{a} & \texttt{a} & \texttt{abc} & \texttt{02} & \texttt{0} & \texttt{-{}1} & \texttt{s} & \texttt{b@1}\\
N598 & F4.1811 & \texttt{T-{}-{}} & 1 & \texttt{a\textbar{}a\textbar{}a\textbar{}b} & \texttt{a} & \texttt{a} & \texttt{ab} & \texttt{012} & \texttt{0} & \texttt{-{}3} & \texttt{s} & \texttt{b@3}\\
N599 & F4.1811 & \texttt{T-{}-{}} & 1 & \texttt{a\textbar{}a\textbar{}b\textbar{}b} & \texttt{ab} & \texttt{ab} & \texttt{ab} & \texttt{0123} & \texttt{0} & \texttt{+} & \texttt{.} & \texttt{a@0}\\
N600 & F4.1811 & \texttt{T-{}-{}} & 1 & \texttt{a\textbar{}b\textbar{}a\textbar{}b} & \texttt{a} & \texttt{a} & \texttt{ab} & \texttt{02} & \texttt{0} & \texttt{-{}1} & \texttt{s} & \texttt{b@1}\\
N601 & F4.1811 & \texttt{T-{}-{}} & 1 & \texttt{a\textbar{}a\textbar{}a\textbar{}a} & \texttt{a} & \texttt{a} & \texttt{a} & \texttt{0123} & \texttt{0} & \texttt{+} & \texttt{S} & \texttt{none}\\
N602 & F4.1812 & \texttt{T-{}-{}} & 1 & \texttt{a\textbar{}b\textbar{}a\textbar{}a} & \texttt{a} & \texttt{a} & \texttt{ab} & \texttt{023} & \texttt{0} & \texttt{-{}1} & \texttt{s} & \texttt{b@1}\\
N603 & F4.1812 & \texttt{T-{}-{}} & 1 & \texttt{a\textbar{}a\textbar{}b\textbar{}b} & \texttt{ab} & \texttt{ab} & \texttt{ab} & \texttt{0123} & \texttt{0} & \texttt{+} & \texttt{.} & \texttt{a@0}\\
N604 & F4.1812 & \texttt{T-{}-{}} & 1 & \texttt{a\textbar{}a\textbar{}a\textbar{}a} & \texttt{a} & \texttt{a} & \texttt{a} & \texttt{0123} & \texttt{0} & \texttt{+} & \texttt{S} & \texttt{none}\\
N605 & F4.1814 & \texttt{T-{}-{}} & 1 & \texttt{a\textbar{}a\textbar{}b\textbar{}b} & \texttt{ab} & \texttt{ab} & \texttt{ab} & \texttt{0123} & \texttt{0} & \texttt{+} & \texttt{.} & \texttt{a@0}\\
N606 & F4.1814 & \texttt{T-{}-{}} & 1 & \texttt{a\textbar{}a\textbar{}a\textbar{}a} & \texttt{a} & \texttt{a} & \texttt{a} & \texttt{0123} & \texttt{0} & \texttt{+} & \texttt{S} & \texttt{none}\\
N607 & F4.2292 & \texttt{T-{}-{}} & 1 & \texttt{a\textbar{}b\textbar{}c\textbar{}c} & \texttt{c} & \texttt{c} & \texttt{abc} & \texttt{23} & \texttt{-{}} & \texttt{-{}0} & \texttt{s} & \texttt{a@0}\\
N608 & F4.2292 & \texttt{T-{}-{}} & 1 & \texttt{a\textbar{}a\textbar{}b\textbar{}b} & \texttt{b} & \texttt{b} & \texttt{ab} & \texttt{23} & \texttt{-{}} & \texttt{-{}0} & \texttt{s} & \texttt{a@0}\\
N609 & F4.2292 & \texttt{T-{}-{}} & 1 & \texttt{a\textbar{}a\textbar{}a\textbar{}a} & \texttt{a} & \texttt{a} & \texttt{a} & \texttt{0123} & \texttt{0} & \texttt{+} & \texttt{S} & \texttt{none}\\
N610 & F4.2293 & \texttt{T-{}4} & 1 & \texttt{a\textbar{}b\textbar{}c\textbar{}c} & \texttt{c} & \texttt{c} & \texttt{abc} & \texttt{23} & \texttt{-{}} & \texttt{-{}0} & \texttt{s} & \texttt{a@0}\\
N611 & F4.2293 & \texttt{T-{}4} & 1 & \texttt{a\textbar{}b\textbar{}b\textbar{}b} & \texttt{b} & \texttt{b} & \texttt{ab} & \texttt{123} & \texttt{-{}} & \texttt{-{}0} & \texttt{s} & \texttt{a@0}\\
N612 & F4.2293 & \texttt{T-{}4} & 1 & \texttt{a\textbar{}a\textbar{}b\textbar{}b} & \texttt{b} & \texttt{b} & \texttt{ab} & \texttt{23} & \texttt{-{}} & \texttt{-{}0} & \texttt{s} & \texttt{a@0}\\
N613 & F4.2293 & \texttt{T-{}4} & 1 & \texttt{a\textbar{}a\textbar{}a\textbar{}a} & \texttt{a} & \texttt{a} & \texttt{a} & \texttt{0123} & \texttt{0} & \texttt{+} & \texttt{S} & \texttt{none}\\
N614 & F4.2294 & \texttt{T-{}-{}} & 1 & \texttt{a\textbar{}a\textbar{}a\textbar{}a} & \texttt{a} & \texttt{a} & \texttt{a} & \texttt{0123} & \texttt{0} & \texttt{+} & \texttt{S} & \texttt{none}\\
N615 & F4.2295 & \texttt{T-{}-{}} & 1 & \texttt{a\textbar{}a\textbar{}a\textbar{}a} & \texttt{a} & \texttt{a} & \texttt{a} & \texttt{0123} & \texttt{0} & \texttt{+} & \texttt{S} & \texttt{none}\\
N616 & F4.2296 & \texttt{T-{}-{}} & 1 & \texttt{a\textbar{}a\textbar{}a\textbar{}a} & \texttt{a} & \texttt{a} & \texttt{a} & \texttt{0123} & \texttt{0} & \texttt{+} & \texttt{S} & \texttt{none}\\
N617 & F4.2297 & \texttt{T-{}-{}} & 1 & \texttt{a\textbar{}a\textbar{}a\textbar{}a} & \texttt{a} & \texttt{a} & \texttt{a} & \texttt{0123} & \texttt{0} & \texttt{+} & \texttt{S} & \texttt{none}\\
N618 & F4.2298 & \texttt{T-{}-{}} & 1 & \texttt{a\textbar{}b\textbar{}a\textbar{}a} & \texttt{a} & \texttt{a} & \texttt{ab} & \texttt{023} & \texttt{0} & \texttt{-{}1} & \texttt{s} & \texttt{b@1}\\
N619 & F4.2298 & \texttt{T-{}-{}} & 1 & \texttt{a\textbar{}a\textbar{}a\textbar{}a} & \texttt{a} & \texttt{a} & \texttt{a} & \texttt{0123} & \texttt{0} & \texttt{+} & \texttt{S} & \texttt{none}\\
N620 & F4.2299 & \texttt{T-{}-{}} & 1 & \texttt{a\textbar{}a\textbar{}a\textbar{}a} & \texttt{a} & \texttt{a} & \texttt{a} & \texttt{0123} & \texttt{0} & \texttt{+} & \texttt{S} & \texttt{none}\\
N621 & F4.2300 & \texttt{T-{}-{}} & 1 & \texttt{a\textbar{}a\textbar{}a\textbar{}a} & \texttt{a} & \texttt{a} & \texttt{a} & \texttt{0123} & \texttt{0} & \texttt{+} & \texttt{S} & \texttt{none}\\
N622 & F4.2301 & \texttt{T-{}-{}} & 1 & \texttt{a\textbar{}a\textbar{}b\textbar{}a} & \texttt{a} & \texttt{a} & \texttt{ab} & \texttt{013} & \texttt{0} & \texttt{-{}2} & \texttt{s} & \texttt{b@2}\\
N623 & F4.2301 & \texttt{T-{}-{}} & 1 & \texttt{a\textbar{}a\textbar{}a\textbar{}a} & \texttt{a} & \texttt{a} & \texttt{a} & \texttt{0123} & \texttt{0} & \texttt{+} & \texttt{S} & \texttt{none}\\
N624 & F4.2302 & \texttt{T-{}-{}} & 1 & \texttt{a\textbar{}b\textbar{}a\textbar{}a} & \texttt{a} & \texttt{a} & \texttt{ab} & \texttt{023} & \texttt{0} & \texttt{-{}1} & \texttt{s} & \texttt{b@1}\\
N625 & F4.2302 & \texttt{T-{}-{}} & 1 & \texttt{a\textbar{}a\textbar{}a\textbar{}a} & \texttt{a} & \texttt{a} & \texttt{a} & \texttt{0123} & \texttt{0} & \texttt{+} & \texttt{S} & \texttt{none}\\
N626 & F4.2303 & \texttt{T-{}-{}} & 1 & \texttt{a\textbar{}a\textbar{}a\textbar{}a} & \texttt{a} & \texttt{a} & \texttt{a} & \texttt{0123} & \texttt{0} & \texttt{+} & \texttt{S} & \texttt{none}\\
N627 & F4.2304 & \texttt{T-{}-{}} & 1 & \texttt{a\textbar{}a\textbar{}a\textbar{}a} & \texttt{a} & \texttt{a} & \texttt{a} & \texttt{0123} & \texttt{0} & \texttt{+} & \texttt{S} & \texttt{none}\\
N628 & F4.2305 & \texttt{T-{}-{}} & 1 & \texttt{a\textbar{}b\textbar{}a\textbar{}a} & \texttt{a} & \texttt{a} & \texttt{ab} & \texttt{023} & \texttt{0} & \texttt{-{}1} & \texttt{s} & \texttt{b@1}\\
N629 & F4.2305 & \texttt{T-{}-{}} & 1 & \texttt{a\textbar{}a\textbar{}a\textbar{}a} & \texttt{a} & \texttt{a} & \texttt{a} & \texttt{0123} & \texttt{0} & \texttt{+} & \texttt{S} & \texttt{none}\\
N630 & F4.2306 & \texttt{T-{}4} & 1 & \texttt{a\textbar{}b\textbar{}a\textbar{}a} & \texttt{a} & \texttt{a} & \texttt{ab} & \texttt{023} & \texttt{0} & \texttt{-{}1} & \texttt{s} & \texttt{b@1}\\
N631 & F4.2306 & \texttt{T-{}4} & 1 & \texttt{a\textbar{}a\textbar{}a\textbar{}a} & \texttt{a} & \texttt{a} & \texttt{a} & \texttt{0123} & \texttt{0} & \texttt{+} & \texttt{S} & \texttt{none}\\
N632 & F4.2307 & \texttt{T-{}-{}} & 1 & \texttt{a\textbar{}b\textbar{}c\textbar{}a} & \texttt{a} & \texttt{a} & \texttt{abc} & \texttt{03} & \texttt{0} & \texttt{-{}1} & \texttt{s} & \texttt{b@1}\\
N633 & F4.2307 & \texttt{T-{}-{}} & 1 & \texttt{a\textbar{}b\textbar{}a\textbar{}a} & \texttt{a} & \texttt{a} & \texttt{ab} & \texttt{023} & \texttt{0} & \texttt{-{}1} & \texttt{s} & \texttt{b@1}\\
N634 & F4.2307 & \texttt{T-{}-{}} & 1 & \texttt{a\textbar{}b\textbar{}b\textbar{}a} & \texttt{a} & \texttt{a} & \texttt{ab} & \texttt{03} & \texttt{0} & \texttt{-{}1} & \texttt{s} & \texttt{b@1}\\
N635 & F4.2307 & \texttt{T-{}-{}} & 1 & \texttt{a\textbar{}a\textbar{}a\textbar{}a} & \texttt{a} & \texttt{a} & \texttt{a} & \texttt{0123} & \texttt{0} & \texttt{+} & \texttt{S} & \texttt{none}\\
N636 & F4.2308 & \texttt{T-{}-{}} & 1 & \texttt{a\textbar{}b\textbar{}c\textbar{}a} & \texttt{a} & \texttt{a} & \texttt{abc} & \texttt{03} & \texttt{0} & \texttt{-{}1} & \texttt{s} & \texttt{b@1}\\
N637 & F4.2308 & \texttt{T-{}-{}} & 1 & \texttt{a\textbar{}b\textbar{}b\textbar{}a} & \texttt{a} & \texttt{a} & \texttt{ab} & \texttt{03} & \texttt{0} & \texttt{-{}1} & \texttt{s} & \texttt{b@1}\\
N638 & F4.2308 & \texttt{T-{}-{}} & 1 & \texttt{a\textbar{}a\textbar{}a\textbar{}a} & \texttt{a} & \texttt{a} & \texttt{a} & \texttt{0123} & \texttt{0} & \texttt{+} & \texttt{S} & \texttt{none}\\
N639 & F4.2309 & \texttt{T-{}-{}} & 1 & \texttt{a\textbar{}b\textbar{}c\textbar{}b} & \texttt{b} & \texttt{b} & \texttt{abc} & \texttt{13} & \texttt{-{}} & \texttt{-{}0} & \texttt{s} & \texttt{a@0}\\
N640 & F4.2309 & \texttt{T-{}-{}} & 1 & \texttt{a\textbar{}a\textbar{}b\textbar{}a} & \texttt{a} & \texttt{a} & \texttt{ab} & \texttt{013} & \texttt{0} & \texttt{-{}2} & \texttt{s} & \texttt{b@2}\\
N641 & F4.2309 & \texttt{T-{}-{}} & 1 & \texttt{a\textbar{}b\textbar{}a\textbar{}b} & \texttt{b} & \texttt{b} & \texttt{ab} & \texttt{13} & \texttt{-{}} & \texttt{-{}0} & \texttt{s} & \texttt{a@0}\\
N642 & F4.2309 & \texttt{T-{}-{}} & 1 & \texttt{a\textbar{}a\textbar{}a\textbar{}a} & \texttt{a} & \texttt{a} & \texttt{a} & \texttt{0123} & \texttt{0} & \texttt{+} & \texttt{S} & \texttt{none}\\
N643 & F4.2310 & \texttt{T-{}-{}} & 1 & \texttt{a\textbar{}b\textbar{}c\textbar{}b} & \texttt{b} & \texttt{b} & \texttt{abc} & \texttt{13} & \texttt{-{}} & \texttt{-{}0} & \texttt{s} & \texttt{a@0}\\
N644 & F4.2310 & \texttt{T-{}-{}} & 1 & \texttt{a\textbar{}b\textbar{}a\textbar{}b} & \texttt{b} & \texttt{b} & \texttt{ab} & \texttt{13} & \texttt{-{}} & \texttt{-{}0} & \texttt{s} & \texttt{a@0}\\
N645 & F4.2310 & \texttt{T-{}-{}} & 1 & \texttt{a\textbar{}a\textbar{}a\textbar{}a} & \texttt{a} & \texttt{a} & \texttt{a} & \texttt{0123} & \texttt{0} & \texttt{+} & \texttt{S} & \texttt{none}\\
N646 & F4.2311 & \texttt{T-{}4} & 1 & \texttt{a\textbar{}b\textbar{}c\textbar{}d} & \texttt{d} & \texttt{d} & \texttt{abcd} & \texttt{3} & \texttt{-{}} & \texttt{-{}0} & \texttt{s} & \texttt{a@0}\\
N647 & F4.2311 & \texttt{T-{}4} & 1 & \texttt{a\textbar{}b\textbar{}c\textbar{}b} & \texttt{b} & \texttt{b} & \texttt{abc} & \texttt{13} & \texttt{-{}} & \texttt{-{}0} & \texttt{s} & \texttt{a@0}\\
N648 & F4.2311 & \texttt{T-{}4} & 1 & \texttt{a\textbar{}b\textbar{}b\textbar{}c} & \texttt{c} & \texttt{c} & \texttt{abc} & \texttt{3} & \texttt{-{}} & \texttt{-{}0} & \texttt{s} & \texttt{a@0}\\
N649 & F4.2311 & \texttt{T-{}4} & 1 & \texttt{a\textbar{}b\textbar{}b\textbar{}b} & \texttt{b} & \texttt{b} & \texttt{ab} & \texttt{123} & \texttt{-{}} & \texttt{-{}0} & \texttt{s} & \texttt{a@0}\\
N650 & F4.2311 & \texttt{T-{}4} & 1 & \texttt{a\textbar{}b\textbar{}a\textbar{}b} & \texttt{b} & \texttt{b} & \texttt{ab} & \texttt{13} & \texttt{-{}} & \texttt{-{}0} & \texttt{s} & \texttt{a@0}\\
N651 & F4.2311 & \texttt{T-{}4} & 1 & \texttt{a\textbar{}a\textbar{}a\textbar{}b} & \texttt{b} & \texttt{b} & \texttt{ab} & \texttt{3} & \texttt{-{}} & \texttt{-{}0} & \texttt{s} & \texttt{a@0}\\
N652 & F4.2311 & \texttt{T-{}4} & 1 & \texttt{a\textbar{}a\textbar{}a\textbar{}a} & \texttt{a} & \texttt{a} & \texttt{a} & \texttt{0123} & \texttt{0} & \texttt{+} & \texttt{S} & \texttt{none}\\
N653 & F4.2312 & \texttt{T-{}4} & 1 & \texttt{a\textbar{}b\textbar{}c\textbar{}d} & \texttt{d} & \texttt{d} & \texttt{abcd} & \texttt{3} & \texttt{-{}} & \texttt{-{}0} & \texttt{s} & \texttt{a@0}\\
N654 & F4.2312 & \texttt{T-{}4} & 1 & \texttt{a\textbar{}b\textbar{}c\textbar{}c} & \texttt{c} & \texttt{c} & \texttt{abc} & \texttt{23} & \texttt{-{}} & \texttt{-{}0} & \texttt{s} & \texttt{a@0}\\
N655 & F4.2312 & \texttt{T-{}4} & 1 & \texttt{a\textbar{}a\textbar{}b\textbar{}c} & \texttt{c} & \texttt{c} & \texttt{abc} & \texttt{3} & \texttt{-{}} & \texttt{-{}0} & \texttt{s} & \texttt{b@2}\\
N656 & F4.2312 & \texttt{T-{}4} & 1 & \texttt{a\textbar{}a\textbar{}b\textbar{}b} & \texttt{b} & \texttt{b} & \texttt{ab} & \texttt{23} & \texttt{-{}} & \texttt{-{}0} & \texttt{s} & \texttt{a@0}\\
N657 & F4.2312 & \texttt{T-{}4} & 1 & \texttt{a\textbar{}a\textbar{}b\textbar{}a} & \texttt{a} & \texttt{a} & \texttt{ab} & \texttt{013} & \texttt{0} & \texttt{-{}2} & \texttt{s} & \texttt{b@2}\\
N658 & F4.2312 & \texttt{T-{}4} & 1 & \texttt{a\textbar{}a\textbar{}a\textbar{}b} & \texttt{b} & \texttt{b} & \texttt{ab} & \texttt{3} & \texttt{-{}} & \texttt{-{}0} & \texttt{s} & \texttt{a@0}\\
N659 & F4.2312 & \texttt{T-{}4} & 1 & \texttt{a\textbar{}a\textbar{}a\textbar{}a} & \texttt{a} & \texttt{a} & \texttt{a} & \texttt{0123} & \texttt{0} & \texttt{+} & \texttt{S} & \texttt{none}\\
N660 & F4.2313 & \texttt{T-{}-{}} & 1 & \texttt{a\textbar{}b\textbar{}c\textbar{}d} & \texttt{d} & \texttt{d} & \texttt{abcd} & \texttt{3} & \texttt{-{}} & \texttt{-{}0} & \texttt{s} & \texttt{a@0}\\
N661 & F4.2313 & \texttt{T-{}-{}} & 1 & \texttt{a\textbar{}b\textbar{}c\textbar{}c} & \texttt{c} & \texttt{c} & \texttt{abc} & \texttt{23} & \texttt{-{}} & \texttt{-{}0} & \texttt{s} & \texttt{a@0}\\
N662 & F4.2313 & \texttt{T-{}-{}} & 1 & \texttt{a\textbar{}b\textbar{}a\textbar{}c} & \texttt{c} & \texttt{c} & \texttt{abc} & \texttt{3} & \texttt{-{}} & \texttt{-{}0} & \texttt{s} & \texttt{b@1}\\
N663 & F4.2313 & \texttt{T-{}-{}} & 1 & \texttt{a\textbar{}b\textbar{}a\textbar{}a} & \texttt{a} & \texttt{a} & \texttt{ab} & \texttt{023} & \texttt{0} & \texttt{-{}1} & \texttt{s} & \texttt{b@1}\\
N664 & F4.2313 & \texttt{T-{}-{}} & 1 & \texttt{a\textbar{}a\textbar{}b\textbar{}b} & \texttt{b} & \texttt{b} & \texttt{ab} & \texttt{23} & \texttt{-{}} & \texttt{-{}0} & \texttt{s} & \texttt{a@0}\\
N665 & F4.2313 & \texttt{T-{}-{}} & 1 & \texttt{a\textbar{}a\textbar{}a\textbar{}b} & \texttt{b} & \texttt{b} & \texttt{ab} & \texttt{3} & \texttt{-{}} & \texttt{-{}0} & \texttt{s} & \texttt{a@0}\\
N666 & F4.2313 & \texttt{T-{}-{}} & 1 & \texttt{a\textbar{}a\textbar{}a\textbar{}a} & \texttt{a} & \texttt{a} & \texttt{a} & \texttt{0123} & \texttt{0} & \texttt{+} & \texttt{S} & \texttt{none}\\
N667 & F4.2314 & \texttt{T-{}4} & 1 & \texttt{a\textbar{}b\textbar{}c\textbar{}d} & \texttt{d} & \texttt{d} & \texttt{abcd} & \texttt{3} & \texttt{-{}} & \texttt{-{}0} & \texttt{s} & \texttt{a@0}\\
N668 & F4.2314 & \texttt{T-{}4} & 1 & \texttt{a\textbar{}b\textbar{}c\textbar{}c} & \texttt{c} & \texttt{c} & \texttt{abc} & \texttt{23} & \texttt{-{}} & \texttt{-{}0} & \texttt{s} & \texttt{a@0}\\
N669 & F4.2314 & \texttt{T-{}4} & 1 & \texttt{a\textbar{}b\textbar{}a\textbar{}c} & \texttt{c} & \texttt{c} & \texttt{abc} & \texttt{3} & \texttt{-{}} & \texttt{-{}0} & \texttt{s} & \texttt{b@1}\\
N670 & F4.2314 & \texttt{T-{}4} & 1 & \texttt{a\textbar{}b\textbar{}a\textbar{}a} & \texttt{a} & \texttt{a} & \texttt{ab} & \texttt{023} & \texttt{0} & \texttt{-{}1} & \texttt{s} & \texttt{b@1}\\
N671 & F4.2314 & \texttt{T-{}4} & 1 & \texttt{a\textbar{}a\textbar{}b\textbar{}c} & \texttt{c} & \texttt{c} & \texttt{abc} & \texttt{3} & \texttt{-{}} & \texttt{-{}0} & \texttt{s} & \texttt{b@2}\\
N672 & F4.2314 & \texttt{T-{}4} & 1 & \texttt{a\textbar{}a\textbar{}b\textbar{}b} & \texttt{b} & \texttt{b} & \texttt{ab} & \texttt{23} & \texttt{-{}} & \texttt{-{}0} & \texttt{s} & \texttt{a@0}\\
N673 & F4.2314 & \texttt{T-{}4} & 1 & \texttt{a\textbar{}a\textbar{}a\textbar{}b} & \texttt{b} & \texttt{b} & \texttt{ab} & \texttt{3} & \texttt{-{}} & \texttt{-{}0} & \texttt{s} & \texttt{a@0}\\
N674 & F4.2314 & \texttt{T-{}4} & 1 & \texttt{a\textbar{}a\textbar{}a\textbar{}a} & \texttt{a} & \texttt{a} & \texttt{a} & \texttt{0123} & \texttt{0} & \texttt{+} & \texttt{S} & \texttt{none}\\
N675 & F4.2315 & \texttt{T-{}-{}} & 1 & \texttt{a\textbar{}b\textbar{}b\textbar{}b} & \texttt{ab} & \texttt{ab} & \texttt{ab} & \texttt{0123} & \texttt{0} & \texttt{+} & \texttt{.} & \texttt{a@0}\\
N676 & F4.2315 & \texttt{T-{}-{}} & 1 & \texttt{a\textbar{}a\textbar{}a\textbar{}b} & \texttt{a} & \texttt{a} & \texttt{ab} & \texttt{012} & \texttt{0} & \texttt{-{}3} & \texttt{s} & \texttt{b@3}\\
N677 & F4.2315 & \texttt{T-{}-{}} & 1 & \texttt{a\textbar{}a\textbar{}a\textbar{}a} & \texttt{a} & \texttt{a} & \texttt{a} & \texttt{0123} & \texttt{0} & \texttt{+} & \texttt{S} & \texttt{none}\\
N678 & F4.2316 & \texttt{T-{}-{}} & 1 & \texttt{a\textbar{}a\textbar{}b\textbar{}b} & \texttt{ab} & \texttt{ab} & \texttt{ab} & \texttt{0123} & \texttt{0} & \texttt{+} & \texttt{.} & \texttt{a@0}\\
N679 & F4.2316 & \texttt{T-{}-{}} & 1 & \texttt{a\textbar{}a\textbar{}a\textbar{}b} & \texttt{a} & \texttt{a} & \texttt{ab} & \texttt{012} & \texttt{0} & \texttt{-{}3} & \texttt{s} & \texttt{b@3}\\
N680 & F4.2316 & \texttt{T-{}-{}} & 1 & \texttt{a\textbar{}a\textbar{}a\textbar{}a} & \texttt{a} & \texttt{a} & \texttt{a} & \texttt{0123} & \texttt{0} & \texttt{+} & \texttt{S} & \texttt{none}\\
N681 & F4.2317 & \texttt{T-{}-{}} & 1 & \texttt{a\textbar{}a\textbar{}b\textbar{}b} & \texttt{ab} & \texttt{ab} & \texttt{ab} & \texttt{0123} & \texttt{0} & \texttt{+} & \texttt{.} & \texttt{a@0}\\
N682 & F4.2317 & \texttt{T-{}-{}} & 1 & \texttt{a\textbar{}a\textbar{}a\textbar{}a} & \texttt{a} & \texttt{a} & \texttt{a} & \texttt{0123} & \texttt{0} & \texttt{+} & \texttt{S} & \texttt{none}\\
N683 & F4.2318 & \texttt{T-{}-{}} & 1 & \texttt{a\textbar{}b\textbar{}a\textbar{}a} & \texttt{a} & \texttt{a} & \texttt{ab} & \texttt{023} & \texttt{0} & \texttt{-{}1} & \texttt{s} & \texttt{b@1}\\
N684 & F4.2318 & \texttt{T-{}-{}} & 1 & \texttt{a\textbar{}a\textbar{}b\textbar{}b} & \texttt{ab} & \texttt{ab} & \texttt{ab} & \texttt{0123} & \texttt{0} & \texttt{+} & \texttt{.} & \texttt{a@0}\\
N685 & F4.2318 & \texttt{T-{}-{}} & 1 & \texttt{a\textbar{}a\textbar{}a\textbar{}a} & \texttt{a} & \texttt{a} & \texttt{a} & \texttt{0123} & \texttt{0} & \texttt{+} & \texttt{S} & \texttt{none}\\
N686 & F4.2319 & \texttt{T-{}-{}} & 1 & \texttt{a\textbar{}b\textbar{}b\textbar{}b} & \texttt{ab} & \texttt{ab} & \texttt{ab} & \texttt{0123} & \texttt{0} & \texttt{+} & \texttt{.} & \texttt{a@0}\\
N687 & F4.2319 & \texttt{T-{}-{}} & 1 & \texttt{a\textbar{}a\textbar{}a\textbar{}b} & \texttt{a} & \texttt{a} & \texttt{ab} & \texttt{012} & \texttt{0} & \texttt{-{}3} & \texttt{s} & \texttt{b@3}\\
N688 & F4.2319 & \texttt{T-{}-{}} & 1 & \texttt{a\textbar{}a\textbar{}a\textbar{}a} & \texttt{a} & \texttt{a} & \texttt{a} & \texttt{0123} & \texttt{0} & \texttt{+} & \texttt{S} & \texttt{none}\\
N689 & F4.2320 & \texttt{T-{}-{}} & 1 & \texttt{a\textbar{}b\textbar{}c\textbar{}c} & \texttt{ac} & \texttt{ac} & \texttt{abc} & \texttt{023} & \texttt{0} & \texttt{-{}1} & \texttt{.} & \texttt{a@0}\\
N690 & F4.2320 & \texttt{T-{}-{}} & 1 & \texttt{a\textbar{}b\textbar{}b\textbar{}b} & \texttt{ab} & \texttt{ab} & \texttt{ab} & \texttt{0123} & \texttt{0} & \texttt{+} & \texttt{.} & \texttt{a@0}\\
N691 & F4.2320 & \texttt{T-{}-{}} & 1 & \texttt{a\textbar{}b\textbar{}a\textbar{}c} & \texttt{a} & \texttt{a} & \texttt{abc} & \texttt{02} & \texttt{0} & \texttt{-{}1} & \texttt{s} & \texttt{b@1}\\
N692 & F4.2320 & \texttt{T-{}-{}} & 1 & \texttt{a\textbar{}b\textbar{}a\textbar{}a} & \texttt{a} & \texttt{a} & \texttt{ab} & \texttt{023} & \texttt{0} & \texttt{-{}1} & \texttt{s} & \texttt{b@1}\\
N693 & F4.2320 & \texttt{T-{}-{}} & 1 & \texttt{a\textbar{}a\textbar{}b\textbar{}b} & \texttt{ab} & \texttt{ab} & \texttt{ab} & \texttt{0123} & \texttt{0} & \texttt{+} & \texttt{.} & \texttt{a@0}\\
N694 & F4.2320 & \texttt{T-{}-{}} & 1 & \texttt{a\textbar{}b\textbar{}a\textbar{}b} & \texttt{a} & \texttt{a} & \texttt{ab} & \texttt{02} & \texttt{0} & \texttt{-{}1} & \texttt{s} & \texttt{b@1}\\
N695 & F4.2320 & \texttt{T-{}-{}} & 1 & \texttt{a\textbar{}a\textbar{}a\textbar{}b} & \texttt{a} & \texttt{a} & \texttt{ab} & \texttt{012} & \texttt{0} & \texttt{-{}3} & \texttt{s} & \texttt{b@3}\\
N696 & F4.2320 & \texttt{T-{}-{}} & 1 & \texttt{a\textbar{}a\textbar{}a\textbar{}a} & \texttt{a} & \texttt{a} & \texttt{a} & \texttt{0123} & \texttt{0} & \texttt{+} & \texttt{S} & \texttt{none}\\
N697 & F4.2321 & \texttt{T-{}-{}} & 1 & \texttt{a\textbar{}b\textbar{}c\textbar{}c} & \texttt{ac} & \texttt{ac} & \texttt{abc} & \texttt{023} & \texttt{0} & \texttt{-{}1} & \texttt{.} & \texttt{a@0}\\
N698 & F4.2321 & \texttt{T-{}-{}} & 1 & \texttt{a\textbar{}b\textbar{}b\textbar{}b} & \texttt{ab} & \texttt{ab} & \texttt{ab} & \texttt{0123} & \texttt{0} & \texttt{+} & \texttt{.} & \texttt{a@0}\\
N699 & F4.2321 & \texttt{T-{}-{}} & 1 & \texttt{a\textbar{}b\textbar{}c\textbar{}a} & \texttt{a} & \texttt{a} & \texttt{abc} & \texttt{03} & \texttt{0} & \texttt{-{}1} & \texttt{s} & \texttt{b@1}\\
N700 & F4.2321 & \texttt{T-{}-{}} & 1 & \texttt{a\textbar{}b\textbar{}a\textbar{}a} & \texttt{a} & \texttt{a} & \texttt{ab} & \texttt{023} & \texttt{0} & \texttt{-{}1} & \texttt{s} & \texttt{b@1}\\
N701 & F4.2321 & \texttt{T-{}-{}} & 1 & \texttt{a\textbar{}a\textbar{}b\textbar{}b} & \texttt{ab} & \texttt{ab} & \texttt{ab} & \texttt{0123} & \texttt{0} & \texttt{+} & \texttt{.} & \texttt{a@0}\\
N702 & F4.2321 & \texttt{T-{}-{}} & 1 & \texttt{a\textbar{}b\textbar{}b\textbar{}a} & \texttt{a} & \texttt{a} & \texttt{ab} & \texttt{03} & \texttt{0} & \texttt{-{}1} & \texttt{s} & \texttt{b@1}\\
N703 & F4.2321 & \texttt{T-{}-{}} & 1 & \texttt{a\textbar{}a\textbar{}a\textbar{}a} & \texttt{a} & \texttt{a} & \texttt{a} & \texttt{0123} & \texttt{0} & \texttt{+} & \texttt{S} & \texttt{none}\\
N704 & F4.2322 & \texttt{T-{}-{}} & 1 & \texttt{a\textbar{}a\textbar{}a\textbar{}a} & \texttt{a} & \texttt{a} & \texttt{a} & \texttt{0123} & \texttt{0} & \texttt{+} & \texttt{S} & \texttt{none}\\
N705 & F4.2323 & \texttt{T-{}-{}} & 1 & \texttt{a\textbar{}b\textbar{}a\textbar{}c} & \texttt{a} & \texttt{a} & \texttt{abc} & \texttt{02} & \texttt{0} & \texttt{-{}1} & \texttt{s} & \texttt{b@1}\\
N706 & F4.2323 & \texttt{T-{}-{}} & 1 & \texttt{a\textbar{}b\textbar{}a\textbar{}a} & \texttt{a} & \texttt{a} & \texttt{ab} & \texttt{023} & \texttt{0} & \texttt{-{}1} & \texttt{s} & \texttt{b@1}\\
N707 & F4.2323 & \texttt{T-{}-{}} & 1 & \texttt{a\textbar{}a\textbar{}b\textbar{}b} & \texttt{ab} & \texttt{ab} & \texttt{ab} & \texttt{0123} & \texttt{0} & \texttt{+} & \texttt{.} & \texttt{a@0}\\
N708 & F4.2323 & \texttt{T-{}-{}} & 1 & \texttt{a\textbar{}b\textbar{}a\textbar{}b} & \texttt{a} & \texttt{a} & \texttt{ab} & \texttt{02} & \texttt{0} & \texttt{-{}1} & \texttt{s} & \texttt{b@1}\\
N709 & F4.2323 & \texttt{T-{}-{}} & 1 & \texttt{a\textbar{}a\textbar{}a\textbar{}a} & \texttt{a} & \texttt{a} & \texttt{a} & \texttt{0123} & \texttt{0} & \texttt{+} & \texttt{S} & \texttt{none}\\
N710 & F4.2324 & \texttt{T-{}-{}} & 1 & \texttt{a\textbar{}b\textbar{}a\textbar{}b} & \texttt{ab} & \texttt{ab} & \texttt{ab} & \texttt{0123} & \texttt{0} & \texttt{+} & \texttt{.} & \texttt{a@0}\\
N711 & F4.2324 & \texttt{T-{}-{}} & 1 & \texttt{a\textbar{}a\textbar{}a\textbar{}b} & \texttt{a} & \texttt{a} & \texttt{ab} & \texttt{012} & \texttt{0} & \texttt{-{}3} & \texttt{s} & \texttt{b@3}\\
N712 & F4.2324 & \texttt{T-{}-{}} & 1 & \texttt{a\textbar{}a\textbar{}a\textbar{}a} & \texttt{a} & \texttt{a} & \texttt{a} & \texttt{0123} & \texttt{0} & \texttt{+} & \texttt{S} & \texttt{none}\\
N713 & F4.2325 & \texttt{T-{}-{}} & 1 & \texttt{a\textbar{}a\textbar{}b\textbar{}b} & \texttt{ab} & \texttt{ab} & \texttt{ab} & \texttt{0123} & \texttt{0} & \texttt{+} & \texttt{.} & \texttt{a@0}\\
N714 & F4.2325 & \texttt{T-{}-{}} & 1 & \texttt{a\textbar{}b\textbar{}a\textbar{}b} & \texttt{ab} & \texttt{ab} & \texttt{ab} & \texttt{0123} & \texttt{0} & \texttt{+} & \texttt{.} & \texttt{a@0}\\
N715 & F4.2325 & \texttt{T-{}-{}} & 1 & \texttt{a\textbar{}a\textbar{}a\textbar{}a} & \texttt{a} & \texttt{a} & \texttt{a} & \texttt{0123} & \texttt{0} & \texttt{+} & \texttt{S} & \texttt{none}\\
N716 & F4.2326 & \texttt{T-{}-{}} & 1 & \texttt{a\textbar{}a\textbar{}b\textbar{}b} & \texttt{ab} & \texttt{ab} & \texttt{ab} & \texttt{0123} & \texttt{0} & \texttt{+} & \texttt{.} & \texttt{a@0}\\
N717 & F4.2326 & \texttt{T-{}-{}} & 1 & \texttt{a\textbar{}a\textbar{}a\textbar{}b} & \texttt{a} & \texttt{a} & \texttt{ab} & \texttt{012} & \texttt{0} & \texttt{-{}3} & \texttt{s} & \texttt{b@3}\\
N718 & F4.2326 & \texttt{T-{}-{}} & 1 & \texttt{a\textbar{}a\textbar{}a\textbar{}a} & \texttt{a} & \texttt{a} & \texttt{a} & \texttt{0123} & \texttt{0} & \texttt{+} & \texttt{S} & \texttt{none}\\
N719 & F4.2327 & \texttt{T-{}-{}} & 1 & \texttt{a\textbar{}a\textbar{}b\textbar{}b} & \texttt{ab} & \texttt{ab} & \texttt{ab} & \texttt{0123} & \texttt{0} & \texttt{+} & \texttt{.} & \texttt{a@0}\\
N720 & F4.2327 & \texttt{T-{}-{}} & 1 & \texttt{a\textbar{}a\textbar{}a\textbar{}a} & \texttt{a} & \texttt{a} & \texttt{a} & \texttt{0123} & \texttt{0} & \texttt{+} & \texttt{S} & \texttt{none}\\
N721 & F4.2328 & \texttt{T-{}-{}} & 1 & \texttt{a\textbar{}a\textbar{}b\textbar{}b} & \texttt{ab} & \texttt{ab} & \texttt{ab} & \texttt{0123} & \texttt{0} & \texttt{+} & \texttt{.} & \texttt{a@0}\\
N722 & F4.2328 & \texttt{T-{}-{}} & 1 & \texttt{a\textbar{}a\textbar{}a\textbar{}a} & \texttt{a} & \texttt{a} & \texttt{a} & \texttt{0123} & \texttt{0} & \texttt{+} & \texttt{S} & \texttt{none}\\
N723 & F4.2329 & \texttt{T-{}-{}} & 1 & \texttt{a\textbar{}a\textbar{}b\textbar{}a} & \texttt{a} & \texttt{a} & \texttt{ab} & \texttt{013} & \texttt{0} & \texttt{-{}2} & \texttt{s} & \texttt{b@2}\\
N724 & F4.2329 & \texttt{T-{}-{}} & 1 & \texttt{a\textbar{}a\textbar{}a\textbar{}a} & \texttt{a} & \texttt{a} & \texttt{a} & \texttt{0123} & \texttt{0} & \texttt{+} & \texttt{S} & \texttt{none}\\
N725 & F4.2655 & \texttt{T-{}4} & 1 & \texttt{a\textbar{}b\textbar{}c\textbar{}c} & \texttt{c} & \texttt{c} & \texttt{abc} & \texttt{23} & \texttt{-{}} & \texttt{-{}0} & \texttt{s} & \texttt{a@0}\\
N726 & F4.2655 & \texttt{T-{}4} & 1 & \texttt{a\textbar{}a\textbar{}b\textbar{}b} & \texttt{b} & \texttt{b} & \texttt{ab} & \texttt{23} & \texttt{-{}} & \texttt{-{}0} & \texttt{s} & \texttt{a@0}\\
N727 & F4.2655 & \texttt{T-{}4} & 1 & \texttt{a\textbar{}a\textbar{}a\textbar{}a} & \texttt{a} & \texttt{a} & \texttt{a} & \texttt{0123} & \texttt{0} & \texttt{+} & \texttt{S} & \texttt{none}\\
N728 & F4.2656 & \texttt{T-{}-{}} & 1 & \texttt{a\textbar{}a\textbar{}a\textbar{}a} & \texttt{a} & \texttt{a} & \texttt{a} & \texttt{0123} & \texttt{0} & \texttt{+} & \texttt{S} & \texttt{none}\\
N729 & F4.2657 & \texttt{T-{}-{}} & 1 & \texttt{a\textbar{}a\textbar{}a\textbar{}a} & \texttt{a} & \texttt{a} & \texttt{a} & \texttt{0123} & \texttt{0} & \texttt{+} & \texttt{S} & \texttt{none}\\
N730 & F4.2658 & \texttt{T-{}-{}} & 1 & \texttt{a\textbar{}a\textbar{}a\textbar{}a} & \texttt{a} & \texttt{a} & \texttt{a} & \texttt{0123} & \texttt{0} & \texttt{+} & \texttt{S} & \texttt{none}\\
N731 & F4.2659 & \texttt{T-{}-{}} & 1 & \texttt{a\textbar{}a\textbar{}a\textbar{}a} & \texttt{a} & \texttt{a} & \texttt{a} & \texttt{0123} & \texttt{0} & \texttt{+} & \texttt{S} & \texttt{none}\\
N732 & F4.2660 & \texttt{T-{}-{}} & 1 & \texttt{a\textbar{}a\textbar{}a\textbar{}a} & \texttt{a} & \texttt{a} & \texttt{a} & \texttt{0123} & \texttt{0} & \texttt{+} & \texttt{S} & \texttt{none}\\
N733 & F4.2661 & \texttt{T-{}-{}} & 1 & \texttt{a\textbar{}b\textbar{}a\textbar{}a} & \texttt{a} & \texttt{a} & \texttt{ab} & \texttt{023} & \texttt{0} & \texttt{-{}1} & \texttt{s} & \texttt{b@1}\\
N734 & F4.2661 & \texttt{T-{}-{}} & 1 & \texttt{a\textbar{}a\textbar{}a\textbar{}a} & \texttt{a} & \texttt{a} & \texttt{a} & \texttt{0123} & \texttt{0} & \texttt{+} & \texttt{S} & \texttt{none}\\
N735 & F4.2663 & \texttt{T-{}-{}} & 1 & \texttt{a\textbar{}a\textbar{}a\textbar{}a} & \texttt{a} & \texttt{a} & \texttt{a} & \texttt{0123} & \texttt{0} & \texttt{+} & \texttt{S} & \texttt{none}\\
N736 & F4.2664 & \texttt{T-{}-{}} & 1 & \texttt{a\textbar{}b\textbar{}c\textbar{}b} & \texttt{b} & \texttt{b} & \texttt{abc} & \texttt{13} & \texttt{-{}} & \texttt{-{}0} & \texttt{s} & \texttt{a@0}\\
N737 & F4.2664 & \texttt{T-{}-{}} & 1 & \texttt{a\textbar{}b\textbar{}a\textbar{}b} & \texttt{b} & \texttt{b} & \texttt{ab} & \texttt{13} & \texttt{-{}} & \texttt{-{}0} & \texttt{s} & \texttt{a@0}\\
N738 & F4.2664 & \texttt{T-{}-{}} & 1 & \texttt{a\textbar{}a\textbar{}a\textbar{}a} & \texttt{a} & \texttt{a} & \texttt{a} & \texttt{0123} & \texttt{0} & \texttt{+} & \texttt{S} & \texttt{none}\\
N739 & F4.2665 & \texttt{T-{}-{}} & 1 & \texttt{a\textbar{}a\textbar{}b\textbar{}a} & \texttt{a} & \texttt{a} & \texttt{ab} & \texttt{013} & \texttt{0} & \texttt{-{}2} & \texttt{s} & \texttt{b@2}\\
N740 & F4.2665 & \texttt{T-{}-{}} & 1 & \texttt{a\textbar{}a\textbar{}a\textbar{}a} & \texttt{a} & \texttt{a} & \texttt{a} & \texttt{0123} & \texttt{0} & \texttt{+} & \texttt{S} & \texttt{none}\\
N741 & F4.2666 & \texttt{T-{}-{}} & 1 & \texttt{a\textbar{}a\textbar{}a\textbar{}a} & \texttt{a} & \texttt{a} & \texttt{a} & \texttt{0123} & \texttt{0} & \texttt{+} & \texttt{S} & \texttt{none}\\
N742 & F4.2667 & \texttt{T-{}-{}} & 1 & \texttt{a\textbar{}a\textbar{}a\textbar{}a} & \texttt{a} & \texttt{a} & \texttt{a} & \texttt{0123} & \texttt{0} & \texttt{+} & \texttt{S} & \texttt{none}\\
N743 & F4.2668 & \texttt{T-{}-{}} & 1 & \texttt{a\textbar{}b\textbar{}c\textbar{}d} & \texttt{d} & \texttt{d} & \texttt{abcd} & \texttt{3} & \texttt{-{}} & \texttt{-{}0} & \texttt{s} & \texttt{a@0}\\
N744 & F4.2668 & \texttt{T-{}-{}} & 1 & \texttt{a\textbar{}b\textbar{}c\textbar{}c} & \texttt{c} & \texttt{c} & \texttt{abc} & \texttt{23} & \texttt{-{}} & \texttt{-{}0} & \texttt{s} & \texttt{a@0}\\
N745 & F4.2668 & \texttt{T-{}-{}} & 1 & \texttt{a\textbar{}a\textbar{}b\textbar{}b} & \texttt{b} & \texttt{b} & \texttt{ab} & \texttt{23} & \texttt{-{}} & \texttt{-{}0} & \texttt{s} & \texttt{a@0}\\
N746 & F4.2668 & \texttt{T-{}-{}} & 1 & \texttt{a\textbar{}a\textbar{}a\textbar{}b} & \texttt{b} & \texttt{b} & \texttt{ab} & \texttt{3} & \texttt{-{}} & \texttt{-{}0} & \texttt{s} & \texttt{a@0}\\
N747 & F4.2668 & \texttt{T-{}-{}} & 1 & \texttt{a\textbar{}a\textbar{}a\textbar{}a} & \texttt{a} & \texttt{a} & \texttt{a} & \texttt{0123} & \texttt{0} & \texttt{+} & \texttt{S} & \texttt{none}\\
N748 & F4.2669 & \texttt{T-{}4} & 1 & \texttt{a\textbar{}b\textbar{}c\textbar{}d} & \texttt{d} & \texttt{d} & \texttt{abcd} & \texttt{3} & \texttt{-{}} & \texttt{-{}0} & \texttt{s} & \texttt{a@0}\\
N749 & F4.2669 & \texttt{T-{}4} & 1 & \texttt{a\textbar{}b\textbar{}c\textbar{}c} & \texttt{c} & \texttt{c} & \texttt{abc} & \texttt{23} & \texttt{-{}} & \texttt{-{}0} & \texttt{s} & \texttt{a@0}\\
N750 & F4.2669 & \texttt{T-{}4} & 1 & \texttt{a\textbar{}b\textbar{}b\textbar{}c} & \texttt{c} & \texttt{c} & \texttt{abc} & \texttt{3} & \texttt{-{}} & \texttt{-{}0} & \texttt{s} & \texttt{a@0}\\
N751 & F4.2669 & \texttt{T-{}4} & 1 & \texttt{a\textbar{}b\textbar{}b\textbar{}b} & \texttt{b} & \texttt{b} & \texttt{ab} & \texttt{123} & \texttt{-{}} & \texttt{-{}0} & \texttt{s} & \texttt{a@0}\\
N752 & F4.2669 & \texttt{T-{}4} & 1 & \texttt{a\textbar{}a\textbar{}b\textbar{}c} & \texttt{c} & \texttt{c} & \texttt{abc} & \texttt{3} & \texttt{-{}} & \texttt{-{}0} & \texttt{s} & \texttt{b@2}\\
N753 & F4.2669 & \texttt{T-{}4} & 1 & \texttt{a\textbar{}a\textbar{}b\textbar{}b} & \texttt{b} & \texttt{b} & \texttt{ab} & \texttt{23} & \texttt{-{}} & \texttt{-{}0} & \texttt{s} & \texttt{a@0}\\
N754 & F4.2669 & \texttt{T-{}4} & 1 & \texttt{a\textbar{}a\textbar{}a\textbar{}b} & \texttt{b} & \texttt{b} & \texttt{ab} & \texttt{3} & \texttt{-{}} & \texttt{-{}0} & \texttt{s} & \texttt{a@0}\\
N755 & F4.2669 & \texttt{T-{}4} & 1 & \texttt{a\textbar{}a\textbar{}a\textbar{}a} & \texttt{a} & \texttt{a} & \texttt{a} & \texttt{0123} & \texttt{0} & \texttt{+} & \texttt{S} & \texttt{none}\\
N756 & F4.2670 & \texttt{T-{}-{}} & 1 & \texttt{a\textbar{}b\textbar{}c\textbar{}d} & \texttt{d} & \texttt{d} & \texttt{abcd} & \texttt{3} & \texttt{-{}} & \texttt{-{}0} & \texttt{s} & \texttt{a@0}\\
N757 & F4.2670 & \texttt{T-{}-{}} & 1 & \texttt{a\textbar{}b\textbar{}b\textbar{}c} & \texttt{c} & \texttt{c} & \texttt{abc} & \texttt{3} & \texttt{-{}} & \texttt{-{}0} & \texttt{s} & \texttt{a@0}\\
N758 & F4.2670 & \texttt{T-{}-{}} & 1 & \texttt{a\textbar{}a\textbar{}b\textbar{}c} & \texttt{c} & \texttt{c} & \texttt{abc} & \texttt{3} & \texttt{-{}} & \texttt{-{}0} & \texttt{s} & \texttt{b@2}\\
N759 & F4.2670 & \texttt{T-{}-{}} & 1 & \texttt{a\textbar{}a\textbar{}b\textbar{}a} & \texttt{a} & \texttt{a} & \texttt{ab} & \texttt{013} & \texttt{0} & \texttt{-{}2} & \texttt{s} & \texttt{b@2}\\
N760 & F4.2670 & \texttt{T-{}-{}} & 1 & \texttt{a\textbar{}a\textbar{}a\textbar{}b} & \texttt{b} & \texttt{b} & \texttt{ab} & \texttt{3} & \texttt{-{}} & \texttt{-{}0} & \texttt{s} & \texttt{a@0}\\
N761 & F4.2670 & \texttt{T-{}-{}} & 1 & \texttt{a\textbar{}a\textbar{}a\textbar{}a} & \texttt{a} & \texttt{a} & \texttt{a} & \texttt{0123} & \texttt{0} & \texttt{+} & \texttt{S} & \texttt{none}\\
N762 & F4.2671 & \texttt{T-{}-{}} & 1 & \texttt{a\textbar{}b\textbar{}c\textbar{}d} & \texttt{d} & \texttt{d} & \texttt{abcd} & \texttt{3} & \texttt{-{}} & \texttt{-{}0} & \texttt{s} & \texttt{a@0}\\
N763 & F4.2671 & \texttt{T-{}-{}} & 1 & \texttt{a\textbar{}a\textbar{}a\textbar{}b} & \texttt{b} & \texttt{b} & \texttt{ab} & \texttt{3} & \texttt{-{}} & \texttt{-{}0} & \texttt{s} & \texttt{a@0}\\
N764 & F4.2671 & \texttt{T-{}-{}} & 1 & \texttt{a\textbar{}a\textbar{}a\textbar{}a} & \texttt{a} & \texttt{a} & \texttt{a} & \texttt{0123} & \texttt{0} & \texttt{+} & \texttt{S} & \texttt{none}\\
N765 & F4.2673 & \texttt{T-{}-{}} & 1 & \texttt{a\textbar{}a\textbar{}b\textbar{}b} & \texttt{ab} & \texttt{ab} & \texttt{ab} & \texttt{0123} & \texttt{0} & \texttt{+} & \texttt{.} & \texttt{a@0}\\
N766 & F4.2673 & \texttt{T-{}-{}} & 1 & \texttt{a\textbar{}a\textbar{}a\textbar{}a} & \texttt{a} & \texttt{a} & \texttt{a} & \texttt{0123} & \texttt{0} & \texttt{+} & \texttt{S} & \texttt{none}\\
N767 & F4.2674 & \texttt{T-{}-{}} & 1 & \texttt{a\textbar{}b\textbar{}b\textbar{}b} & \texttt{ab} & \texttt{ab} & \texttt{ab} & \texttt{0123} & \texttt{0} & \texttt{+} & \texttt{.} & \texttt{a@0}\\
N768 & F4.2674 & \texttt{T-{}-{}} & 1 & \texttt{a\textbar{}a\textbar{}b\textbar{}b} & \texttt{ab} & \texttt{ab} & \texttt{ab} & \texttt{0123} & \texttt{0} & \texttt{+} & \texttt{.} & \texttt{a@0}\\
N769 & F4.2674 & \texttt{T-{}-{}} & 1 & \texttt{a\textbar{}a\textbar{}a\textbar{}b} & \texttt{a} & \texttt{a} & \texttt{ab} & \texttt{012} & \texttt{0} & \texttt{-{}3} & \texttt{s} & \texttt{b@3}\\
N770 & F4.2674 & \texttt{T-{}-{}} & 1 & \texttt{a\textbar{}a\textbar{}a\textbar{}a} & \texttt{a} & \texttt{a} & \texttt{a} & \texttt{0123} & \texttt{0} & \texttt{+} & \texttt{S} & \texttt{none}\\
N771 & F4.2675 & \texttt{T-{}-{}} & 1 & \texttt{a\textbar{}b\textbar{}b\textbar{}b} & \texttt{ab} & \texttt{ab} & \texttt{ab} & \texttt{0123} & \texttt{0} & \texttt{+} & \texttt{.} & \texttt{a@0}\\
N772 & F4.2675 & \texttt{T-{}-{}} & 1 & \texttt{a\textbar{}a\textbar{}b\textbar{}b} & \texttt{ab} & \texttt{ab} & \texttt{ab} & \texttt{0123} & \texttt{0} & \texttt{+} & \texttt{.} & \texttt{a@0}\\
N773 & F4.2675 & \texttt{T-{}-{}} & 1 & \texttt{a\textbar{}a\textbar{}a\textbar{}a} & \texttt{a} & \texttt{a} & \texttt{a} & \texttt{0123} & \texttt{0} & \texttt{+} & \texttt{S} & \texttt{none}\\
N774 & F4.2676 & \texttt{T-{}-{}} & 1 & \texttt{a\textbar{}b\textbar{}c\textbar{}c} & \texttt{ac} & \texttt{ac} & \texttt{abc} & \texttt{023} & \texttt{0} & \texttt{-{}1} & \texttt{.} & \texttt{a@0}\\
N775 & F4.2676 & \texttt{T-{}-{}} & 1 & \texttt{a\textbar{}b\textbar{}b\textbar{}b} & \texttt{ab} & \texttt{ab} & \texttt{ab} & \texttt{0123} & \texttt{0} & \texttt{+} & \texttt{.} & \texttt{a@0}\\
N776 & F4.2676 & \texttt{T-{}-{}} & 1 & \texttt{a\textbar{}b\textbar{}a\textbar{}a} & \texttt{a} & \texttt{a} & \texttt{ab} & \texttt{023} & \texttt{0} & \texttt{-{}1} & \texttt{s} & \texttt{b@1}\\
N777 & F4.2676 & \texttt{T-{}-{}} & 1 & \texttt{a\textbar{}a\textbar{}b\textbar{}b} & \texttt{ab} & \texttt{ab} & \texttt{ab} & \texttt{0123} & \texttt{0} & \texttt{+} & \texttt{.} & \texttt{a@0}\\
N778 & F4.2676 & \texttt{T-{}-{}} & 1 & \texttt{a\textbar{}a\textbar{}a\textbar{}a} & \texttt{a} & \texttt{a} & \texttt{a} & \texttt{0123} & \texttt{0} & \texttt{+} & \texttt{S} & \texttt{none}\\
N779 & F4.2677 & \texttt{T-{}-{}} & 1 & \texttt{a\textbar{}a\textbar{}b\textbar{}b} & \texttt{ab} & \texttt{ab} & \texttt{ab} & \texttt{0123} & \texttt{0} & \texttt{+} & \texttt{.} & \texttt{a@0}\\
N780 & F4.2677 & \texttt{T-{}-{}} & 1 & \texttt{a\textbar{}a\textbar{}a\textbar{}a} & \texttt{a} & \texttt{a} & \texttt{a} & \texttt{0123} & \texttt{0} & \texttt{+} & \texttt{S} & \texttt{none}\\
N781 & F4.2678 & \texttt{T-{}-{}} & 1 & \texttt{a\textbar{}b\textbar{}a\textbar{}a} & \texttt{a} & \texttt{a} & \texttt{ab} & \texttt{023} & \texttt{0} & \texttt{-{}1} & \texttt{s} & \texttt{b@1}\\
N782 & F4.2678 & \texttt{T-{}-{}} & 1 & \texttt{a\textbar{}a\textbar{}b\textbar{}b} & \texttt{ab} & \texttt{ab} & \texttt{ab} & \texttt{0123} & \texttt{0} & \texttt{+} & \texttt{.} & \texttt{a@0}\\
N783 & F4.2678 & \texttt{T-{}-{}} & 1 & \texttt{a\textbar{}a\textbar{}a\textbar{}a} & \texttt{a} & \texttt{a} & \texttt{a} & \texttt{0123} & \texttt{0} & \texttt{+} & \texttt{S} & \texttt{none}\\
N784 & F4.2679 & \texttt{T-{}-{}} & 1 & \texttt{a\textbar{}b\textbar{}b\textbar{}b} & \texttt{ab} & \texttt{ab} & \texttt{ab} & \texttt{0123} & \texttt{0} & \texttt{+} & \texttt{.} & \texttt{a@0}\\
N785 & F4.2679 & \texttt{T-{}-{}} & 1 & \texttt{a\textbar{}a\textbar{}a\textbar{}a} & \texttt{a} & \texttt{a} & \texttt{a} & \texttt{0123} & \texttt{0} & \texttt{+} & \texttt{S} & \texttt{none}\\
N786 & F4.2680 & \texttt{T-{}-{}} & 1 & \texttt{a\textbar{}b\textbar{}c\textbar{}c} & \texttt{ac} & \texttt{ac} & \texttt{abc} & \texttt{023} & \texttt{0} & \texttt{-{}1} & \texttt{.} & \texttt{a@0}\\
N787 & F4.2680 & \texttt{T-{}-{}} & 1 & \texttt{a\textbar{}b\textbar{}b\textbar{}b} & \texttt{ab} & \texttt{ab} & \texttt{ab} & \texttt{0123} & \texttt{0} & \texttt{+} & \texttt{.} & \texttt{a@0}\\
N788 & F4.2680 & \texttt{T-{}-{}} & 1 & \texttt{a\textbar{}b\textbar{}a\textbar{}c} & \texttt{a} & \texttt{a} & \texttt{abc} & \texttt{02} & \texttt{0} & \texttt{-{}1} & \texttt{s} & \texttt{b@1}\\
N789 & F4.2680 & \texttt{T-{}-{}} & 1 & \texttt{a\textbar{}b\textbar{}a\textbar{}a} & \texttt{a} & \texttt{a} & \texttt{ab} & \texttt{023} & \texttt{0} & \texttt{-{}1} & \texttt{s} & \texttt{b@1}\\
N790 & F4.2680 & \texttt{T-{}-{}} & 1 & \texttt{a\textbar{}a\textbar{}b\textbar{}b} & \texttt{ab} & \texttt{ab} & \texttt{ab} & \texttt{0123} & \texttt{0} & \texttt{+} & \texttt{.} & \texttt{a@0}\\
N791 & F4.2680 & \texttt{T-{}-{}} & 1 & \texttt{a\textbar{}b\textbar{}a\textbar{}b} & \texttt{a} & \texttt{a} & \texttt{ab} & \texttt{02} & \texttt{0} & \texttt{-{}1} & \texttt{s} & \texttt{b@1}\\
N792 & F4.2680 & \texttt{T-{}-{}} & 1 & \texttt{a\textbar{}a\textbar{}a\textbar{}a} & \texttt{a} & \texttt{a} & \texttt{a} & \texttt{0123} & \texttt{0} & \texttt{+} & \texttt{S} & \texttt{none}\\
N793 & F4.2681 & \texttt{T-{}-{}} & 1 & \texttt{a\textbar{}b\textbar{}a\textbar{}b} & \texttt{ab} & \texttt{ab} & \texttt{ab} & \texttt{0123} & \texttt{0} & \texttt{+} & \texttt{.} & \texttt{a@0}\\
N794 & F4.2681 & \texttt{T-{}-{}} & 1 & \texttt{a\textbar{}a\textbar{}a\textbar{}a} & \texttt{a} & \texttt{a} & \texttt{a} & \texttt{0123} & \texttt{0} & \texttt{+} & \texttt{S} & \texttt{none}\\
N795 & F4.2682 & \texttt{T-{}-{}} & 1 & \texttt{a\textbar{}a\textbar{}b\textbar{}b} & \texttt{ab} & \texttt{ab} & \texttt{ab} & \texttt{0123} & \texttt{0} & \texttt{+} & \texttt{.} & \texttt{a@0}\\
N796 & F4.2682 & \texttt{T-{}-{}} & 1 & \texttt{a\textbar{}b\textbar{}a\textbar{}b} & \texttt{ab} & \texttt{ab} & \texttt{ab} & \texttt{0123} & \texttt{0} & \texttt{+} & \texttt{.} & \texttt{a@0}\\
N797 & F4.2682 & \texttt{T-{}-{}} & 1 & \texttt{a\textbar{}a\textbar{}a\textbar{}b} & \texttt{a} & \texttt{a} & \texttt{ab} & \texttt{012} & \texttt{0} & \texttt{-{}3} & \texttt{s} & \texttt{b@3}\\
N798 & F4.2682 & \texttt{T-{}-{}} & 1 & \texttt{a\textbar{}a\textbar{}a\textbar{}a} & \texttt{a} & \texttt{a} & \texttt{a} & \texttt{0123} & \texttt{0} & \texttt{+} & \texttt{S} & \texttt{none}\\
N799 & F4.2684 & \texttt{T-{}-{}} & 1 & \texttt{a\textbar{}a\textbar{}b\textbar{}b} & \texttt{ab} & \texttt{ab} & \texttt{ab} & \texttt{0123} & \texttt{0} & \texttt{+} & \texttt{.} & \texttt{a@0}\\
N800 & F4.2684 & \texttt{T-{}-{}} & 1 & \texttt{a\textbar{}b\textbar{}a\textbar{}b} & \texttt{ab} & \texttt{ab} & \texttt{ab} & \texttt{0123} & \texttt{0} & \texttt{+} & \texttt{.} & \texttt{a@0}\\
N801 & F4.2684 & \texttt{T-{}-{}} & 1 & \texttt{a\textbar{}a\textbar{}a\textbar{}a} & \texttt{a} & \texttt{a} & \texttt{a} & \texttt{0123} & \texttt{0} & \texttt{+} & \texttt{S} & \texttt{none}\\
N802 & F4.2685 & \texttt{T-{}-{}} & 1 & \texttt{a\textbar{}a\textbar{}b\textbar{}b} & \texttt{ab} & \texttt{ab} & \texttt{ab} & \texttt{0123} & \texttt{0} & \texttt{+} & \texttt{.} & \texttt{a@0}\\
N803 & F4.2685 & \texttt{T-{}-{}} & 1 & \texttt{a\textbar{}a\textbar{}a\textbar{}a} & \texttt{a} & \texttt{a} & \texttt{a} & \texttt{0123} & \texttt{0} & \texttt{+} & \texttt{S} & \texttt{none}\\
N804 & F4.2686 & \texttt{T-{}-{}} & 1 & \texttt{a\textbar{}a\textbar{}b\textbar{}b} & \texttt{ab} & \texttt{ab} & \texttt{ab} & \texttt{0123} & \texttt{0} & \texttt{+} & \texttt{.} & \texttt{a@0}\\
N805 & F4.2686 & \texttt{T-{}-{}} & 1 & \texttt{a\textbar{}b\textbar{}a\textbar{}b} & \texttt{ab} & \texttt{ab} & \texttt{ab} & \texttt{0123} & \texttt{0} & \texttt{+} & \texttt{.} & \texttt{a@0}\\
N806 & F4.2686 & \texttt{T-{}-{}} & 1 & \texttt{a\textbar{}a\textbar{}a\textbar{}a} & \texttt{a} & \texttt{a} & \texttt{a} & \texttt{0123} & \texttt{0} & \texttt{+} & \texttt{S} & \texttt{none}\\
N807 & F4.2687 & \texttt{T-{}-{}} & 1 & \texttt{a\textbar{}a\textbar{}b\textbar{}b} & \texttt{ab} & \texttt{ab} & \texttt{ab} & \texttt{0123} & \texttt{0} & \texttt{+} & \texttt{.} & \texttt{a@0}\\
N808 & F4.2687 & \texttt{T-{}-{}} & 1 & \texttt{a\textbar{}b\textbar{}a\textbar{}b} & \texttt{ab} & \texttt{ab} & \texttt{ab} & \texttt{0123} & \texttt{0} & \texttt{+} & \texttt{.} & \texttt{a@0}\\
N809 & F4.2687 & \texttt{T-{}-{}} & 1 & \texttt{a\textbar{}a\textbar{}a\textbar{}a} & \texttt{a} & \texttt{a} & \texttt{a} & \texttt{0123} & \texttt{0} & \texttt{+} & \texttt{S} & \texttt{none}\\
N810 & F4.2688 & \texttt{T-{}-{}} & 1 & \texttt{a\textbar{}a\textbar{}b\textbar{}b} & \texttt{ab} & \texttt{ab} & \texttt{ab} & \texttt{0123} & \texttt{0} & \texttt{+} & \texttt{.} & \texttt{a@0}\\
N811 & F4.2688 & \texttt{T-{}-{}} & 1 & \texttt{a\textbar{}b\textbar{}a\textbar{}b} & \texttt{ab} & \texttt{ab} & \texttt{ab} & \texttt{0123} & \texttt{0} & \texttt{+} & \texttt{.} & \texttt{a@0}\\
N812 & F4.2688 & \texttt{T-{}-{}} & 1 & \texttt{a\textbar{}a\textbar{}a\textbar{}a} & \texttt{a} & \texttt{a} & \texttt{a} & \texttt{0123} & \texttt{0} & \texttt{+} & \texttt{S} & \texttt{none}\\
N813 & F4.2689 & \texttt{T-{}-{}} & 1 & \texttt{a\textbar{}b\textbar{}b\textbar{}b} & \texttt{ab} & \texttt{ab} & \texttt{ab} & \texttt{0123} & \texttt{0} & \texttt{+} & \texttt{.} & \texttt{a@0}\\
N814 & F4.2689 & \texttt{T-{}-{}} & 1 & \texttt{a\textbar{}b\textbar{}a\textbar{}b} & \texttt{ab} & \texttt{ab} & \texttt{ab} & \texttt{0123} & \texttt{0} & \texttt{+} & \texttt{.} & \texttt{a@0}\\
N815 & F4.2689 & \texttt{T-{}-{}} & 1 & \texttt{a\textbar{}a\textbar{}a\textbar{}b} & \texttt{a} & \texttt{a} & \texttt{ab} & \texttt{012} & \texttt{0} & \texttt{-{}3} & \texttt{s} & \texttt{b@3}\\
N816 & F4.2689 & \texttt{T-{}-{}} & 1 & \texttt{a\textbar{}a\textbar{}a\textbar{}a} & \texttt{a} & \texttt{a} & \texttt{a} & \texttt{0123} & \texttt{0} & \texttt{+} & \texttt{S} & \texttt{none}\\
N817 & F4.2690 & \texttt{T-{}-{}} & 1 & \texttt{a\textbar{}b\textbar{}b\textbar{}b} & \texttt{ab} & \texttt{ab} & \texttt{ab} & \texttt{0123} & \texttt{0} & \texttt{+} & \texttt{.} & \texttt{a@0}\\
N818 & F4.2690 & \texttt{T-{}-{}} & 1 & \texttt{a\textbar{}b\textbar{}c\textbar{}a} & \texttt{a} & \texttt{a} & \texttt{abc} & \texttt{03} & \texttt{0} & \texttt{-{}1} & \texttt{s} & \texttt{b@1}\\
N819 & F4.2690 & \texttt{T-{}-{}} & 1 & \texttt{a\textbar{}b\textbar{}a\textbar{}b} & \texttt{ab} & \texttt{ab} & \texttt{ab} & \texttt{0123} & \texttt{0} & \texttt{+} & \texttt{.} & \texttt{a@0}\\
N820 & F4.2690 & \texttt{T-{}-{}} & 1 & \texttt{a\textbar{}b\textbar{}b\textbar{}a} & \texttt{a} & \texttt{a} & \texttt{ab} & \texttt{03} & \texttt{0} & \texttt{-{}1} & \texttt{s} & \texttt{b@1}\\
N821 & F4.2690 & \texttt{T-{}-{}} & 1 & \texttt{a\textbar{}a\textbar{}a\textbar{}a} & \texttt{a} & \texttt{a} & \texttt{a} & \texttt{0123} & \texttt{0} & \texttt{+} & \texttt{S} & \texttt{none}\\
N822 & F4.2691 & \texttt{T-{}-{}} & 1 & \texttt{a\textbar{}a\textbar{}b\textbar{}b} & \texttt{ab} & \texttt{ab} & \texttt{ab} & \texttt{0123} & \texttt{0} & \texttt{+} & \texttt{.} & \texttt{a@0}\\
N823 & F4.2691 & \texttt{T-{}-{}} & 1 & \texttt{a\textbar{}a\textbar{}a\textbar{}a} & \texttt{a} & \texttt{a} & \texttt{a} & \texttt{0123} & \texttt{0} & \texttt{+} & \texttt{S} & \texttt{none}\\
N824 & F4.2692 & \texttt{T-{}-{}} & 1 & \texttt{a\textbar{}a\textbar{}b\textbar{}b} & \texttt{ab} & \texttt{ab} & \texttt{ab} & \texttt{0123} & \texttt{0} & \texttt{+} & \texttt{.} & \texttt{a@0}\\
N825 & F4.2692 & \texttt{T-{}-{}} & 1 & \texttt{a\textbar{}a\textbar{}b\textbar{}a} & \texttt{a} & \texttt{a} & \texttt{ab} & \texttt{013} & \texttt{0} & \texttt{-{}2} & \texttt{s} & \texttt{b@2}\\
N826 & F4.2692 & \texttt{T-{}-{}} & 1 & \texttt{a\textbar{}a\textbar{}a\textbar{}a} & \texttt{a} & \texttt{a} & \texttt{a} & \texttt{0123} & \texttt{0} & \texttt{+} & \texttt{S} & \texttt{none}\\
N827 & F4.2693 & \texttt{T-{}-{}} & 1 & \texttt{a\textbar{}a\textbar{}b\textbar{}b} & \texttt{ab} & \texttt{ab} & \texttt{ab} & \texttt{0123} & \texttt{0} & \texttt{+} & \texttt{.} & \texttt{a@0}\\
N828 & F4.2693 & \texttt{T-{}-{}} & 1 & \texttt{a\textbar{}a\textbar{}a\textbar{}a} & \texttt{a} & \texttt{a} & \texttt{a} & \texttt{0123} & \texttt{0} & \texttt{+} & \texttt{S} & \texttt{none}\\
N829 & F4.2694 & \texttt{T-{}-{}} & 1 & \texttt{a\textbar{}b\textbar{}c\textbar{}b} & \texttt{ab} & \texttt{ab} & \texttt{abc} & \texttt{013} & \texttt{0} & \texttt{-{}2} & \texttt{.} & \texttt{a@0}\\
N830 & F4.2694 & \texttt{T-{}-{}} & 1 & \texttt{a\textbar{}b\textbar{}b\textbar{}b} & \texttt{ab} & \texttt{ab} & \texttt{ab} & \texttt{0123} & \texttt{0} & \texttt{+} & \texttt{.} & \texttt{a@0}\\
N831 & F4.2694 & \texttt{T-{}-{}} & 1 & \texttt{a\textbar{}a\textbar{}b\textbar{}a} & \texttt{a} & \texttt{a} & \texttt{ab} & \texttt{013} & \texttt{0} & \texttt{-{}2} & \texttt{s} & \texttt{b@2}\\
N832 & F4.2694 & \texttt{T-{}-{}} & 1 & \texttt{a\textbar{}b\textbar{}a\textbar{}b} & \texttt{ab} & \texttt{ab} & \texttt{ab} & \texttt{0123} & \texttt{0} & \texttt{+} & \texttt{.} & \texttt{a@0}\\
N833 & F4.2694 & \texttt{T-{}-{}} & 1 & \texttt{a\textbar{}a\textbar{}a\textbar{}a} & \texttt{a} & \texttt{a} & \texttt{a} & \texttt{0123} & \texttt{0} & \texttt{+} & \texttt{S} & \texttt{none}\\
N834 & F4.2695 & \texttt{T-{}-{}} & 1 & \texttt{a\textbar{}b\textbar{}b\textbar{}b} & \texttt{ab} & \texttt{ab} & \texttt{ab} & \texttt{0123} & \texttt{0} & \texttt{+} & \texttt{.} & \texttt{a@0}\\
N835 & F4.2695 & \texttt{T-{}-{}} & 1 & \texttt{a\textbar{}a\textbar{}b\textbar{}b} & \texttt{ab} & \texttt{ab} & \texttt{ab} & \texttt{0123} & \texttt{0} & \texttt{+} & \texttt{.} & \texttt{a@0}\\
N836 & F4.2695 & \texttt{T-{}-{}} & 1 & \texttt{a\textbar{}a\textbar{}a\textbar{}a} & \texttt{a} & \texttt{a} & \texttt{a} & \texttt{0123} & \texttt{0} & \texttt{+} & \texttt{S} & \texttt{none}\\
N837 & F4.2696 & \texttt{T-{}-{}} & 1 & \texttt{a\textbar{}b\textbar{}c\textbar{}b} & \texttt{ab} & \texttt{ab} & \texttt{abc} & \texttt{013} & \texttt{0} & \texttt{-{}2} & \texttt{.} & \texttt{a@0}\\
N838 & F4.2696 & \texttt{T-{}-{}} & 1 & \texttt{a\textbar{}b\textbar{}b\textbar{}b} & \texttt{ab} & \texttt{ab} & \texttt{ab} & \texttt{0123} & \texttt{0} & \texttt{+} & \texttt{.} & \texttt{a@0}\\
N839 & F4.2696 & \texttt{T-{}-{}} & 1 & \texttt{a\textbar{}b\textbar{}c\textbar{}a} & \texttt{a} & \texttt{a} & \texttt{abc} & \texttt{03} & \texttt{0} & \texttt{-{}1} & \texttt{s} & \texttt{b@1}\\
N840 & F4.2696 & \texttt{T-{}-{}} & 1 & \texttt{a\textbar{}a\textbar{}b\textbar{}b} & \texttt{ab} & \texttt{ab} & \texttt{ab} & \texttt{0123} & \texttt{0} & \texttt{+} & \texttt{.} & \texttt{a@0}\\
N841 & F4.2696 & \texttt{T-{}-{}} & 1 & \texttt{a\textbar{}a\textbar{}b\textbar{}a} & \texttt{a} & \texttt{a} & \texttt{ab} & \texttt{013} & \texttt{0} & \texttt{-{}2} & \texttt{s} & \texttt{b@2}\\
N842 & F4.2696 & \texttt{T-{}-{}} & 1 & \texttt{a\textbar{}b\textbar{}a\textbar{}b} & \texttt{ab} & \texttt{ab} & \texttt{ab} & \texttt{0123} & \texttt{0} & \texttt{+} & \texttt{.} & \texttt{a@0}\\
N843 & F4.2696 & \texttt{T-{}-{}} & 1 & \texttt{a\textbar{}b\textbar{}b\textbar{}a} & \texttt{a} & \texttt{a} & \texttt{ab} & \texttt{03} & \texttt{0} & \texttt{-{}1} & \texttt{s} & \texttt{b@1}\\
N844 & F4.2696 & \texttt{T-{}-{}} & 1 & \texttt{a\textbar{}a\textbar{}a\textbar{}a} & \texttt{a} & \texttt{a} & \texttt{a} & \texttt{0123} & \texttt{0} & \texttt{+} & \texttt{S} & \texttt{none}\\
N845 & F4.2697 & \texttt{T-{}-{}} & 1 & \texttt{a\textbar{}b\textbar{}a\textbar{}b} & \texttt{ab} & \texttt{ab} & \texttt{ab} & \texttt{0123} & \texttt{0} & \texttt{+} & \texttt{.} & \texttt{a@0}\\
N846 & F4.2697 & \texttt{T-{}-{}} & 1 & \texttt{a\textbar{}a\textbar{}a\textbar{}a} & \texttt{a} & \texttt{a} & \texttt{a} & \texttt{0123} & \texttt{0} & \texttt{+} & \texttt{S} & \texttt{none}\\
N847 & F4.2698 & \texttt{T-{}-{}} & 1 & \texttt{a\textbar{}b\textbar{}c\textbar{}a} & \texttt{a} & \texttt{a} & \texttt{abc} & \texttt{03} & \texttt{0} & \texttt{-{}1} & \texttt{s} & \texttt{b@1}\\
N848 & F4.2698 & \texttt{T-{}-{}} & 1 & \texttt{a\textbar{}a\textbar{}b\textbar{}b} & \texttt{ab} & \texttt{ab} & \texttt{ab} & \texttt{0123} & \texttt{0} & \texttt{+} & \texttt{.} & \texttt{a@0}\\
N849 & F4.2698 & \texttt{T-{}-{}} & 1 & \texttt{a\textbar{}b\textbar{}a\textbar{}b} & \texttt{ab} & \texttt{ab} & \texttt{ab} & \texttt{0123} & \texttt{0} & \texttt{+} & \texttt{.} & \texttt{a@0}\\
N850 & F4.2698 & \texttt{T-{}-{}} & 1 & \texttt{a\textbar{}b\textbar{}b\textbar{}a} & \texttt{a} & \texttt{a} & \texttt{ab} & \texttt{03} & \texttt{0} & \texttt{-{}1} & \texttt{s} & \texttt{b@1}\\
N851 & F4.2698 & \texttt{T-{}-{}} & 1 & \texttt{a\textbar{}a\textbar{}a\textbar{}a} & \texttt{a} & \texttt{a} & \texttt{a} & \texttt{0123} & \texttt{0} & \texttt{+} & \texttt{S} & \texttt{none}\\
N852 & F4.2699 & \texttt{T-{}-{}} & 1 & \texttt{a\textbar{}b\textbar{}a\textbar{}b} & \texttt{ab} & \texttt{ab} & \texttt{ab} & \texttt{0123} & \texttt{0} & \texttt{+} & \texttt{.} & \texttt{a@0}\\
N853 & F4.2699 & \texttt{T-{}-{}} & 1 & \texttt{a\textbar{}a\textbar{}a\textbar{}a} & \texttt{a} & \texttt{a} & \texttt{a} & \texttt{0123} & \texttt{0} & \texttt{+} & \texttt{S} & \texttt{none}\\
N854 & F4.2700 & \texttt{T-{}-{}} & 1 & \texttt{a\textbar{}b\textbar{}a\textbar{}c} & \texttt{ac} & \texttt{ac} & \texttt{abc} & \texttt{023} & \texttt{0} & \texttt{-{}1} & \texttt{.} & \texttt{b@1}\\
N855 & F4.2700 & \texttt{T-{}-{}} & 1 & \texttt{a\textbar{}b\textbar{}a\textbar{}a} & \texttt{a} & \texttt{a} & \texttt{ab} & \texttt{023} & \texttt{0} & \texttt{-{}1} & \texttt{s} & \texttt{b@1}\\
N856 & F4.2700 & \texttt{T-{}-{}} & 1 & \texttt{a\textbar{}b\textbar{}a\textbar{}b} & \texttt{ab} & \texttt{ab} & \texttt{ab} & \texttt{0123} & \texttt{0} & \texttt{+} & \texttt{.} & \texttt{a@0}\\
N857 & F4.2700 & \texttt{T-{}-{}} & 1 & \texttt{a\textbar{}a\textbar{}a\textbar{}b} & \texttt{ab} & \texttt{ab} & \texttt{ab} & \texttt{0123} & \texttt{0} & \texttt{+} & \texttt{.} & \texttt{b@3}\\
N858 & F4.2700 & \texttt{T-{}-{}} & 1 & \texttt{a\textbar{}a\textbar{}a\textbar{}a} & \texttt{a} & \texttt{a} & \texttt{a} & \texttt{0123} & \texttt{0} & \texttt{+} & \texttt{S} & \texttt{none}\\
N859 & F4.2701 & \texttt{T-{}-{}} & 1 & \texttt{a\textbar{}b\textbar{}a\textbar{}a} & \texttt{a} & \texttt{a} & \texttt{ab} & \texttt{023} & \texttt{0} & \texttt{-{}1} & \texttt{s} & \texttt{b@1}\\
N860 & F4.2701 & \texttt{T-{}-{}} & 1 & \texttt{a\textbar{}a\textbar{}a\textbar{}a} & \texttt{a} & \texttt{a} & \texttt{a} & \texttt{0123} & \texttt{0} & \texttt{+} & \texttt{S} & \texttt{none}\\
N861 & F4.2702 & \texttt{T-{}4} & 1 & \texttt{a\textbar{}b\textbar{}c\textbar{}d} & \texttt{ad} & \texttt{ad} & \texttt{abcd} & \texttt{03} & \texttt{0} & \texttt{-{}1} & \texttt{.} & \texttt{a@0}\\
N862 & F4.2702 & \texttt{T-{}4} & 1 & \texttt{a\textbar{}a\textbar{}b\textbar{}c} & \texttt{ac} & \texttt{ac} & \texttt{abc} & \texttt{013} & \texttt{0} & \texttt{-{}2} & \texttt{.} & \texttt{b@2}\\
N863 & F4.2702 & \texttt{T-{}4} & 1 & \texttt{a\textbar{}b\textbar{}b\textbar{}c} & \texttt{ac} & \texttt{ac} & \texttt{abc} & \texttt{03} & \texttt{0} & \texttt{-{}1} & \texttt{.} & \texttt{a@0}\\
N864 & F4.2702 & \texttt{T-{}4} & 1 & \texttt{a\textbar{}a\textbar{}a\textbar{}b} & \texttt{ab} & \texttt{ab} & \texttt{ab} & \texttt{0123} & \texttt{0} & \texttt{+} & \texttt{.} & \texttt{b@3}\\
N865 & F4.2702 & \texttt{T-{}4} & 1 & \texttt{a\textbar{}b\textbar{}c\textbar{}a} & \texttt{a} & \texttt{a} & \texttt{abc} & \texttt{03} & \texttt{0} & \texttt{-{}1} & \texttt{s} & \texttt{b@1}\\
N866 & F4.2702 & \texttt{T-{}4} & 1 & \texttt{a\textbar{}a\textbar{}b\textbar{}a} & \texttt{a} & \texttt{a} & \texttt{ab} & \texttt{013} & \texttt{0} & \texttt{-{}2} & \texttt{s} & \texttt{b@2}\\
N867 & F4.2702 & \texttt{T-{}4} & 1 & \texttt{a\textbar{}b\textbar{}a\textbar{}b} & \texttt{ab} & \texttt{ab} & \texttt{ab} & \texttt{0123} & \texttt{0} & \texttt{+} & \texttt{.} & \texttt{a@0}\\
N868 & F4.2702 & \texttt{T-{}4} & 1 & \texttt{a\textbar{}b\textbar{}b\textbar{}a} & \texttt{a} & \texttt{a} & \texttt{ab} & \texttt{03} & \texttt{0} & \texttt{-{}1} & \texttt{s} & \texttt{b@1}\\
N869 & F4.2702 & \texttt{T-{}4} & 1 & \texttt{a\textbar{}a\textbar{}a\textbar{}a} & \texttt{a} & \texttt{a} & \texttt{a} & \texttt{0123} & \texttt{0} & \texttt{+} & \texttt{S} & \texttt{none}\\
N870 & F4.2868 & \texttt{T-{}-{}} & 1 & \texttt{a\textbar{}a\textbar{}a\textbar{}a} & \texttt{a} & \texttt{a} & \texttt{a} & \texttt{0123} & \texttt{0} & \texttt{+} & \texttt{S} & \texttt{none}\\
N871 & F4.2869 & \texttt{T-{}-{}} & 1 & \texttt{a\textbar{}a\textbar{}a\textbar{}a} & \texttt{a} & \texttt{a} & \texttt{a} & \texttt{0123} & \texttt{0} & \texttt{+} & \texttt{S} & \texttt{none}\\
N872 & F4.2870 & \texttt{T-{}-{}} & 1 & \texttt{a\textbar{}a\textbar{}a\textbar{}a} & \texttt{a} & \texttt{a} & \texttt{a} & \texttt{0123} & \texttt{0} & \texttt{+} & \texttt{S} & \texttt{none}\\
N873 & F4.2871 & \texttt{T-{}-{}} & 1 & \texttt{a\textbar{}a\textbar{}a\textbar{}a} & \texttt{a} & \texttt{a} & \texttt{a} & \texttt{0123} & \texttt{0} & \texttt{+} & \texttt{S} & \texttt{none}\\
N874 & F4.2872 & \texttt{T-{}4} & 1 & \texttt{a\textbar{}b\textbar{}c\textbar{}d} & \texttt{d} & \texttt{d} & \texttt{abcd} & \texttt{3} & \texttt{-{}} & \texttt{-{}0} & \texttt{s} & \texttt{a@0}\\
N875 & F4.2872 & \texttt{T-{}4} & 1 & \texttt{a\textbar{}b\textbar{}c\textbar{}c} & \texttt{c} & \texttt{c} & \texttt{abc} & \texttt{23} & \texttt{-{}} & \texttt{-{}0} & \texttt{s} & \texttt{a@0}\\
N876 & F4.2872 & \texttt{T-{}4} & 1 & \texttt{a\textbar{}a\textbar{}b\textbar{}c} & \texttt{c} & \texttt{c} & \texttt{abc} & \texttt{3} & \texttt{-{}} & \texttt{-{}0} & \texttt{s} & \texttt{b@2}\\
N877 & F4.2872 & \texttt{T-{}4} & 1 & \texttt{a\textbar{}a\textbar{}b\textbar{}b} & \texttt{b} & \texttt{b} & \texttt{ab} & \texttt{23} & \texttt{-{}} & \texttt{-{}0} & \texttt{s} & \texttt{a@0}\\
N878 & F4.2872 & \texttt{T-{}4} & 1 & \texttt{a\textbar{}a\textbar{}a\textbar{}b} & \texttt{b} & \texttt{b} & \texttt{ab} & \texttt{3} & \texttt{-{}} & \texttt{-{}0} & \texttt{s} & \texttt{a@0}\\
N879 & F4.2872 & \texttt{T-{}4} & 1 & \texttt{a\textbar{}a\textbar{}a\textbar{}a} & \texttt{a} & \texttt{a} & \texttt{a} & \texttt{0123} & \texttt{0} & \texttt{+} & \texttt{S} & \texttt{none}\\
N880 & F4.2873 & \texttt{T-{}-{}} & 1 & \texttt{a\textbar{}b\textbar{}c\textbar{}d} & \texttt{d} & \texttt{d} & \texttt{abcd} & \texttt{3} & \texttt{-{}} & \texttt{-{}0} & \texttt{s} & \texttt{a@0}\\
N881 & F4.2873 & \texttt{T-{}-{}} & 1 & \texttt{a\textbar{}b\textbar{}b\textbar{}c} & \texttt{c} & \texttt{c} & \texttt{abc} & \texttt{3} & \texttt{-{}} & \texttt{-{}0} & \texttt{s} & \texttt{a@0}\\
N882 & F4.2873 & \texttt{T-{}-{}} & 1 & \texttt{a\textbar{}a\textbar{}a\textbar{}b} & \texttt{b} & \texttt{b} & \texttt{ab} & \texttt{3} & \texttt{-{}} & \texttt{-{}0} & \texttt{s} & \texttt{a@0}\\
N883 & F4.2873 & \texttt{T-{}-{}} & 1 & \texttt{a\textbar{}a\textbar{}a\textbar{}a} & \texttt{a} & \texttt{a} & \texttt{a} & \texttt{0123} & \texttt{0} & \texttt{+} & \texttt{S} & \texttt{none}\\
N884 & F4.2874 & \texttt{T-{}-{}} & 1 & \texttt{a\textbar{}b\textbar{}c\textbar{}d} & \texttt{d} & \texttt{d} & \texttt{abcd} & \texttt{3} & \texttt{-{}} & \texttt{-{}0} & \texttt{s} & \texttt{a@0}\\
N885 & F4.2874 & \texttt{T-{}-{}} & 1 & \texttt{a\textbar{}a\textbar{}b\textbar{}c} & \texttt{c} & \texttt{c} & \texttt{abc} & \texttt{3} & \texttt{-{}} & \texttt{-{}0} & \texttt{s} & \texttt{b@2}\\
N886 & F4.2874 & \texttt{T-{}-{}} & 1 & \texttt{a\textbar{}a\textbar{}a\textbar{}b} & \texttt{b} & \texttt{b} & \texttt{ab} & \texttt{3} & \texttt{-{}} & \texttt{-{}0} & \texttt{s} & \texttt{a@0}\\
N887 & F4.2874 & \texttt{T-{}-{}} & 1 & \texttt{a\textbar{}a\textbar{}a\textbar{}a} & \texttt{a} & \texttt{a} & \texttt{a} & \texttt{0123} & \texttt{0} & \texttt{+} & \texttt{S} & \texttt{none}\\
N888 & F4.2875 & \texttt{T-{}4} & 1 & \texttt{a\textbar{}b\textbar{}c\textbar{}d} & \texttt{d} & \texttt{d} & \texttt{abcd} & \texttt{3} & \texttt{-{}} & \texttt{-{}0} & \texttt{s} & \texttt{a@0}\\
N889 & F4.2875 & \texttt{T-{}4} & 1 & \texttt{a\textbar{}a\textbar{}b\textbar{}c} & \texttt{c} & \texttt{c} & \texttt{abc} & \texttt{3} & \texttt{-{}} & \texttt{-{}0} & \texttt{s} & \texttt{b@2}\\
N890 & F4.2875 & \texttt{T-{}4} & 1 & \texttt{a\textbar{}b\textbar{}a\textbar{}c} & \texttt{c} & \texttt{c} & \texttt{abc} & \texttt{3} & \texttt{-{}} & \texttt{-{}0} & \texttt{s} & \texttt{b@1}\\
N891 & F4.2875 & \texttt{T-{}4} & 1 & \texttt{a\textbar{}b\textbar{}a\textbar{}a} & \texttt{a} & \texttt{a} & \texttt{ab} & \texttt{023} & \texttt{0} & \texttt{-{}1} & \texttt{s} & \texttt{b@1}\\
N892 & F4.2875 & \texttt{T-{}4} & 1 & \texttt{a\textbar{}a\textbar{}a\textbar{}b} & \texttt{b} & \texttt{b} & \texttt{ab} & \texttt{3} & \texttt{-{}} & \texttt{-{}0} & \texttt{s} & \texttt{a@0}\\
N893 & F4.2875 & \texttt{T-{}4} & 1 & \texttt{a\textbar{}a\textbar{}a\textbar{}a} & \texttt{a} & \texttt{a} & \texttt{a} & \texttt{0123} & \texttt{0} & \texttt{+} & \texttt{S} & \texttt{none}\\
N894 & F4.2876 & \texttt{T-{}-{}} & 1 & \texttt{a\textbar{}b\textbar{}b\textbar{}b} & \texttt{ab} & \texttt{ab} & \texttt{ab} & \texttt{0123} & \texttt{0} & \texttt{+} & \texttt{.} & \texttt{a@0}\\
N895 & F4.2876 & \texttt{T-{}-{}} & 1 & \texttt{a\textbar{}a\textbar{}b\textbar{}b} & \texttt{ab} & \texttt{ab} & \texttt{ab} & \texttt{0123} & \texttt{0} & \texttt{+} & \texttt{.} & \texttt{a@0}\\
N896 & F4.2876 & \texttt{T-{}-{}} & 1 & \texttt{a\textbar{}a\textbar{}a\textbar{}a} & \texttt{a} & \texttt{a} & \texttt{a} & \texttt{0123} & \texttt{0} & \texttt{+} & \texttt{S} & \texttt{none}\\
N897 & F4.2877 & \texttt{T-{}-{}} & 1 & \texttt{a\textbar{}a\textbar{}b\textbar{}b} & \texttt{ab} & \texttt{ab} & \texttt{ab} & \texttt{0123} & \texttt{0} & \texttt{+} & \texttt{.} & \texttt{a@0}\\
N898 & F4.2877 & \texttt{T-{}-{}} & 1 & \texttt{a\textbar{}a\textbar{}a\textbar{}a} & \texttt{a} & \texttt{a} & \texttt{a} & \texttt{0123} & \texttt{0} & \texttt{+} & \texttt{S} & \texttt{none}\\
N899 & F4.2878 & \texttt{T-{}-{}} & 1 & \texttt{a\textbar{}b\textbar{}b\textbar{}b} & \texttt{ab} & \texttt{ab} & \texttt{ab} & \texttt{0123} & \texttt{0} & \texttt{+} & \texttt{.} & \texttt{a@0}\\
N900 & F4.2878 & \texttt{T-{}-{}} & 1 & \texttt{a\textbar{}a\textbar{}b\textbar{}b} & \texttt{ab} & \texttt{ab} & \texttt{ab} & \texttt{0123} & \texttt{0} & \texttt{+} & \texttt{.} & \texttt{a@0}\\
N901 & F4.2878 & \texttt{T-{}-{}} & 1 & \texttt{a\textbar{}a\textbar{}a\textbar{}a} & \texttt{a} & \texttt{a} & \texttt{a} & \texttt{0123} & \texttt{0} & \texttt{+} & \texttt{S} & \texttt{none}\\
N902 & F4.2879 & \texttt{T-{}-{}} & 1 & \texttt{a\textbar{}b\textbar{}c\textbar{}c} & \texttt{ac} & \texttt{ac} & \texttt{abc} & \texttt{023} & \texttt{0} & \texttt{-{}1} & \texttt{.} & \texttt{a@0}\\
N903 & F4.2879 & \texttt{T-{}-{}} & 1 & \texttt{a\textbar{}b\textbar{}b\textbar{}b} & \texttt{ab} & \texttt{ab} & \texttt{ab} & \texttt{0123} & \texttt{0} & \texttt{+} & \texttt{.} & \texttt{a@0}\\
N904 & F4.2879 & \texttt{T-{}-{}} & 1 & \texttt{a\textbar{}b\textbar{}a\textbar{}a} & \texttt{a} & \texttt{a} & \texttt{ab} & \texttt{023} & \texttt{0} & \texttt{-{}1} & \texttt{s} & \texttt{b@1}\\
N905 & F4.2879 & \texttt{T-{}-{}} & 1 & \texttt{a\textbar{}a\textbar{}b\textbar{}b} & \texttt{ab} & \texttt{ab} & \texttt{ab} & \texttt{0123} & \texttt{0} & \texttt{+} & \texttt{.} & \texttt{a@0}\\
N906 & F4.2879 & \texttt{T-{}-{}} & 1 & \texttt{a\textbar{}a\textbar{}a\textbar{}a} & \texttt{a} & \texttt{a} & \texttt{a} & \texttt{0123} & \texttt{0} & \texttt{+} & \texttt{S} & \texttt{none}\\
N907 & F4.2880 & \texttt{T-{}-{}} & 1 & \texttt{a\textbar{}a\textbar{}b\textbar{}b} & \texttt{ab} & \texttt{ab} & \texttt{ab} & \texttt{0123} & \texttt{0} & \texttt{+} & \texttt{.} & \texttt{a@0}\\
N908 & F4.2880 & \texttt{T-{}-{}} & 1 & \texttt{a\textbar{}b\textbar{}a\textbar{}b} & \texttt{ab} & \texttt{ab} & \texttt{ab} & \texttt{0123} & \texttt{0} & \texttt{+} & \texttt{.} & \texttt{a@0}\\
N909 & F4.2880 & \texttt{T-{}-{}} & 1 & \texttt{a\textbar{}a\textbar{}a\textbar{}a} & \texttt{a} & \texttt{a} & \texttt{a} & \texttt{0123} & \texttt{0} & \texttt{+} & \texttt{S} & \texttt{none}\\
N910 & F4.2881 & \texttt{T-{}-{}} & 1 & \texttt{a\textbar{}a\textbar{}b\textbar{}b} & \texttt{ab} & \texttt{ab} & \texttt{ab} & \texttt{0123} & \texttt{0} & \texttt{+} & \texttt{.} & \texttt{a@0}\\
N911 & F4.2881 & \texttt{T-{}-{}} & 1 & \texttt{a\textbar{}b\textbar{}a\textbar{}b} & \texttt{ab} & \texttt{ab} & \texttt{ab} & \texttt{0123} & \texttt{0} & \texttt{+} & \texttt{.} & \texttt{a@0}\\
N912 & F4.2881 & \texttt{T-{}-{}} & 1 & \texttt{a\textbar{}a\textbar{}a\textbar{}a} & \texttt{a} & \texttt{a} & \texttt{a} & \texttt{0123} & \texttt{0} & \texttt{+} & \texttt{S} & \texttt{none}\\
N913 & F4.2882 & \texttt{T-{}-{}} & 1 & \texttt{a\textbar{}a\textbar{}b\textbar{}b} & \texttt{ab} & \texttt{ab} & \texttt{ab} & \texttt{0123} & \texttt{0} & \texttt{+} & \texttt{.} & \texttt{a@0}\\
N914 & F4.2882 & \texttt{T-{}-{}} & 1 & \texttt{a\textbar{}b\textbar{}a\textbar{}b} & \texttt{ab} & \texttt{ab} & \texttt{ab} & \texttt{0123} & \texttt{0} & \texttt{+} & \texttt{.} & \texttt{a@0}\\
N915 & F4.2882 & \texttt{T-{}-{}} & 1 & \texttt{a\textbar{}a\textbar{}a\textbar{}a} & \texttt{a} & \texttt{a} & \texttt{a} & \texttt{0123} & \texttt{0} & \texttt{+} & \texttt{S} & \texttt{none}\\
N916 & F4.2883 & \texttt{T-{}-{}} & 1 & \texttt{a\textbar{}b\textbar{}b\textbar{}b} & \texttt{ab} & \texttt{ab} & \texttt{ab} & \texttt{0123} & \texttt{0} & \texttt{+} & \texttt{.} & \texttt{a@0}\\
N917 & F4.2883 & \texttt{T-{}-{}} & 1 & \texttt{a\textbar{}b\textbar{}a\textbar{}b} & \texttt{ab} & \texttt{ab} & \texttt{ab} & \texttt{0123} & \texttt{0} & \texttt{+} & \texttt{.} & \texttt{a@0}\\
N918 & F4.2883 & \texttt{T-{}-{}} & 1 & \texttt{a\textbar{}a\textbar{}a\textbar{}b} & \texttt{a} & \texttt{a} & \texttt{ab} & \texttt{012} & \texttt{0} & \texttt{-{}3} & \texttt{s} & \texttt{b@3}\\
N919 & F4.2883 & \texttt{T-{}-{}} & 1 & \texttt{a\textbar{}a\textbar{}a\textbar{}a} & \texttt{a} & \texttt{a} & \texttt{a} & \texttt{0123} & \texttt{0} & \texttt{+} & \texttt{S} & \texttt{none}\\
N920 & F4.2884 & \texttt{T-{}-{}} & 1 & \texttt{a\textbar{}b\textbar{}b\textbar{}b} & \texttt{ab} & \texttt{ab} & \texttt{ab} & \texttt{0123} & \texttt{0} & \texttt{+} & \texttt{.} & \texttt{a@0}\\
N921 & F4.2884 & \texttt{T-{}-{}} & 1 & \texttt{a\textbar{}a\textbar{}b\textbar{}b} & \texttt{ab} & \texttt{ab} & \texttt{ab} & \texttt{0123} & \texttt{0} & \texttt{+} & \texttt{.} & \texttt{a@0}\\
N922 & F4.2884 & \texttt{T-{}-{}} & 1 & \texttt{a\textbar{}b\textbar{}a\textbar{}b} & \texttt{ab} & \texttt{ab} & \texttt{ab} & \texttt{0123} & \texttt{0} & \texttt{+} & \texttt{.} & \texttt{a@0}\\
N923 & F4.2884 & \texttt{T-{}-{}} & 1 & \texttt{a\textbar{}a\textbar{}a\textbar{}a} & \texttt{a} & \texttt{a} & \texttt{a} & \texttt{0123} & \texttt{0} & \texttt{+} & \texttt{S} & \texttt{none}\\
N924 & F4.2885 & \texttt{T-{}-{}} & 1 & \texttt{a\textbar{}a\textbar{}b\textbar{}b} & \texttt{ab} & \texttt{ab} & \texttt{ab} & \texttt{0123} & \texttt{0} & \texttt{+} & \texttt{.} & \texttt{a@0}\\
N925 & F4.2885 & \texttt{T-{}-{}} & 1 & \texttt{a\textbar{}a\textbar{}a\textbar{}a} & \texttt{a} & \texttt{a} & \texttt{a} & \texttt{0123} & \texttt{0} & \texttt{+} & \texttt{S} & \texttt{none}\\
N926 & F4.2886 & \texttt{T-{}-{}} & 1 & \texttt{a\textbar{}b\textbar{}b\textbar{}b} & \texttt{ab} & \texttt{ab} & \texttt{ab} & \texttt{0123} & \texttt{0} & \texttt{+} & \texttt{.} & \texttt{a@0}\\
N927 & F4.2886 & \texttt{T-{}-{}} & 1 & \texttt{a\textbar{}a\textbar{}b\textbar{}b} & \texttt{ab} & \texttt{ab} & \texttt{ab} & \texttt{0123} & \texttt{0} & \texttt{+} & \texttt{.} & \texttt{a@0}\\
N928 & F4.2886 & \texttt{T-{}-{}} & 1 & \texttt{a\textbar{}b\textbar{}a\textbar{}b} & \texttt{ab} & \texttt{ab} & \texttt{ab} & \texttt{0123} & \texttt{0} & \texttt{+} & \texttt{.} & \texttt{a@0}\\
N929 & F4.2886 & \texttt{T-{}-{}} & 1 & \texttt{a\textbar{}a\textbar{}a\textbar{}a} & \texttt{a} & \texttt{a} & \texttt{a} & \texttt{0123} & \texttt{0} & \texttt{+} & \texttt{S} & \texttt{none}\\
N930 & F4.2887 & \texttt{T-{}-{}} & 1 & \texttt{a\textbar{}b\textbar{}c\textbar{}b} & \texttt{ab} & \texttt{ab} & \texttt{abc} & \texttt{013} & \texttt{0} & \texttt{-{}2} & \texttt{.} & \texttt{a@0}\\
N931 & F4.2887 & \texttt{T-{}-{}} & 1 & \texttt{a\textbar{}b\textbar{}b\textbar{}b} & \texttt{ab} & \texttt{ab} & \texttt{ab} & \texttt{0123} & \texttt{0} & \texttt{+} & \texttt{.} & \texttt{a@0}\\
N932 & F4.2887 & \texttt{T-{}-{}} & 1 & \texttt{a\textbar{}a\textbar{}b\textbar{}b} & \texttt{ab} & \texttt{ab} & \texttt{ab} & \texttt{0123} & \texttt{0} & \texttt{+} & \texttt{.} & \texttt{a@0}\\
N933 & F4.2887 & \texttt{T-{}-{}} & 1 & \texttt{a\textbar{}a\textbar{}b\textbar{}a} & \texttt{a} & \texttt{a} & \texttt{ab} & \texttt{013} & \texttt{0} & \texttt{-{}2} & \texttt{s} & \texttt{b@2}\\
N934 & F4.2887 & \texttt{T-{}-{}} & 1 & \texttt{a\textbar{}b\textbar{}a\textbar{}b} & \texttt{ab} & \texttt{ab} & \texttt{ab} & \texttt{0123} & \texttt{0} & \texttt{+} & \texttt{.} & \texttt{a@0}\\
N935 & F4.2887 & \texttt{T-{}-{}} & 1 & \texttt{a\textbar{}a\textbar{}a\textbar{}a} & \texttt{a} & \texttt{a} & \texttt{a} & \texttt{0123} & \texttt{0} & \texttt{+} & \texttt{S} & \texttt{none}\\
N936 & F4.2888 & \texttt{T-{}-{}} & 1 & \texttt{a\textbar{}b\textbar{}b\textbar{}b} & \texttt{ab} & \texttt{ab} & \texttt{ab} & \texttt{0123} & \texttt{0} & \texttt{+} & \texttt{.} & \texttt{a@0}\\
N937 & F4.2888 & \texttt{T-{}-{}} & 1 & \texttt{a\textbar{}a\textbar{}b\textbar{}b} & \texttt{ab} & \texttt{ab} & \texttt{ab} & \texttt{0123} & \texttt{0} & \texttt{+} & \texttt{.} & \texttt{a@0}\\
N938 & F4.2888 & \texttt{T-{}-{}} & 1 & \texttt{a\textbar{}b\textbar{}a\textbar{}b} & \texttt{ab} & \texttt{ab} & \texttt{ab} & \texttt{0123} & \texttt{0} & \texttt{+} & \texttt{.} & \texttt{a@0}\\
N939 & F4.2888 & \texttt{T-{}-{}} & 1 & \texttt{a\textbar{}a\textbar{}a\textbar{}a} & \texttt{a} & \texttt{a} & \texttt{a} & \texttt{0123} & \texttt{0} & \texttt{+} & \texttt{S} & \texttt{none}\\
N940 & F4.2889 & \texttt{T-{}-{}} & 1 & \texttt{a\textbar{}a\textbar{}b\textbar{}b} & \texttt{ab} & \texttt{ab} & \texttt{ab} & \texttt{0123} & \texttt{0} & \texttt{+} & \texttt{.} & \texttt{a@0}\\
N941 & F4.2889 & \texttt{T-{}-{}} & 1 & \texttt{a\textbar{}b\textbar{}a\textbar{}b} & \texttt{ab} & \texttt{ab} & \texttt{ab} & \texttt{0123} & \texttt{0} & \texttt{+} & \texttt{.} & \texttt{a@0}\\
N942 & F4.2889 & \texttt{T-{}-{}} & 1 & \texttt{a\textbar{}a\textbar{}a\textbar{}a} & \texttt{a} & \texttt{a} & \texttt{a} & \texttt{0123} & \texttt{0} & \texttt{+} & \texttt{S} & \texttt{none}\\
N943 & F4.2890 & \texttt{T-{}-{}} & 1 & \texttt{a\textbar{}a\textbar{}b\textbar{}b} & \texttt{ab} & \texttt{ab} & \texttt{ab} & \texttt{0123} & \texttt{0} & \texttt{+} & \texttt{.} & \texttt{a@0}\\
N944 & F4.2890 & \texttt{T-{}-{}} & 1 & \texttt{a\textbar{}b\textbar{}a\textbar{}b} & \texttt{ab} & \texttt{ab} & \texttt{ab} & \texttt{0123} & \texttt{0} & \texttt{+} & \texttt{.} & \texttt{a@0}\\
N945 & F4.2890 & \texttt{T-{}-{}} & 1 & \texttt{a\textbar{}a\textbar{}a\textbar{}a} & \texttt{a} & \texttt{a} & \texttt{a} & \texttt{0123} & \texttt{0} & \texttt{+} & \texttt{S} & \texttt{none}\\
N946 & F4.2891 & \texttt{T-{}-{}} & 1 & \texttt{a\textbar{}a\textbar{}b\textbar{}b} & \texttt{ab} & \texttt{ab} & \texttt{ab} & \texttt{0123} & \texttt{0} & \texttt{+} & \texttt{.} & \texttt{a@0}\\
N947 & F4.2891 & \texttt{T-{}-{}} & 1 & \texttt{a\textbar{}b\textbar{}a\textbar{}b} & \texttt{ab} & \texttt{ab} & \texttt{ab} & \texttt{0123} & \texttt{0} & \texttt{+} & \texttt{.} & \texttt{a@0}\\
N948 & F4.2891 & \texttt{T-{}-{}} & 1 & \texttt{a\textbar{}a\textbar{}a\textbar{}a} & \texttt{a} & \texttt{a} & \texttt{a} & \texttt{0123} & \texttt{0} & \texttt{+} & \texttt{S} & \texttt{none}\\
N949 & F4.2892 & \texttt{T-{}-{}} & 1 & \texttt{a\textbar{}b\textbar{}b\textbar{}b} & \texttt{ab} & \texttt{ab} & \texttt{ab} & \texttt{0123} & \texttt{0} & \texttt{+} & \texttt{.} & \texttt{a@0}\\
N950 & F4.2892 & \texttt{T-{}-{}} & 1 & \texttt{a\textbar{}b\textbar{}a\textbar{}b} & \texttt{ab} & \texttt{ab} & \texttt{ab} & \texttt{0123} & \texttt{0} & \texttt{+} & \texttt{.} & \texttt{a@0}\\
N951 & F4.2892 & \texttt{T-{}-{}} & 1 & \texttt{a\textbar{}a\textbar{}a\textbar{}a} & \texttt{a} & \texttt{a} & \texttt{a} & \texttt{0123} & \texttt{0} & \texttt{+} & \texttt{S} & \texttt{none}\\
N952 & F4.2893 & \texttt{T-{}-{}} & 1 & \texttt{a\textbar{}b\textbar{}b\textbar{}b} & \texttt{ab} & \texttt{ab} & \texttt{ab} & \texttt{0123} & \texttt{0} & \texttt{+} & \texttt{.} & \texttt{a@0}\\
N953 & F4.2893 & \texttt{T-{}-{}} & 1 & \texttt{a\textbar{}a\textbar{}b\textbar{}b} & \texttt{ab} & \texttt{ab} & \texttt{ab} & \texttt{0123} & \texttt{0} & \texttt{+} & \texttt{.} & \texttt{a@0}\\
N954 & F4.2893 & \texttt{T-{}-{}} & 1 & \texttt{a\textbar{}a\textbar{}a\textbar{}a} & \texttt{a} & \texttt{a} & \texttt{a} & \texttt{0123} & \texttt{0} & \texttt{+} & \texttt{S} & \texttt{none}\\
N955 & F4.2894 & \texttt{T-{}-{}} & 1 & \texttt{a\textbar{}b\textbar{}b\textbar{}b} & \texttt{ab} & \texttt{ab} & \texttt{ab} & \texttt{0123} & \texttt{0} & \texttt{+} & \texttt{.} & \texttt{a@0}\\
N956 & F4.2894 & \texttt{T-{}-{}} & 1 & \texttt{a\textbar{}b\textbar{}c\textbar{}a} & \texttt{a} & \texttt{a} & \texttt{abc} & \texttt{03} & \texttt{0} & \texttt{-{}1} & \texttt{s} & \texttt{b@1}\\
N957 & F4.2894 & \texttt{T-{}-{}} & 1 & \texttt{a\textbar{}a\textbar{}b\textbar{}b} & \texttt{ab} & \texttt{ab} & \texttt{ab} & \texttt{0123} & \texttt{0} & \texttt{+} & \texttt{.} & \texttt{a@0}\\
N958 & F4.2894 & \texttt{T-{}-{}} & 1 & \texttt{a\textbar{}b\textbar{}a\textbar{}b} & \texttt{ab} & \texttt{ab} & \texttt{ab} & \texttt{0123} & \texttt{0} & \texttt{+} & \texttt{.} & \texttt{a@0}\\
N959 & F4.2894 & \texttt{T-{}-{}} & 1 & \texttt{a\textbar{}b\textbar{}b\textbar{}a} & \texttt{a} & \texttt{a} & \texttt{ab} & \texttt{03} & \texttt{0} & \texttt{-{}1} & \texttt{s} & \texttt{b@1}\\
N960 & F4.2894 & \texttt{T-{}-{}} & 1 & \texttt{a\textbar{}a\textbar{}a\textbar{}a} & \texttt{a} & \texttt{a} & \texttt{a} & \texttt{0123} & \texttt{0} & \texttt{+} & \texttt{S} & \texttt{none}\\
N961 & F4.2895 & \texttt{T-{}-{}} & 1 & \texttt{a\textbar{}b\textbar{}a\textbar{}b} & \texttt{ab} & \texttt{ab} & \texttt{ab} & \texttt{0123} & \texttt{0} & \texttt{+} & \texttt{.} & \texttt{a@0}\\
N962 & F4.2895 & \texttt{T-{}-{}} & 1 & \texttt{a\textbar{}a\textbar{}a\textbar{}b} & \texttt{ab} & \texttt{ab} & \texttt{ab} & \texttt{0123} & \texttt{0} & \texttt{+} & \texttt{.} & \texttt{b@3}\\
N963 & F4.2895 & \texttt{T-{}-{}} & 1 & \texttt{a\textbar{}a\textbar{}a\textbar{}a} & \texttt{a} & \texttt{a} & \texttt{a} & \texttt{0123} & \texttt{0} & \texttt{+} & \texttt{S} & \texttt{none}\\
N964 & F4.2896 & \texttt{T-{}-{}} & 1 & \texttt{a\textbar{}b\textbar{}a\textbar{}b} & \texttt{ab} & \texttt{ab} & \texttt{ab} & \texttt{0123} & \texttt{0} & \texttt{+} & \texttt{.} & \texttt{a@0}\\
N965 & F4.2896 & \texttt{T-{}-{}} & 1 & \texttt{a\textbar{}a\textbar{}a\textbar{}a} & \texttt{a} & \texttt{a} & \texttt{a} & \texttt{0123} & \texttt{0} & \texttt{+} & \texttt{S} & \texttt{none}\\
N966 & F4.2897 & \texttt{T-{}-{}} & 1 & \texttt{a\textbar{}b\textbar{}a\textbar{}c} & \texttt{ac} & \texttt{ac} & \texttt{abc} & \texttt{023} & \texttt{0} & \texttt{-{}1} & \texttt{.} & \texttt{b@1}\\
N967 & F4.2897 & \texttt{T-{}-{}} & 1 & \texttt{a\textbar{}b\textbar{}a\textbar{}a} & \texttt{a} & \texttt{a} & \texttt{ab} & \texttt{023} & \texttt{0} & \texttt{-{}1} & \texttt{s} & \texttt{b@1}\\
N968 & F4.2897 & \texttt{T-{}-{}} & 1 & \texttt{a\textbar{}b\textbar{}a\textbar{}b} & \texttt{ab} & \texttt{ab} & \texttt{ab} & \texttt{0123} & \texttt{0} & \texttt{+} & \texttt{.} & \texttt{a@0}\\
N969 & F4.2897 & \texttt{T-{}-{}} & 1 & \texttt{a\textbar{}a\textbar{}a\textbar{}b} & \texttt{ab} & \texttt{ab} & \texttt{ab} & \texttt{0123} & \texttt{0} & \texttt{+} & \texttt{.} & \texttt{b@3}\\
N970 & F4.2897 & \texttt{T-{}-{}} & 1 & \texttt{a\textbar{}a\textbar{}a\textbar{}a} & \texttt{a} & \texttt{a} & \texttt{a} & \texttt{0123} & \texttt{0} & \texttt{+} & \texttt{S} & \texttt{none}\\
N971 & F4.2898 & \texttt{T-{}-{}} & 1 & \texttt{a\textbar{}b\textbar{}b\textbar{}b} & \texttt{ab} & \texttt{ab} & \texttt{ab} & \texttt{0123} & \texttt{0} & \texttt{+} & \texttt{.} & \texttt{a@0}\\
N972 & F4.2898 & \texttt{T-{}-{}} & 1 & \texttt{a\textbar{}b\textbar{}a\textbar{}b} & \texttt{ab} & \texttt{ab} & \texttt{ab} & \texttt{0123} & \texttt{0} & \texttt{+} & \texttt{.} & \texttt{a@0}\\
N973 & F4.2898 & \texttt{T-{}-{}} & 1 & \texttt{a\textbar{}a\textbar{}a\textbar{}a} & \texttt{a} & \texttt{a} & \texttt{a} & \texttt{0123} & \texttt{0} & \texttt{+} & \texttt{S} & \texttt{none}\\
N974 & F4.2899 & \texttt{T-{}-{}} & 1 & \texttt{a\textbar{}b\textbar{}c\textbar{}b} & \texttt{ab} & \texttt{ab} & \texttt{abc} & \texttt{013} & \texttt{0} & \texttt{-{}2} & \texttt{.} & \texttt{a@0}\\
N975 & F4.2899 & \texttt{T-{}-{}} & 1 & \texttt{a\textbar{}b\textbar{}b\textbar{}b} & \texttt{ab} & \texttt{ab} & \texttt{ab} & \texttt{0123} & \texttt{0} & \texttt{+} & \texttt{.} & \texttt{a@0}\\
N976 & F4.2899 & \texttt{T-{}-{}} & 1 & \texttt{a\textbar{}a\textbar{}b\textbar{}c} & \texttt{ac} & \texttt{ac} & \texttt{abc} & \texttt{013} & \texttt{0} & \texttt{-{}2} & \texttt{.} & \texttt{b@2}\\
N977 & F4.2899 & \texttt{T-{}-{}} & 1 & \texttt{a\textbar{}a\textbar{}b\textbar{}b} & \texttt{ab} & \texttt{ab} & \texttt{ab} & \texttt{0123} & \texttt{0} & \texttt{+} & \texttt{.} & \texttt{a@0}\\
N978 & F4.2899 & \texttt{T-{}-{}} & 1 & \texttt{a\textbar{}a\textbar{}b\textbar{}a} & \texttt{a} & \texttt{a} & \texttt{ab} & \texttt{013} & \texttt{0} & \texttt{-{}2} & \texttt{s} & \texttt{b@2}\\
N979 & F4.2899 & \texttt{T-{}-{}} & 1 & \texttt{a\textbar{}b\textbar{}a\textbar{}b} & \texttt{ab} & \texttt{ab} & \texttt{ab} & \texttt{0123} & \texttt{0} & \texttt{+} & \texttt{.} & \texttt{a@0}\\
N980 & F4.2899 & \texttt{T-{}-{}} & 1 & \texttt{a\textbar{}a\textbar{}a\textbar{}b} & \texttt{ab} & \texttt{ab} & \texttt{ab} & \texttt{0123} & \texttt{0} & \texttt{+} & \texttt{.} & \texttt{b@3}\\
N981 & F4.2899 & \texttt{T-{}-{}} & 1 & \texttt{a\textbar{}a\textbar{}a\textbar{}a} & \texttt{a} & \texttt{a} & \texttt{a} & \texttt{0123} & \texttt{0} & \texttt{+} & \texttt{S} & \texttt{none}\\
N982 & F4.2900 & \texttt{T-{}-{}} & 1 & \texttt{a\textbar{}a\textbar{}b\textbar{}b} & \texttt{ab} & \texttt{ab} & \texttt{ab} & \texttt{0123} & \texttt{0} & \texttt{+} & \texttt{.} & \texttt{a@0}\\
N983 & F4.2900 & \texttt{T-{}-{}} & 1 & \texttt{a\textbar{}b\textbar{}a\textbar{}b} & \texttt{ab} & \texttt{ab} & \texttt{ab} & \texttt{0123} & \texttt{0} & \texttt{+} & \texttt{.} & \texttt{a@0}\\
N984 & F4.2900 & \texttt{T-{}-{}} & 1 & \texttt{a\textbar{}a\textbar{}a\textbar{}a} & \texttt{a} & \texttt{a} & \texttt{a} & \texttt{0123} & \texttt{0} & \texttt{+} & \texttt{S} & \texttt{none}\\
N985 & F4.2901 & \texttt{T-{}-{}} & 1 & \texttt{a\textbar{}a\textbar{}b\textbar{}b} & \texttt{ab} & \texttt{ab} & \texttt{ab} & \texttt{0123} & \texttt{0} & \texttt{+} & \texttt{.} & \texttt{a@0}\\
N986 & F4.2901 & \texttt{T-{}-{}} & 1 & \texttt{a\textbar{}b\textbar{}a\textbar{}b} & \texttt{ab} & \texttt{ab} & \texttt{ab} & \texttt{0123} & \texttt{0} & \texttt{+} & \texttt{.} & \texttt{a@0}\\
N987 & F4.2901 & \texttt{T-{}-{}} & 1 & \texttt{a\textbar{}a\textbar{}a\textbar{}b} & \texttt{ab} & \texttt{ab} & \texttt{ab} & \texttt{0123} & \texttt{0} & \texttt{+} & \texttt{.} & \texttt{b@3}\\
N988 & F4.2901 & \texttt{T-{}-{}} & 1 & \texttt{a\textbar{}a\textbar{}a\textbar{}a} & \texttt{a} & \texttt{a} & \texttt{a} & \texttt{0123} & \texttt{0} & \texttt{+} & \texttt{S} & \texttt{none}\\
N989 & F4.2902 & \texttt{T-{}-{}} & 1 & \texttt{a\textbar{}b\textbar{}a\textbar{}c} & \texttt{ac} & \texttt{ac} & \texttt{abc} & \texttt{023} & \texttt{0} & \texttt{-{}1} & \texttt{.} & \texttt{b@1}\\
N990 & F4.2902 & \texttt{T-{}-{}} & 1 & \texttt{a\textbar{}b\textbar{}a\textbar{}a} & \texttt{a} & \texttt{a} & \texttt{ab} & \texttt{023} & \texttt{0} & \texttt{-{}1} & \texttt{s} & \texttt{b@1}\\
N991 & F4.2902 & \texttt{T-{}-{}} & 1 & \texttt{a\textbar{}a\textbar{}b\textbar{}b} & \texttt{ab} & \texttt{ab} & \texttt{ab} & \texttt{0123} & \texttt{0} & \texttt{+} & \texttt{.} & \texttt{a@0}\\
N992 & F4.2902 & \texttt{T-{}-{}} & 1 & \texttt{a\textbar{}b\textbar{}a\textbar{}b} & \texttt{ab} & \texttt{ab} & \texttt{ab} & \texttt{0123} & \texttt{0} & \texttt{+} & \texttt{.} & \texttt{a@0}\\
N993 & F4.2902 & \texttt{T-{}-{}} & 1 & \texttt{a\textbar{}a\textbar{}a\textbar{}b} & \texttt{ab} & \texttt{ab} & \texttt{ab} & \texttt{0123} & \texttt{0} & \texttt{+} & \texttt{.} & \texttt{b@3}\\
N994 & F4.2902 & \texttt{T-{}-{}} & 1 & \texttt{a\textbar{}a\textbar{}a\textbar{}a} & \texttt{a} & \texttt{a} & \texttt{a} & \texttt{0123} & \texttt{0} & \texttt{+} & \texttt{S} & \texttt{none}\\
N995 & F4.2903 & \texttt{T-{}4} & 1 & \texttt{a\textbar{}b\textbar{}c\textbar{}d} & \texttt{ad} & \texttt{ad} & \texttt{abcd} & \texttt{03} & \texttt{0} & \texttt{-{}1} & \texttt{.} & \texttt{a@0}\\
N996 & F4.2903 & \texttt{T-{}4} & 1 & \texttt{a\textbar{}b\textbar{}a\textbar{}c} & \texttt{ac} & \texttt{ac} & \texttt{abc} & \texttt{023} & \texttt{0} & \texttt{-{}1} & \texttt{.} & \texttt{b@1}\\
N997 & F4.2903 & \texttt{T-{}4} & 1 & \texttt{a\textbar{}b\textbar{}b\textbar{}c} & \texttt{ac} & \texttt{ac} & \texttt{abc} & \texttt{03} & \texttt{0} & \texttt{-{}1} & \texttt{.} & \texttt{a@0}\\
N998 & F4.2903 & \texttt{T-{}4} & 1 & \texttt{a\textbar{}a\textbar{}a\textbar{}b} & \texttt{ab} & \texttt{ab} & \texttt{ab} & \texttt{0123} & \texttt{0} & \texttt{+} & \texttt{.} & \texttt{b@3}\\
N999 & F4.2903 & \texttt{T-{}4} & 1 & \texttt{a\textbar{}b\textbar{}c\textbar{}a} & \texttt{a} & \texttt{a} & \texttt{abc} & \texttt{03} & \texttt{0} & \texttt{-{}1} & \texttt{s} & \texttt{b@1}\\
N1000 & F4.2903 & \texttt{T-{}4} & 1 & \texttt{a\textbar{}b\textbar{}a\textbar{}a} & \texttt{a} & \texttt{a} & \texttt{ab} & \texttt{023} & \texttt{0} & \texttt{-{}1} & \texttt{s} & \texttt{b@1}\\
N1001 & F4.2903 & \texttt{T-{}4} & 1 & \texttt{a\textbar{}b\textbar{}a\textbar{}b} & \texttt{ab} & \texttt{ab} & \texttt{ab} & \texttt{0123} & \texttt{0} & \texttt{+} & \texttt{.} & \texttt{a@0}\\
N1002 & F4.2903 & \texttt{T-{}4} & 1 & \texttt{a\textbar{}b\textbar{}b\textbar{}a} & \texttt{a} & \texttt{a} & \texttt{ab} & \texttt{03} & \texttt{0} & \texttt{-{}1} & \texttt{s} & \texttt{b@1}\\
N1003 & F4.2903 & \texttt{T-{}4} & 1 & \texttt{a\textbar{}a\textbar{}a\textbar{}a} & \texttt{a} & \texttt{a} & \texttt{a} & \texttt{0123} & \texttt{0} & \texttt{+} & \texttt{S} & \texttt{none}\\
N1004 & F4.2904 & \texttt{T-{}-{}} & 1 & \texttt{a\textbar{}b\textbar{}a\textbar{}b} & \texttt{ab} & \texttt{ab} & \texttt{ab} & \texttt{0123} & \texttt{0} & \texttt{+} & \texttt{.} & \texttt{a@0}\\
N1005 & F4.2904 & \texttt{T-{}-{}} & 1 & \texttt{a\textbar{}a\textbar{}a\textbar{}a} & \texttt{a} & \texttt{a} & \texttt{a} & \texttt{0123} & \texttt{0} & \texttt{+} & \texttt{S} & \texttt{none}\\
N1006 & F4.2905 & \texttt{T-{}-{}} & 1 & \texttt{a\textbar{}a\textbar{}b\textbar{}b} & \texttt{ab} & \texttt{ab} & \texttt{ab} & \texttt{0123} & \texttt{0} & \texttt{+} & \texttt{.} & \texttt{a@0}\\
N1007 & F4.2905 & \texttt{T-{}-{}} & 1 & \texttt{a\textbar{}b\textbar{}a\textbar{}b} & \texttt{ab} & \texttt{ab} & \texttt{ab} & \texttt{0123} & \texttt{0} & \texttt{+} & \texttt{.} & \texttt{a@0}\\
N1008 & F4.2905 & \texttt{T-{}-{}} & 1 & \texttt{a\textbar{}a\textbar{}a\textbar{}b} & \texttt{ab} & \texttt{ab} & \texttt{ab} & \texttt{0123} & \texttt{0} & \texttt{+} & \texttt{.} & \texttt{b@3}\\
N1009 & F4.2905 & \texttt{T-{}-{}} & 1 & \texttt{a\textbar{}a\textbar{}a\textbar{}a} & \texttt{a} & \texttt{a} & \texttt{a} & \texttt{0123} & \texttt{0} & \texttt{+} & \texttt{S} & \texttt{none}\\
N1010 & F4.2974 & \texttt{T-{}-{}} & 1 & \texttt{a\textbar{}a\textbar{}a\textbar{}a} & \texttt{a} & \texttt{a} & \texttt{a} & \texttt{0123} & \texttt{0} & \texttt{+} & \texttt{S} & \texttt{none}\\
N1011 & F4.2975 & \texttt{T-{}-{}} & 1 & \texttt{a\textbar{}b\textbar{}c\textbar{}d} & \texttt{d} & \texttt{d} & \texttt{abcd} & \texttt{3} & \texttt{-{}} & \texttt{-{}0} & \texttt{s} & \texttt{a@0}\\
N1012 & F4.2975 & \texttt{T-{}-{}} & 1 & \texttt{a\textbar{}b\textbar{}b\textbar{}c} & \texttt{c} & \texttt{c} & \texttt{abc} & \texttt{3} & \texttt{-{}} & \texttt{-{}0} & \texttt{s} & \texttt{a@0}\\
N1013 & F4.2975 & \texttt{T-{}-{}} & 1 & \texttt{a\textbar{}a\textbar{}b\textbar{}c} & \texttt{c} & \texttt{c} & \texttt{abc} & \texttt{3} & \texttt{-{}} & \texttt{-{}0} & \texttt{s} & \texttt{b@2}\\
N1014 & F4.2975 & \texttt{T-{}-{}} & 1 & \texttt{a\textbar{}a\textbar{}a\textbar{}b} & \texttt{b} & \texttt{b} & \texttt{ab} & \texttt{3} & \texttt{-{}} & \texttt{-{}0} & \texttt{s} & \texttt{a@0}\\
N1015 & F4.2975 & \texttt{T-{}-{}} & 1 & \texttt{a\textbar{}a\textbar{}a\textbar{}a} & \texttt{a} & \texttt{a} & \texttt{a} & \texttt{0123} & \texttt{0} & \texttt{+} & \texttt{S} & \texttt{none}\\
N1016 & F4.2976 & \texttt{T-{}-{}} & 1 & \texttt{a\textbar{}b\textbar{}b\textbar{}b} & \texttt{ab} & \texttt{ab} & \texttt{ab} & \texttt{0123} & \texttt{0} & \texttt{+} & \texttt{.} & \texttt{a@0}\\
N1017 & F4.2976 & \texttt{T-{}-{}} & 1 & \texttt{a\textbar{}a\textbar{}b\textbar{}b} & \texttt{ab} & \texttt{ab} & \texttt{ab} & \texttt{0123} & \texttt{0} & \texttt{+} & \texttt{.} & \texttt{a@0}\\
N1018 & F4.2976 & \texttt{T-{}-{}} & 1 & \texttt{a\textbar{}a\textbar{}a\textbar{}a} & \texttt{a} & \texttt{a} & \texttt{a} & \texttt{0123} & \texttt{0} & \texttt{+} & \texttt{S} & \texttt{none}\\
N1019 & F4.2977 & \texttt{T-{}-{}} & 1 & \texttt{a\textbar{}a\textbar{}b\textbar{}b} & \texttt{ab} & \texttt{ab} & \texttt{ab} & \texttt{0123} & \texttt{0} & \texttt{+} & \texttt{.} & \texttt{a@0}\\
N1020 & F4.2977 & \texttt{T-{}-{}} & 1 & \texttt{a\textbar{}b\textbar{}a\textbar{}b} & \texttt{ab} & \texttt{ab} & \texttt{ab} & \texttt{0123} & \texttt{0} & \texttt{+} & \texttt{.} & \texttt{a@0}\\
N1021 & F4.2977 & \texttt{T-{}-{}} & 1 & \texttt{a\textbar{}a\textbar{}a\textbar{}a} & \texttt{a} & \texttt{a} & \texttt{a} & \texttt{0123} & \texttt{0} & \texttt{+} & \texttt{S} & \texttt{none}\\
N1022 & F4.2979 & \texttt{T-{}-{}} & 1 & \texttt{a\textbar{}b\textbar{}b\textbar{}b} & \texttt{ab} & \texttt{ab} & \texttt{ab} & \texttt{0123} & \texttt{0} & \texttt{+} & \texttt{.} & \texttt{a@0}\\
N1023 & F4.2979 & \texttt{T-{}-{}} & 1 & \texttt{a\textbar{}a\textbar{}b\textbar{}b} & \texttt{ab} & \texttt{ab} & \texttt{ab} & \texttt{0123} & \texttt{0} & \texttt{+} & \texttt{.} & \texttt{a@0}\\
N1024 & F4.2979 & \texttt{T-{}-{}} & 1 & \texttt{a\textbar{}b\textbar{}a\textbar{}b} & \texttt{ab} & \texttt{ab} & \texttt{ab} & \texttt{0123} & \texttt{0} & \texttt{+} & \texttt{.} & \texttt{a@0}\\
N1025 & F4.2979 & \texttt{T-{}-{}} & 1 & \texttt{a\textbar{}a\textbar{}a\textbar{}a} & \texttt{a} & \texttt{a} & \texttt{a} & \texttt{0123} & \texttt{0} & \texttt{+} & \texttt{S} & \texttt{none}\\
N1026 & F4.2980 & \texttt{T-{}-{}} & 1 & \texttt{a\textbar{}a\textbar{}b\textbar{}b} & \texttt{ab} & \texttt{ab} & \texttt{ab} & \texttt{0123} & \texttt{0} & \texttt{+} & \texttt{.} & \texttt{a@0}\\
N1027 & F4.2980 & \texttt{T-{}-{}} & 1 & \texttt{a\textbar{}b\textbar{}a\textbar{}b} & \texttt{ab} & \texttt{ab} & \texttt{ab} & \texttt{0123} & \texttt{0} & \texttt{+} & \texttt{.} & \texttt{a@0}\\
N1028 & F4.2980 & \texttt{T-{}-{}} & 1 & \texttt{a\textbar{}a\textbar{}a\textbar{}a} & \texttt{a} & \texttt{a} & \texttt{a} & \texttt{0123} & \texttt{0} & \texttt{+} & \texttt{S} & \texttt{none}\\
N1029 & F4.2981 & \texttt{T-{}-{}} & 1 & \texttt{a\textbar{}b\textbar{}b\textbar{}b} & \texttt{ab} & \texttt{ab} & \texttt{ab} & \texttt{0123} & \texttt{0} & \texttt{+} & \texttt{.} & \texttt{a@0}\\
N1030 & F4.2981 & \texttt{T-{}-{}} & 1 & \texttt{a\textbar{}a\textbar{}b\textbar{}b} & \texttt{ab} & \texttt{ab} & \texttt{ab} & \texttt{0123} & \texttt{0} & \texttt{+} & \texttt{.} & \texttt{a@0}\\
N1031 & F4.2981 & \texttt{T-{}-{}} & 1 & \texttt{a\textbar{}b\textbar{}a\textbar{}b} & \texttt{ab} & \texttt{ab} & \texttt{ab} & \texttt{0123} & \texttt{0} & \texttt{+} & \texttt{.} & \texttt{a@0}\\
N1032 & F4.2981 & \texttt{T-{}-{}} & 1 & \texttt{a\textbar{}a\textbar{}a\textbar{}a} & \texttt{a} & \texttt{a} & \texttt{a} & \texttt{0123} & \texttt{0} & \texttt{+} & \texttt{S} & \texttt{none}\\
N1033 & F4.2982 & \texttt{T-{}-{}} & 1 & \texttt{a\textbar{}b\textbar{}b\textbar{}b} & \texttt{ab} & \texttt{ab} & \texttt{ab} & \texttt{0123} & \texttt{0} & \texttt{+} & \texttt{.} & \texttt{a@0}\\
N1034 & F4.2982 & \texttt{T-{}-{}} & 1 & \texttt{a\textbar{}a\textbar{}b\textbar{}b} & \texttt{ab} & \texttt{ab} & \texttt{ab} & \texttt{0123} & \texttt{0} & \texttt{+} & \texttt{.} & \texttt{a@0}\\
N1035 & F4.2982 & \texttt{T-{}-{}} & 1 & \texttt{a\textbar{}b\textbar{}a\textbar{}b} & \texttt{ab} & \texttt{ab} & \texttt{ab} & \texttt{0123} & \texttt{0} & \texttt{+} & \texttt{.} & \texttt{a@0}\\
N1036 & F4.2982 & \texttt{T-{}-{}} & 1 & \texttt{a\textbar{}a\textbar{}a\textbar{}a} & \texttt{a} & \texttt{a} & \texttt{a} & \texttt{0123} & \texttt{0} & \texttt{+} & \texttt{S} & \texttt{none}\\
N1037 & F4.2983 & \texttt{T-{}-{}} & 1 & \texttt{a\textbar{}b\textbar{}b\textbar{}b} & \texttt{ab} & \texttt{ab} & \texttt{ab} & \texttt{0123} & \texttt{0} & \texttt{+} & \texttt{.} & \texttt{a@0}\\
N1038 & F4.2983 & \texttt{T-{}-{}} & 1 & \texttt{a\textbar{}a\textbar{}b\textbar{}b} & \texttt{ab} & \texttt{ab} & \texttt{ab} & \texttt{0123} & \texttt{0} & \texttt{+} & \texttt{.} & \texttt{a@0}\\
N1039 & F4.2983 & \texttt{T-{}-{}} & 1 & \texttt{a\textbar{}b\textbar{}a\textbar{}b} & \texttt{ab} & \texttt{ab} & \texttt{ab} & \texttt{0123} & \texttt{0} & \texttt{+} & \texttt{.} & \texttt{a@0}\\
N1040 & F4.2983 & \texttt{T-{}-{}} & 1 & \texttt{a\textbar{}a\textbar{}a\textbar{}a} & \texttt{a} & \texttt{a} & \texttt{a} & \texttt{0123} & \texttt{0} & \texttt{+} & \texttt{S} & \texttt{none}\\
N1041 & F4.2984 & \texttt{T-{}-{}} & 1 & \texttt{a\textbar{}b\textbar{}a\textbar{}b} & \texttt{ab} & \texttt{ab} & \texttt{ab} & \texttt{0123} & \texttt{0} & \texttt{+} & \texttt{.} & \texttt{a@0}\\
N1042 & F4.2984 & \texttt{T-{}-{}} & 1 & \texttt{a\textbar{}a\textbar{}a\textbar{}b} & \texttt{ab} & \texttt{ab} & \texttt{ab} & \texttt{0123} & \texttt{0} & \texttt{+} & \texttt{.} & \texttt{b@3}\\
N1043 & F4.2984 & \texttt{T-{}-{}} & 1 & \texttt{a\textbar{}a\textbar{}a\textbar{}a} & \texttt{a} & \texttt{a} & \texttt{a} & \texttt{0123} & \texttt{0} & \texttt{+} & \texttt{S} & \texttt{none}\\
N1044 & F4.2985 & \texttt{T-{}-{}} & 1 & \texttt{a\textbar{}b\textbar{}b\textbar{}b} & \texttt{ab} & \texttt{ab} & \texttt{ab} & \texttt{0123} & \texttt{0} & \texttt{+} & \texttt{.} & \texttt{a@0}\\
N1045 & F4.2985 & \texttt{T-{}-{}} & 1 & \texttt{a\textbar{}b\textbar{}a\textbar{}b} & \texttt{ab} & \texttt{ab} & \texttt{ab} & \texttt{0123} & \texttt{0} & \texttt{+} & \texttt{.} & \texttt{a@0}\\
N1046 & F4.2985 & \texttt{T-{}-{}} & 1 & \texttt{a\textbar{}a\textbar{}a\textbar{}b} & \texttt{ab} & \texttt{ab} & \texttt{ab} & \texttt{0123} & \texttt{0} & \texttt{+} & \texttt{.} & \texttt{b@3}\\
N1047 & F4.2985 & \texttt{T-{}-{}} & 1 & \texttt{a\textbar{}a\textbar{}a\textbar{}a} & \texttt{a} & \texttt{a} & \texttt{a} & \texttt{0123} & \texttt{0} & \texttt{+} & \texttt{S} & \texttt{none}\\
N1048 & F4.2986 & \texttt{T-{}-{}} & 1 & \texttt{a\textbar{}a\textbar{}b\textbar{}b} & \texttt{ab} & \texttt{ab} & \texttt{ab} & \texttt{0123} & \texttt{0} & \texttt{+} & \texttt{.} & \texttt{a@0}\\
N1049 & F4.2986 & \texttt{T-{}-{}} & 1 & \texttt{a\textbar{}b\textbar{}a\textbar{}b} & \texttt{ab} & \texttt{ab} & \texttt{ab} & \texttt{0123} & \texttt{0} & \texttt{+} & \texttt{.} & \texttt{a@0}\\
N1050 & F4.2986 & \texttt{T-{}-{}} & 1 & \texttt{a\textbar{}a\textbar{}a\textbar{}b} & \texttt{ab} & \texttt{ab} & \texttt{ab} & \texttt{0123} & \texttt{0} & \texttt{+} & \texttt{.} & \texttt{b@3}\\
N1051 & F4.2986 & \texttt{T-{}-{}} & 1 & \texttt{a\textbar{}a\textbar{}a\textbar{}a} & \texttt{a} & \texttt{a} & \texttt{a} & \texttt{0123} & \texttt{0} & \texttt{+} & \texttt{S} & \texttt{none}\\
N1052 & F4.2987 & \texttt{T-{}-{}} & 1 & \texttt{a\textbar{}b\textbar{}b\textbar{}b} & \texttt{ab} & \texttt{ab} & \texttt{ab} & \texttt{0123} & \texttt{0} & \texttt{+} & \texttt{.} & \texttt{a@0}\\
N1053 & F4.2987 & \texttt{T-{}-{}} & 1 & \texttt{a\textbar{}a\textbar{}b\textbar{}b} & \texttt{ab} & \texttt{ab} & \texttt{ab} & \texttt{0123} & \texttt{0} & \texttt{+} & \texttt{.} & \texttt{a@0}\\
N1054 & F4.2987 & \texttt{T-{}-{}} & 1 & \texttt{a\textbar{}b\textbar{}a\textbar{}b} & \texttt{ab} & \texttt{ab} & \texttt{ab} & \texttt{0123} & \texttt{0} & \texttt{+} & \texttt{.} & \texttt{a@0}\\
N1055 & F4.2987 & \texttt{T-{}-{}} & 1 & \texttt{a\textbar{}a\textbar{}a\textbar{}a} & \texttt{a} & \texttt{a} & \texttt{a} & \texttt{0123} & \texttt{0} & \texttt{+} & \texttt{S} & \texttt{none}\\
N1056 & F4.2988 & \texttt{T-{}-{}} & 1 & \texttt{a\textbar{}b\textbar{}b\textbar{}b} & \texttt{ab} & \texttt{ab} & \texttt{ab} & \texttt{0123} & \texttt{0} & \texttt{+} & \texttt{.} & \texttt{a@0}\\
N1057 & F4.2988 & \texttt{T-{}-{}} & 1 & \texttt{a\textbar{}a\textbar{}b\textbar{}b} & \texttt{ab} & \texttt{ab} & \texttt{ab} & \texttt{0123} & \texttt{0} & \texttt{+} & \texttt{.} & \texttt{a@0}\\
N1058 & F4.2988 & \texttt{T-{}-{}} & 1 & \texttt{a\textbar{}b\textbar{}a\textbar{}b} & \texttt{ab} & \texttt{ab} & \texttt{ab} & \texttt{0123} & \texttt{0} & \texttt{+} & \texttt{.} & \texttt{a@0}\\
N1059 & F4.2988 & \texttt{T-{}-{}} & 1 & \texttt{a\textbar{}a\textbar{}a\textbar{}b} & \texttt{ab} & \texttt{ab} & \texttt{ab} & \texttt{0123} & \texttt{0} & \texttt{+} & \texttt{.} & \texttt{b@3}\\
N1060 & F4.2988 & \texttt{T-{}-{}} & 1 & \texttt{a\textbar{}a\textbar{}a\textbar{}a} & \texttt{a} & \texttt{a} & \texttt{a} & \texttt{0123} & \texttt{0} & \texttt{+} & \texttt{S} & \texttt{none}\\
N1061 & F4.2989 & \texttt{T-{}-{}} & 1 & \texttt{a\textbar{}b\textbar{}c\textbar{}c} & \texttt{ac} & \texttt{ac} & \texttt{abc} & \texttt{023} & \texttt{0} & \texttt{-{}1} & \texttt{.} & \texttt{a@0}\\
N1062 & F4.2989 & \texttt{T-{}-{}} & 1 & \texttt{a\textbar{}b\textbar{}b\textbar{}b} & \texttt{ab} & \texttt{ab} & \texttt{ab} & \texttt{0123} & \texttt{0} & \texttt{+} & \texttt{.} & \texttt{a@0}\\
N1063 & F4.2989 & \texttt{T-{}-{}} & 1 & \texttt{a\textbar{}b\textbar{}a\textbar{}c} & \texttt{ac} & \texttt{ac} & \texttt{abc} & \texttt{023} & \texttt{0} & \texttt{-{}1} & \texttt{.} & \texttt{b@1}\\
N1064 & F4.2989 & \texttt{T-{}-{}} & 1 & \texttt{a\textbar{}b\textbar{}a\textbar{}a} & \texttt{a} & \texttt{a} & \texttt{ab} & \texttt{023} & \texttt{0} & \texttt{-{}1} & \texttt{s} & \texttt{b@1}\\
N1065 & F4.2989 & \texttt{T-{}-{}} & 1 & \texttt{a\textbar{}a\textbar{}b\textbar{}b} & \texttt{ab} & \texttt{ab} & \texttt{ab} & \texttt{0123} & \texttt{0} & \texttt{+} & \texttt{.} & \texttt{a@0}\\
N1066 & F4.2989 & \texttt{T-{}-{}} & 1 & \texttt{a\textbar{}b\textbar{}a\textbar{}b} & \texttt{ab} & \texttt{ab} & \texttt{ab} & \texttt{0123} & \texttt{0} & \texttt{+} & \texttt{.} & \texttt{a@0}\\
N1067 & F4.2989 & \texttt{T-{}-{}} & 1 & \texttt{a\textbar{}a\textbar{}a\textbar{}b} & \texttt{ab} & \texttt{ab} & \texttt{ab} & \texttt{0123} & \texttt{0} & \texttt{+} & \texttt{.} & \texttt{b@3}\\
N1068 & F4.2989 & \texttt{T-{}-{}} & 1 & \texttt{a\textbar{}a\textbar{}a\textbar{}a} & \texttt{a} & \texttt{a} & \texttt{a} & \texttt{0123} & \texttt{0} & \texttt{+} & \texttt{S} & \texttt{none}\\
N1069 & F4.2990 & \texttt{T-{}-{}} & 1 & \texttt{a\textbar{}b\textbar{}a\textbar{}b} & \texttt{ab} & \texttt{ab} & \texttt{ab} & \texttt{0123} & \texttt{0} & \texttt{+} & \texttt{.} & \texttt{a@0}\\
N1070 & F4.2990 & \texttt{T-{}-{}} & 1 & \texttt{a\textbar{}a\textbar{}a\textbar{}b} & \texttt{ab} & \texttt{ab} & \texttt{ab} & \texttt{0123} & \texttt{0} & \texttt{+} & \texttt{.} & \texttt{b@3}\\
N1071 & F4.2990 & \texttt{T-{}-{}} & 1 & \texttt{a\textbar{}a\textbar{}a\textbar{}a} & \texttt{a} & \texttt{a} & \texttt{a} & \texttt{0123} & \texttt{0} & \texttt{+} & \texttt{S} & \texttt{none}\\
N1072 & F4.2991 & \texttt{T-{}-{}} & 1 & \texttt{a\textbar{}a\textbar{}b\textbar{}b} & \texttt{ab} & \texttt{ab} & \texttt{ab} & \texttt{0123} & \texttt{0} & \texttt{+} & \texttt{.} & \texttt{a@0}\\
N1073 & F4.2991 & \texttt{T-{}-{}} & 1 & \texttt{a\textbar{}b\textbar{}a\textbar{}b} & \texttt{ab} & \texttt{ab} & \texttt{ab} & \texttt{0123} & \texttt{0} & \texttt{+} & \texttt{.} & \texttt{a@0}\\
N1074 & F4.2991 & \texttt{T-{}-{}} & 1 & \texttt{a\textbar{}a\textbar{}a\textbar{}a} & \texttt{a} & \texttt{a} & \texttt{a} & \texttt{0123} & \texttt{0} & \texttt{+} & \texttt{S} & \texttt{none}\\
N1075 & F4.2992 & \texttt{T-{}-{}} & 1 & \texttt{a\textbar{}a\textbar{}b\textbar{}b} & \texttt{ab} & \texttt{ab} & \texttt{ab} & \texttt{0123} & \texttt{0} & \texttt{+} & \texttt{.} & \texttt{a@0}\\
N1076 & F4.2992 & \texttt{T-{}-{}} & 1 & \texttt{a\textbar{}b\textbar{}a\textbar{}b} & \texttt{ab} & \texttt{ab} & \texttt{ab} & \texttt{0123} & \texttt{0} & \texttt{+} & \texttt{.} & \texttt{a@0}\\
N1077 & F4.2992 & \texttt{T-{}-{}} & 1 & \texttt{a\textbar{}a\textbar{}a\textbar{}b} & \texttt{ab} & \texttt{ab} & \texttt{ab} & \texttt{0123} & \texttt{0} & \texttt{+} & \texttt{.} & \texttt{b@3}\\
N1078 & F4.2992 & \texttt{T-{}-{}} & 1 & \texttt{a\textbar{}a\textbar{}a\textbar{}a} & \texttt{a} & \texttt{a} & \texttt{a} & \texttt{0123} & \texttt{0} & \texttt{+} & \texttt{S} & \texttt{none}\\
N1079 & F4.2993 & \texttt{T-{}-{}} & 1 & \texttt{a\textbar{}a\textbar{}b\textbar{}b} & \texttt{ab} & \texttt{ab} & \texttt{ab} & \texttt{0123} & \texttt{0} & \texttt{+} & \texttt{.} & \texttt{a@0}\\
N1080 & F4.2993 & \texttt{T-{}-{}} & 1 & \texttt{a\textbar{}b\textbar{}a\textbar{}b} & \texttt{ab} & \texttt{ab} & \texttt{ab} & \texttt{0123} & \texttt{0} & \texttt{+} & \texttt{.} & \texttt{a@0}\\
N1081 & F4.2993 & \texttt{T-{}-{}} & 1 & \texttt{a\textbar{}a\textbar{}a\textbar{}b} & \texttt{ab} & \texttt{ab} & \texttt{ab} & \texttt{0123} & \texttt{0} & \texttt{+} & \texttt{.} & \texttt{b@3}\\
N1082 & F4.2993 & \texttt{T-{}-{}} & 1 & \texttt{a\textbar{}a\textbar{}a\textbar{}a} & \texttt{a} & \texttt{a} & \texttt{a} & \texttt{0123} & \texttt{0} & \texttt{+} & \texttt{S} & \texttt{none}\\
N1083 & F4.2994 & \texttt{T-{}4} & 1 & \texttt{a\textbar{}b\textbar{}c\textbar{}d} & \texttt{ad} & \texttt{ad} & \texttt{abcd} & \texttt{03} & \texttt{0} & \texttt{-{}1} & \texttt{.} & \texttt{a@0}\\
N1084 & F4.2994 & \texttt{T-{}4} & 1 & \texttt{a\textbar{}b\textbar{}b\textbar{}c} & \texttt{ac} & \texttt{ac} & \texttt{abc} & \texttt{03} & \texttt{0} & \texttt{-{}1} & \texttt{.} & \texttt{a@0}\\
N1085 & F4.2994 & \texttt{T-{}4} & 1 & \texttt{a\textbar{}a\textbar{}a\textbar{}b} & \texttt{ab} & \texttt{ab} & \texttt{ab} & \texttt{0123} & \texttt{0} & \texttt{+} & \texttt{.} & \texttt{b@3}\\
N1086 & F4.2994 & \texttt{T-{}4} & 1 & \texttt{a\textbar{}b\textbar{}c\textbar{}a} & \texttt{a} & \texttt{a} & \texttt{abc} & \texttt{03} & \texttt{0} & \texttt{-{}1} & \texttt{s} & \texttt{b@1}\\
N1087 & F4.2994 & \texttt{T-{}4} & 1 & \texttt{a\textbar{}b\textbar{}a\textbar{}b} & \texttt{ab} & \texttt{ab} & \texttt{ab} & \texttt{0123} & \texttt{0} & \texttt{+} & \texttt{.} & \texttt{a@0}\\
N1088 & F4.2994 & \texttt{T-{}4} & 1 & \texttt{a\textbar{}b\textbar{}b\textbar{}a} & \texttt{a} & \texttt{a} & \texttt{ab} & \texttt{03} & \texttt{0} & \texttt{-{}1} & \texttt{s} & \texttt{b@1}\\
N1089 & F4.2994 & \texttt{T-{}4} & 1 & \texttt{a\textbar{}a\textbar{}a\textbar{}a} & \texttt{a} & \texttt{a} & \texttt{a} & \texttt{0123} & \texttt{0} & \texttt{+} & \texttt{S} & \texttt{none}\\
N1090 & F4.2995 & \texttt{T-{}-{}} & 1 & \texttt{a\textbar{}b\textbar{}c\textbar{}c} & \texttt{abc} & \texttt{abc} & \texttt{abc} & \texttt{0123} & \texttt{0} & \texttt{+} & \texttt{.} & \texttt{a@0}\\
N1091 & F4.2995 & \texttt{T-{}-{}} & 1 & \texttt{a\textbar{}b\textbar{}a\textbar{}c} & \texttt{ab} & \texttt{ab} & \texttt{abc} & \texttt{012} & \texttt{0} & \texttt{-{}3} & \texttt{.} & \texttt{b@1}\\
N1092 & F4.2995 & \texttt{T-{}-{}} & 1 & \texttt{a\textbar{}b\textbar{}a\textbar{}a} & \texttt{ab} & \texttt{ab} & \texttt{ab} & \texttt{0123} & \texttt{0} & \texttt{+} & \texttt{.} & \texttt{b@1}\\
N1093 & F4.2995 & \texttt{T-{}-{}} & 1 & \texttt{a\textbar{}a\textbar{}b\textbar{}b} & \texttt{ab} & \texttt{ab} & \texttt{ab} & \texttt{0123} & \texttt{0} & \texttt{+} & \texttt{.} & \texttt{a@0}\\
N1094 & F4.2995 & \texttt{T-{}-{}} & 1 & \texttt{a\textbar{}b\textbar{}b\textbar{}a} & \texttt{ab} & \texttt{ab} & \texttt{ab} & \texttt{0123} & \texttt{0} & \texttt{+} & \texttt{.} & \texttt{b@1}\\
N1095 & F4.2995 & \texttt{T-{}-{}} & 1 & \texttt{a\textbar{}a\textbar{}a\textbar{}b} & \texttt{a} & \texttt{a} & \texttt{ab} & \texttt{012} & \texttt{0} & \texttt{-{}3} & \texttt{s} & \texttt{b@3}\\
N1096 & F4.2995 & \texttt{T-{}-{}} & 1 & \texttt{a\textbar{}a\textbar{}a\textbar{}a} & \texttt{a} & \texttt{a} & \texttt{a} & \texttt{0123} & \texttt{0} & \texttt{+} & \texttt{S} & \texttt{none}\\
N1097 & F4.2996 & \texttt{T-{}-{}} & 1 & \texttt{a\textbar{}b\textbar{}c\textbar{}c} & \texttt{abc} & \texttt{abc} & \texttt{abc} & \texttt{0123} & \texttt{0} & \texttt{+} & \texttt{.} & \texttt{a@0}\\
N1098 & F4.2996 & \texttt{T-{}-{}} & 1 & \texttt{a\textbar{}b\textbar{}a\textbar{}a} & \texttt{ab} & \texttt{ab} & \texttt{ab} & \texttt{0123} & \texttt{0} & \texttt{+} & \texttt{.} & \texttt{b@1}\\
N1099 & F4.2996 & \texttt{T-{}-{}} & 1 & \texttt{a\textbar{}a\textbar{}b\textbar{}b} & \texttt{ab} & \texttt{ab} & \texttt{ab} & \texttt{0123} & \texttt{0} & \texttt{+} & \texttt{.} & \texttt{a@0}\\
N1100 & F4.2996 & \texttt{T-{}-{}} & 1 & \texttt{a\textbar{}b\textbar{}a\textbar{}b} & \texttt{ab} & \texttt{ab} & \texttt{ab} & \texttt{0123} & \texttt{0} & \texttt{+} & \texttt{.} & \texttt{a@0}\\
N1101 & F4.2996 & \texttt{T-{}-{}} & 1 & \texttt{a\textbar{}b\textbar{}b\textbar{}a} & \texttt{ab} & \texttt{ab} & \texttt{ab} & \texttt{0123} & \texttt{0} & \texttt{+} & \texttt{.} & \texttt{b@1}\\
N1102 & F4.2996 & \texttt{T-{}-{}} & 1 & \texttt{a\textbar{}a\textbar{}a\textbar{}a} & \texttt{a} & \texttt{a} & \texttt{a} & \texttt{0123} & \texttt{0} & \texttt{+} & \texttt{S} & \texttt{none}\\
N1103 & F4.2997 & \texttt{T-{}-{}} & 1 & \texttt{a\textbar{}b\textbar{}b\textbar{}b} & \texttt{ab} & \texttt{ab} & \texttt{ab} & \texttt{0123} & \texttt{0} & \texttt{+} & \texttt{.} & \texttt{a@0}\\
N1104 & F4.2997 & \texttt{T-{}-{}} & 1 & \texttt{a\textbar{}a\textbar{}b\textbar{}b} & \texttt{ab} & \texttt{ab} & \texttt{ab} & \texttt{0123} & \texttt{0} & \texttt{+} & \texttt{.} & \texttt{a@0}\\
N1105 & F4.2997 & \texttt{T-{}-{}} & 1 & \texttt{a\textbar{}b\textbar{}a\textbar{}b} & \texttt{ab} & \texttt{ab} & \texttt{ab} & \texttt{0123} & \texttt{0} & \texttt{+} & \texttt{.} & \texttt{a@0}\\
N1106 & F4.2997 & \texttt{T-{}-{}} & 1 & \texttt{a\textbar{}b\textbar{}b\textbar{}a} & \texttt{ab} & \texttt{ab} & \texttt{ab} & \texttt{0123} & \texttt{0} & \texttt{+} & \texttt{.} & \texttt{b@1}\\
N1107 & F4.2997 & \texttt{T-{}-{}} & 1 & \texttt{a\textbar{}a\textbar{}a\textbar{}a} & \texttt{a} & \texttt{a} & \texttt{a} & \texttt{0123} & \texttt{0} & \texttt{+} & \texttt{S} & \texttt{none}\\
N1108 & F4.2999 & \texttt{T-{}-{}} & 1 & \texttt{a\textbar{}b\textbar{}b\textbar{}b} & \texttt{ab} & \texttt{ab} & \texttt{ab} & \texttt{0123} & \texttt{0} & \texttt{+} & \texttt{.} & \texttt{a@0}\\
N1109 & F4.2999 & \texttt{T-{}-{}} & 1 & \texttt{a\textbar{}a\textbar{}b\textbar{}b} & \texttt{ab} & \texttt{ab} & \texttt{ab} & \texttt{0123} & \texttt{0} & \texttt{+} & \texttt{.} & \texttt{a@0}\\
N1110 & F4.2999 & \texttt{T-{}-{}} & 1 & \texttt{a\textbar{}b\textbar{}a\textbar{}b} & \texttt{ab} & \texttt{ab} & \texttt{ab} & \texttt{0123} & \texttt{0} & \texttt{+} & \texttt{.} & \texttt{a@0}\\
N1111 & F4.2999 & \texttt{T-{}-{}} & 1 & \texttt{a\textbar{}b\textbar{}b\textbar{}a} & \texttt{ab} & \texttt{ab} & \texttt{ab} & \texttt{0123} & \texttt{0} & \texttt{+} & \texttt{.} & \texttt{b@1}\\
N1112 & F4.2999 & \texttt{T-{}-{}} & 1 & \texttt{a\textbar{}a\textbar{}a\textbar{}a} & \texttt{a} & \texttt{a} & \texttt{a} & \texttt{0123} & \texttt{0} & \texttt{+} & \texttt{S} & \texttt{none}\\
N1113 & F4.3000 & \texttt{T-{}-{}} & 1 & \texttt{a\textbar{}b\textbar{}b\textbar{}a} & \texttt{ab} & \texttt{ab} & \texttt{ab} & \texttt{0123} & \texttt{0} & \texttt{+} & \texttt{.} & \texttt{b@1}\\
N1114 & F4.3000 & \texttt{T-{}-{}} & 1 & \texttt{a\textbar{}b\textbar{}a\textbar{}b} & \texttt{ab} & \texttt{ab} & \texttt{ab} & \texttt{0123} & \texttt{0} & \texttt{+} & \texttt{.} & \texttt{a@0}\\
N1115 & F4.3000 & \texttt{T-{}-{}} & 1 & \texttt{a\textbar{}a\textbar{}a\textbar{}a} & \texttt{a} & \texttt{a} & \texttt{a} & \texttt{0123} & \texttt{0} & \texttt{+} & \texttt{S} & \texttt{none}\\
N1116 & F4.3020 & \texttt{T-{}4} & 1 & \texttt{a\textbar{}b\textbar{}c\textbar{}d} & \texttt{d} & \texttt{d} & \texttt{abcd} & \texttt{3} & \texttt{-{}} & \texttt{-{}0} & \texttt{s} & \texttt{a@0}\\
N1117 & F4.3020 & \texttt{T-{}4} & 1 & \texttt{a\textbar{}a\textbar{}b\textbar{}c} & \texttt{c} & \texttt{c} & \texttt{abc} & \texttt{3} & \texttt{-{}} & \texttt{-{}0} & \texttt{s} & \texttt{b@2}\\
N1118 & F4.3020 & \texttt{T-{}4} & 1 & \texttt{a\textbar{}a\textbar{}a\textbar{}b} & \texttt{b} & \texttt{b} & \texttt{ab} & \texttt{3} & \texttt{-{}} & \texttt{-{}0} & \texttt{s} & \texttt{a@0}\\
N1119 & F4.3020 & \texttt{T-{}4} & 1 & \texttt{a\textbar{}a\textbar{}a\textbar{}a} & \texttt{a} & \texttt{a} & \texttt{a} & \texttt{0123} & \texttt{0} & \texttt{+} & \texttt{S} & \texttt{none}\\
N1120 & F4.3021 & \texttt{T-{}-{}} & 1 & \texttt{a\textbar{}b\textbar{}b\textbar{}b} & \texttt{ab} & \texttt{ab} & \texttt{ab} & \texttt{0123} & \texttt{0} & \texttt{+} & \texttt{.} & \texttt{a@0}\\
N1121 & F4.3021 & \texttt{T-{}-{}} & 1 & \texttt{a\textbar{}a\textbar{}b\textbar{}b} & \texttt{ab} & \texttt{ab} & \texttt{ab} & \texttt{0123} & \texttt{0} & \texttt{+} & \texttt{.} & \texttt{a@0}\\
N1122 & F4.3021 & \texttt{T-{}-{}} & 1 & \texttt{a\textbar{}b\textbar{}a\textbar{}b} & \texttt{ab} & \texttt{ab} & \texttt{ab} & \texttt{0123} & \texttt{0} & \texttt{+} & \texttt{.} & \texttt{a@0}\\
N1123 & F4.3021 & \texttt{T-{}-{}} & 1 & \texttt{a\textbar{}a\textbar{}a\textbar{}a} & \texttt{a} & \texttt{a} & \texttt{a} & \texttt{0123} & \texttt{0} & \texttt{+} & \texttt{S} & \texttt{none}\\
N1124 & F4.3022 & \texttt{T-{}-{}} & 1 & \texttt{a\textbar{}b\textbar{}b\textbar{}b} & \texttt{ab} & \texttt{ab} & \texttt{ab} & \texttt{0123} & \texttt{0} & \texttt{+} & \texttt{.} & \texttt{a@0}\\
N1125 & F4.3022 & \texttt{T-{}-{}} & 1 & \texttt{a\textbar{}a\textbar{}b\textbar{}b} & \texttt{ab} & \texttt{ab} & \texttt{ab} & \texttt{0123} & \texttt{0} & \texttt{+} & \texttt{.} & \texttt{a@0}\\
N1126 & F4.3022 & \texttt{T-{}-{}} & 1 & \texttt{a\textbar{}b\textbar{}a\textbar{}b} & \texttt{ab} & \texttt{ab} & \texttt{ab} & \texttt{0123} & \texttt{0} & \texttt{+} & \texttt{.} & \texttt{a@0}\\
N1127 & F4.3022 & \texttt{T-{}-{}} & 1 & \texttt{a\textbar{}a\textbar{}a\textbar{}b} & \texttt{ab} & \texttt{ab} & \texttt{ab} & \texttt{0123} & \texttt{0} & \texttt{+} & \texttt{.} & \texttt{b@3}\\
N1128 & F4.3022 & \texttt{T-{}-{}} & 1 & \texttt{a\textbar{}a\textbar{}a\textbar{}a} & \texttt{a} & \texttt{a} & \texttt{a} & \texttt{0123} & \texttt{0} & \texttt{+} & \texttt{S} & \texttt{none}\\
N1129 & F4.3023 & \texttt{T-{}-{}} & 1 & \texttt{a\textbar{}a\textbar{}b\textbar{}b} & \texttt{ab} & \texttt{ab} & \texttt{ab} & \texttt{0123} & \texttt{0} & \texttt{+} & \texttt{.} & \texttt{a@0}\\
N1130 & F4.3023 & \texttt{T-{}-{}} & 1 & \texttt{a\textbar{}b\textbar{}a\textbar{}b} & \texttt{ab} & \texttt{ab} & \texttt{ab} & \texttt{0123} & \texttt{0} & \texttt{+} & \texttt{.} & \texttt{a@0}\\
N1131 & F4.3023 & \texttt{T-{}-{}} & 1 & \texttt{a\textbar{}a\textbar{}a\textbar{}b} & \texttt{ab} & \texttt{ab} & \texttt{ab} & \texttt{0123} & \texttt{0} & \texttt{+} & \texttt{.} & \texttt{b@3}\\
N1132 & F4.3023 & \texttt{T-{}-{}} & 1 & \texttt{a\textbar{}a\textbar{}a\textbar{}a} & \texttt{a} & \texttt{a} & \texttt{a} & \texttt{0123} & \texttt{0} & \texttt{+} & \texttt{S} & \texttt{none}\\
N1133 & F4.3024 & \texttt{T-{}-{}} & 1 & \texttt{a\textbar{}b\textbar{}b\textbar{}b} & \texttt{ab} & \texttt{ab} & \texttt{ab} & \texttt{0123} & \texttt{0} & \texttt{+} & \texttt{.} & \texttt{a@0}\\
N1134 & F4.3024 & \texttt{T-{}-{}} & 1 & \texttt{a\textbar{}b\textbar{}a\textbar{}b} & \texttt{ab} & \texttt{ab} & \texttt{ab} & \texttt{0123} & \texttt{0} & \texttt{+} & \texttt{.} & \texttt{a@0}\\
N1135 & F4.3024 & \texttt{T-{}-{}} & 1 & \texttt{a\textbar{}a\textbar{}a\textbar{}a} & \texttt{a} & \texttt{a} & \texttt{a} & \texttt{0123} & \texttt{0} & \texttt{+} & \texttt{S} & \texttt{none}\\
N1136 & F4.3025 & \texttt{T-{}-{}} & 1 & \texttt{a\textbar{}b\textbar{}b\textbar{}b} & \texttt{ab} & \texttt{ab} & \texttt{ab} & \texttt{0123} & \texttt{0} & \texttt{+} & \texttt{.} & \texttt{a@0}\\
N1137 & F4.3025 & \texttt{T-{}-{}} & 1 & \texttt{a\textbar{}a\textbar{}b\textbar{}b} & \texttt{ab} & \texttt{ab} & \texttt{ab} & \texttt{0123} & \texttt{0} & \texttt{+} & \texttt{.} & \texttt{a@0}\\
N1138 & F4.3025 & \texttt{T-{}-{}} & 1 & \texttt{a\textbar{}b\textbar{}a\textbar{}b} & \texttt{ab} & \texttt{ab} & \texttt{ab} & \texttt{0123} & \texttt{0} & \texttt{+} & \texttt{.} & \texttt{a@0}\\
N1139 & F4.3025 & \texttt{T-{}-{}} & 1 & \texttt{a\textbar{}a\textbar{}a\textbar{}b} & \texttt{ab} & \texttt{ab} & \texttt{ab} & \texttt{0123} & \texttt{0} & \texttt{+} & \texttt{.} & \texttt{b@3}\\
N1140 & F4.3025 & \texttt{T-{}-{}} & 1 & \texttt{a\textbar{}a\textbar{}a\textbar{}a} & \texttt{a} & \texttt{a} & \texttt{a} & \texttt{0123} & \texttt{0} & \texttt{+} & \texttt{S} & \texttt{none}\\
N1141 & F4.3026 & \texttt{T-{}-{}} & 1 & \texttt{a\textbar{}b\textbar{}c\textbar{}c} & \texttt{abc} & \texttt{abc} & \texttt{abc} & \texttt{0123} & \texttt{0} & \texttt{+} & \texttt{.} & \texttt{a@0}\\
N1142 & F4.3026 & \texttt{T-{}-{}} & 1 & \texttt{a\textbar{}b\textbar{}b\textbar{}b} & \texttt{ab} & \texttt{ab} & \texttt{ab} & \texttt{0123} & \texttt{0} & \texttt{+} & \texttt{.} & \texttt{a@0}\\
N1143 & F4.3026 & \texttt{T-{}-{}} & 1 & \texttt{a\textbar{}b\textbar{}a\textbar{}a} & \texttt{ab} & \texttt{ab} & \texttt{ab} & \texttt{0123} & \texttt{0} & \texttt{+} & \texttt{.} & \texttt{b@1}\\
N1144 & F4.3026 & \texttt{T-{}-{}} & 1 & \texttt{a\textbar{}a\textbar{}b\textbar{}b} & \texttt{ab} & \texttt{ab} & \texttt{ab} & \texttt{0123} & \texttt{0} & \texttt{+} & \texttt{.} & \texttt{a@0}\\
N1145 & F4.3026 & \texttt{T-{}-{}} & 1 & \texttt{a\textbar{}b\textbar{}a\textbar{}b} & \texttt{ab} & \texttt{ab} & \texttt{ab} & \texttt{0123} & \texttt{0} & \texttt{+} & \texttt{.} & \texttt{a@0}\\
N1146 & F4.3026 & \texttt{T-{}-{}} & 1 & \texttt{a\textbar{}b\textbar{}b\textbar{}a} & \texttt{ab} & \texttt{ab} & \texttt{ab} & \texttt{0123} & \texttt{0} & \texttt{+} & \texttt{.} & \texttt{b@1}\\
N1147 & F4.3026 & \texttt{T-{}-{}} & 1 & \texttt{a\textbar{}a\textbar{}a\textbar{}a} & \texttt{a} & \texttt{a} & \texttt{a} & \texttt{0123} & \texttt{0} & \texttt{+} & \texttt{S} & \texttt{none}\\
N1148 & F4.3027 & \texttt{T-{}-{}} & 1 & \texttt{a\textbar{}b\textbar{}a\textbar{}a} & \texttt{ab} & \texttt{ab} & \texttt{ab} & \texttt{0123} & \texttt{0} & \texttt{+} & \texttt{.} & \texttt{b@1}\\
N1149 & F4.3027 & \texttt{T-{}-{}} & 1 & \texttt{a\textbar{}a\textbar{}b\textbar{}b} & \texttt{ab} & \texttt{ab} & \texttt{ab} & \texttt{0123} & \texttt{0} & \texttt{+} & \texttt{.} & \texttt{a@0}\\
N1150 & F4.3027 & \texttt{T-{}-{}} & 1 & \texttt{a\textbar{}b\textbar{}b\textbar{}a} & \texttt{ab} & \texttt{ab} & \texttt{ab} & \texttt{0123} & \texttt{0} & \texttt{+} & \texttt{.} & \texttt{b@1}\\
N1151 & F4.3027 & \texttt{T-{}-{}} & 1 & \texttt{a\textbar{}a\textbar{}a\textbar{}a} & \texttt{a} & \texttt{a} & \texttt{a} & \texttt{0123} & \texttt{0} & \texttt{+} & \texttt{S} & \texttt{none}\\
N1152 & F4.3028 & \texttt{T-{}-{}} & 1 & \texttt{a\textbar{}b\textbar{}b\textbar{}b} & \texttt{ab} & \texttt{ab} & \texttt{ab} & \texttt{0123} & \texttt{0} & \texttt{+} & \texttt{.} & \texttt{a@0}\\
N1153 & F4.3028 & \texttt{T-{}-{}} & 1 & \texttt{a\textbar{}b\textbar{}b\textbar{}a} & \texttt{ab} & \texttt{ab} & \texttt{ab} & \texttt{0123} & \texttt{0} & \texttt{+} & \texttt{.} & \texttt{b@1}\\
N1154 & F4.3028 & \texttt{T-{}-{}} & 1 & \texttt{a\textbar{}a\textbar{}a\textbar{}a} & \texttt{a} & \texttt{a} & \texttt{a} & \texttt{0123} & \texttt{0} & \texttt{+} & \texttt{S} & \texttt{none}\\
N1155 & F4.3029 & \texttt{T-{}-{}} & 1 & \texttt{a\textbar{}a\textbar{}b\textbar{}b} & \texttt{ab} & \texttt{ab} & \texttt{ab} & \texttt{0123} & \texttt{0} & \texttt{+} & \texttt{.} & \texttt{a@0}\\
N1156 & F4.3029 & \texttt{T-{}-{}} & 1 & \texttt{a\textbar{}b\textbar{}a\textbar{}b} & \texttt{ab} & \texttt{ab} & \texttt{ab} & \texttt{0123} & \texttt{0} & \texttt{+} & \texttt{.} & \texttt{a@0}\\
N1157 & F4.3029 & \texttt{T-{}-{}} & 1 & \texttt{a\textbar{}b\textbar{}b\textbar{}a} & \texttt{ab} & \texttt{ab} & \texttt{ab} & \texttt{0123} & \texttt{0} & \texttt{+} & \texttt{.} & \texttt{b@1}\\
N1158 & F4.3029 & \texttt{T-{}-{}} & 1 & \texttt{a\textbar{}a\textbar{}a\textbar{}b} & \texttt{ab} & \texttt{ab} & \texttt{ab} & \texttt{0123} & \texttt{0} & \texttt{+} & \texttt{.} & \texttt{b@3}\\
N1159 & F4.3029 & \texttt{T-{}-{}} & 1 & \texttt{a\textbar{}a\textbar{}a\textbar{}a} & \texttt{a} & \texttt{a} & \texttt{a} & \texttt{0123} & \texttt{0} & \texttt{+} & \texttt{S} & \texttt{none}\\
N1160 & F4.3030 & \texttt{T-{}-{}} & 1 & \texttt{a\textbar{}b\textbar{}b\textbar{}c} & \texttt{abc} & \texttt{abc} & \texttt{abc} & \texttt{0123} & \texttt{0} & \texttt{+} & \texttt{.} & \texttt{a@0}\\
N1161 & F4.3030 & \texttt{T-{}-{}} & 1 & \texttt{a\textbar{}a\textbar{}a\textbar{}b} & \texttt{ab} & \texttt{ab} & \texttt{ab} & \texttt{0123} & \texttt{0} & \texttt{+} & \texttt{.} & \texttt{b@3}\\
N1162 & F4.3030 & \texttt{T-{}-{}} & 1 & \texttt{a\textbar{}b\textbar{}a\textbar{}b} & \texttt{ab} & \texttt{ab} & \texttt{ab} & \texttt{0123} & \texttt{0} & \texttt{+} & \texttt{.} & \texttt{a@0}\\
N1163 & F4.3030 & \texttt{T-{}-{}} & 1 & \texttt{a\textbar{}b\textbar{}b\textbar{}a} & \texttt{ab} & \texttt{ab} & \texttt{ab} & \texttt{0123} & \texttt{0} & \texttt{+} & \texttt{.} & \texttt{b@1}\\
N1164 & F4.3030 & \texttt{T-{}-{}} & 1 & \texttt{a\textbar{}a\textbar{}a\textbar{}a} & \texttt{a} & \texttt{a} & \texttt{a} & \texttt{0123} & \texttt{0} & \texttt{+} & \texttt{S} & \texttt{none}\\
N1165 & F4.3031 & \texttt{T-{}-{}} & 1 & \texttt{a\textbar{}b\textbar{}b\textbar{}c} & \texttt{abc} & \texttt{abc} & \texttt{abc} & \texttt{0123} & \texttt{0} & \texttt{+} & \texttt{.} & \texttt{a@0}\\
N1166 & F4.3031 & \texttt{T-{}-{}} & 1 & \texttt{a\textbar{}b\textbar{}b\textbar{}b} & \texttt{ab} & \texttt{ab} & \texttt{ab} & \texttt{0123} & \texttt{0} & \texttt{+} & \texttt{.} & \texttt{a@0}\\
N1167 & F4.3031 & \texttt{T-{}-{}} & 1 & \texttt{a\textbar{}a\textbar{}b\textbar{}b} & \texttt{ab} & \texttt{ab} & \texttt{ab} & \texttt{0123} & \texttt{0} & \texttt{+} & \texttt{.} & \texttt{a@0}\\
N1168 & F4.3031 & \texttt{T-{}-{}} & 1 & \texttt{a\textbar{}b\textbar{}a\textbar{}b} & \texttt{ab} & \texttt{ab} & \texttt{ab} & \texttt{0123} & \texttt{0} & \texttt{+} & \texttt{.} & \texttt{a@0}\\
N1169 & F4.3031 & \texttt{T-{}-{}} & 1 & \texttt{a\textbar{}b\textbar{}b\textbar{}a} & \texttt{ab} & \texttt{ab} & \texttt{ab} & \texttt{0123} & \texttt{0} & \texttt{+} & \texttt{.} & \texttt{b@1}\\
N1170 & F4.3031 & \texttt{T-{}-{}} & 1 & \texttt{a\textbar{}a\textbar{}a\textbar{}b} & \texttt{ab} & \texttt{ab} & \texttt{ab} & \texttt{0123} & \texttt{0} & \texttt{+} & \texttt{.} & \texttt{b@3}\\
N1171 & F4.3031 & \texttt{T-{}-{}} & 1 & \texttt{a\textbar{}a\textbar{}a\textbar{}a} & \texttt{a} & \texttt{a} & \texttt{a} & \texttt{0123} & \texttt{0} & \texttt{+} & \texttt{S} & \texttt{none}\\
N1172 & F4.3032 & \texttt{T-{}4} & 1 & \texttt{a\textbar{}b\textbar{}c\textbar{}d} & \texttt{abd} & \texttt{abd} & \texttt{abcd} & \texttt{013} & \texttt{0} & \texttt{-{}2} & \texttt{.} & \texttt{a@0}\\
N1173 & F4.3032 & \texttt{T-{}4} & 1 & \texttt{a\textbar{}b\textbar{}a\textbar{}c} & \texttt{abc} & \texttt{abc} & \texttt{abc} & \texttt{0123} & \texttt{0} & \texttt{+} & \texttt{.} & \texttt{b@1}\\
N1174 & F4.3032 & \texttt{T-{}4} & 1 & \texttt{a\textbar{}a\textbar{}b\textbar{}c} & \texttt{ac} & \texttt{ac} & \texttt{abc} & \texttt{013} & \texttt{0} & \texttt{-{}2} & \texttt{.} & \texttt{b@2}\\
N1175 & F4.3032 & \texttt{T-{}4} & 1 & \texttt{a\textbar{}a\textbar{}a\textbar{}b} & \texttt{ab} & \texttt{ab} & \texttt{ab} & \texttt{0123} & \texttt{0} & \texttt{+} & \texttt{.} & \texttt{b@3}\\
N1176 & F4.3032 & \texttt{T-{}4} & 1 & \texttt{a\textbar{}b\textbar{}b\textbar{}a} & \texttt{ab} & \texttt{ab} & \texttt{ab} & \texttt{0123} & \texttt{0} & \texttt{+} & \texttt{.} & \texttt{b@1}\\
N1177 & F4.3032 & \texttt{T-{}4} & 1 & \texttt{a\textbar{}a\textbar{}b\textbar{}a} & \texttt{a} & \texttt{a} & \texttt{ab} & \texttt{013} & \texttt{0} & \texttt{-{}2} & \texttt{s} & \texttt{b@2}\\
N1178 & F4.3032 & \texttt{T-{}4} & 1 & \texttt{a\textbar{}a\textbar{}a\textbar{}a} & \texttt{a} & \texttt{a} & \texttt{a} & \texttt{0123} & \texttt{0} & \texttt{+} & \texttt{S} & \texttt{none}\\
N1179 & F4.3037 & \texttt{T-{}-{}} & 1 & \texttt{a\textbar{}b\textbar{}b\textbar{}b} & \texttt{ab} & \texttt{ab} & \texttt{ab} & \texttt{0123} & \texttt{0} & \texttt{+} & \texttt{.} & \texttt{a@0}\\
N1180 & F4.3037 & \texttt{T-{}-{}} & 1 & \texttt{a\textbar{}b\textbar{}a\textbar{}b} & \texttt{ab} & \texttt{ab} & \texttt{ab} & \texttt{0123} & \texttt{0} & \texttt{+} & \texttt{.} & \texttt{a@0}\\
N1181 & F4.3037 & \texttt{T-{}-{}} & 1 & \texttt{a\textbar{}a\textbar{}a\textbar{}b} & \texttt{ab} & \texttt{ab} & \texttt{ab} & \texttt{0123} & \texttt{0} & \texttt{+} & \texttt{.} & \texttt{b@3}\\
N1182 & F4.3037 & \texttt{T-{}-{}} & 1 & \texttt{a\textbar{}a\textbar{}a\textbar{}a} & \texttt{a} & \texttt{a} & \texttt{a} & \texttt{0123} & \texttt{0} & \texttt{+} & \texttt{S} & \texttt{none}\\
N1183 & F4.3039 & \texttt{T-{}-{}} & 1 & \texttt{a\textbar{}b\textbar{}b\textbar{}c} & \texttt{abc} & \texttt{abc} & \texttt{abc} & \texttt{0123} & \texttt{0} & \texttt{+} & \texttt{.} & \texttt{a@0}\\
N1184 & F4.3039 & \texttt{T-{}-{}} & 1 & \texttt{a\textbar{}b\textbar{}b\textbar{}b} & \texttt{ab} & \texttt{ab} & \texttt{ab} & \texttt{0123} & \texttt{0} & \texttt{+} & \texttt{.} & \texttt{a@0}\\
N1185 & F4.3039 & \texttt{T-{}-{}} & 1 & \texttt{a\textbar{}a\textbar{}b\textbar{}b} & \texttt{ab} & \texttt{ab} & \texttt{ab} & \texttt{0123} & \texttt{0} & \texttt{+} & \texttt{.} & \texttt{a@0}\\
N1186 & F4.3039 & \texttt{T-{}-{}} & 1 & \texttt{a\textbar{}b\textbar{}a\textbar{}b} & \texttt{ab} & \texttt{ab} & \texttt{ab} & \texttt{0123} & \texttt{0} & \texttt{+} & \texttt{.} & \texttt{a@0}\\
N1187 & F4.3039 & \texttt{T-{}-{}} & 1 & \texttt{a\textbar{}b\textbar{}b\textbar{}a} & \texttt{ab} & \texttt{ab} & \texttt{ab} & \texttt{0123} & \texttt{0} & \texttt{+} & \texttt{.} & \texttt{b@1}\\
N1188 & F4.3039 & \texttt{T-{}-{}} & 1 & \texttt{a\textbar{}a\textbar{}a\textbar{}b} & \texttt{ab} & \texttt{ab} & \texttt{ab} & \texttt{0123} & \texttt{0} & \texttt{+} & \texttt{.} & \texttt{b@3}\\
N1189 & F4.3039 & \texttt{T-{}-{}} & 1 & \texttt{a\textbar{}a\textbar{}a\textbar{}a} & \texttt{a} & \texttt{a} & \texttt{a} & \texttt{0123} & \texttt{0} & \texttt{+} & \texttt{S} & \texttt{none}\\
N1190 & F4.3040 & \texttt{T-{}-{}} & 1 & \texttt{a\textbar{}b\textbar{}b\textbar{}c} & \texttt{abc} & \texttt{abc} & \texttt{abc} & \texttt{0123} & \texttt{0} & \texttt{+} & \texttt{.} & \texttt{a@0}\\
N1191 & F4.3040 & \texttt{T-{}-{}} & 1 & \texttt{a\textbar{}b\textbar{}b\textbar{}b} & \texttt{ab} & \texttt{ab} & \texttt{ab} & \texttt{0123} & \texttt{0} & \texttt{+} & \texttt{.} & \texttt{a@0}\\
N1192 & F4.3040 & \texttt{T-{}-{}} & 1 & \texttt{a\textbar{}b\textbar{}a\textbar{}c} & \texttt{abc} & \texttt{abc} & \texttt{abc} & \texttt{0123} & \texttt{0} & \texttt{+} & \texttt{.} & \texttt{b@1}\\
N1193 & F4.3040 & \texttt{T-{}-{}} & 1 & \texttt{a\textbar{}a\textbar{}a\textbar{}b} & \texttt{ab} & \texttt{ab} & \texttt{ab} & \texttt{0123} & \texttt{0} & \texttt{+} & \texttt{.} & \texttt{b@3}\\
N1194 & F4.3040 & \texttt{T-{}-{}} & 1 & \texttt{a\textbar{}b\textbar{}b\textbar{}a} & \texttt{ab} & \texttt{ab} & \texttt{ab} & \texttt{0123} & \texttt{0} & \texttt{+} & \texttt{.} & \texttt{b@1}\\
N1195 & F4.3040 & \texttt{T-{}-{}} & 1 & \texttt{a\textbar{}b\textbar{}a\textbar{}b} & \texttt{ab} & \texttt{ab} & \texttt{ab} & \texttt{0123} & \texttt{0} & \texttt{+} & \texttt{.} & \texttt{a@0}\\
N1196 & F4.3040 & \texttt{T-{}-{}} & 1 & \texttt{a\textbar{}a\textbar{}a\textbar{}a} & \texttt{a} & \texttt{a} & \texttt{a} & \texttt{0123} & \texttt{0} & \texttt{+} & \texttt{S} & \texttt{none}\\
N1197 & F4.3041 & \texttt{T-{}-{}} & 1 & \texttt{a\textbar{}b\textbar{}a\textbar{}c} & \texttt{abc} & \texttt{abc} & \texttt{abc} & \texttt{0123} & \texttt{0} & \texttt{+} & \texttt{.} & \texttt{b@1}\\
N1198 & F4.3041 & \texttt{T-{}-{}} & 1 & \texttt{a\textbar{}a\textbar{}a\textbar{}b} & \texttt{ab} & \texttt{ab} & \texttt{ab} & \texttt{0123} & \texttt{0} & \texttt{+} & \texttt{.} & \texttt{b@3}\\
N1199 & F4.3041 & \texttt{T-{}-{}} & 1 & \texttt{a\textbar{}a\textbar{}b\textbar{}b} & \texttt{ab} & \texttt{ab} & \texttt{ab} & \texttt{0123} & \texttt{0} & \texttt{+} & \texttt{.} & \texttt{a@0}\\
N1200 & F4.3041 & \texttt{T-{}-{}} & 1 & \texttt{a\textbar{}b\textbar{}a\textbar{}b} & \texttt{ab} & \texttt{ab} & \texttt{ab} & \texttt{0123} & \texttt{0} & \texttt{+} & \texttt{.} & \texttt{a@0}\\
N1201 & F4.3041 & \texttt{T-{}-{}} & 1 & \texttt{a\textbar{}b\textbar{}b\textbar{}a} & \texttt{ab} & \texttt{ab} & \texttt{ab} & \texttt{0123} & \texttt{0} & \texttt{+} & \texttt{.} & \texttt{b@1}\\
N1202 & F4.3041 & \texttt{T-{}-{}} & 1 & \texttt{a\textbar{}a\textbar{}a\textbar{}a} & \texttt{a} & \texttt{a} & \texttt{a} & \texttt{0123} & \texttt{0} & \texttt{+} & \texttt{S} & \texttt{none}\\
N1203 & F4.3043 & \texttt{T-{}-{}} & 1 & \texttt{a\textbar{}b\textbar{}c\textbar{}c} & \texttt{abc} & \texttt{abc} & \texttt{abc} & \texttt{0123} & \texttt{0} & \texttt{+} & \texttt{.} & \texttt{a@0}\\
N1204 & F4.3043 & \texttt{T-{}-{}} & 1 & \texttt{a\textbar{}b\textbar{}a\textbar{}c} & \texttt{abc} & \texttt{abc} & \texttt{abc} & \texttt{0123} & \texttt{0} & \texttt{+} & \texttt{.} & \texttt{b@1}\\
N1205 & F4.3043 & \texttt{T-{}-{}} & 1 & \texttt{a\textbar{}b\textbar{}a\textbar{}a} & \texttt{ab} & \texttt{ab} & \texttt{ab} & \texttt{0123} & \texttt{0} & \texttt{+} & \texttt{.} & \texttt{b@1}\\
N1206 & F4.3043 & \texttt{T-{}-{}} & 1 & \texttt{a\textbar{}a\textbar{}b\textbar{}b} & \texttt{ab} & \texttt{ab} & \texttt{ab} & \texttt{0123} & \texttt{0} & \texttt{+} & \texttt{.} & \texttt{a@0}\\
N1207 & F4.3043 & \texttt{T-{}-{}} & 1 & \texttt{a\textbar{}b\textbar{}b\textbar{}a} & \texttt{ab} & \texttt{ab} & \texttt{ab} & \texttt{0123} & \texttt{0} & \texttt{+} & \texttt{.} & \texttt{b@1}\\
N1208 & F4.3043 & \texttt{T-{}-{}} & 1 & \texttt{a\textbar{}a\textbar{}a\textbar{}b} & \texttt{ab} & \texttt{ab} & \texttt{ab} & \texttt{0123} & \texttt{0} & \texttt{+} & \texttt{.} & \texttt{b@3}\\
N1209 & F4.3043 & \texttt{T-{}-{}} & 1 & \texttt{a\textbar{}a\textbar{}a\textbar{}a} & \texttt{a} & \texttt{a} & \texttt{a} & \texttt{0123} & \texttt{0} & \texttt{+} & \texttt{S} & \texttt{none}\\
\bottomrule\end{longtable}}


\clearpage
\section{Part B: checked derivations}\label{app:derivations}

\subsection{The system}

Derivations are written in a labelled natural-deduction system for second-order modal logic \citep{Vigano2000}, in the dependency style of \citet{Lemmon1965}. A line is either a labelled formula $w:\varphi$, read \enquote{$\varphi$ holds at the world $w$}, or an accessibility statement $w\,R\,v$. Every line depends on a set of assumption lines, which the checker computes; a derivation proves its last line if that line depends on nothing. The rules are the classical introduction and elimination rules for the connectives and for first- and second-order quantifiers (eigenvariable conditions checked against the dependencies), reductio and ex falso, and the modal rules: $\nec$I discharges an assumption $w\,R\,v$ and infers $w:\nec\varphi$ from $v:\varphi$, provided $v$ is fresh; $\nec$E infers $v:\varphi$ from $w:\nec\varphi$ and $w\,R\,v$; $\pos$I and $\pos$E are dual. Def replaces a formula by one that is equal to it after unfolding Definitions~\ref{d:godlike}--\ref{d:ne} (or Scott's definitions) and $\beta$-reduction. Ax introduces an axiom at any label. Rel restates a line that depends on nothing at another label. Sym infers $v\,R\,w$ from $w\,R\,v$; it is the only relational rule used. Second-order instantiation ($\forall^2$E, $\exists^2$I) is restricted to property terms admissible in the models of Part~A: property variables, the defined properties, and $\lambda$-terms built from atoms by Boolean connectives, modal operators and individual quantifiers, with at most one free individual variable in each subformula.

The derivations are written as data in \texttt{gen/derivations.py} and checked by \texttt{gen/nd.py}, which verifies for each line that the cited lines exist and precede it, that the rule applies to them, that eigenvariable and freshness conditions hold, that discharged lines are assumptions, and that instantiation terms are admissible; it computes the dependency set shown in the second column. The program was also run on corrupted copies of the derivations (a symmetry step replaced by an assumption, a violated eigenvariable condition, a box introduction on a non-fresh world, an inadmissible comprehension term), and rejected each. The tables below are typeset from the checked data.

\subsection{Derivation B.1: the main theorem of AF}

\csLinesAf{} lines; open assumptions: none; relational rules: Sym; axioms cited: F2, F3, A3, A5. Lines 1--19 prove Lemma~\ref{l:poss} at $@$; lines 22--26 and 28--32 are Lemma~\ref{l:rigid}, at $v$ and at $@$; lines 37--66 prove Lemma~\ref{l:essence}; lines 67--70 prove Lemma~\ref{l:ne}; the last line is Theorem~\ref{t:main} at $@$.

% generated by gen/generate.py from gen/derivations.py; every line checked by gen/nd.py
{\footnotesize\setlength{\tabcolsep}{4pt}
\begin{longtable}{@{}r>{\raggedright}p{0.12\textwidth}p{0.55\textwidth}>{\raggedright\arraybackslash}p{0.21\textwidth}@{}}\toprule
no. & depends on & line & rule\\\midrule\endhead
1 & -- & $@:\ \mathcal{P}(G)$ \hfill{\itshape G is positive} & Ax A3\\
2 & 2 & $@:\ \neg \Diamond \exists x\, G(x)$ \hfill{\itshape for reductio} & A\\
3 & 3 & $@\,R\,v$ & A\\
4 & 4 & $v:\ G(a)$ & A\\
5 & 4 & $v:\ \exists x\, G(x)$ & $\exists$I 4 $[a]$\\
6 & 3,4 & $@:\ \Diamond \exists x\, G(x)$ & $\Diamond$I 3,5\\
7 & 2,3,4 & $@:\ \bot$ & $\neg$E 6,2\\
8 & 2,3,4 & $v:\ \bot$ & EFQ 7\\
9 & 2,3 & $v:\ G(a) \to \bot$ & $\to$I 4,8\\
10 & 2,3 & $v:\ \forall x\, (G(x) \to \bot)$ & $\forall$I 9 $(a)$\\
11 & 2 & $@:\ \Box \forall x\, (G(x) \to \bot)$ & $\Box$I 3,10\\
12 & -- & $@:\ \forall Y\, \forall Z\, (\mathcal{P}(Y) \wedge \Box \forall x\, (Y(x) \to Z(x)) \to \mathcal{P}(Z))$ & Ax F2\\
13 & -- & $@:\ \forall Z\, (\mathcal{P}(G) \wedge \Box \forall x\, (G(x) \to Z(x)) \to \mathcal{P}(Z))$ & $\forall^2$E 12 $[G]$\\
14 & -- & $@:\ \mathcal{P}(G) \wedge \Box \forall x\, (G(x) \to \bot) \to \mathcal{P}(\lambda x.\,\bot)$ & $\forall^2$E 13 $[\lambda x.\,\bot]$\\
15 & 2 & $@:\ \mathcal{P}(G) \wedge \Box \forall x\, (G(x) \to \bot)$ & $\wedge$I 1,11\\
16 & 2 & $@:\ \mathcal{P}(\lambda x.\,\bot)$ & $\to$E 14,15\\
17 & -- & $@:\ \neg \mathcal{P}(\lambda x.\,\bot)$ & Ax F3\\
18 & 2 & $@:\ \bot$ & $\neg$E 16,17\\
19 & -- & $@:\ \Diamond \exists x\, G(x)$ \hfill{\itshape Lemma 4.1} & RAA 2,18\\
20 & 20 & $@\,R\,v$ & A\\
21 & 21 & $v:\ \exists x\, G(x)$ & A\\
22 & 22 & $v:\ G(g)$ & A\\
23 & 22 & $v:\ \forall Y\, (\mathcal{P}(Y) \leftrightarrow \Box Y(g))$ & Def 22\\
24 & 22 & $v:\ \mathcal{P}(G) \leftrightarrow \Box G(g)$ & $\forall^2$E 23 $[G]$\\
25 & -- & $v:\ \mathcal{P}(G)$ & Ax A3\\
26 & 22 & $v:\ \Box G(g)$ & $\leftrightarrow$E 24,25\\
27 & 20 & $v\,R\,@$ \hfill{\itshape symmetry} & Sym 20\\
28 & 20,22 & $@:\ G(g)$ & $\Box$E 26,27\\
29 & 20,22 & $@:\ \forall Y\, (\mathcal{P}(Y) \leftrightarrow \Box Y(g))$ & Def 28\\
30 & 20,22 & $@:\ \mathcal{P}(G) \leftrightarrow \Box G(g)$ & $\forall^2$E 29 $[G]$\\
31 & -- & $@:\ \mathcal{P}(G)$ & Ax A3\\
32 & 20,22 & $@:\ \Box G(g)$ & $\leftrightarrow$E 30,31\\
33 & 33 & $@\,R\,u$ & A\\
34 & 20,22,33 & $u:\ G(g)$ & $\Box$E 32,33\\
35 & 20,22,33 & $u:\ \exists x\, G(x)$ & $\exists$I 34 $[g]$\\
36 & 20,22 & $@:\ \Box \exists x\, G(x)$ & $\Box$I 33,35\\
37 & 37 & $@:\ \Box Z(g)$ & A\\
38 & 20,22 & $@:\ \mathcal{P}(Z) \leftrightarrow \Box Z(g)$ & $\forall^2$E 29 $[Z]$\\
39 & 20,22,37 & $@:\ \mathcal{P}(Z)$ & $\leftrightarrow$E 38,37\\
40 & 40 & $@\,R\,u$ & A\\
41 & 41 & $u:\ G(y)$ & A\\
42 & 41 & $u:\ \forall Y\, (\mathcal{P}(Y) \leftrightarrow \Box Y(y))$ & Def 41\\
43 & 41 & $u:\ \mathcal{P}(G) \leftrightarrow \Box G(y)$ & $\forall^2$E 42 $[G]$\\
44 & -- & $u:\ \mathcal{P}(G)$ & Ax A3\\
45 & 41 & $u:\ \Box G(y)$ & $\leftrightarrow$E 43,44\\
46 & 40 & $u\,R\,@$ & Sym 40\\
47 & 40,41 & $@:\ G(y)$ & $\Box$E 45,46\\
48 & 40,41 & $@:\ \forall Y\, (\mathcal{P}(Y) \leftrightarrow \Box Y(y))$ & Def 47\\
49 & 40,41 & $@:\ \mathcal{P}(Z) \leftrightarrow \Box Z(y)$ & $\forall^2$E 48 $[Z]$\\
50 & 20,22,37,40,41 & $@:\ \Box Z(y)$ & $\leftrightarrow$E 49,39\\
51 & 20,22,37,40,41 & $u:\ Z(y)$ & $\Box$E 50,40\\
52 & 20,22,37,40 & $u:\ G(y) \to Z(y)$ & $\to$I 41,51\\
53 & 20,22,37,40 & $u:\ \forall y\, (G(y) \to Z(y))$ & $\forall$I 52 $(y)$\\
54 & 20,22,37 & $@:\ \Box \forall y\, (G(y) \to Z(y))$ & $\Box$I 40,53\\
55 & 20,22 & $@:\ \Box Z(g) \to \Box \forall y\, (G(y) \to Z(y))$ & $\to$I 37,54\\
56 & 56 & $@:\ \Box \forall y\, (G(y) \to Z(y))$ & A\\
57 & 57 & $@\,R\,u$ & A\\
58 & 56,57 & $u:\ \forall y\, (G(y) \to Z(y))$ & $\Box$E 56,57\\
59 & 56,57 & $u:\ G(g) \to Z(g)$ & $\forall$E 58 $[g]$\\
60 & 20,22,57 & $u:\ G(g)$ & $\Box$E 32,57\\
61 & 20,22,56,57 & $u:\ Z(g)$ & $\to$E 59,60\\
62 & 20,22,56 & $@:\ \Box Z(g)$ & $\Box$I 57,61\\
63 & 20,22 & $@:\ \Box \forall y\, (G(y) \to Z(y)) \to \Box Z(g)$ & $\to$I 56,62\\
64 & 20,22 & $@:\ \Box Z(g) \leftrightarrow \Box \forall y\, (G(y) \to Z(y))$ & $\leftrightarrow$I 55,63\\
65 & 20,22 & $@:\ \forall Z\, (\Box Z(g) \leftrightarrow \Box \forall y\, (G(y) \to Z(y)))$ & $\forall^2$I 64 $(Z)$\\
66 & 20,22 & $@:\ \mathrm{Ess}(G,g)$ \hfill{\itshape Lemma 4.3} & Def 65\\
67 & 22 & $v:\ \mathcal{P}(\mathit{NE}) \leftrightarrow \Box \mathit{NE}(g)$ & $\forall^2$E 23 $[\mathit{NE}]$\\
68 & -- & $v:\ \mathcal{P}(\mathit{NE})$ & Ax A5\\
69 & 22 & $v:\ \Box \mathit{NE}(g)$ & $\leftrightarrow$E 67,68\\
70 & 20,22 & $@:\ \mathit{NE}(g)$ \hfill{\itshape Lemma 4.4} & $\Box$E 69,27\\
71 & 20,22 & $@:\ G(g) \wedge \mathrm{Ess}(G,g)$ & $\wedge$I 28,66\\
72 & 20,22 & $@:\ (G(g) \wedge \mathrm{Ess}(G,g)) \wedge \mathit{NE}(g)$ & $\wedge$I 71,70\\
73 & 20,22 & $@:\ \exists x\, ((G(x) \wedge \mathrm{Ess}(G,x)) \wedge \mathit{NE}(x))$ & $\exists$I 72 $[g]$\\
74 & 20,22 & $@:\ \Box \exists x\, G(x) \wedge \exists x\, ((G(x) \wedge \mathrm{Ess}(G,x)) \wedge \mathit{NE}(x))$ & $\wedge$I 36,73\\
75 & 20,21 & $@:\ \Box \exists x\, G(x) \wedge \exists x\, ((G(x) \wedge \mathrm{Ess}(G,x)) \wedge \mathit{NE}(x))$ & $\exists$E 21,22,74 $(g)$\\
76 & -- & $@:\ \Box \exists x\, G(x) \wedge \exists x\, ((G(x) \wedge \mathrm{Ess}(G,x)) \wedge \mathit{NE}(x))$ \hfill{\itshape Theorem 4.5} & $\Diamond$E 19,20,21,75\\
\bottomrule\end{longtable}}


\subsection{Derivation B.2: modal collapse in Scott's system}

\csLinesScott{} lines; open assumptions: none; relational rules: Sym; axioms cited: S1, S2 (written F2), S3, S4 (written A4), S5. Lines 1--21 prove that possibly something is Scott-godlike; lines 22--50 that Scott-godlikeness is the Scott essence of any Scott-godlike being; lines 51--62 that a Scott-godlike being has necessary existence; lines 63--72 that something is Scott-godlike, at an arbitrary world; lines 73--92 derive the collapse of every proposition of the form $\exists x\,Y(x)$, the comprehension instance at line~80 being $\lambda z.\,\exists x\,Y(x)$.

% generated by gen/generate.py from gen/derivations.py; every line checked by gen/nd.py
{\footnotesize\setlength{\tabcolsep}{4pt}
\begin{longtable}{@{}r>{\raggedright}p{0.12\textwidth}p{0.55\textwidth}>{\raggedright\arraybackslash}p{0.21\textwidth}@{}}\toprule
no. & depends on & line & rule\\\midrule\endhead
1 & -- & $w:\ \mathcal{P}(G_{\!S})$ & Ax S3\\
2 & 2 & $w:\ \neg \Diamond \exists x\, G_{\!S}(x)$ & A\\
3 & 3 & $w\,R\,v$ & A\\
4 & 4 & $v:\ G_{\!S}(a)$ & A\\
5 & 4 & $v:\ \exists x\, G_{\!S}(x)$ & $\exists$I 4 $[a]$\\
6 & 3,4 & $w:\ \Diamond \exists x\, G_{\!S}(x)$ & $\Diamond$I 3,5\\
7 & 2,3,4 & $w:\ \bot$ & $\neg$E 6,2\\
8 & 2,3,4 & $v:\ \neg G_{\!S}(a)$ & EFQ 7\\
9 & 2,3 & $v:\ G_{\!S}(a) \to \neg G_{\!S}(a)$ & $\to$I 4,8\\
10 & 2,3 & $v:\ \forall x\, (G_{\!S}(x) \to \neg G_{\!S}(x))$ & $\forall$I 9 $(a)$\\
11 & 2 & $w:\ \Box \forall x\, (G_{\!S}(x) \to \neg G_{\!S}(x))$ & $\Box$I 3,10\\
12 & -- & $w:\ \forall Y\, \forall Z\, (\mathcal{P}(Y) \wedge \Box \forall x\, (Y(x) \to Z(x)) \to \mathcal{P}(Z))$ & Ax F2\\
13 & -- & $w:\ \forall Z\, (\mathcal{P}(G_{\!S}) \wedge \Box \forall x\, (G_{\!S}(x) \to Z(x)) \to \mathcal{P}(Z))$ & $\forall^2$E 12 $[G_{\!S}]$\\
14 & -- & $w:\ \mathcal{P}(G_{\!S}) \wedge \Box \forall x\, (G_{\!S}(x) \to \neg G_{\!S}(x)) \to \mathcal{P}(\lambda x.\,\neg G_{\!S}(x))$ & $\forall^2$E 13 $[\lambda x.\,\neg G_{\!S}(x)]$\\
15 & 2 & $w:\ \mathcal{P}(G_{\!S}) \wedge \Box \forall x\, (G_{\!S}(x) \to \neg G_{\!S}(x))$ & $\wedge$I 1,11\\
16 & 2 & $w:\ \mathcal{P}(\lambda x.\,\neg G_{\!S}(x))$ & $\to$E 14,15\\
17 & -- & $w:\ \forall Y\, (\mathcal{P}(\lambda x.\,\neg Y(x)) \leftrightarrow \neg \mathcal{P}(Y))$ & Ax S1\\
18 & -- & $w:\ \mathcal{P}(\lambda x.\,\neg G_{\!S}(x)) \leftrightarrow \neg \mathcal{P}(G_{\!S})$ & $\forall^2$E 17 $[G_{\!S}]$\\
19 & 2 & $w:\ \neg \mathcal{P}(G_{\!S})$ & $\leftrightarrow$E 18,16\\
20 & 2 & $w:\ \bot$ & $\neg$E 1,19\\
21 & -- & $w:\ \Diamond \exists x\, G_{\!S}(x)$ \hfill{\itshape possibility} & RAA 2,20\\
22 & 22 & $w:\ G_{\!S}(g)$ & A\\
23 & 22 & $w:\ \forall Y\, (\mathcal{P}(Y) \to Y(g))$ & Def 22\\
24 & 24 & $w:\ Z(g)$ & A\\
25 & 25 & $w:\ \neg \mathcal{P}(Z)$ & A\\
26 & -- & $w:\ \forall Y\, (\mathcal{P}(\lambda x.\,\neg Y(x)) \leftrightarrow \neg \mathcal{P}(Y))$ & Ax S1\\
27 & -- & $w:\ \mathcal{P}(\lambda x.\,\neg Z(x)) \leftrightarrow \neg \mathcal{P}(Z)$ & $\forall^2$E 26 $[Z]$\\
28 & 25 & $w:\ \mathcal{P}(\lambda x.\,\neg Z(x))$ & $\leftrightarrow$E 27,25\\
29 & 22 & $w:\ \mathcal{P}(\lambda x.\,\neg Z(x)) \to \neg Z(g)$ & $\forall^2$E 23 $[\lambda x.\,\neg Z(x)]$\\
30 & 22,25 & $w:\ \neg Z(g)$ & $\to$E 29,28\\
31 & 22,24,25 & $w:\ \bot$ & $\neg$E 24,30\\
32 & 22,24 & $w:\ \mathcal{P}(Z)$ & RAA 25,31\\
33 & -- & $w:\ \forall Y\, (\mathcal{P}(Y) \to \Box \mathcal{P}(Y))$ & Ax A4\\
34 & -- & $w:\ \mathcal{P}(Z) \to \Box \mathcal{P}(Z)$ & $\forall^2$E 33 $[Z]$\\
35 & 22,24 & $w:\ \Box \mathcal{P}(Z)$ & $\to$E 34,32\\
36 & 36 & $w\,R\,u$ & A\\
37 & 22,24,36 & $u:\ \mathcal{P}(Z)$ & $\Box$E 35,36\\
38 & 38 & $u:\ G_{\!S}(y)$ & A\\
39 & 38 & $u:\ \forall Y\, (\mathcal{P}(Y) \to Y(y))$ & Def 38\\
40 & 38 & $u:\ \mathcal{P}(Z) \to Z(y)$ & $\forall^2$E 39 $[Z]$\\
41 & 22,24,36,38 & $u:\ Z(y)$ & $\to$E 40,37\\
42 & 22,24,36 & $u:\ G_{\!S}(y) \to Z(y)$ & $\to$I 38,41\\
43 & 22,24,36 & $u:\ \forall y\, (G_{\!S}(y) \to Z(y))$ & $\forall$I 42 $(y)$\\
44 & 22,24 & $w:\ \Box \forall y\, (G_{\!S}(y) \to Z(y))$ & $\Box$I 36,43\\
45 & 22 & $w:\ Z(g) \to \Box \forall y\, (G_{\!S}(y) \to Z(y))$ & $\to$I 24,44\\
46 & 22 & $w:\ \forall Z\, (Z(g) \to \Box \forall y\, (G_{\!S}(y) \to Z(y)))$ & $\forall^2$I 45 $(Z)$\\
47 & 22 & $w:\ G_{\!S}(g) \wedge \forall Z\, (Z(g) \to \Box \forall y\, (G_{\!S}(y) \to Z(y)))$ & $\wedge$I 22,46\\
48 & 22 & $w:\ \mathrm{Ess}_S(G_{\!S},g)$ & Def 47\\
49 & -- & $w:\ G_{\!S}(g) \to \mathrm{Ess}_S(G_{\!S},g)$ & $\to$I 22,48\\
50 & -- & $w:\ \forall x\, (G_{\!S}(x) \to \mathrm{Ess}_S(G_{\!S},x))$ \hfill{\itshape essence lemma} & $\forall$I 49 $(g)$\\
51 & 51 & $w:\ G_{\!S}(g)$ & A\\
52 & 51 & $w:\ \forall Y\, (\mathcal{P}(Y) \to Y(g))$ & Def 51\\
53 & 51 & $w:\ \mathcal{P}(\mathit{NE}_{\!S}) \to \mathit{NE}_{\!S}(g)$ & $\forall^2$E 52 $[\mathit{NE}_{\!S}]$\\
54 & -- & $w:\ \mathcal{P}(\mathit{NE}_{\!S})$ & Ax S5\\
55 & 51 & $w:\ \mathit{NE}_{\!S}(g)$ & $\to$E 53,54\\
56 & 51 & $w:\ \forall Y\, (\mathrm{Ess}_S(Y,g) \to \Box \exists y\, Y(y))$ & Def 55\\
57 & 51 & $w:\ \mathrm{Ess}_S(G_{\!S},g) \to \Box \exists y\, G_{\!S}(y)$ & $\forall^2$E 56 $[G_{\!S}]$\\
58 & -- & $w:\ G_{\!S}(g) \to \mathrm{Ess}_S(G_{\!S},g)$ & $\forall$E 50 $[g]$\\
59 & 51 & $w:\ \mathrm{Ess}_S(G_{\!S},g)$ & $\to$E 58,51\\
60 & 51 & $w:\ \Box \exists y\, G_{\!S}(y)$ & $\to$E 57,59\\
61 & -- & $w:\ G_{\!S}(g) \to \Box \exists y\, G_{\!S}(y)$ & $\to$I 51,60\\
62 & -- & $w:\ \forall x\, (G_{\!S}(x) \to \Box \exists y\, G_{\!S}(y))$ \hfill{\itshape necessary existence lemma} & $\forall$I 61 $(g)$\\
63 & 63 & $w\,R\,v$ & A\\
64 & 64 & $v:\ \exists x\, G_{\!S}(x)$ & A\\
65 & 65 & $v:\ G_{\!S}(h)$ & A\\
66 & -- & $v:\ \forall x\, (G_{\!S}(x) \to \Box \exists y\, G_{\!S}(y))$ & Rel 62\\
67 & -- & $v:\ G_{\!S}(h) \to \Box \exists y\, G_{\!S}(y)$ & $\forall$E 66 $[h]$\\
68 & 65 & $v:\ \Box \exists y\, G_{\!S}(y)$ & $\to$E 67,65\\
69 & 63 & $v\,R\,w$ \hfill{\itshape symmetry} & Sym 63\\
70 & 63,65 & $w:\ \exists y\, G_{\!S}(y)$ & $\Box$E 68,69\\
71 & 63,64 & $w:\ \exists y\, G_{\!S}(y)$ & $\exists$E 64,65,70 $(h)$\\
72 & -- & $w:\ \exists y\, G_{\!S}(y)$ \hfill{\itshape theorem (Scott)} & $\Diamond$E 21,63,64,71\\
73 & 73 & $w:\ \exists x\, Y(x)$ & A\\
74 & 74 & $w:\ G_{\!S}(h)$ & A\\
75 & -- & $w:\ \forall x\, (G_{\!S}(x) \to \mathrm{Ess}_S(G_{\!S},x))$ & Rel 50\\
76 & -- & $w:\ G_{\!S}(h) \to \mathrm{Ess}_S(G_{\!S},h)$ & $\forall$E 75 $[h]$\\
77 & 74 & $w:\ \mathrm{Ess}_S(G_{\!S},h)$ & $\to$E 76,74\\
78 & 74 & $w:\ G_{\!S}(h) \wedge \forall Z\, (Z(h) \to \Box \forall y\, (G_{\!S}(y) \to Z(y)))$ & Def 77\\
79 & 74 & $w:\ \forall Z\, (Z(h) \to \Box \forall y\, (G_{\!S}(y) \to Z(y)))$ & $\wedge$E 78\\
80 & 74 & $w:\ \exists x\, Y(x) \to \Box \forall y\, (G_{\!S}(y) \to \exists x\, Y(x))$ \hfill{\itshape comprehension} & $\forall^2$E 79 $[\lambda z.\,\exists x\, Y(x)]$\\
81 & 73,74 & $w:\ \Box \forall y\, (G_{\!S}(y) \to \exists x\, Y(x))$ & $\to$E 80,73\\
82 & 73 & $w:\ \Box \forall y\, (G_{\!S}(y) \to \exists x\, Y(x))$ & $\exists$E 72,74,81 $(h)$\\
83 & 83 & $w\,R\,u$ & A\\
84 & 73,83 & $u:\ \forall y\, (G_{\!S}(y) \to \exists x\, Y(x))$ & $\Box$E 82,83\\
85 & -- & $u:\ \exists y\, G_{\!S}(y)$ & Rel 72\\
86 & 86 & $u:\ G_{\!S}(k)$ & A\\
87 & 73,83 & $u:\ G_{\!S}(k) \to \exists x\, Y(x)$ & $\forall$E 84 $[k]$\\
88 & 73,83,86 & $u:\ \exists x\, Y(x)$ & $\to$E 87,86\\
89 & 73,83 & $u:\ \exists x\, Y(x)$ & $\exists$E 85,86,88 $(k)$\\
90 & 73 & $w:\ \Box \exists x\, Y(x)$ & $\Box$I 83,89\\
91 & -- & $w:\ \exists x\, Y(x) \to \Box \exists x\, Y(x)$ & $\to$I 73,90\\
92 & -- & $w:\ \forall Y\, (\exists x\, Y(x) \to \Box \exists x\, Y(x))$ \hfill{\itshape collapse} & $\forall^2$I 91 $(Y)$\\
\bottomrule\end{longtable}}


\subsection{Derivation B.3: positive properties are possibly instantiated}

\csLinesPositive{} lines; open assumptions: none; relational rules: none; axioms cited: F2, F3. This is the general form of Lemma~\ref{l:poss}, and the form of Lemma~2417 that REL-01-f states.

% generated by gen/generate.py from gen/derivations.py; every line checked by gen/nd.py
{\footnotesize\setlength{\tabcolsep}{4pt}
\begin{longtable}{@{}r>{\raggedright}p{0.12\textwidth}p{0.55\textwidth}>{\raggedright\arraybackslash}p{0.21\textwidth}@{}}\toprule
no. & depends on & line & rule\\\midrule\endhead
1 & 1 & $w:\ \mathcal{P}(Y)$ & A\\
2 & 2 & $w:\ \neg \Diamond \exists x\, Y(x)$ & A\\
3 & 3 & $w\,R\,v$ & A\\
4 & 4 & $v:\ Y(a)$ & A\\
5 & 4 & $v:\ \exists x\, Y(x)$ & $\exists$I 4 $[a]$\\
6 & 3,4 & $w:\ \Diamond \exists x\, Y(x)$ & $\Diamond$I 3,5\\
7 & 2,3,4 & $w:\ \bot$ & $\neg$E 6,2\\
8 & 2,3,4 & $v:\ \bot$ & EFQ 7\\
9 & 2,3 & $v:\ Y(a) \to \bot$ & $\to$I 4,8\\
10 & 2,3 & $v:\ \forall x\, (Y(x) \to \bot)$ & $\forall$I 9 $(a)$\\
11 & 2 & $w:\ \Box \forall x\, (Y(x) \to \bot)$ & $\Box$I 3,10\\
12 & -- & $w:\ \forall Y\, \forall Z\, (\mathcal{P}(Y) \wedge \Box \forall x\, (Y(x) \to Z(x)) \to \mathcal{P}(Z))$ & Ax F2\\
13 & -- & $w:\ \forall Z\, (\mathcal{P}(Y) \wedge \Box \forall x\, (Y(x) \to Z(x)) \to \mathcal{P}(Z))$ & $\forall^2$E 12 $[Y]$\\
14 & -- & $w:\ \mathcal{P}(Y) \wedge \Box \forall x\, (Y(x) \to \bot) \to \mathcal{P}(\lambda x.\,\bot)$ & $\forall^2$E 13 $[\lambda x.\,\bot]$\\
15 & 1,2 & $w:\ \mathcal{P}(Y) \wedge \Box \forall x\, (Y(x) \to \bot)$ & $\wedge$I 1,11\\
16 & 1,2 & $w:\ \mathcal{P}(\lambda x.\,\bot)$ & $\to$E 14,15\\
17 & -- & $w:\ \neg \mathcal{P}(\lambda x.\,\bot)$ & Ax F3\\
18 & 1,2 & $w:\ \bot$ & $\neg$E 16,17\\
19 & 1 & $w:\ \Diamond \exists x\, Y(x)$ & RAA 2,18\\
20 & -- & $w:\ \mathcal{P}(Y) \to \Diamond \exists x\, Y(x)$ & $\to$I 1,19\\
21 & -- & $w:\ \forall Y\, (\mathcal{P}(Y) \to \Diamond \exists x\, Y(x))$ & $\forall^2$I 20 $(Y)$\\
\bottomrule\end{longtable}}


\subsection{Derivation B.4: Anderson's first axiom from the filter axioms}

\csLinesAnderson{} lines; open assumptions: none; relational rules: none; axioms cited: F1, F2, F3.

% generated by gen/generate.py from gen/derivations.py; every line checked by gen/nd.py
{\footnotesize\setlength{\tabcolsep}{4pt}
\begin{longtable}{@{}r>{\raggedright}p{0.12\textwidth}p{0.55\textwidth}>{\raggedright\arraybackslash}p{0.21\textwidth}@{}}\toprule
no. & depends on & line & rule\\\midrule\endhead
1 & 1 & $w:\ \mathcal{P}(Y)$ & A\\
2 & 2 & $w:\ \mathcal{P}(\lambda x.\,\neg Y(x))$ & A\\
3 & -- & $w:\ \forall Y\, \forall Z\, (\mathcal{P}(Y) \wedge \mathcal{P}(Z) \to \mathcal{P}(\lambda x.\,Y(x) \wedge Z(x)))$ & Ax F1\\
4 & -- & $w:\ \forall Z\, (\mathcal{P}(Y) \wedge \mathcal{P}(Z) \to \mathcal{P}(\lambda x.\,Y(x) \wedge Z(x)))$ & $\forall^2$E 3 $[Y]$\\
5 & -- & $w:\ \mathcal{P}(Y) \wedge \mathcal{P}(\lambda x.\,\neg Y(x)) \to \mathcal{P}(\lambda x.\,Y(x) \wedge \neg Y(x))$ & $\forall^2$E 4 $[\lambda x.\,\neg Y(x)]$\\
6 & 1,2 & $w:\ \mathcal{P}(Y) \wedge \mathcal{P}(\lambda x.\,\neg Y(x))$ & $\wedge$I 1,2\\
7 & 1,2 & $w:\ \mathcal{P}(\lambda x.\,Y(x) \wedge \neg Y(x))$ & $\to$E 5,6\\
8 & 8 & $w\,R\,v$ & A\\
9 & 9 & $v:\ Y(a) \wedge \neg Y(a)$ & A\\
10 & 9 & $v:\ Y(a)$ & $\wedge$E 9\\
11 & 9 & $v:\ \neg Y(a)$ & $\wedge$E 9\\
12 & 9 & $v:\ \bot$ & $\neg$E 10,11\\
13 & -- & $v:\ Y(a) \wedge \neg Y(a) \to \bot$ & $\to$I 9,12\\
14 & -- & $v:\ \forall x\, (Y(x) \wedge \neg Y(x) \to \bot)$ & $\forall$I 13 $(a)$\\
15 & -- & $w:\ \Box \forall x\, (Y(x) \wedge \neg Y(x) \to \bot)$ & $\Box$I 8,14\\
16 & -- & $w:\ \forall Y\, \forall Z\, (\mathcal{P}(Y) \wedge \Box \forall x\, (Y(x) \to Z(x)) \to \mathcal{P}(Z))$ & Ax F2\\
17 & -- & $w:\ \forall Z\, (\mathcal{P}(\lambda x.\,Y(x) \wedge \neg Y(x)) \wedge \Box \forall x\, (Y(x) \wedge \neg Y(x) \to Z(x)) \to \mathcal{P}(Z))$ & $\forall^2$E 16 $[\lambda x.\,Y(x) \wedge \neg Y(x)]$\\
18 & -- & $w:\ \mathcal{P}(\lambda x.\,Y(x) \wedge \neg Y(x)) \wedge \Box \forall x\, (Y(x) \wedge \neg Y(x) \to \bot) \to \mathcal{P}(\lambda x.\,\bot)$ & $\forall^2$E 17 $[\lambda x.\,\bot]$\\
19 & 1,2 & $w:\ \mathcal{P}(\lambda x.\,Y(x) \wedge \neg Y(x)) \wedge \Box \forall x\, (Y(x) \wedge \neg Y(x) \to \bot)$ & $\wedge$I 7,15\\
20 & 1,2 & $w:\ \mathcal{P}(\lambda x.\,\bot)$ & $\to$E 18,19\\
21 & -- & $w:\ \neg \mathcal{P}(\lambda x.\,\bot)$ & Ax F3\\
22 & 1,2 & $w:\ \bot$ & $\neg$E 20,21\\
23 & 1 & $w:\ \neg \mathcal{P}(\lambda x.\,\neg Y(x))$ & $\neg$I 2,22\\
24 & -- & $w:\ \mathcal{P}(Y) \to \neg \mathcal{P}(\lambda x.\,\neg Y(x))$ & $\to$I 1,23\\
25 & -- & $w:\ \forall Y\, (\mathcal{P}(Y) \to \neg \mathcal{P}(\lambda x.\,\neg Y(x)))$ & $\forall^2$I 24 $(Y)$\\
\bottomrule\end{longtable}}


\subsection{Derivation B.5: godlike beings share their necessary properties}

\csLinesIndisc{} lines; open assumptions: none; relational rules: none; axioms cited: none. The result follows from Definition~\ref{d:godlike} alone.

% generated by gen/generate.py from gen/derivations.py; every line checked by gen/nd.py
{\footnotesize\setlength{\tabcolsep}{4pt}
\begin{longtable}{@{}r>{\raggedright}p{0.12\textwidth}p{0.55\textwidth}>{\raggedright\arraybackslash}p{0.21\textwidth}@{}}\toprule
no. & depends on & line & rule\\\midrule\endhead
1 & 1 & $w:\ G(b)$ & A\\
2 & 2 & $w:\ G(c)$ & A\\
3 & 1 & $w:\ \forall Y\, (\mathcal{P}(Y) \leftrightarrow \Box Y(b))$ & Def 1\\
4 & 2 & $w:\ \forall Y\, (\mathcal{P}(Y) \leftrightarrow \Box Y(c))$ & Def 2\\
5 & 1 & $w:\ \mathcal{P}(Z) \leftrightarrow \Box Z(b)$ & $\forall^2$E 3 $[Z]$\\
6 & 2 & $w:\ \mathcal{P}(Z) \leftrightarrow \Box Z(c)$ & $\forall^2$E 4 $[Z]$\\
7 & 7 & $w:\ \Box Z(b)$ & A\\
8 & 1,7 & $w:\ \mathcal{P}(Z)$ & $\leftrightarrow$E 5,7\\
9 & 1,2,7 & $w:\ \Box Z(c)$ & $\leftrightarrow$E 6,8\\
10 & 1,2 & $w:\ \Box Z(b) \to \Box Z(c)$ & $\to$I 7,9\\
11 & 11 & $w:\ \Box Z(c)$ & A\\
12 & 2,11 & $w:\ \mathcal{P}(Z)$ & $\leftrightarrow$E 6,11\\
13 & 1,2,11 & $w:\ \Box Z(b)$ & $\leftrightarrow$E 5,12\\
14 & 1,2 & $w:\ \Box Z(c) \to \Box Z(b)$ & $\to$I 11,13\\
15 & 1,2 & $w:\ \Box Z(b) \leftrightarrow \Box Z(c)$ & $\leftrightarrow$I 10,14\\
16 & 1,2 & $w:\ \forall Z\, (\Box Z(b) \leftrightarrow \Box Z(c))$ & $\forall^2$I 15 $(Z)$\\
17 & 1 & $w:\ G(c) \to \forall Z\, (\Box Z(b) \leftrightarrow \Box Z(c))$ & $\to$I 2,16\\
18 & -- & $w:\ G(b) \to G(c) \to \forall Z\, (\Box Z(b) \leftrightarrow \Box Z(c))$ & $\to$I 1,17\\
19 & -- & $w:\ \forall y\, (G(b) \to G(y) \to \forall Z\, (\Box Z(b) \leftrightarrow \Box Z(y)))$ & $\forall$I 18 $(c)$\\
20 & -- & $w:\ \forall x\, \forall y\, (G(x) \to G(y) \to \forall Z\, (\Box Z(x) \leftrightarrow \Box Z(y)))$ & $\forall$I 19 $(b)$\\
\bottomrule\end{longtable}}


\subsection{Derivation B.6: the converse argument}

\csLinesConverse{} lines; open assumptions: none; relational rules: none; axioms cited: A3. If it is possible that nothing is godlike, nothing is godlike, on every frame.

% generated by gen/generate.py from gen/derivations.py; every line checked by gen/nd.py
{\footnotesize\setlength{\tabcolsep}{4pt}
\begin{longtable}{@{}r>{\raggedright}p{0.12\textwidth}p{0.55\textwidth}>{\raggedright\arraybackslash}p{0.21\textwidth}@{}}\toprule
no. & depends on & line & rule\\\midrule\endhead
1 & 1 & $w:\ \Diamond \neg \exists x\, G(x)$ & A\\
2 & 2 & $w:\ \exists x\, G(x)$ & A\\
3 & 3 & $w:\ G(b)$ & A\\
4 & 3 & $w:\ \forall Y\, (\mathcal{P}(Y) \leftrightarrow \Box Y(b))$ & Def 3\\
5 & 3 & $w:\ \mathcal{P}(G) \leftrightarrow \Box G(b)$ & $\forall^2$E 4 $[G]$\\
6 & -- & $w:\ \mathcal{P}(G)$ & Ax A3\\
7 & 3 & $w:\ \Box G(b)$ & $\leftrightarrow$E 5,6\\
8 & 8 & $w\,R\,v$ & A\\
9 & 9 & $v:\ \neg \exists x\, G(x)$ & A\\
10 & 3,8 & $v:\ G(b)$ & $\Box$E 7,8\\
11 & 3,8 & $v:\ \exists x\, G(x)$ & $\exists$I 10 $[b]$\\
12 & 3,8,9 & $v:\ \bot$ & $\neg$E 11,9\\
13 & 3,8,9 & $w:\ \bot$ & EFQ 12\\
14 & 1,3 & $w:\ \bot$ & $\Diamond$E 1,8,9,13\\
15 & 1,2 & $w:\ \bot$ & $\exists$E 2,3,14 $(b)$\\
16 & 1 & $w:\ \neg \exists x\, G(x)$ & $\neg$I 2,15\\
17 & -- & $w:\ \Diamond \neg \exists x\, G(x) \to \neg \exists x\, G(x)$ & $\to$I 1,16\\
\bottomrule\end{longtable}}


\end{document}
